%%%%%%%%%%%%%%%%%%%%%%%%%%%%%%%%%%%%%%%%%%%%%%%%%%%%%%%%%%%%%%%%%
\chapter{Unit 10: Polar Coordinates}
\label{chap:unit10}
%%%%%%%%%%%%%%%%%%%%%%%%%%%%%%%%%%%%%%%%%%%%%%%%%%%%%%%%%%%%%%%%%

%%===============================================================
\section{Polar Coordinates}
\label{sec:polar-coordinates}
%%===============================================================

% [Image: A polar grid with several angles labeled in

degrees]



Figure: A polar grid with several angles labeled in degrees



The \textbf{polar coordinate system} is a two-dimensional coordinate system in which each point on a plane is determined by an angle and a distance. The polar coordinate system is especially useful in situations where the relationship between two points is most easily expressed in terms of angles and distance; in the more familiar Cartesian coordinate system or rectangular coordinate system, such a relationship can only be found through trigonometric formulae.



As the coordinate system is two-dimensional, each point is determined by two polar coordinates: the radial coordinate and the angular coordinate. The radial coordinate (usually denoted as $r$) denotes the point's distance from a central point known as the \textit{pole} (equivalent to the \textit{origin} in the Cartesian system). The angular coordinate (also known as the polar angle or the azimuth angle, and usually denoted by $\theta$ or $t$) denotes the positive or anticlockwise (counterclockwise) angle required to reach the point from the 0° ray or \textit{polar axis} (which is equivalent to the positive $x$-axis in the Cartesian coordinate plane).



\paragraph{Plotting points with polar coordinates}



% [Image: The points (3,60°) and (4,210°) on a polar coordinate

system]



Figure: The points (3,60°) and (4,210°) on a polar coordinate system



Each point in the polar coordinate system can be described with the two polar coordinates, which are usually called $r$ (the radial coordinate) and $\theta$ (the angular coordinate, polar angle, or azimuth angle, sometimes represented as $\phi$ or $t$). The $r$ coordinate represents the radial distance from the pole, and the $\theta$ coordinate represents the anticlockwise (counterclockwise) angle from the $0^\circ$ ray (sometimes called the polar axis), known as the positive $x$-axis on the Cartesian coordinate plane.



For example, the polar coordinates $(3,60^\circ)$ would be plotted as a point 3 units from the pole on the $60^\circ$ ray. The coordinates $(-3,240^\circ)$ would also be plotted at this point because a negative radial distance is measured as a positive distance on the opposite ray (the ray reflected about the origin, which differs from the original ray by $180^\circ$).



One important aspect of the polar coordinate system, not present in the Cartesian coordinate system, is that a single point can be expressed with an infinite number of different coordinates. This is because any number of multiple revolutions can be made around the central pole without affecting the actual location of the point plotted. In general, the point $(r,\theta)$ can be represented as $(r,\theta\pm360^\circ k)$ or $(-r,\theta\pm(2k+1)180^\circ)$ , where $k$ is any integer.



The arbitrary coordinates $(0,\theta)$ are conventionally used to represent the pole, as regardless of the $\theta$ coordinate, a point with radius 0 will always be on the pole. To get a unique representation of a point, it is usual to limit $r$ to negative and non-negative numbers $r\ge0$ and $\theta$ to the interval $[0,360^\circ)$ or $(-180^\circ,180^\circ]$ (or, in radian measure, $[0,2\pi)$ or $(-\pi,\pi]$).



Angles in polar notation are generally expressed in either degrees or radians, using the conversion $2\pi\ \text{rad}=360^\circ$ . The choice depends largely on the context. Navigation applications use degree measure, while some physics applications (specifically rotational mechanics) and almost all mathematical literature on calculus use radian measure.



\paragraph{Converting between polar and Cartesian coordinates}



% [Image: A diagram illustrating the conversion

formulae]



Figure: A diagram illustrating the conversion formulae



The two polar coordinates $r,\theta$ can be converted to the Cartesian coordinates $x,y$ by using the trigonometric functions sine and cosine:



$$x=r\cos(\theta)$$



$$y=r\sin(\theta)$$



while the two Cartesian coordinates $x,y$ can be converted to polar coordinate $r$ by



$$r=\sqrt{x^2+y^2}$$ (by a simple application of the Pythagorean theorem).



To determine the angular coordinate $\theta$ , the following two ideas must be considered:



\begin{itemize}
    \item For $r=0$ , $\theta$ can be set to any real value.
    \item For $r\ne0$ , to get a unique representation for $\theta$ , it must be limited to an interval of size $2\pi$ . Conventional choices for such an interval are $[0,2\pi)$ and $(-\pi,\pi]$ .
\end{itemize}

To obtain $\theta$ in the interval $[0,2\pi)$ , the following may be used ($\arctan$ denotes the inverse of the tangent function):



$$\theta=

\begin{cases}

\arctan\left(\frac{y}{x}\right) & \text{if }x>0\text{ and }y \ge 0 \\\ \arctan\left(\frac{y}{x}\right)+2\pi & \text{if }x>0\text{ and }y < 0 \\\ \arctan \left(\frac{y}{x} \right)+\pi & \text{if } x < 0 \\\ \frac{\pi}{2} & \text{if }x=0\text{ and }y > 0 \\\ \frac{3\pi}{2} & \text{if }x=0\text{ and }y < 0

\end{cases}$$



To obtain $\theta$ in the interval $(-\pi,\pi]$ , the following may be used:



$$\theta=

\begin{cases}

\arctan\left(\frac{y}{x}\right) \& \text{if }x>0 \\\ \arctan\left(\frac{y}{x}\right)+\pi \& \text{if }x 0 \\\ -\frac{\pi}{2} \& \text{if }x=0\text{ and }y0$ and one for $\theta1$ , this equation defines a hyperbola.

2.  If $\epsilon=1$ , it defines a parabola.

3.  If $0<\epsilon<1$ , it defines an ellipse. The special case $\epsilon=0$ of the latter results in a circle of radius $\ell$ .



\paragraph{Resources}



\begin{itemize}
    \item \href{https://mathresearch.utsa.edu/wikiFiles/MAT1093/Polar\%20Coordinates/Esparza\%201093\%20Notes\%205.1.pdf}{Polar Coordinates}. Written notes created by Professor Esparza, UTSA.
    \item \href{https://www.youtube.com/watch?v=aSdaT62ndYE}{Polar Coordinates Basic Introduction, Conversion to Rectangular, How to Plot Points, Negative R Value} Video by The Organic Chemistry Tutor 2017
    \item \href{https://www.youtube.com/watch?v=r0fv9V9GHdo\&t=69s}{Polar Coordinates - The Basics} Video by patrickJMT 2018
\end{itemize}

\paragraph{Licensing}



Content obtained and/or adapted from:



\begin{itemize}
    \item \href{https://en.wikibooks.org/wiki/Calculus/Polar_Introduction}{Polar Introduction, Wikibooks} under a CC BY-SA license
\end{itemize}

\subsection{Learning Objectives}

\begin{enumerate}
    \item \textbf{Plot Points Using Polar Coordinates:}
    \item \textbf{Outcome:} Students will be able to plot points on a polar coordinate system using given radial and angular coordinates.
    \item \textbf{Details:} This includes understanding the relationship between the radial distance and the angle from the polar axis and being able to plot points such as $(3, 60^\circ)$ and $(-3, 240^\circ)$ on the polar grid.
    \item \textbf{Convert Between Polar and Cartesian Coordinates:}
    \item \textbf{Outcome:} Students will be able to convert coordinates between polar and Cartesian systems.
    \item \textbf{Details:} This includes using the formulas $x = r\cos(\theta)$ and $y = r\sin(\theta)$ for converting from polar to Cartesian coordinates, and $r = \sqrt{x^2 + y^2}$ and $\theta = \arctan\left(\frac{y}{x}\right)$ for converting from Cartesian to polar coordinates.
    \item \textbf{Interpret Polar Equations:}
    \item \textbf{Outcome:} Students will be able to interpret and graph equations in polar coordinates.
    \item \textbf{Details:} This includes recognizing and graphing common polar curves such as circles, lines, and roses by analyzing equations like $r = a$, $\theta = \phi$, and $r = a\cos(k\theta + \phi_0)$. Students will also be able to determine symmetries and special characteristics of these curves.
\end{enumerate}

\subsection{Assessment Instrument}

\begin{enumerate}
    \item \textbf{Assessment Instrument: Polar Coordinates Quiz} \textbf{Format:}
    \item \textbf{Type:} Mixed-format quiz (multiple-choice, short answer, and problem-solving questions)
    \item \textbf{Duration:} 45 minutes
    \item \textbf{Total Marks:} 30 points \textbf{Sections:}
    \item \textbf{Plot Points Using Polar Coordinates (10 points):}
    \item \textbf{Question 1:} Plot the following points on the polar coordinate system:
    \item a. $(4, 45^\circ)$
    \item b. $(-2, 135^\circ)$
    \item c. $(3, 210^\circ)$
    \item d. $(-5, 300^\circ)$
    \item \textbf{Convert Between Polar and Cartesian Coordinates (10 points):}
    \item \textbf{Question 2:} Convert the following polar coordinates to Cartesian coordinates:
    \item a. $(5, 60^\circ)$
    \item b. $(7, 225^\circ)$
    \item c. $(2, 120^\circ)$
    \item d. $(6, 330^\circ)$
    \item \textbf{Question 3:} Convert the following Cartesian coordinates to polar coordinates:
    \item a. $(3, 3\sqrt{3})$
    \item b. $(-4, -4)$
    \item \textbf{Interpret Polar Equations (10 points):}
    \item \textbf{Question 4:} Identify the type of curve represented by the following polar equations and sketch them:
    \item a. $r = 3$
    \item b. $\theta = 45^\circ$
    \item c. $r = 2\sin(3\theta)$
    \item d. $r = 1 + 2\cos(\theta)$ \textbf{Scoring Rubric:}
    \item \textbf{Plot Points Using Polar Coordinates:} 2.5 points for each correctly plotted point (total 10 points)
    \item \textbf{Convert Between Polar and Cartesian Coordinates:} 2 points for each correctly converted coordinate (total 10 points)
    \item \textbf{Interpret Polar Equations:} 2.5 points for each correctly identified and sketched curve (total 10 points) The quiz will help assess the students' understanding and application of the polar coordinate system. It ensures that they can correctly plot points, convert between coordinate systems, and interpret polar equations accurately.
\end{enumerate}

%%===============================================================
\section{Three-Dimensional Coordinate System}
\label{sec:three-dimensional-coordinate-system}
%%===============================================================

\paragraph{Three-Dimensional Coordinates and Vectors}



\paragraph{Basic definition}



Two-dimensional Cartesian coordinates as we've discussed so far can be easily extended to three-dimensions by adding one more value: $z$ . If the standard $(x,y)$ coordinate axes are drawn on a sheet of paper, the $z$-axis would extend upwards off of the paper.



% [Image: 3D coordinate axes.]



Figure: 3D coordinate axes.



Similar to the two coordinate axes in two-dimensional coordinates, there are three \textbf{coordinate planes} in space. These are the $xy$\textbf{-plane}, the $yz$\textbf{-plane}, and the $xz$\textbf{-plane}. Each plane is the "sheet of paper" that contains both axes the name mentions. For instance, the $yz$-plane contains both the $y$ and $z$ axes and is perpendicular to the $x$-axis.



% [Image: Coordinate planes in space.]



Figure: Coordinate planes in space.



Therefore, vectors can be extended to three dimensions by simply adding the $z$ value.





$$\mathbf{u}=\begin{pmatrix}x \\\\ y \\\\ z\end{pmatrix}$$





To facilitate standard form notation, we add another standard unit vector:





$\mathbf{k}=\begin{pmatrix}0 \\\\ 0 \\\\ 1\end{pmatrix}$





Again, both forms (component and standard) are equivalent.





$\begin{pmatrix}1 \\\\ 2 \\\\ 3\end{pmatrix}=1\mathbf{i}+2\mathbf{j}+3\mathbf{k}$





\textbf{Magnitude}: Magnitude in three dimensions is the same as in two dimensions, with the addition of a $z$ term in the radicand.





$$|\mathbf{u}|=\sqrt{u _x^2+u _y^2+u _z^2}$$





\paragraph{Three dimensions}



The polar coordinate system is extended into three dimensions with two different coordinate systems, the cylindrical and spherical coordinate systems, both of which include two-dimensional or planar polar coordinates as a subset. In essence, the cylindrical coordinate system extends polar coordinates by adding an additional distance coordinate, while the spherical system instead adds an additional angular coordinate.



\paragraph{Cylindrical coordinates}



% [Image: A point plotted with cylindrical coordinates]



Figure: A point plotted with cylindrical coordinates



The \textit{cylindrical coordinate system} is a coordinate system that essentially extends the two-dimensional polar coordinate system by adding a third coordinate measuring the height of a point above the plane, similar to the way in which the Cartesian coordinate system is extended into three dimensions. The third coordinate is usually denoted $h$ , making the three cylindrical coordinates $(r,\theta,h)$ .



The three cylindrical coordinates can be converted to Cartesian coordinates by



$$\begin{align}x & =r\cos(\theta) \\\\ y & =r\sin(\theta) \\\\ z & =h\end{align}$$



\paragraph{Spherical coordinates}



% [Image: A point plotted using spherical coordinates]



Figure: A point plotted using spherical coordinates



Polar coordinates can also be extended into three dimensions using the coordinates $(\rho,\phi,\theta)$ , where $\rho$ is the distance from the origin, $\phi$ is the angle from the $z$-axis (called the colatitude or zenith and measured from 0 to 180°) and $\theta$ is the angle from the $x$-axis (as in the polar coordinates). This coordinate system, called the \textit{spherical coordinate system}, is similar to the latitude and longitude system used for Earth, with the origin in the centre of Earth, the latitude $\delta$ being the complement of $\phi$ , determined by $\delta=90^\circ-\phi$ , and the longitude $l$ being measured by $l=\theta-180^\circ$ .



The three spherical coordinates are converted to Cartesian coordinates

by



$$\begin{align}x & =\rho\sin(\phi)\cos(\theta) \\\\ y & =\rho\sin(\phi)\sin(\theta) \\\\ z & =\rho\cos(\phi)\end{align}$$



$$\rho=\sqrt{x^2+y^2+z^2}$$



$$\theta=\arctan\left(\frac{y}{x}\right)$$



$$\phi=\arccos\left(\frac{z}{\rho}\right)$$



\paragraph{Cross Product}



The cross product of two vectors is a determinant:



$$\mathbf{u}\times\mathbf{v}=

\begin{vmatrix}

  \mathbf{i} & \mathbf{j} & \mathbf{k} \\\  u _x & u _y & u _z \\\  v _x & v _y & v _z

\end{vmatrix}$$



and is also a pseudovector.



The cross product of two vectors is orthogonal to both vectors. The magnitude of the cross product is the product of the magnitude of the vectors and the sin of the angle between them.



$$|\mathbf{u}\times\mathbf{v}|=|\mathbf{u}||\mathbf{v}|\sin(\theta)$$

This magnitude is the area of the parallelogram defined by the two vectors.



The cross product is \textit{linear} and \textit{anticommutative}. For any numbers $a,b$ ,



$$\mathbf{u}\times\left(a\mathbf{v}+b\mathbf{w}\right)=a\mathbf{u}\times\mathbf{v}+b\mathbf{u}\times\mathbf{w}\qquad\mathbf{u}\times\mathbf{v}=-\mathbf{v}\times\mathbf{u}$$



If both vectors point in the same direction, their cross product is 0.



\paragraph{Deriving the Cross Product}



% [Image: A depiction of the cross product of vectors $\mathbf{u}$ and $\mathbf{v}$.]



Figure: A depiction of the cross product of vectors $\mathbf{u}$ and $\mathbf{v}$.



Start with the following definition for the cross product: $\mathbf{u}\times\mathbf{v} = (|\mathbf{u}||\mathbf{v}|\sin\theta)\hat{\mathbf{n}}$. Vector $\hat{\mathbf{n}}$ is a unit length vector that is perpendicular to both $\mathbf{u}$ and $\mathbf{v}$ and oriented according to the "right hand rule". Angle $\theta$ is the counterclockwise angle from $\mathbf{u}$ to $\mathbf{v}$ within the plane that is orthogonal to $\hat{\mathbf{n}}$.



The formula

$\mathbf{u}\times\mathbf{v} = (u _yv _z-u _zv _y)\mathbf{i} + (u _zv _x-u _xv _z)\mathbf{j} + (u _xv _y-u _yv _x)\mathbf{k}$ can be derived by exploiting the bilinearity of the cross product. To establish the cross product as a bilinear operator, the following must be established:



1.  Holding $\mathbf{u}$ constant, the cross product $\mathbf{u}\times\mathbf{v}$ must be linear with respect to $\mathbf{v}$.

2.  Holding $\mathbf{v}$ constant, the cross product $\mathbf{u}\times\mathbf{v}$ must be linear with respect to $\mathbf{u}$.



From the definition

$\mathbf{u}\times\mathbf{v} = (|\mathbf{u}||\mathbf{v}|\sin\theta)\hat{\mathbf{n}}$, the anticommutative property of the cross product can be inferred: $\mathbf{u}\times\mathbf{v} = -\mathbf{v}\times\mathbf{u}$ (vector $\hat{\mathbf{n}}$ reverses direction when $\mathbf{u}$ and $\mathbf{v}$ are swapped). This means that linearity with respect to $\mathbf{v}$ implies linearity with respect to $\mathbf{u}$. It is hence only necessary to establish that the cross product is linear with respect to $\mathbf{v}$ to establish bilinearity.



% [Image: The perpendicular component of $\mathbf{v}$ relative to the line $L(\mathbf{u})$.]



Figure: The perpendicular component of $\mathbf{v}$ relative to the line $L(\mathbf{u})$.



$\mathbf{u}\times\mathbf{v} = (|\mathbf{u}||\mathbf{v}|\sin\theta)\hat{\mathbf{n}} = |\mathbf{u}|\textbf{rotate}(\textbf{perp}(\mathbf{v}|\mathbf{u})|\mathbf{u})$ where:



\begin{itemize}
    \item $\textbf{perp}(\mathbf{v}|\mathbf{u})$ is the component of $\mathbf{v}$ that is perpendicular to a line $L(\mathbf{u})$ whose direction is the direction of $\mathbf{u}$. $\textbf{perp}(\mathbf{v}|\mathbf{u}) = (|\mathbf{v}|\sin\theta)\hat{\mathbf{m}}$ where vector $\hat{\mathbf{m}}$ is a unit length vector that is parallel to $\textbf{perp}(\mathbf{v}|\mathbf{u})$.
    \item $\textbf{rotate}(\mathbf{w}|\mathbf{u})$ is a 90 degree counterclockwise rotation of $\mathbf{w}$ around $L(\mathbf{u})$. Note that $\textbf{rotate}(\hat{\mathbf{m}}|\mathbf{u}) = \hat{\mathbf{n}}$.
\end{itemize}

It can easily be observed that $\textbf{perp}(\mathbf{v}|\mathbf{u})$ is linear with respect to $\mathbf{v}$ with $\mathbf{u}$ held constant, and that $\textbf{rotate}(\mathbf{w}|\mathbf{u})$ is linear with respect to $\mathbf{w}$ with $\mathbf{u}$ held constant. The cross product $\mathbf{u}\times\mathbf{v} = (|\mathbf{u}||\mathbf{v}|\sin\theta)\hat{\mathbf{n}} = |\mathbf{u}|\textbf{rotate}(\textbf{perp}(\mathbf{v}|\mathbf{u})|\mathbf{u})$ is linear with respect to $\mathbf{v}$, and therefore the cross product is a bilinear operator.



The bilinearity of the cross product now enables the derivation:



$\mathbf{u}\times\mathbf{v} = (u _x\mathbf{i}+u _y\mathbf{j}+u _z\mathbf{k})\times(v _x\mathbf{i}+v _y\mathbf{j}+v _z\mathbf{k})$



$= v _x((u _x\mathbf{i}+u _y\mathbf{j}+u _z\mathbf{k})\times\mathbf{i}) + v _y((u _x\mathbf{i}+u _y\mathbf{j}+u _z\mathbf{k})\times\mathbf{j}) + v _z((u _x\mathbf{i}+u _y\mathbf{j}+u _z\mathbf{k})\times\mathbf{k})$



$= u _xv _x(\mathbf{i}\times\mathbf{i}) + u _yv _x(\mathbf{j}\times\mathbf{i}) + u _zv _x(\mathbf{k}\times\mathbf{i}) + u _xv _y(\mathbf{i}\times\mathbf{j}) + u _yv _y(\mathbf{j}\times\mathbf{j}) + u _zv _y(\mathbf{k}\times\mathbf{j}) + u _xv _z(\mathbf{i}\times\mathbf{k}) + u _yv _z(\mathbf{j}\times\mathbf{k}) + u _zv _z(\mathbf{k}\times\mathbf{k})$



$= u _xv _y\mathbf{0} + u _yv _x(-\mathbf{k}) + u _zv _x\mathbf{j} + u _xv _y\mathbf{k} + u _yv _y\mathbf{0} + u _zv _y(-\mathbf{i}) + u _xv _z(-\mathbf{j}) + u _yv _z\mathbf{i} + u _zv _z\mathbf{0}$



$= (u _yv _z-u _zv _y)\mathbf{i} + (u _zv _x-u _xv _z)\mathbf{j} + (u _xv _y-u _yv _x)\mathbf{k}$



Therefore

$\mathbf{u}\times\mathbf{v} = (u _yv _z-u _zv _y)\mathbf{i} + (u _zv _x-u _xv _z)\mathbf{j} + (u _xv _y-u _yv _x)\mathbf{k}$.



\paragraph{Triple Products}



If we have three vectors we can combine them in two ways, a triple scalar product,



$$\mathbf{u}\cdot(\mathbf{v}\times\mathbf{w})$$



and a triple vector product



$$\mathbf{u}\times(\mathbf{v}\times\mathbf{w})$$





The triple scalar product is a determinant



$$\mathbf{u}\cdot(\mathbf{v}\times\mathbf{w})=\begin{vmatrix}u _x&u _y&u _z \\\\ v _x&v _y&v _z \\\\ w _x&w _y&w _z\end{vmatrix}$$



If the three vectors are listed clockwise, looking from the origin, the sign of this product is positive. If they are listed anticlockwise the sign is negative.



The order of the cross and dot products doesn't matter.



$$\mathbf{u}\cdot(\mathbf{v}\times\mathbf{w})=(\mathbf{u}\times\mathbf{v})\cdot\mathbf{w}$$



Either way, the absolute value of this product is the volume of the parallelepiped defined by the three vectors $\mathbf{u},\mathbf{v},\mathbf{w}$



The triple vector product can be simplified



$$\mathbf{u}\times(\mathbf{v}\times\mathbf{w})=(\mathbf{u}\cdot\mathbf{w})\mathbf{v}-(\mathbf{u}\cdot\mathbf{v})\mathbf{w}$$



This form is easier to do calculations with.



The triple vector product is \textbf{not} associative.



$$\mathbf{u}\times(\mathbf{v}\times\mathbf{w})\ne(\mathbf{u}\times\mathbf{v})\times\mathbf{w}$$



There are special cases where the two sides are equal, but in general the brackets matter. They must not be omitted.



\paragraph{Three-Dimensional Lines and Planes}



We will use $\mathbf r$ to denote the position of a point.



The multiples of a vector, $\mathbf a$ all lie on a line through the origin. Adding a constant vector $\mathbf b$ will shift the line, but leave it straight, so the equation of a line is,



$$\mathbf{r}=s\mathbf{a}+\mathbf{b}$$



This is a \textit{parametric equation}. The position is specified in terms of the parameter $s$ .



Any linear combination of two vectors, $\mathbf{a},\mathbf{b}$ lies on a single plane through the origin, provided the two vectors are not colinear. We can shift this plane by a constant vector again and write



$$\mathbf{r}=s\mathbf{a}+t\mathbf{b}+\mathbf{c}$$



If we choose $\mathbf{a},\mathbf{b}$ to be \textit{orthonormal} vectors in the plane (i.e. unit vectors at right angles) then $s,t$ are Cartesian coordinates for points in the plane.



These parametric equations can be extended to higher dimensions.



Instead of giving parametric equations for the line and plane, we could use constraints. E.g., for any point in the $xy$ plane $z=0$



For a plane through the origin, the single vector normal to the plane, $\mathbf n$ , is at right angle with every vector in the plane, by definition, so



$$\mathbf{r}\cdot\mathbf{n}=0$$ is a plane through the origin, normal to

$\mathbf n$ .



For planes not through the origin we get



$$(\mathbf{r}-\mathbf{a})\cdot\mathbf{n}=0\qquad\mathbf{r}\cdot\mathbf{n}=\mathbf{a}\cdot\mathbf{n}$$



A line lies on the intersection of two planes, so it must obey the constraint for both planes, i.e.



$$\mathbf{r}\cdot\mathbf{n}=a\qquad\mathbf{r}\cdot\mathbf{m}=b$$



These constraint equations can also be extended to higher dimensions.



\paragraph{Resources}



\begin{itemize}
    \item \href{https://www.youtube.com/watch?v=EzJP9uwV3ms}{Plotting Points In a Three Dimensional Coordinate System} Video by The Organic Chemistry Tutor 2017
\end{itemize}

\paragraph{Licensing}



Content obtained and/or adapted from:



\begin{itemize}
    \item \href{https://en.wikibooks.org/wiki/Calculus/Vectors}{Vectors, Wikibooks: Calculus} under a CC BY-SA license
\end{itemize}

\subsection{Learning Objectives}

\begin{enumerate}
    \item Student Learning Objectives (SLOs) for Three-Dimensional Coordinates and Vectors Basic Definition
    \item \textbf{Identify and Plot Three-Dimensional Coordinates:}
    \item \textbf{Outcome:} Students will be able to identify and plot points in three-dimensional space using Cartesian coordinates.
    \item \textbf{Details:} This includes understanding the $x$, $y$, and $z$ axes, the three coordinate planes ($xy$-plane, $yz$-plane, $xz$-plane), and correctly plotting points based on given coordinates.
    \item \textbf{Define and Use Standard Unit Vectors in Three Dimensions:}
    \item \textbf{Outcome:} Students will be able to define the standard unit vectors $\mathbf{i}$, $\mathbf{j}$, and $\mathbf{k}$ and use them to express vectors in three-dimensional space.
    \item \textbf{Details:} This includes understanding the notation of standard unit vectors and converting between component form and standard form of vectors.
    \item \textbf{Calculate Magnitude of Three-Dimensional Vectors:}
    \item \textbf{Outcome:} Students will be able to calculate the magnitude of vectors in three-dimensional space.
    \item \textbf{Details:} This includes using the formula $|\mathbf{u}|=\sqrt{u_x^2+u_y^2+u_z^2}$ and applying it to find the length of a given vector. Cylindrical Coordinates
    \item \textbf{Convert Between Cartesian and Cylindrical Coordinates:}
    \item \textbf{Outcome:} Students will be able to convert points from Cartesian coordinates to cylindrical coordinates and vice versa.
    \item \textbf{Details:} This includes using the formulas $x = r\cos(\theta)$, $y = r\sin(\theta)$, and $z = h$ to switch between coordinate systems and understanding the significance of each coordinate in both systems.
    \item \textbf{Plot Points Using Cylindrical Coordinates:}
    \item \textbf{Outcome:} Students will be able to plot points using cylindrical coordinates $(r, \theta, h)$.
    \item \textbf{Details:} This includes correctly identifying and plotting points based on radial distance, angular displacement, and height from the plane. Spherical Coordinates
    \item \textbf{Convert Between Cartesian and Spherical Coordinates:}
    \item \textbf{Outcome:} Students will be able to convert points from Cartesian coordinates to spherical coordinates and vice versa.
    \item \textbf{Details:} This includes using the formulas $x = \rho\sin(\phi)\cos(\theta)$, $y = \rho\sin(\phi)\sin(\theta)$, $z = \rho\cos(\phi)$ and $\rho = \sqrt{x^2+y^2+z^2}$, $\theta = \arctan\left(\frac{y}{x}\right)$, $\phi = \arccos\left(\frac{z}{\rho}\right)$ to switch between coordinate systems and understanding the significance of each coordinate in both systems.
    \item \textbf{Plot Points Using Spherical Coordinates:}
    \item \textbf{Outcome:} Students will be able to plot points using spherical coordinates $(\rho, \phi, \theta)$.
    \item \textbf{Details:} This includes correctly identifying and plotting points based on radial distance from the origin, angle from the z-axis (colatitude), and angle from the x-axis (longitude).
\end{enumerate}

\subsection{Assessment Instrument}

\begin{enumerate}
    \item \textbf{Assessment Instrument: 3D Coordinate Systems Quiz} \textbf{Format:}
    \item \textbf{Type:} Mixed-format quiz (multiple-choice, short answer, and problem-solving questions)
    \item \textbf{Duration:} 45 minutes
    \item \textbf{Total Marks:} 30 points \textbf{Sections:}
    \item \textbf{Basic Definitions (10 points):}
    \item \textbf{Question 1:} Identify the coordinate plane for each of the following descriptions:
    \item a. The plane containing the x and y axes.
    \item b. The plane perpendicular to the x-axis containing the y and z axes.
    \item c. The plane perpendicular to the y-axis containing the x and z axes.
    \item \textbf{Question 2:} Given the point (3, -2, 5), specify which coordinate corresponds to the height above the $xy$-plane.
    \item \textbf{Question 3:} Plot the point (2, 3, -1) in the three-dimensional coordinate system.
    \item \textbf{Vectors and Unit Vectors (10 points):}
    \item \textbf{Question 4:} Express the vector $\mathbf{u} = \begin{pmatrix} 4 \\ 5 \\ 6 \end{pmatrix}$ in standard form notation using $\mathbf{i}$, $\mathbf{j}$, and $\mathbf{k}$.
    \item \textbf{Question 5:} Calculate the magnitude of the vector $\mathbf{v} = \begin{pmatrix} 3 \\ -4 \\ 12 \end{pmatrix}$.
    \item \textbf{Question 6:} Identify the standard unit vector that corresponds to the z-axis.
    \item \textbf{Cylindrical and Spherical Coordinates (10 points):}
    \item \textbf{Question 7:} Convert the cylindrical coordinates (5, $\frac{\pi}{4}$, 3) to Cartesian coordinates.
    \item \textbf{Question 8:} Convert the Cartesian coordinates (2, 2, 1) to spherical coordinates.
    \item \textbf{Question 9:} Plot the point given by the cylindrical coordinates (3, $\frac{\pi}{2}$, 4) on the three-dimensional coordinate system. \textbf{Scoring Rubric:}
    \item \textbf{Basic Definitions:}
    \item \textbf{Question 1:} 1 point for each correct plane identification (3 points)
    \item \textbf{Question 2:} 2 points for correct identification of the coordinate
    \item \textbf{Question 3:} 5 points for correctly plotting the point
    \item \textbf{Vectors and Unit Vectors:}
    \item \textbf{Question 4:} 3 points for correctly expressing the vector in standard form
    \item \textbf{Question 5:} 5 points for correctly calculating the magnitude
    \item \textbf{Question 6:} 2 points for correctly identifying the unit vector
    \item \textbf{Cylindrical and Spherical Coordinates:}
    \item \textbf{Question 7:} 4 points for correctly converting to Cartesian coordinates
    \item \textbf{Question 8:} 4 points for correctly converting to spherical coordinates
    \item \textbf{Question 9:} 2 points for correctly plotting the point This quiz will help assess students' understanding of three-dimensional coordinate systems, vectors, and conversions between different coordinate systems. It ensures that they can identify, plot, and convert coordinates accurately, as well as work with vectors and their magnitudes.
\end{enumerate}

%%===============================================================
\section{Vectors}
\label{sec:vectors}
%%===============================================================

"Higher-dimensional geometry" sounds exotic. It is exotic — interesting and eye-opening. But it isn't distant or unreachable.



We begin by defining one-dimensional space to be the set $\mathbb{R}^1$. To see that definition is reasonable, draw a one-dimensional space



% [Image: ]



and make the usual correspondence with $\mathbb{R}$: pick a point to label $0$ and another to label $1$.



% [Image: ]



Now, with a scale and a direction, finding the point corresponding to, say $+2.17$, is easy — start at $0$ and head in the direction of $1$ (i.e., the positive direction), but don't stop there, go $2.17$ times as far.



The basic idea here, combining magnitude with direction, is the key to extending to higher dimensions.



An object comprised of a magnitude and a direction is a \textbf{vector} (we will use the same word as in the previous section because we shall show below how to describe such an object with a column vector). We can draw a vector as having some length, and pointing somewhere.



% [Image: ]



There is a subtlety here — these vectors



% [Image: ]



are equal, even though they start in different places, because they have equal lengths and equal directions. Again: those vectors are not just alike, they are equal.



How can things that are in different places be equal? Think of a vector as representing a displacement ("vector" is Latin for "carrier" or "traveler"). These squares undergo the same displacement, despite that those displacements start in different places.



% [Image: ]



\paragraph{Free Vectors}



Sometimes, to emphasize this property vectors have of not being anchored, they are referred to as \textbf{free} vectors. Thus, these free vectors are equal as each is a displacement of one over and two up.



% [Image: ]



More generally, vectors in the plane are the same if and only if they have the same change in first components and the same change in second components: the vector extending from $(a _1,a _2)$ to $(b _1,b _2)$ equals the vector from $(c _1,c _2)$ to $(d _1,d _2)$ if and only if $b _1-a _1=d _1-c _1$ and $b _2-a _2=d _2-c _2$.



An expression like "the vector that, were it to start at $(a _1,a _2)$, would extend to $(b _1,b _2)$" is awkward. We instead describe such a vector as



$$\begin{pmatrix} b _1-a _1 \\\\ b _2-a _2 \end{pmatrix}$$



so that, for instance, the "one over and two up" arrows shown above

picture this vector.



$$\begin{pmatrix} 1 \\\\ 2 \end{pmatrix}$$



\paragraph{Canonical Position}



We often draw the arrow as starting at the origin, and we then say it is in the \textbf{canonical position} (or \textbf{natural position}). When the vector



$$\begin{pmatrix} b _1-a _1 \\\\ b _2-a _2 \end{pmatrix}$$



is in its canonical position then it extends to the endpoint $(b _1-a _1,b _2-a _2)$.



We typically just refer to "the point



$$\begin{pmatrix} 1 \\\\ 2 \end{pmatrix}$$"



rather than "the endpoint of the canonical position of" that vector.



Thus, we will call both of these sets $\mathbb{R}^2$.



$$\{(x _1,x _2)\,\big|\, x _1,x _2\in\mathbb{R}\}

\qquad

\{\begin{pmatrix} x _1 \\\\ x _2 \end{pmatrix}\,\big|\, x _1,x _2\in\mathbb{R}\}$$



\paragraph{Vector Operations}



In the prior section we defined vectors and vector operations with an algebraic motivation;



$$r\cdot\begin{pmatrix} v _1 \\\\ v _2 \end{pmatrix} = \begin{pmatrix} rv _1 \\\\ rv _2 \end{pmatrix} \qquad \begin{pmatrix} v _1 \\\\  v _2 \end{pmatrix} + \begin{pmatrix} w _1 \\\\ w _2 \end{pmatrix} = \begin{pmatrix} v _1+w _1 \\\\ v _2+w _2 \end{pmatrix}$$



we can now interpret those operations geometrically. For instance, if $\vec{v}$ represents a displacement then $3\vec{v}\,$ represents a displacement in the same direction but three times as far, and $-1\vec{v}\,$ represents a displacement of the same distance as $\vec{v}\,$ but in the opposite direction.



% [Image: ]



And, where $\vec{v}$ and $\vec{w}$ represent displacements, $\vec{v}+\vec{w}$ represents those displacements combined.



% [Image: ]



The long arrow is the combined displacement in this sense: if, in one minute, a ship's motion gives it the displacement relative to the earth of $\vec{v}$ and a passenger's motion gives a displacement relative to the ship's deck of $\vec{w}$, then $\vec{v}+\vec{w}$ is the displacement of the passenger relative to the earth.



\paragraph{Parallelogram Rule}



Another way to understand the vector sum is with the \textbf{parallelogram rule}. Draw the parallelogram formed by the vectors

$\vec{v} _1,\vec{v} _2$ and then the sum $\vec{v} _1+\vec{v} _2$ extends along the diagonal to the far corner.



% [Image: ]



The above drawings show how vectors and vector operations behave in $\mathbb{R}^2$. We can extend to $\mathbb{R}^3$, or to even higher-dimensional spaces where we have no pictures, with the obvious generalization: the free vector that, if it starts at $(a _1,\ldots,a _n)$, ends at $(b _1,\ldots,b _n)$, is represented by this column



$$\begin{pmatrix} b _1-a _1 \\\\ \vdots \\\\ b _n-a _n \end{pmatrix}$$



(vectors are equal if they have the same representation), we aren't too careful to distinguish between a point and the vector whose canonical representation ends at that point,



$$\mathbb{R}^n= \{\begin{pmatrix} v _1 \\\\ \vdots \\\\ v _n \end{pmatrix}\,\big|\, v _1,\ldots,v _n\in\mathbb{R}\}$$



and addition and scalar multiplication are component-wise.



Having considered points, we now turn to the lines.



In $\mathbb{R}^2$, the line through $(1,2)$ and $(3,1)$ s comprised of (the endpoints of) the vectors in this set



$$\{ \begin{pmatrix} 1 \\\\ 2 \end{pmatrix}+t\cdot\begin{pmatrix} 2 \\\\ -1 \end{pmatrix}\,\big|\, t\in\mathbb{R}\}$$



That description expresses this picture.



% [Image: ]



The vector associated with the parameter $t$ has its whole body in the line — it is a \textbf{direction vector} for the line. Note that points on the line to the left of $x=1$ are described using negative values of $t$.



In $\mathbb{R}^3$, the line through $(1,2,1)$ and $(2,3,2)$ is the set of (endpoints of) vectors of this form



% [Image: ]



and lines in even higher-dimensional spaces work in the same way.



If a line uses one parameter, so that there is freedom to move back and forth in one dimension, then a plane must involve two. For example, the plane through the points $(1,0,5)$, $(2,1,-3)$, and $(-2,4,0.5)$ consists of (endpoints of) the vectors in



$$\\{ \begin{pmatrix} 1 \\\\ 0 \\\\ 5 \end{pmatrix} +t\cdot\begin{pmatrix} 1 \\\\ 1 \\\\ -8 \end{pmatrix} +s\cdot\begin{pmatrix} -3 \\\\ 4 \\\\ -4.5 \end{pmatrix} \,\big|\, t,s\in\mathbb{R}      \\}$$



(the column vectors associated with the parameters



$$\begin{pmatrix} 1 \\\\ 1 \\\\ -8 \end{pmatrix} = \begin{pmatrix} 2 \\\\ 1 \\\\ -3 \end{pmatrix} - \begin{pmatrix} 1 \\\\ 0 \\\\ 5 \end{pmatrix} \qquad \begin{pmatrix} -3 \\\\ 4 \\\\ -4.5 \end{pmatrix} = \begin{pmatrix} -2 \\\\ 4 \\\\ 0.5 \end{pmatrix} - \begin{pmatrix} 1 \\\\ 0 \\\\ 5 \end{pmatrix}$$



are two vectors whose whole bodies lie in the plane). As with the line, note that some points in this plane are described with negative $t$'s or negative $s$'s or both.



A description of planes that is often encountered in algebra and calculus uses a single equation as the condition that describes the relationship among the first, second, and third coordinates of points in a plane.



% [Image: ]



The translation from such a description to the vector description that we favor in this book is to think of the condition as a one-equation linear system and parametrize $x=(1/2)(4-y-z)$.



% [Image: ]



\paragraph{Linear Surface}



Generalizing from lines and planes, we define a \textbf{$k$-dimensional linear surface} (or \textbf{$k$-flat}) in $\mathbb{R}^n$ to be $\{\vec{p}+t _1\vec{v} _1+t _2\vec{v} _2+\cdots+t _k\vec{v} _k \,\big|\, t _1,\ldots ,t _k\in\mathbb{R}\}$ where $\vec{v} _1,\ldots,\vec{v} _k\in\mathbb{R}^n$. For example, in $\mathbb{R}^4$,



$$\\{\begin{pmatrix} 2 \\\\ \pi \\\\ 3 \\\\ -0.5 \end{pmatrix} +t\begin{pmatrix} 1 \\\\ 0 \\\\ 0 \\\\ 0 \end{pmatrix} \,\big|\, t\in\mathbb{R}\\}$$



is a line,



$$\\{ \begin{pmatrix} 0 \\\\ 0 \\\\ 0 \\\\ 0 \end{pmatrix} +t\begin{pmatrix} 1 \\\\ 1 \\\\ 0 \\\\ -1 \end{pmatrix} +s\begin{pmatrix} 2 \\\\ 0 \\\\ 1 \\\\ 0 \end{pmatrix} \,\big|\, t,s\in\mathbb{R}\\}$$



is a plane, and



$$\\{ \begin{pmatrix} 3 \\\\ 1 \\\\ -2 \\\\ 0.5 \end{pmatrix} +r\begin{pmatrix} 0 \\\\ 0 \\\\ 0 \\\\ -1 \end{pmatrix} +s\begin{pmatrix} 1 \\\\ 0 \\\\ 1 \\\\ 0 \end{pmatrix} +t\begin{pmatrix} 2 \\\\ 0 \\\\ 1 \\\\ 0 \end{pmatrix} \,\big|\, r,s,t\in\mathbb{R}\\}$$



is a three-dimensional linear surface. Again, the intuition is that a line permits motion in one direction, a plane permits motion in combinations of two directions, etc.



A linear surface description can be misleading about the dimension — this



$$L=\\{ \begin{pmatrix} 1 \\\\ 0 \\\\ -1 \\\\ -2 \end{pmatrix} +t\begin{pmatrix} 1 \\\\ 1 \\\\ 0 \\\\ -1 \end{pmatrix} +s\begin{pmatrix} 2 \\\\ 2 \\\\ 0 \\\\ -2 \end{pmatrix} \,\big|\, t,s\in\mathbb{R}\\}$$



is a \textbf{degenerate} plane because it is actually a line.



$$L=\\{ \begin{pmatrix} 1 \\\\ 0 \\\\ -1 \\\\ -2 \end{pmatrix} +r\begin{pmatrix} 1 \\\\ 1 \\\\ 0 \\\\ -1 \end{pmatrix} \,\big|\, r\in\mathbb{R}\\}$$



We shall see in the Linear Independence section of Chapter Two what relationships among vectors causes the linear surface they generate to be degenerate.



We finish this subsection by restating our conclusions from the first section in geometric terms. First, the solution set of a linear system with $n$ unknowns is a linear surface in $\mathbb{R}^n$. Specifically, it is a $k$-dimensional linear surface, where $k$ is the number of free variables in an echelon form version of the system. Second, the solution set of a homogeneous linear system is a linear surface passing through the origin. Finally, we can view the general solution set of any linear system as being the solution set of its associated homogeneous system offset from the origin by a vector, namely by any particular solution.



\paragraph{Practice Problems}



\href{https://en.wikibooks.org/wiki/Linear_Algebra/Vectors_in_Space/Solutions}{Vectors in Space (With Solutions)}



\paragraph{Resources}



\begin{itemize}
    \item \href{https://math.libretexts.org/Courses/Monroe_Community_College/MTH_212_Calculus_III/Chapter_11\%3A_Vectors_and_the_Geometry_of_Space/11.1\%3A_Vectors_in_the_Plane}{Vectors in the Plane}, Mathematics LibreTexts
\end{itemize}

\paragraph{Licensing}



Content obtained and/or adapted from:



\begin{itemize}
    \item \href{https://en.wikibooks.org/wiki/Linear_Algebra/Vectors_in_Space}{Vectors in Space, Wikibooks: Linear Algebra} under a CC BY-SA license
\end{itemize}

\subsection{Learning Objectives}

\begin{enumerate}
    \item \textbf{Understanding One-Dimensional Space:}
    \item \textbf{Outcome:} Students will be able to define one-dimensional space as the set $\mathbb{R}^1$ and understand its representation on a number line.
    \item \textbf{Details:} This includes labeling key points such as 0 and 1, and locating specific points based on their magnitude and direction from the origin.
    \item \textbf{Identifying and Representing Vectors:}
    \item \textbf{Outcome:} Students will be able to identify vectors as objects with both magnitude and direction and represent them graphically.
    \item \textbf{Details:} This involves understanding free vectors, canonical position, and using column vectors to represent displacements.
    \item \textbf{Vector Operations and Their Geometric Interpretations:}
    \item \textbf{Outcome:} Students will be able to perform vector operations such as addition and scalar multiplication and interpret these operations geometrically.
    \item \textbf{Details:} This includes understanding how to combine vectors using the parallelogram rule and interpreting scalar multiplication as changing the magnitude and direction of vectors.
\end{enumerate}

\subsection{Assessment Instrument}

\begin{enumerate}
    \item \textbf{Assessment Instrument: Vectors in The Plane, Space Quiz} \textbf{Format:}
    \item \textbf{Type:} Mixed-format quiz (multiple-choice, short answer, and problem-solving questions)
    \item \textbf{Duration:} 45 minutes
    \item \textbf{Total Marks:} 30 points \textbf{Sections:}
    \item \textbf{Understanding One-Dimensional Space (10 points):}
    \item \textbf{Question 1:}
    \item a. Define one-dimensional space and explain how it can be represented on a number line. (2 points)
    \item b. Given a number line, label the points corresponding to $-1$, $0$, and $2.5$. (3 points)
    \item c. How would you find the point corresponding to $+3.5$ starting from $0$ on the number line? (2 points)
    \item d. Why is the definition of one-dimensional space using $\mathbb{R}^1$ reasonable? (3 points)
    \item \textbf{Identifying and Representing Vectors (10 points):}
    \item \textbf{Question 2:}
    \item a. Explain the concept of a free vector and how it differs from a point in space. (2 points)
    \item b. Draw the vector $\begin{pmatrix} 3 \\ 2 \end{pmatrix}$ starting from the origin on a coordinate plane. (3 points)
    \item c. If a vector $\vec{v}$ represents a displacement of "one over and two up," write it in column vector form and draw it starting from $(1,1)$. (3 points)
    \item d. Why are vectors with equal lengths and directions considered equal even if they start from different points? (2 points)
    \item \textbf{Vector Operations and Their Geometric Interpretations (10 points):}
    \item \textbf{Question 3:}
    \item a. Perform the following vector operations and interpret them geometrically:
    \item i. $2 \cdot \begin{pmatrix} 1 \\ 3 \end{pmatrix}$ (2 points)
    \item ii. $\begin{pmatrix} 1 \\ 2 \end{pmatrix} + \begin{pmatrix} 3 \\ -1 \end{pmatrix}$ (3 points)
    \item b. Use the parallelogram rule to find the sum of $\begin{pmatrix} 2 \\ 1 \end{pmatrix}$ and $\begin{pmatrix} -1 \\ 4 \end{pmatrix}$. (3 points)
    \item c. Explain the significance of scalar multiplication and vector addition in the context of vector operations. (2 points) \textbf{Scoring Rubric:}
    \item \textbf{Understanding One-Dimensional Space:} Points distributed as per each sub-question (total 10 points)
    \item \textbf{Identifying and Representing Vectors:} Points distributed as per each sub-question (total 10 points)
    \item \textbf{Vector Operations and Their Geometric Interpretations:} Points distributed as per each sub-question (total 10 points) This quiz will assess students' understanding and application of higher-dimensional geometry concepts, including one-dimensional space, vectors, and vector operations. It ensures that they can correctly define, represent, and manipulate vectors both algebraically and geometrically.
\end{enumerate}

%%===============================================================
\section{The Dot Product}
\label{sec:the-dot-product}
%%===============================================================

The dot product is a way of multiplying two vectors to produce a scalar value. Because it combines the components of two vectors to form a \textit{scalar}, it is sometimes called a scalar product. If you were asked to take the 'dot product of two rectangular vectors' you would do the following:



$$\mathbf{u}\cdot\mathbf{v}=u _xv _x+u _yv _y$$



It is very important to note that the dot product of two vectors does \textbf{not} result in another vector, it gives you a scalar, just a numerical value.



Another common pitfall may arise if your vectors are not in rectangular ('cartesian') format. Sometimes, vectors are instead expressed in polar coordinates, where the first component is the vector's magnitude (length) and the second is the angle from the $x$-axis at which the vector should be oriented. Dot products cannot be performed using the conventional method on these sorts of vectors; vectors in polar format must be converted to their equivalent rectangular form before you can work with them using the formula given above. A common way to convert to rectangular coordinates is to imagine that the vector was projected horizontally and vertically to form a right triangle. You could then use properties of sin and cos to find the length of the two legs the right triangle. The horizontal length would then be the x-component of the rectangular expression of the vector and the vertical length would be the y-component. Remember that if the vector is pointing down or to the left, the corresponding components would have to be negative to indicate that.



With some rearrangement and trigonometric manipulation, we can see that the number that results from the dot product of two vectors is a surprising and useful identity:



$$\mathbf{u}\cdot\mathbf{v}=|\mathbf{u}||\mathbf{v}|\cos(\theta)$$



where $\theta$ is the angle between the two vectors.



% [Image: Calculating bond angles of a symmetrical tetrahedral molecule such as methane using a dot product]



Figure: Calculating bond angles of a symmetrical tetrahedral molecule such as methane using a dot product



This provides a convenient way of finding the angle between two vectors:



$$\cos(\theta)=\frac{\mathbf{u}\cdot\mathbf{v}}{|\mathbf{u}||\mathbf{v}|}$$



Notice that the dot product is 'commutative', that is:



$$\mathbf{u}\cdot\mathbf{v}=\mathbf{v}\cdot\mathbf{u}$$



Also, the dot product of two vectors will be the length of the vector squared:



$$\mathbf{u}\cdot\mathbf{u}=u _xu _x+u _yu _y=(u _x)^2+(u _y)^2$$



and by the Pythagorean theorem,



$$(u _x)^2+(u _y)^2=\|u\|^2$$



The dot product can be visualized as the length of a projection of one vector on to the other. In other words, the dot product asks 'how much magnitude of this vector is going in the direction of that vector?'



\paragraph{Resources}



\begin{itemize}
    \item \href{https://www.youtube.com/watch?v=0iNrGpwZwog}{The Vector Dot Product} Video by Professor Dave Explains 2018
    \item \href{https://en.wikibooks.org/wiki/Calculus/Vectors}{Vectors}, WikiBooks: Calculus
\end{itemize}

\paragraph{Licensing}



Content obtained and/or adapted from:



\begin{itemize}
    \item \href{https://en.wikibooks.org/wiki/Calculus/Vectors}{Vectors, WikiBooks: Calculus} under a CC BY-SA license
\end{itemize}

\subsection{Learning Objectives}

\begin{enumerate}
    \item \textbf{Calculate the Dot Product of Two Vectors in Rectangular Coordinates:}
    \item \textbf{Outcome:} Students will be able to calculate the dot product of two vectors given in rectangular coordinates.
    \item \textbf{Details:} This includes understanding the formula for the dot product $$\mathbf{u}\cdot\mathbf{v}=u_xv_x + u_yv_y$$, correctly applying it to pairs of vectors, and interpreting the resulting scalar value.
    \item \textbf{Convert Polar Coordinates to Rectangular Coordinates:}
    \item \textbf{Outcome:} Students will be able to convert vectors from polar coordinates to rectangular coordinates.
    \item \textbf{Details:} This includes using trigonometric functions to determine the x and y components of a vector from its magnitude and angle, ensuring proper handling of direction (signs of components), and verifying the conversion through examples.
    \item \textbf{Determine the Angle Between Two Vectors Using the Dot Product:}
    \item \textbf{Outcome:} Students will be able to determine the angle between two vectors using the dot product.
    \item \textbf{Details:} This includes applying the identity $$\cos(\theta)=\frac{\mathbf{u}\cdot\mathbf{v}}{|\mathbf{u}||\mathbf{v}|}$$ to find the cosine of the angle between two vectors, using the inverse cosine function to find the angle, and interpreting the angle in the context of the vectors' orientations.
\end{enumerate}

\subsection{Assessment Instrument}

\begin{enumerate}
    \item \textbf{Assessment Instrument: The Dot Product Quiz} \textbf{Format:}
    \item \textbf{Type:} Mixed-format quiz (multiple-choice, short answer, and problem-solving questions)
    \item \textbf{Duration:} 45 minutes
    \item \textbf{Total Marks:} 30 points \textbf{Sections:}
    \item \textbf{Calculating the Dot Product (10 points):}
    \item \textbf{Question 1:} Calculate the dot product of the following vectors given in rectangular coordinates:
    \item a. $\mathbf{u} = (3, 4)$ and $\mathbf{v} = (2, -1)$
    \item b. $\mathbf{u} = (-1, 5)$ and $\mathbf{v} = (4, 3)$
    \item \textbf{Question 2:} Given $\mathbf{u} = (6, 8)$, find $\mathbf{u} \cdot \mathbf{u}$ and interpret the result.
    \item \textbf{Converting Polar Coordinates to Rectangular Coordinates (10 points):}
    \item \textbf{Question 3:} Convert the following vectors from polar to rectangular coordinates:
    \item a. $|\mathbf{u}| = 5$, $\theta = 53^\circ$
    \item b. $|\mathbf{v}| = 7$, $\theta = 120^\circ$
    \item \textbf{Question 4:} Explain the process of converting a vector from polar to rectangular coordinates and why it is necessary for calculating the dot product.
    \item \textbf{Determining the Angle Between Vectors (10 points):}
    \item \textbf{Question 5:} Find the angle between the following pairs of vectors using the dot product:
    \item a. $\mathbf{u} = (1, 2)$ and $\mathbf{v} = (3, 4)$
    \item b. $\mathbf{u} = (2, 3)$ and $\mathbf{v} = (-1, 5)$
    \item \textbf{Question 6:} Given $\mathbf{u} = (4, 4)$ and $\mathbf{v} = (4, -4)$, find the cosine of the angle between them and interpret the result. \textbf{Scoring Rubric:}
    \item \textbf{Calculating the Dot Product:} 3.3 points for each correct calculation and interpretation (total 10 points)
    \item \textbf{Converting Polar Coordinates to Rectangular Coordinates:} 3.3 points for each correct conversion and explanation (total 10 points)
    \item \textbf{Determining the Angle Between Vectors:} 3.3 points for each correct angle calculation and interpretation (total 10 points) This quiz will help assess the students' understanding and application of the dot product, the process of converting polar coordinates to rectangular coordinates, and determining the angle between vectors. It ensures that they can correctly perform these operations and understand their significance in vector mathematics.
\end{enumerate}

%%===============================================================
\section{The Cross-Product}
\label{sec:the-cross-product}
%%===============================================================

In mathematics, the \textbf{cross product} or \textbf{vector product} (occasionally \textbf{directed area product}, to emphasize its geometric significance) is a binary operation on two vectors in a three-dimensional oriented Euclidean vector space (named here $E$), and is denoted by the symbol $\times$. Given two linearly independent vectors $\mathbf{a}$ and $\mathbf{b}$, the cross product, $\mathbf{a} \times \mathbf{b}$ (read "a cross b"), is a vector that is perpendicular to both $\mathbf{a}$ and $\mathbf{b}$, and thus normal to the plane containing them. It has many applications in mathematics, physics, engineering, and computer programming. It should not be confused with the dot product (projection product).



If two vectors have the same direction or have the exact opposite direction from each other (that is, they are \textit{not} linearly independent), or if either one has zero length, then their cross product is zero. More generally, the magnitude of the product equals the area of a parallelogram with the vectors for sides; in particular, the magnitude of the product of two perpendicular vectors is the product of their lengths.



The cross product is anticommutative (that is, $\mathbf{a} \times \mathbf{b} = - \mathbf{b} \times \mathbf{a}$) and is distributive over addition (that is, $\mathbf{a} \times \left( \mathbf{b + \mathbf{c}} \right) = \mathbf{a} \times \mathbf{b} + \mathbf{a} \times \mathbf{c}$). The space $E$ together with the cross product is an algebra over the real numbers, which is neither commutative nor associative, but is a Lie algebra with the cross product being the Lie bracket.



Like the dot product, it depends on the metric of Euclidean space, but unlike the dot product, it also depends on a choice of orientation (or "handedness") of the space (it's why an oriented space is needed). In connection with the cross product, the exterior product of vectors can be used in arbitrary dimensions (with a bivector or 2-form result) and is independent of the orientation of the space.



The product can be generalized in various ways, using the orientation and metric structure just as for the traditional 3-dimensional cross product, one can, in $n$ dimensions, take the product of $n-1$ vectors to produce a vector perpendicular to all of them. But if the product is limited to non-trivial binary products with vector results, it exists only in three and seven dimensions.



% [Image: The cross product with respect to a right-handed coordinate system]



Figure: The cross product with respect to a right-handed coordinate system.



\paragraph{Definition}



% [Image: Finding the direction of the cross product by the right-hand rule.]



Figure: Finding the direction of the cross product by the right-hand rule.



The cross product of two vectors $\mathbf{a}$ and $\mathbf{b}$ is defined only in three-dimensional space and is denoted by $\mathbf{a} \times \mathbf{b}$. In physics and applied mathematics, the wedge notation $\mathbf{a} \wedge \mathbf{b}$ is often used (in conjunction with the name \textit{vector product}), although in pure mathematics such notation is usually reserved for just the exterior product, an abstraction of the vector product to $n$ dimensions.



The cross product $\mathbf{a} \times \mathbf{b}$ is defined as a vector $\mathbf{c}$ that is perpendicular (orthogonal) to both $\mathbf{a}$ and $\mathbf{b}$, with a direction given by the right-hand rule and a magnitude equal to the area of the parallelogram that the vectors span.



The cross product is defined by the formula



$$\mathbf{a} \times \mathbf{b} = \left\| \mathbf{a} \right\| \left\| \mathbf{b} \right\| \sin (\theta) \ \mathbf{n}$$



where:



\begin{itemize}
    \item $\theta$ is the angle between $\mathbf{a}$ and $\mathbf{b}$ in the plane containing them (hence, it is between 0° and 180°)
    \item $\| \| \mathbf{a} \| \|$ and $\| \| \mathbf{b} \| \|$ are the magnitudes of vectors $\mathbf{a}$ and $\mathbf{b}$
    \item and $\mathbf{n}$ is a unit vector perpendicular to the plane containing $\mathbf{a}$ and $\mathbf{b}$, in the direction given by the right-hand rule (illustrated).
\end{itemize}

If the vectors $\mathbf{a}$ and $\mathbf{b}$ are parallel (that is, the angle $\theta$ between them is either 0° or 180°), by the above formula, the cross product of $\mathbf{a}$ and $\mathbf{b}$ is the zero vector \textbf{0}.



\paragraph{Direction}



% [Image: Animated cross product]



Figure: The cross product $\mathbf{a} \times \mathbf{b}$ (vertical, in purple) changes as the angle between the vectors $\mathbf{a}$ (blue) and $\mathbf{b}$ (red) changes. The cross product is always orthogonal to both vectors, and has magnitude zero when the vectors are parallel and maximum magnitude $\| \| \mathbf{a} \| \| \text{ } \| \| \mathbf{b} \| \|$ when they are orthogonal.



By convention, the direction of the vector $\mathbf{n}$ is given by the right-hand rule, where one simply points the forefinger of the right hand in the direction of $\mathbf{a}$ and the middle finger in the direction of $\mathbf{b}$. Then, the vector $\mathbf{n}$ is coming out of the thumb (see the adjacent picture). Using this rule implies that the cross product is anti-commutative; that is, $\mathbf{b} \times \mathbf{a} = -(\mathbf{a} \times \mathbf{b})$). By pointing the forefinger toward $\mathbf{b}$ first, and then pointing the middle finger toward $\mathbf{a}$, the thumb will be forced in the opposite direction, reversing the sign of the product vector.



As the cross product operator depends on the orientation of the space (as explicit in the definition above), the cross product of two vectors is not a "true" vector, but a \textit{pseudovector}.



\paragraph{Names}



% [Image: Sarrus's rule]



Figure: According to Sarrus's rule, the determinant of a 3×3 matrix involves multiplications between matrix elements identified by crossed diagonals.



In 1881, Josiah Willard Gibbs, and independently Oliver Heaviside, introduced both the dot product and the cross product using a period (\textbf{a} . \textbf{b}) and an "x" (\textbf{a} x \textbf{b}), respectively, to denote them.



In 1877, to emphasize the fact that the result of a dot product is a scalar while the result of a cross product is a vector, William Kingdon Clifford coined the alternative names \textbf{scalar product} and \textbf{vector product} for the two operations. These alternative names are still widely used in the literature.



Both the cross notation ($\mathbf{a} \times \mathbf{b}$) and the name \textbf{cross product} were possibly inspired by the fact that each scalar component of $\mathbf{a} \times \mathbf{b}$ is computed by multiplying non-corresponding components of $\mathbf{a}$ and $\mathbf{b}$. Conversely, a dot product $\mathbf{a} \cdot \mathbf{b}$ involves multiplications between corresponding components of $\mathbf{a}$ and $\mathbf{b}$. As explained below, the cross product can be expressed in the form of a determinant of a special $3 \times 3$ matrix. According to Sarrus's rule, this involves multiplications between matrix elements identified by crossed diagonals.



\paragraph{Computing}



\paragraph{Coordinate notation}



% [Image: Standard basis vectors and vector components]



Figure: Standard basis vectors ($\mathbf{i}$, $\mathbf{j}$, $\mathbf{k}$, also denoted $\mathbf{e} _1$, $\mathbf{e} _2$, $\mathbf{e} _3$) and vector components of $\mathbf{a}$ ($\mathbf{a} _x$, $\mathbf{a} _y$, $\mathbf{a} _z$, also denoted $\mathbf{a} _1$, $\mathbf{a} _2$, $\mathbf{a} _3$)



If $(\mathbf{i}, \mathbf{j}, \mathbf{k})$ is a positively oriented orthonormal basis, the basis vectors satisfy the following equalities



$$\begin{alignat}{2}

  \mathbf{\color{blue}{i}}&\times\mathbf{\color{red}{j}} && = \mathbf{\color{green}{k}} \\\  \mathbf{\color{red}{j}}&\times\mathbf{\color{green}{k}} && = \mathbf{\color{blue}{i}} \\\  \mathbf{\color{green}{k}}&\times\mathbf{\color{blue}{i}} && = \mathbf{\color{red}{j}}

\end{alignat}$$



which imply, by the anticommutativity of the cross product, that



$$\begin{alignat}{2}

  \mathbf{\color{red}{j}} & \times\mathbf{\color{blue}{i}} && = -\mathbf{\color{green}{k}} \\\  \mathbf{\color{green}{k}} & \times\mathbf{\color{red}{j}} && = -\mathbf{\color{blue}{i}} \\\  \mathbf{\color{blue}{i}} & \times\mathbf{\color{green}{k}} && = -\mathbf{\color{red}{j}}

\end{alignat}$$



The anticommutativity of the cross product (and the obvious lack of

linear independence) also implies that



$$\mathbf{\color{blue}{i}}\times\mathbf{\color{blue}{i}} = \mathbf{\color{red}{j}}\times\mathbf{\color{red}{j}} = \mathbf{\color{green}{k}}\times\mathbf{\color{green}{k}} = \mathbf{0}$$

(the zero vector).



These equalities, together with the distributivity and linearity of the

cross product (though neither follows easily from the definition given

above), are sufficient to determine the cross product of any two vectors

$\mathbf{a}$ and $\mathbf{b}$. Each vector can be defined as the sum of three

orthogonal components parallel to the standard basis vectors:



$$\begin{alignat}{3}

  \mathbf{a} &= a _1\mathbf{\color{blue}{i}} &&+ a _2\mathbf{\color{red}{j}} &&+ a _3\mathbf{\color{green}{k}} \\\  \mathbf{b} &= b _1\mathbf{\color{blue}{i}} &&+ b _2\mathbf{\color{red}{j}} &&+ b _3\mathbf{\color{green}{k}}

\end{alignat}$$



Their cross product $\mathbf{a} \times \mathbf{b}$ can be expanded using distributivity:



$$\begin{align} \mathbf{a} \times \mathbf{b} = {} & (a _1\mathbf{\color{blue}{i}} + a _2\mathbf{\color{red}{j}} + a _3\mathbf{\color{green}{k}}) \times (b _1\mathbf{\color{blue}{i}} + b _2\mathbf{\color{red}{j}} + b _3\mathbf{\color{green}{k}}) \\\\ = {} & a _1b _1(\mathbf{\color{blue}{i}} \times \mathbf{\color{blue}{i}}) + a _1b _2(\mathbf{\color{blue}{i}} \times \mathbf{\color{red}{j}}) + a _1b _3(\mathbf{\color{blue}{i}} \times \mathbf{\color{green}{k}}) + {} \\\\ & a _2b _1(\mathbf{\color{red}{j}} \times \mathbf{\color{blue}{i}}) + a _2b _2(\mathbf{\color{red}{j}} \times \mathbf{\color{red}{j}}) + a _2b _3(\mathbf{\color{red}{j}} \times \mathbf{\color{green}{k}}) + {} \\\\ & a _3b _1(\mathbf{\color{green}{k}} \times \mathbf{\color{blue}{i}}) + a _3b _2(\mathbf{\color{green}{k}} \times \mathbf{\color{red}{j}}) + a _3b _3(\mathbf{\color{green}{k}} \times \mathbf{\color{green}{k}}) \end{align}$$



This can be interpreted as the decomposition of $\mathbf{a} \times \mathbf{b}$ into the sum of nine simpler cross products involving vectors aligned with $\mathbf{i}$, $\mathbf{j}$, or $\mathbf{k}$. Each one of these nine cross products operates on two vectors that are easy to handle as they are either parallel or orthogonal to each other. From this decomposition, by using the above-mentioned equalities and collecting similar terms, we obtain:



$$\begin{align}

 \mathbf{a}\times\mathbf{b} = {} &\quad\ a _1b _1\mathbf{0} + a _1b _2\mathbf{\color{green}{k}} - a _1b _3\mathbf{\color{red}{j}} \\\\                              &- a _2b _1\mathbf{\color{green}{k}} + a _2b _2\mathbf{0} + a _2b _3\mathbf{\color{blue}{i}} \\\\                              &+ a _3b _1\mathbf{\color{red}{j}}\ - a _3b _2\mathbf{\color{blue}{i}}\ + a _3b _3\mathbf{0} \\\\                         = {} &(a _2b _3 - a _3b _2)\mathbf{\color{blue}{i}} + (a _3b _1 - a _1b _3)\mathbf{\color{red}{j}} + (a _1b _2 - a _2b _1)\mathbf{\color{green}{k}} \end{align}$$



meaning that the three scalar components of the resulting vector $\mathbf{s} = s _1 \mathbf{i} + s _2 \mathbf{j} + s _3 \mathbf{k} = \mathbf{a} \times \mathbf{b}$ are



$$\begin{align}

  s _1 & = a _2b _3-a _3b _2 \\\  s _2 & = a _3b _1-a _1b _3 \\\  s _3 & = a _1b _2-a _2b _1

\end{align}$$



Using column vectors, we can represent the same result as follows:



$$\begin{bmatrix}s _1 \\\\ s _2 \\\\ s _3\end{bmatrix}=\begin{bmatrix}a _2b _3-a _3b _2 \\\\ a _3b _1-a _1b _3 \\\\ a _1b _2-a _2b _1\end{bmatrix}$$



\paragraph{Matrix notation}



% [Image: Use of Sarrus's rule to find the cross product of \textbf{a} and \textbf{b}]



Figure: Use of Sarrus's rule to find the cross product of $\mathbf{a}$ and $\mathbf{b}$



The cross product can also be expressed as the formal determinant:



$$\mathbf{a\times b} = \begin{vmatrix}   \mathbf{i} & \mathbf{j} & \mathbf{k} \\\\   a _1 & a _2 & a _3 \\\\   b _1 & b _2 & b _3 \\\\ \end{vmatrix}$$



This determinant can be computed using Sarrus's rule or cofactor

expansion. Using Sarrus's rule, it expands to



$$\begin{align}

  \mathbf{a\times b} & =(a _2b _3\mathbf{i}+a _3b _1\mathbf{j}+a _1b _2\mathbf{k}) - (a _3b _2\mathbf{i}+a _1b _3\mathbf{j}+a _2b _1\mathbf{k}) \\\\ &=(a _2b _3 - a _3b _2)\mathbf{i} +(a _3b _1 - a _1b _3)\mathbf{j} +(a _1b _2 - a _2b _1)\mathbf{k}.

\end{align}$$



Using cofactor expansion along the first row instead, it expands to



$$\begin{align}

\mathbf{a\times b} &=

  \begin{vmatrix} a _2 & a _3 \\\\ b _2&b _3

  \end{vmatrix}\mathbf{i} -

  \begin{vmatrix}  a _1&a _3 \\\\  b _1&b _3

  \end{vmatrix}\mathbf{j} +

  \begin{vmatrix} a _1&a _2 \\\\ b _1 & b _2

  \end{vmatrix}\mathbf{k} \\\ &=(a _2b _3 - a _3b _2)\mathbf{i} -(a _1b _3 - a _3b _1)\mathbf{j} +(a _1b _2 - a _2b _1)\mathbf{k},

\end{align}$$ which gives the components of the resulting vector

directly.



\paragraph{Using Levi-Civita tensors}



\begin{itemize}
    \item In any basis, the cross-product $a \times b$ is given by the tensorial formula $E _{ijk}a^ib^j$ where $E _{ijk}$ is the covariant Levi-Civita tensor (we note the position of the indices). That corresponds to the intrinsic formula given here.
    \item In an orthonormal basis \textbf{having the same orientation as the space}, $a \times b$ is given by the pseudo-tensorial formula $\varepsilon _{ijk}a^ib^j$ where $\varepsilon _{ijk}$ is the Levi-Civita symbol (which is a pseudo-tensor). That's the formula used for everyday physics but it works only for this special choice of basis.
    \item In any orthonormal basis, $a \times b$ is given by the pseudo-tensorial formula $(-1)^B\varepsilon _{ijk}a^ib^j$ where $(-1)^B = \pm 1$ indicates whether the basis has the same orientation as the space or not.
\end{itemize}

The latter formula avoids having to change the orientation of the space

when we inverse an orthonormal basis.



\paragraph{Properties}



\paragraph{Geometric meaning}



% [Image: Figure 1. The area of a parallelogram as the magnitude of a cross product]



Figure 1: The area of a parallelogram as the magnitude of a cross product



% [Image: Figure 2. Three vectors defining a parallelepiped]



Figure 2: Three vectors defining a parallelepiped



The magnitude of the cross product can be interpreted as the positive area of the parallelogram having $\mathbf{a}$ and $\mathbf{b}$ as sides (see Figure 1):



$$\left\| \mathbf{a} \times \mathbf{b} \right\| = \left\| \mathbf{a} \right\| \left\| \mathbf{b} \right\| | \sin \theta | .$$



Indeed, one can also compute the volume \textit{V} of a parallelepiped having $\mathbf{a}$, $\mathbf{b}$ and $\mathbf{c}$ as edges by using a combination of a cross product and a dot product, called scalar triple product (see Figure 2):



$$\mathbf{a}\cdot(\mathbf{b}\times \mathbf{c})=

  \mathbf{b}\cdot(\mathbf{c}\times \mathbf{a})=

  \mathbf{c}\cdot(\mathbf{a}\times \mathbf{b}).$$



Since the result of the scalar triple product may be negative, the

volume of the parallelepiped is given by its absolute value:



$$V = |\mathbf{a} \cdot (\mathbf{b} \times \mathbf{c})|.$$



Because the magnitude of the cross product goes by the sine of the angle

between its arguments, the cross product can be thought of as a measure

of \textit{perpendicularity} in the same way that the dot product is a measure

of \textit{parallelism}. Given two unit vectors, their cross product has a

magnitude of 1 if the two are perpendicular and a magnitude of zero if

the two are parallel. The dot product of two unit vectors behaves just

oppositely: it is zero when the unit vectors are perpendicular and 1 if

the unit vectors are parallel.



Unit vectors enable two convenient identities: the dot product of two

unit vectors yields the cosine (which may be positive or negative) of

the angle between the two unit vectors. The magnitude of the cross

product of the two unit vectors yields the sine (which will always be

positive).



\paragraph{Algebraic properties}



% [Image: Cross product scalar multiplication. \textbf{Left:} Decomposition of \textbf{b} into components parallel and perpendicular to \textbf{a}. Right: Scaling of the perpendicular components by a positive real number \textit{r} (if negative, \textbf{b} and the cross product are reversed).]



Figure: Cross product scalar multiplication. \textbf{Left:} Decomposition of \textbf{b} into components parallel and perpendicular to \textbf{a}. Right: Scaling of the perpendicular components by a positive real number \textit{r} (if negative, \textbf{b} and the cross product are reversed).



% [Image: Cross product distributivity over vector addition. \textbf{Left:} The vectors \textbf{b} and \textbf{c} are resolved into parallel and perpendicular components to \textbf{a}. \textbf{Right:} The parallel components vanish in the cross product, only the perpendicular components shown in the plane perpendicular to \textbf{a} remain.]



Figure: Cross product distributivity over vector addition. \textbf{Left:} The vectors \textbf{b} and \textbf{c} are resolved into parallel and perpendicular components to \textbf{a}. \textbf{Right:} The parallel components vanish in the cross product, only the perpendicular components shown in the plane perpendicular to \textbf{a} remain.



% [Image: The two nonequivalent triple cross products of three vectors \textbf{a}, \textbf{b}, \textbf{c}. In each case, two vectors define a plane, the other is out of the plane and can be split into parallel and perpendicular components to the cross product of the vectors defining the plane. These components can be found by vector projection and rejection. The triple product is in the plane and is rotated as shown.]



Figure: The two nonequivalent triple cross products of three vectors $\mathbf{a}$, $\mathbf{b}$, $\mathbf{c}$. In each case, two vectors define a plane, the other is out of the plane and can be split into parallel and perpendicular components to the cross product of the vectors defining the plane. These components can be found by vector projection and rejection. The triple product is in the plane and is rotated as shown.



If the cross product of two vectors is the zero vector (that is, $\mathbf{a} \times \mathbf{b} = \mathbf{0}$), then either one or both of the inputs is the zero vector, ($\mathbf{a} = \mathbf{0}$ or $\mathbf{b} = \mathbf{0}$) or else they are parallel or antiparallel ($\mathbf{a}$ $\| \|$ $\mathbf{a}$) so that the sine of the angle between them is zero ($\theta$ = 0° or $\theta$ = 180° and sin $\theta$ = 0).



The self cross product of a vector is the zero vector:



$$\mathbf{a} \times \mathbf{a} = \mathbf{0}.$$ The cross product is anticommutative,



$$\mathbf{a} \times \mathbf{b} = -(\mathbf{b} \times \mathbf{a}),$$ distributive over addition,



$\quad$ $\mathbf{a} \times (\mathbf{b} + \mathbf{c}) = (\mathbf{a} \times \mathbf{b}) + (\mathbf{a} \times \mathbf{c}),$



and compatible with scalar multiplication so that



$$(r\,\mathbf{a}) \times \mathbf{b} = \mathbf{a} \times (r\,\mathbf{b}) = r\,(\mathbf{a} \times \mathbf{b}).$$



It is not associative, but satisfies the Jacobi identity:



$$\mathbf{a} \times (\mathbf{b} \times \mathbf{c}) + \mathbf{b} \times (\mathbf{c} \times \mathbf{a}) + \mathbf{c} \times (\mathbf{a} \times \mathbf{b}) = \mathbf{0}.$$

Distributivity, linearity and Jacobi identity show that the $\mathbf{R}^3$ vector space together with vector addition and the cross product forms a Lie algebra, the Lie algebra of the real orthogonal group in 3 dimensions, SO(3). The cross product does not obey the cancellation law; that is, $\mathbf{a} \times \mathbf{b} = \mathbf{a} \times \mathbf{c}$ with $\mathbf{a} \ne \mathbf{0}$ does not imply $\mathbf{b} = \mathbf{c}$, but only that:



$$\begin{align} \mathbf{0} & = (\mathbf{a} \times \mathbf{b}) - (\mathbf{a} \times \mathbf{c}) \\\\ & = \mathbf{a} \times (\mathbf{b} - \mathbf{c}). \\\\ \end{align}$$



This can be the case where $\mathbf{b}$ and $\mathbf{c}$ cancel, but additionally where $\mathbf{a}$ and $\mathbf{b} - \mathbf{c}$ are parallel; that is, they are related by a scale factor \textit{t}, leading to:



$$\mathbf{c} = \mathbf{b} + t\,\mathbf{a},$$ for some scalar \textit{t}.



If, in addition to $\mathbf{a} \times \mathbf{b} = \mathbf{a} \times \mathbf{c}$ and $\mathbf{a} \ne \mathbf{0}$ as above, it is the case that $\mathbf{a} \cdot \mathbf{b} = \mathbf{a} \cdot \mathbf{c}$ then



$$\begin{align}   \mathbf{a} \times (\mathbf{b} - \mathbf{c}) & = \mathbf{0} \\\\    \mathbf{a} \cdot (\mathbf{b} - \mathbf{c}) & = 0, \end{align}$$

As $\mathbf{b} - \mathbf{c}$ cannot be simultaneously parallel (for the cross product to be \textbf{0}) and perpendicular (for the dot product to be 0) to $\mathbf{a}$, it must be the case that $\mathbf{b}$ and $\mathbf{c}$ cancel: $\mathbf{b} = \mathbf{c}$.



From the geometrical definition, the cross product is invariant under proper rotations about the axis defined by $\mathbf{a} \times \mathbf{b}$. In formulae:



$$(R\mathbf{a}) \times (R\mathbf{b}) = R(\mathbf{a} \times \mathbf{b})$$, where $R$ is a rotation matrix with $\det(R)=1$.



More generally, the cross product obeys the following identity under matrix transformations:



$$(M\mathbf{a}) \times (M\mathbf{b}) = (\det M) \left(M^{-1}\right)^\mathrm{T}(\mathbf{a} \times \mathbf{b}) = \operatorname{cof} M (\mathbf{a} \times \mathbf{b})$$

where $M$ is a 3-by-3 matrix and $\left(M^{-1}\right)^\mathrm{T}$ is the transpose of the inverse and $\operatorname{cof}$ is the cofactor matrix. It can be readily seen how this formula reduces to the former one if $M$ is a rotation matrix.



The cross product of two vectors lies in the null space of the 2 × 3 matrix with the vectors as rows:



$$\mathbf{a} \times \mathbf{b} \in NS\left(\begin{bmatrix}\mathbf{a} \\\\ \mathbf{b}\end{bmatrix}\right).$$

For the sum of two cross products, the following identity holds:



$$\mathbf{a} \times \mathbf{b} + \mathbf{c} \times \mathbf{d} = (\mathbf{a} - \mathbf{c}) \times (\mathbf{b} - \mathbf{d}) + \mathbf{a} \times \mathbf{d} + \mathbf{c} \times \mathbf{b}.$$



\paragraph{Differentiation}



The product rule of differential calculus applies to any bilinear operation, and therefore also to the cross product:



$$\frac{d}{dt}(\mathbf{a} \times \mathbf{b}) = \frac{d\mathbf{a}}{dt} \times \mathbf{b} + \mathbf{a} \times \frac{d\mathbf{b}}{dt} ,$$



where $\mathbf{a}$ and $\mathbf{b}$ are vectors that depend on the real variable \textit{t}.



\paragraph{Triple product expansion}



The cross product is used in both forms of the triple product. The scalar triple product of three vectors is defined as



$$\mathbf{a} \cdot (\mathbf{b} \times \mathbf{c}),$$



It is the signed volume of the parallelepiped with edges $\mathbf{a}$, $\mathbf{b}$ and $\mathbf{c}$ and as such the vectors can be used in any order that's an even permutation of the above ordering. The following therefore are equal:



$$\mathbf{a} \cdot (\mathbf{b} \times \mathbf{c}) = \mathbf{b} \cdot (\mathbf{c} \times \mathbf{a}) = \mathbf{c} \cdot (\mathbf{a} \times \mathbf{b}),$$



The vector triple product is the cross product of a vector with the result of another cross product, and is related to the dot product by the following formula



$$\mathbf{a} \times (\mathbf{b} \times \mathbf{c}) = \mathbf{b}(\mathbf{a} \cdot \mathbf{c}) - \mathbf{c}(\mathbf{a} \cdot \mathbf{b}).$$



The mnemonic "BAC minus CAB" is used to remember the order of the vectors in the right hand member. This formula is used in physics to simplify vector calculations. A special case, regarding gradients and useful in vector calculus, is



$$\begin{align} \nabla \times (\nabla \times \mathbf{f}) & = \nabla (\nabla \cdot \mathbf{f} ) - (\nabla \cdot \nabla) \mathbf{f} \\\\ & = \nabla (\nabla \cdot \mathbf{f} ) - \nabla^2 \mathbf{f}, \\\\ \end{align}$$



where $\nabla ^2$ is the vector Laplacian operator.



Other identities relate the cross product to the scalar triple product:



$$\begin{align} (\mathbf{a}\times \mathbf{b})\times (\mathbf{a}\times \mathbf{c}) & = (\mathbf{a}\cdot(\mathbf{b}\times \mathbf{c})) \mathbf{a} \\\\ (\mathbf{a}\times \mathbf{b})\cdot(\mathbf{c}\times \mathbf{d}) & = \mathbf{b}^\mathrm{T} \left( \left( \mathbf{c}^\mathrm{T} \mathbf{a}\right)I - \mathbf{c} \mathbf{a}^\mathrm{T} \right) \mathbf{d} \\\\ & = (\mathbf{a}\cdot \mathbf{c})(\mathbf{b}\cdot \mathbf{d})-(\mathbf{a}\cdot \mathbf{d}) (\mathbf{b}\cdot \mathbf{c}) \end{align}$$



where \textit{I} is the identity matrix.



\paragraph{Alternative formulation}



The cross product and the dot product are related by:



$$\left\| \mathbf{a} \times \mathbf{b} \right\| ^2  = \left\| \mathbf{a}\right\|^2  \left\|\mathbf{b}\right\|^2 - (\mathbf{a} \cdot \mathbf{b})^2 .$$



The right-hand side is the Gram determinant of $\mathbf{a}$ and $\mathbf{b}$, the square of the area of the parallelogram defined by the vectors. This condition determines the magnitude of the cross product. Namely, since the dot product is defined, in terms of the angle $\theta$ between the two vectors, as:



$$\mathbf{a \cdot b} = \left\| \mathbf a \right\| \left\| \mathbf b \right\| \cos \theta ,$$



the above given relationship can be rewritten as follows:



$$\left\| \mathbf{a \times b} \right\|^2 = \left\| \mathbf{a} \right\| ^2 \left\| \mathbf{b}\right \| ^2 \left(1-\cos^2 \theta \right) .$$



Invoking the Pythagorean trigonometric identity one obtains:



$$\left\| \mathbf{a} \times \mathbf{b} \right\| = \left\| \mathbf{a} \right\| \left\| \mathbf{b} \right\| \left| \sin \theta \right| ,$$



which is the magnitude of the cross product expressed in terms of $\theta$, equal to the area of the parallelogram defined by $\mathbf{a}$ and $\mathbf{b}$.



The combination of this requirement and the property that the cross product be orthogonal to its constituents $\mathbf{a}$ and $\mathbf{b}$ provides an alternative definition of the cross product.



\paragraph{Lagrange's identity}



The relation:



$$\left\| \mathbf{a} \times \mathbf{b} \right\|^2 \equiv

  \det \begin{bmatrix} \mathbf{a} \cdot \mathbf{a} & \mathbf{a} \cdot \mathbf{b} \\\\ \mathbf{a} \cdot \mathbf{b}  & \mathbf{b} \cdot \mathbf{b}\\\  \end{bmatrix} \equiv

  \left\| \mathbf{a} \right\| ^2  \left\| \mathbf{b} \right\| ^2 - (\mathbf{a} \cdot \mathbf{b})^2 .$$



can be compared with another relation involving the right-hand side, namely Lagrange's identity expressed as:



$$\sum _{1 \le i < j \le n} \left( a _ib _j - a _jb _i \right)^2 \equiv

  \left\| \mathbf a \right\|^2 \left\| \mathbf b \right\|^2 - ( \mathbf{a \cdot b } )^2\ ,$$



where $\mathbf{a}$ and $\mathbf{b}$ may be \textit{n}-dimensional vectors. This also shows that the Riemannian volume form for surfaces is exactly the surface element from vector calculus. In the case where \textit{n} = 3, combining these two equations results in the expression for the magnitude of the cross product in terms of its components:



$$\begin{align} & \left\|\mathbf{a} \times \mathbf{b}\right\|^2 \equiv \sum _{1 \le i < j \le 3}\left(a _ib _j - a _jb _i \right)^2 \\\\ \equiv{} &\left(a _1b _2 - b _1a _2\right)^2 + \left(a _2b _3 - a _3b _2\right)^2 + \left(a _3b _1 - a _1b _3\right)^2 \ . \end{align}$$



The same result is found directly using the components of the cross product found from:



$$\mathbf{a} \times \mathbf{b} \equiv \det \begin{bmatrix}   \hat{\mathbf{i}} & \hat{\mathbf{j}} & \hat{\mathbf{k}} \\\\   a _1 & a _2 & a _3 \\\\   b _1 & b _2 & b _3 \end{bmatrix}.$$



In $\mathbf{R}^3$, Lagrange's equation is a special case of the multiplicativity \|\textbf{vw}\| = \|\textbf{v}\|\|\textbf{w}\| of the norm in the quaternion algebra.



It is a special case of another formula, also sometimes called Lagrange's identity, which is the three dimensional case of the Binet–Cauchy identity:



$$(\mathbf{a} \times \mathbf{b}) \cdot (\mathbf{c} \times \mathbf{d}) \equiv

  (\mathbf{a} \cdot \mathbf{c})(\mathbf{b} \cdot \mathbf{d}) - (\mathbf{a} \cdot \mathbf{d})(\mathbf{b} \cdot \mathbf{c}).$$



If $\mathbf{a} = \mathbf{c}$ and $\mathbf{b} = \mathbf{d}$ this simplifies to the formula above.



\paragraph{Infinitesimal generators of rotations}



The cross product conveniently describes the infinitesimal generators of rotations in $\mathbf{R}^3$. Specifically, if $\mathbf{n}$ is a unit vector in $\mathbf{R}^3$ and $R(\varphi, \mathbf{n})$ denotes a rotation about the axis through the origin specified by $\mathbf{n}$, with angle $\varphi$ (measured in radians, counterclockwise when viewed from the tip of $\mathbf{n}$), then



$$\left.{d\over d\phi} \right| _{\phi=0} R(\phi,\boldsymbol{n}) \boldsymbol{x} = \boldsymbol{n} \times \boldsymbol{x}$$

for every vector $\mathbf{x}$ in $\mathbf{R}^3$. The cross product with $\mathbf{n}$ therefore describes the infinitesimal generator of the rotations about $\mathbf{n}$. These infinitesimal generators form the Lie algebra \textbf{so}(3) of the rotation group SO(3), and we obtain the result that the Lie algebra $\mathbf{R}^3$ with cross product is isomorphic to the Lie algebra \textbf{so}(3).



\paragraph{Alternative ways to compute}



\paragraph{Conversion to matrix multiplication}



The vector cross product also can be expressed as the product of a skew-symmetric matrix and a vector:



$$\begin{align} \mathbf{a} \times \mathbf{b} = [\mathbf{a}] _{\times} \mathbf{b} &= \begin{bmatrix} 0 & -a _3 & a _2 \\\\ a _3 & 0 & -a _1 \\\\ -a _2 & a _1 & 0 \end{bmatrix} \begin{bmatrix} b _1 \\\\ b _2 \\\\ b _3 \end{bmatrix} \\\\ \mathbf{a} \times \mathbf{b} = {[\mathbf{b}] _\times}^\mathrm{T} \mathbf{a} &= \begin{bmatrix} 0 & b _3 & -b _2 \\\\ -b _3 & 0 & b _1 \\\\ b _2 & -b _1 & 0 \end{bmatrix} \begin{bmatrix} a _1 \\\\ a _2 \\\\ a _3 \end{bmatrix} \end{align}$$



where superscript $^T$ refers to the transpose operation, and $[\mathbf{a}] _{\times}$ is defined by:



$$[\mathbf{a}] _{\times} \stackrel{\rm def}{=} \begin{bmatrix} 0 & -a _3 & a _2 \\\\ a _3 & 0 & -a _1 \\\\ -a _2 & a _1 & 0 \end{bmatrix}.$$



The columns $[\mathbf{a}] _{\times, i}$ of the skew-symmetric matrix for a vector $\mathbf{a}$ can be also obtained by calculating the cross product with unit vectors. That is,



$$[\mathbf{a}] _{\times, i} = \mathbf{a} \times \mathbf{\hat{e} _i}, \; i\in \\{1,2,3\\}$$



or



$$[\mathbf{a}] _{\times} = \sum _{i=1}^3\left(\mathbf{a} \times \mathbf{\hat{e} _i}\right)\otimes\mathbf{\hat{e} _i},$$



where $\otimes$ is the outer product operator.



Also, if $\mathbf{a}$ is itself expressed as a cross product:



$$\mathbf{a} = \mathbf{c} \times \mathbf{d}$$



then



$$[\mathbf{a}]_{\times} = \mathbf{d}\mathbf{c}^\mathrm{T} - \mathbf{c}\mathbf{d}^\mathrm{T} .$$



> \textbf{Proof by substitution}

> Evaluation of the cross product gives

> $$ \mathbf{a} = \mathbf{c} \times \mathbf{d} = \begin{pmatrix} c _2 d _3 - c _3 d _2 \\\\ c _3 d _1 - c _1 d _3 \\\\ c _1 d _2 - c _2 d _1 \end{pmatrix} $$

>

> Hence, the left-hand side equals

> $$

> [\mathbf{a}] _{\times} = \begin{bmatrix} 0 & c _2 d _1 - c _1 d _2 & c _3 d _1 - c _1 d _3 \\\\ c _1 d _2 - c _2 d _1 & 0 & c _3 d _2 - c _2 d _3 \\\\ c _1 d _3 - c _3 d _1 & c _2 d _3 - c _3 d _2 & 0 \end{bmatrix} $$

>

> Now, for the right-hand side,

> $$ \mathbf{c} \mathbf{d}^{\mathrm{T}} = \begin{bmatrix} c _1 d _1 & c _1 d _2 & c _1 d _3 \\\\ c _2 d _1 & c _2 d _2 & c _2 d _3 \\\\ c _3 d _1 & c _3 d _2 & c _3 d _3 \end{bmatrix} $$

>

> And its transpose is

> $$ \mathbf{d} \mathbf{c}^{\mathrm{T}} = \begin{bmatrix} c _1 d _1 & c _2 d _1 & c _3 d _1 \\\\ c _1 d _2 & c _2 d _2 & c _3 d _2 \\\\ c _1 d _3 & c _2 d _3 & c _3 d _3 \end{bmatrix} $$

>

> Evaluation of the right-hand side gives

> $$ \mathbf{d} \mathbf{c}^{\mathrm{T}} - \mathbf{c} \mathbf{d}^{\mathrm{T}} = \begin{bmatrix} 0 & c _2 d _1 - c _1 d _2 & c _3 d _1 - c _1 d _3 \\\\ c _1 d _2 - c _2 d _1 & 0 & c _3 d _2 - c _2 d _3 \\\\ c _1 d _3 - c _3 d _1 & c _2 d _3 - c _3 d _2 & 0 \end{bmatrix} $$

>

> Comparison shows that the left-hand side equals the right-hand side.



This result can be generalized to higher dimensions using geometric algebra. In particular in any dimension bivectors can be identified with skew-symmetric matrices, so the product between a skew-symmetric matrix and vector is equivalent to the grade-1 part of the product of a bivector and vector. In three dimensions bivectors are dual to vectors so the product is equivalent to the cross product, with the bivector instead of its vector dual. In higher dimensions the product can still be calculated but bivectors have more degrees of freedom and are not equivalent to vectors.



This notation is also often much easier to work with, for example, in epipolar geometry.



From the general properties of the cross product follows immediately that



$$[\mathbf{a}] _{\times} \, \mathbf{a} = \mathbf{0}$$

and $\mathbf{a}^\mathrm T \, [\mathbf{a}] _{\times} = \mathbf{0}$



and from fact that $[\mathbf{a}] _{\times}$ is skew-symmetric it follows that



$$\mathbf{b}^\mathrm T \, [\mathbf{a}] _{\times} \, \mathbf{b} = 0.$$



The above-mentioned triple product expansion (bac–cab rule) can be easily proven using this notation.



As mentioned above, the Lie algebra $\mathbf{R}^3$ with cross product is isomorphic to the Lie algebra \textbf{so(3)}, whose elements can be identified with the 3×3 skew-symmetric matrices. The map $\mathbf{a} \Rightarrow [\mathbf{a}] _{\times}$ provides an isomorphism between $\mathbf{R}^3$ and \textbf{so(3)}. Under this map, the cross product of 3-vectors corresponds to the commutator of 3x3 skew-symmetric matrices.



> \textbf{Matrix conversion for cross product with canonical base vectors}

>

> Denoting with $\mathbf{e} _i \in \mathbf{R}^{3 \times 1}$ the $i$-th

> canonical base vector, the cross product of a generic vector

> $\mathbf{v} \in \mathbf{R}^{3 \times 1}$ with $\mathbf{e} _i$ is given

> by: $\mathbf{v} \times \mathbf{e} _i = \mathbf{C} _i \mathbf{v}$, where

>

> $\mathbf{C} _1 = \begin{bmatrix} 0 \& 0 \& 0 \\\\ 0 \& 0 \& 1 \\\\ 0 \& -1 \& 0 \end{bmatrix}, \quad

> \mathbf{C} _2 = \begin{bmatrix} 0 \& 0 \& -1 \\\\ 0 \& 0 \& 0 \\\\ 1 \& 0 \& 0 \end{bmatrix}, \quad

> \mathbf{C} _3 = \begin{bmatrix} 0 \& 1 \& 0 \\\\ -1 \& 0 \& 0 \\\\ 0 \& 0 \& 0 \end{bmatrix}$

>

> These matrices share the following properties:

>

> -   $\mathbf{C} _i^\textrm{T} = -\mathbf{C} _i$ (skew-symmetric);

> -   Both trace and determinant are zero;

> -   $\text{rank}(\mathbf{C} _i) = 2$;

> -   $\mathbf{C} _i \mathbf{C} _i^\textrm{T} = \mathbf{P}^{ ^\perp} _{\mathbf{e} _i}$

>     (see below);

>

> The orthogonal projection matrix of a vector

> $\mathbf{v} \neq \mathbf{0}$ is given by

> $\mathbf{P} _{\mathbf{v}} = \mathbf{v}\left(\mathbf{v}^\textrm{T} \mathbf{v}\right)^{-1} \mathbf{v}^T$.

> The projection matrix onto the orthogonal complement is given by

> $\mathbf{P}^{ ^\perp} _{\mathbf{v}} = \mathbf{I} - \mathbf{P} _{\mathbf{v}}$,

> where $\mathbf{I}$ is the identity matrix. For the special case of

> $\mathbf{v} = \mathbf{e} _i$, it can be verified that

>

> $\mathbf{P}^{^\perp} _{\mathbf{e} _1} = \begin{bmatrix} 0 \& 0 \& 0 \\\\ 0 \& 1 \& 0 \\\\ 0 \& 0 \& 1 \end{bmatrix}, \quad

> \mathbf{P}^{ ^\perp} _{\mathbf{e} _2} = \begin{bmatrix} 1 \& 0 \& 0 \\\\ 0 \& 0 \& 0 \\\\ 0 \& 0 \& 1 \end{bmatrix}, \quad

> \mathbf{P}^{ ^\perp} _{\mathbf{e} _3} = \begin{bmatrix} 1 \& 0 \& 0 \\\\ 0 \& 1 \& 0 \\\\ 0 \& 0 \& 0 \end{bmatrix}$

>

> For other properties of orthogonal projection matrices, see projection

> (linear algebra).



\paragraph{Index notation for tensors}



The cross product can alternatively be defined in terms of the Levi-Civita tensor $E _{ijk}$ and a dot product $\eta^{mi}$, which are useful in converting vector notation for tensor applications:



$$\mathbf{c} = \mathbf{a \times b} \Leftrightarrow\ c^m = \sum _{i=1}^3 \sum _{j=1}^3 \sum _{k=1}^3 \eta^{mi} E _{ijk} a^j b^k$$



where the indices $i,j,k$ correspond to vector components. This characterization of the cross product is often expressed more compactly using the Einstein summation convention as



$$\mathbf{c} = \mathbf{a \times b} \Leftrightarrow\ c^m = \eta^{mi} E _{ijk} a^j b^k$$



in which repeated indices are summed over the values 1 to 3.



In a positively-oriented orthonormal basis $\eta^{mi} = \delta^{mi}$ (the Kronecker delta) and $E _{ijk} = \varepsilon _{ijk}$ (the Levi-Civita symbol). In that case, this representation is another form of the skew-symmetric representation of the cross product:



$$[\varepsilon _{ijk} a^j] = [\mathbf{a}] _\times.$$



In classical mechanics: representing the cross product by using the Levi-Civita symbol can cause mechanical symmetries to be obvious when physical systems are isotropic. (An example: consider a particle in a Hooke's Law potential in three-space, free to oscillate in three dimensions; none of these dimensions are "special" in any sense, so symmetries lie in the cross-product-represented angular momentum, which are made clear by the abovementioned Levi-Civita representation).



\paragraph{Mnemonic}



% [Image: Mnemonic to calculate a cross product in vector form]



Figure: Mnemonic to calculate a cross product in vector form



The word "xyzzy" can be used to remember the definition of the cross product.



If



$$\mathbf{a} = \mathbf{b} \times \mathbf{c}$$



where:



$$\mathbf{a} = \begin{bmatrix}a _x \\\\ a _y \\\\ a _z\end{bmatrix},  \mathbf{b} = \begin{bmatrix}b _x \\\\ b _y \\\\ b _z\end{bmatrix},  \mathbf{c} = \begin{bmatrix}c _x \\\\ c _y \\\\ c _z\end{bmatrix}$$



then:



$$a _x = b _y c _z - b _z c _y$$



$$a _y = b _z c _x - b _x c _z$$



$$a _z = b _x c _y - b _y c _x.$$



The second and third equations can be obtained from the first by simply vertically rotating the subscripts, $x \Rightarrow y \Rightarrow z \Rightarrow x$. The problem, of course, is how to remember the first equation, and two options are available for this purpose: either to remember the relevant two diagonals of Sarrus's scheme (those containing \textit{$\mathbf{i}$}), or to remember the xyzzy sequence.



Since the first diagonal in Sarrus's scheme is just the main diagonal of the above-mentioned 3×3 matrix, the first three letters of the word xyzzy can be very easily remembered.



\paragraph{Cross visualization}



Similarly to the mnemonic device above, a "cross" or X can be visualized between the two vectors in the equation. This may be helpful for remembering the correct cross product formula.



If



$$\mathbf{a} = \mathbf{b} \times \mathbf{c}$$



then:



$$\mathbf{a} = \begin{bmatrix}b _x \\\\ b _y \\\\ b _z \end{bmatrix} \times \begin{bmatrix}c _x \\\\ c _y \\\\ c _z\end{bmatrix}.$$



If we want to obtain the formula for $a _x$ we simply drop the $b _x$ and $c _x$ from the formula, and take the next two components down:



$$a _x = \begin{bmatrix}b _y \\\\ b _z\end{bmatrix} \times \begin{bmatrix}c _y \\\\ c _z\end{bmatrix}.$$



When doing this for $a _y$ the next two elements down should "wrap around" the matrix so that after the z component comes the x component. For clarity, when performing this operation for $a _y$, the next two components should be z and x (in that order). While for $a _z$ the next two components should be taken as x and y.



$$a _y = \begin{bmatrix}b _z \\\\ b _x\end{bmatrix} \times \begin{bmatrix}c _z \\\\ c_x \end{bmatrix},  a _z = \begin{bmatrix}b _x \\\\ b _y\end{bmatrix} \times \begin{bmatrix}c _x \\\\ c _y \end{bmatrix}$$



For $a _x$ then, if we visualize the cross operator as pointing from an element on the left to an element on the right, we can take the first element on the left and simply multiply by the element that the cross points to in the right hand matrix. We then subtract the next element down on the left, multiplied by the element that the cross points to here as well. This results in our $a _x$ formula –



$$a _x = b _y c _z - b _z c _y.$$



We can do this in the same way for $a _y$ and $a _z$ to construct their associated formulas.



\paragraph{Applications}



The cross product has applications in various contexts. For example, it is used in computational geometry, physics and engineering. A non-exhaustive list of examples follows.



\paragraph{Computational geometry}



The cross product appears in the calculation of the distance of two skew lines (lines not in the same plane) from each other in three-dimensional space.



The cross product can be used to calculate the normal for a triangle or polygon, an operation frequently performed in computer graphics. For example, the winding of a polygon (clockwise or anticlockwise) about a point within the polygon can be calculated by triangulating the polygon (like spoking a wheel) and summing the angles (between the spokes) using the cross product to keep track of the sign of each angle.



In computational geometry of the plane, the cross product is used to determine the sign of the acute angle defined by three points $p _1=(x _1,y _1), p _2=(x _2,y _2)$ and $p _3=(x _3,y _3)$. It corresponds to the direction (upward or downward) of the cross product of the two coplanar vectors defined by the two pairs of points $(p _1, p _2)$ and $(p _1, p _3)$. The sign of the acute angle is the sign of the expression



$$P = (x _2-x _1)(y _3-y _1)-(y _2-y _1)(x _3-x _1),$$ which is the signed length of the cross product of the two vectors.



In the "right-handed" coordinate system, if the result is 0, the points are collinear; if it is positive, the three points constitute a positive angle of rotation around $p _1$ from $p _2$ to $p _3$, otherwise a negative angle. From another point of view, the sign of $P$ tells whether $p _3$ lies to the left or to the right of line $p _1, p _2.$



The cross product is used in calculating the volume of a polyhedron such as a tetrahedron or parallelepiped.



\paragraph{Angular momentum and torque}



The angular momentum $\mathbf{L}$ of a particle about a given origin is defined as:



$\quad$ $\mathbf{L} = \mathbf{r} \times \mathbf{p},$



where $\mathbf{r}$ is the position vector of the particle relative to the origin, $\mathbf{p}$ is the linear momentum of the particle.



In the same way, the moment $\mathbf{M}$ of a force $\mathbf{F} _B$ applied at point B around point A is given as:



$\quad$ $\mathbf{M} _\mathrm{A} = \mathbf{r} _\mathrm{AB} \times \mathbf{F} _\mathrm{B}\,$



In mechanics the \textit{moment of a force} is also called \textit{torque} and written as $\mathbf{\tau}$



Since position $\mathbf{r}$, linear momentum $\mathbf{p}$ and force $\mathbf{F}$ are all \textit{true} vectors, both the angular momentum $\mathbf{L}$ and the moment of a force $\mathbf{M}$ are \textit{pseudovectors} or \textit{axial vectors}.



\paragraph{Rigid body}



The cross product frequently appears in the description of rigid motions. Two points \textit{P} and \textit{Q} on a rigid body can be related by:



$\quad$ $\mathbf{v} _P - \mathbf{v} _Q = \boldsymbol\omega \times \left( \mathbf{r} _P - \mathbf{r} _Q \right)\,$



where $\mathbf{r}$ is the point's position, $\mathbf{v}$ is its velocity and $\boldsymbol\omega$ is the body's angular velocity.



Since position $\mathbf{r}$ and velocity $\mathbf{v}$ are \textit{true} vectors, the angular velocity $\boldsymbol\omega$ is a \textit{pseudovector} or \textit{axial vector}.



\paragraph{Lorentz force}



The cross product is used to describe the Lorentz force experienced by a moving electric charge $q _e$:



$\quad$ $\mathbf{F} = q_e \left( \mathbf{E}+ \mathbf{v} \times \mathbf{B} \right)$



Since velocity $\mathbf{v}$, force $\mathbf{F}$ and electric field $\mathbf{E}$ are all \textit{true} vectors, the magnetic field $\mathbf{B}$ is a \textit{pseudovector}.



\paragraph{Other}



In vector calculus, the cross product is used to define the formula for the vector operator curl.



The trick of rewriting a cross product in terms of a matrix multiplication appears frequently in epipolar and multi-view geometry, in particular when deriving matching constraints.



\paragraph{As an external product}



% [Image: The cross product in relation to the exterior product. In red are the orthogonal unit vector, and the "parallel" unit bivector.]



Figure: The cross product in relation to the exterior product. In red are the orthogonal unit vector, and the "parallel" unit bivector.



The cross product can be defined in terms of the exterior product. It can be generalized to an external product in other than three dimensions. This view allows for a natural geometric interpretation of the cross product. In exterior algebra the exterior product of two vectors is a bivector. A bivector is an oriented plane element, in much the same way that a vector is an oriented line element. Given two vectors \textit{a} and \textit{b}, one can view the bivector \textit{a} ? \textit{b} as the oriented parallelogram spanned by $a$ and $b$. The cross product is then obtained by taking the Hodge star of the bivector \textit{a} ? \textit{b}, mapping 2-vectors to vectors:



$$a \times b = \star (a \wedge b) \,.$$



This can be thought of as the oriented multi-dimensional element "perpendicular" to the bivector. Only in three dimensions is the result an oriented one-dimensional element – a vector – whereas, for example, in four dimensions the Hodge dual of a bivector is two-dimensional – a bivector. So, only in three dimensions can a vector cross product of \textit{a} and \textit{b} be defined as the vector dual to the bivector \textit{a} ? \textit{b}: it is perpendicular to the bivector, with orientation dependent on the coordinate system's handedness, and has the same magnitude relative to the unit normal vector as \textit{a} ? \textit{b} has relative to the unit bivector; precisely the properties described above.



\paragraph{Handedness}



\paragraph{Consistency}



When physics laws are written as equations, it is possible to make an arbitrary choice of the coordinate system, including handedness. One should be careful to never write down an equation where the two sides do not behave equally under all transformations that need to be considered. For example, if one side of the equation is a cross product of two polar vectors, one must take into account that the result is an axial vector. Therefore, for consistency, the other side must also be an axial vector. More generally, the result of a cross product may be either a polar vector or an axial vector, depending on the type of its operands (polar vectors or axial vectors). Namely, polar vectors and axial vectors are interrelated in the following ways under application of the cross product:



\begin{itemize}
    \item polar vector × polar vector = axial vector
    \item axial vector × axial vector = axial vector
    \item polar vector × axial vector = polar vector
    \item axial vector × polar vector = polar vector
\end{itemize}

or symbolically



\begin{itemize}
    \item polar × polar = axial
    \item axial × axial = axial
    \item polar × axial = polar
    \item axial × polar = polar
\end{itemize}

Because the cross product may also be a polar vector, it may not change direction with a mirror image transformation. This happens, according to the above relationships, if one of the operands is a polar vector and the other one is an axial vector (e.g., the cross product of two polar vectors). For instance, a vector triple product involving three polar vectors is a polar vector.



A handedness-free approach is possible using exterior algebra.



\paragraph{The paradox of the orthonormal basis}



Let $\left( \mathbf{i}, \mathbf{j}, \mathbf{k} \right)$ be an orthonormal basis. The vectors $\mathbf{i}$, $\mathbf{j}$, and $\mathbf{k}$ don't depend on the orientation of the space. They can even be defined in the absence of any orientation. They can not therefore be axial vectors. But if $\mathbf{i}$ and $\mathbf{j}$ are polar vectors then $\mathbf{k}$ is an axial vector for $\mathbf{i} \times \mathbf{j} = \mathbf{k}$ or $\mathbf{j} \times \mathbf{i} = \mathbf{k}$. This is a paradox.



"Axial" and "polar" are \textit{physical} qualifiers for \textit{physical} vectors; that is, vectors which represent \textit{physical} quantities such as the velocity or the magnetic field. The vectors $\mathbf{i}$, $\mathbf{j}$, and $\mathbf{k}$ are mathematical vectors, neither axial nor polar. In mathematics, the cross-product of two vectors is a vector. There is no contradiction.



\paragraph{Resources}



\begin{itemize}
    \item \href{https://openstax.org/books/calculus-volume-3/pages/2-4-the-cross-product}{The Cross Product}, OpenStax
    \item \href{https://omega0.xyz/omega8008/calc3/cross-product-dir/cornell-lecture.html}{Cross Product}, Cornell University
\end{itemize}

\paragraph{Licensing}



Content obtained and/or adapted from:



\begin{itemize}
    \item \href{https://en.wikipedia.org/wiki/Cross_product}{Cross product, Wikipedia} under a CC BY-SA license
\end{itemize}

\subsection{Learning Objectives}

\begin{enumerate}
    \item \textbf{Calculate the Cross Product:}
    \item \textbf{Outcome:} Students will be able to calculate the cross product of two vectors in three-dimensional space.
    \item \textbf{Details:} This includes applying the cross product formula, using the right-hand rule to determine the direction of the resulting vector, and interpreting the geometric significance of the cross product, such as the area of the parallelogram formed by the vectors.
    \item \textbf{Apply Properties of the Cross Product:}
    \item \textbf{Outcome:} Students will understand and apply the properties of the cross product, including anticommutativity and distributivity over addition.
    \item \textbf{Details:} This includes proving and using the properties in various mathematical contexts, such as solving problems that require the use of these properties and verifying them through vector algebra.
    \item \textbf{Utilize the Cross Product in Physics and Engineering:}
    \item \textbf{Outcome:} Students will be able to apply the cross product to solve problems in physics and engineering, such as calculating torque, angular momentum, and the Lorentz force.
    \item \textbf{Details:} This includes setting up and solving problems involving the cross product in the context of physical quantities, interpreting the physical meaning of the results, and using appropriate units and directions in vector calculations.
\end{enumerate}

\subsection{Assessment Instrument}

\begin{enumerate}
    \item \textbf{Assessment Instrument: Cross Product Quiz} \textbf{Format:}
    \item \textbf{Type:} Mixed-format quiz (multiple-choice, short answer, and problem-solving questions)
    \item \textbf{Duration:} 45 minutes
    \item \textbf{Total Marks:} 30 points \textbf{Sections:}
    \item \textbf{Definition and Calculation of the Cross Product (10 points):}
    \item \textbf{Question 1:} (Multiple-choice) Which of the following is true about the cross product of two vectors in three-dimensional space?
    \item a. The cross product is a scalar.
    \item b. The cross product is a vector perpendicular to the plane containing the original vectors.
    \item c. The cross product is always zero.
    \item d. The cross product is a vector parallel to the original vectors.
    \item \textbf{Question 2:} (Short answer) Given the vectors \(\mathbf{a} = \begin{bmatrix} 2 \\ 3 \\ 4 \end{bmatrix}\) and \(\mathbf{b} = \begin{bmatrix} 5 \\ 6 \\ 7 \end{bmatrix}\), calculate the cross product \(\mathbf{a} \times \mathbf{b}\).
    \item \textbf{Properties of the Cross Product (10 points):}
    \item \textbf{Question 3:} (Multiple-choice) Which property of the cross product is demonstrated by the equation \(\mathbf{a} \times \mathbf{b} = -(\mathbf{b} \times \mathbf{a})\)?
    \item a. Distributivity
    \item b. Anticommutativity
    \item c. Associativity
    \item d. Identity
    \item \textbf{Question 4:} (Short answer) Given the vectors \(\mathbf{u} = \begin{bmatrix} 1 \\ 0 \\ 0 \end{bmatrix}\), \(\mathbf{v} = \begin{bmatrix} 0 \\ 1 \\ 0 \end{bmatrix}\), and \(\mathbf{w} = \begin{bmatrix} 0 \\ 0 \\ 1 \end{bmatrix}\), verify the distributive property of the cross product by calculating \(\mathbf{u} \times (\mathbf{v} + \mathbf{w})\) and \(\mathbf{u} \times \mathbf{v} + \mathbf{u} \times \mathbf{w}\).
    \item \textbf{Applications of the Cross Product (10 points):}
    \item \textbf{Question 5:} (Problem-solving) A force vector \(\mathbf{F} = \begin{bmatrix} 3 \\ -2 \\ 5 \end{bmatrix}\) is applied at a point with position vector \(\mathbf{r} = \begin{bmatrix} 2 \\ 3 \\ -1 \end{bmatrix}\). Calculate the torque \(\mathbf{\tau}\) about the origin using the cross product \(\mathbf{\tau} = \mathbf{r} \times \mathbf{F}\).
    \item \textbf{Question 6:} (Problem-solving) Find the area of the parallelogram formed by the vectors \(\mathbf{p} = \begin{bmatrix} 1 \\ 2 \\ 3 \end{bmatrix}\) and \(\mathbf{q} = \begin{bmatrix} 4 \\ 5 \\ 6 \end{bmatrix}\) using the cross product. \textbf{Scoring Rubric:}
    \item \textbf{Definition and Calculation of the Cross Product:}
    \item Question 1: 2 points for correct answer.
    \item Question 2: 8 points for correct calculation and vector result.
    \item \textbf{Properties of the Cross Product:}
    \item Question 3: 2 points for correct answer.
    \item Question 4: 8 points for correct verification of distributive property.
    \item \textbf{Applications of the Cross Product:}
    \item Question 5: 5 points for correct torque calculation.
    \item Question 6: 5 points for correct area calculation. This quiz will help assess students' understanding of the definition, properties, and applications of the cross product, ensuring they can perform necessary calculations and understand the geometric significance of the cross product in various contexts.
\end{enumerate}

%%===============================================================
\section{Lines and Planes in Space}
\label{sec:lines-and-planes-in-space}
%%===============================================================

For many practical applications you have to work with the mathematical descriptions of \textbf{lines}, \textbf{planes, curves, and surfaces} in 3-dimensional space. This requires some knowledge on vectors and the ability to construct 3-dimensional graphs without using calculators.



\paragraph{Lines in 3-dimensional space}



Although the equation for lines is discussed in previous chapters, this chapter will explain more in detail about the properties and important aspects of lines, as well as the expansion into general curves in 3-dimensional space.



\paragraph{Review of parametric equations}



Parametric equations use a different variable to express the relation between two variables. For example, look at the following equation of a circle:



> $x^2+y^2=1$



If we express variables $x$ and $y$ with a new variable $t$, and we know that $\cos ^2 t+\sin^2 t=1$, we can rewrite the original equation into:



> $x=\cos t$

>

> $y=\sin t$



The equation above is the parametric form of a circle with a radius of $1$.



Now, let's talk about lines in 3-dimensional space.



\paragraph{Equations for lines in 3-dimensional space}



A line in space is defined by two points in space, which I will call $P _1$ and $P _2$. Let $\mathbf{r} _0$ be the vector from the origin to $P _1$, and $\mathbf{r} _2$ the vector from the origin to $P _2$. Given these two points, every other point $P$ on the line can be reached by



$$\mathbf{r} = \mathbf{r} _0 + \mathbf{v}t$$



where $\mathbf{v}$ is the vector from $P _1$ and $P _2$:



$\mathbf{v} = \mathbf{r} _0 - \mathbf{r} _2$



% [Image: Line in 3D Space.]



Figure: Line in 3D Space.



To intuitively understand the equation of a line, imagine that there is a line going through the end point of the vector $\mathbf{r} _0$ and stretches in the direction of the vector $\mathbf{v}$. Just view it as a vector version of the point-slope form. For example, assume there is a line in 3-dimensional space with the equation:



> $\mathbf{r}=\langle2,5,4\rangle+\langle1,1,1\rangle t$



We can graph the line by first finding the point $(2,5,4)$. Then, we stretch the line in the direction parallel with the vector $\langle1,1,1\rangle$. The final graph is the graph of the line $\mathbf{r}=\langle2,5,4\rangle+\langle1,1,1\rangle t$. Sometimes the direction vector $\mathbf{v}$ is unknown. However, it can be easily solved by finding another point on the line. In this case, the point $(5,8,7)$ is on the line. By calculating the vector between the two points, we can find out that the direction vector is $\langle3,3,3\rangle$. Thus, the equation for the line can be written as



> $\mathbf{r}=\langle2,5,4\rangle+\langle3,3,3\rangle t _1$



which is equivalent with the original equation $\mathbf{r}=\langle2,5,4\rangle+\langle1,1,1\rangle t$ because if we let $t=3t _1$, the original equation will turn into the equation above. Therefore, there are infinitely many ways to write an equation of a line in 3-dimensional space using the vector form.



Now, there is a parametric form of expressing a line. Recall that there is another way to write down vectors: $a\mathbf{i}+b\mathbf{j}+c\mathbf{k}$. So, we can rewrite the original equation $\mathbf{r}=\langle2,5,4\rangle+\langle1,1,1\rangle t$ into:



> $\mathbf{r}=\langle2+t,5+t,4+t\rangle=(2+t)\mathbf{i}+(5+t)\mathbf{j}+(4+t)\mathbf{k}$



Then, we assign the respective parts to the corresponding $x,y,z$ axes



> $x=2+t$

>

> $y=5+t$

>

> $z=4+t$



This is the parametric form of the equation for lines.



The final common way to notate lines is the symmetric equations, which is just another slight transformation from the parametric form:



> $x-2=y-5=z-4$



To summarize, there are basically three ways to write the equation for a line that passes through the point $(x _0,y _0,z _0)$ with the direction $\langle a,b,c\rangle$.



$$\begin{array}{|c|c|c|} \hline \textbf{Vector Equation} & \textbf{Parametric Equations} & \textbf{Symmetric Equations} \\\\ \hline \mathbf{r} = \langle x _0, y _0, z _0 \rangle + t \langle a, b, c \rangle & x = x _0 + at & \frac{x - x _0}{a} = \frac{y - y _0}{b} = \frac{z - z _0}{c} \\\\  & y = y _0 + bt & \\\\  & z = z _0 + ct & \\\\ \hline \end{array}$$



\paragraph{Relations between two lines}



There are lines which intersect each other and those which do not. There can be parallel, perpendicular, and skew lines, which will be discussed in this part.



\textbf{Intersection}



Assume there are two lines with the equations:



> $L _1:\langle2,5,4\rangle+\langle1,1,1\rangle t$ or in parametric form

> $L _1: \begin{cases} x=2+t  \\\\ y=5+t  \\\\ z=4+t \end{cases}$



> $L _2:\langle5,-1,0\rangle+\langle -\frac{1}{2},4,3\rangle s$ or in

> parametric form $L _2: \begin{cases} x=5-\frac{1}{2}s  \\\\ y=-1+4s  \\\\ z=3s \end{cases}$



To find out if they intersect or not, we just need to solve the system

of equations



> $\begin{cases} 2+t=5-\frac{1}{2}s  \\\\ 5+t=-1+4s \\\\ 4+t=3s \end{cases}$



If the system of equations has only one solution, then the two lines intersect at one point. If the system of equations has infinitely many solutions, then the two lines are the same. If the system of equations does not have a solution, then the two lines do not intersect at all. In this case, there is a solution:



> $\begin{cases} t=2  \\\\ s=2 \end{cases}$



Thus the two lines intersect at point $(4,7,6)$.



If we want to further know the angle between the two lines, we can apply the dot product formula. The angle between the two lines should be:



> $\theta=\arccos(\frac{\mathbf{v _1}\,\cdot\,\mathbf{v _2}}{\|\mathbf{v _1}\|\|\mathbf{v _2}\|})$



\textbf{Parallel}



To spot two parallel lines in 3-dimensional space, we just need to look at the direction vector $\mathbf{v}$. If the direction vectors of the two lines $\mathbf{v _1},\mathbf{v _2}$ have such relationship that $\mathbf{v _2}=k\mathbf{v _1}, k\in\mathbb{R}$, then the two lines are parallel with each other. For example, the two lines $L _1:\langle2,5,4\rangle+\langle1,1,1\rangle t$ and $L _2:\langle5,-1,0\rangle+\langle -3,-3,-3\rangle s$ are parallel with each other because $\frac{\langle-3,-3,-3\rangle}{-3}=\langle1,1,1\rangle$.



\textbf{Perpendicular}



To have two lines perpendicular with each other, those two lines must intersect first. If they do intersect, recall the dot product of vectors. The dot product states that:



> $\mathbf{a}\,\cdot\,\mathbf{b}=\|\mathbf{a}\|\|\mathbf{b}\|\cos\theta$



If two lines are perpendicular with each other, $\theta=\frac{\pi}{2}$, which results in $\mathbf{a}\,\cdot\,\mathbf{b}=\|\mathbf{a}\|\|\mathbf{b}\|\cos\frac{\pi}{2}=0$.



So, if we continue this flow of thought, we can find that if we choose two vectors on each line, have a dot product between them, and the result is zero, then we can safely say that the two lines are perpendicular with each other. However, there is a more convenient way to simplify the process. Instead of finding two vectors on each line, we just need to apply the dot product on the two direction vectors because the direction vectors are calculated from points on the respective lines.



So if we have two lines:



> $L _1:\langle2,5,4\rangle+\langle1,1,1\rangle t$ and

> $L _2:\langle6,6,5\rangle+\langle -2,1,1\rangle s$



They are perpendicular with each other because



Here is the corrected Markdown with LaTeX to match the formatting shown in the image:



> They are perpendicular with each other because

>

> $$ \begin{cases} 2 + t = 6 - 2s \\\\ 5 + t = 6 + s \\\\ 4 + t = 5 + s \end{cases} $$

>

> has one and only solution:

>

> $$ \begin{cases} t = 2 \\\\ s = 1 \end{cases} $$

>

> which means they intersect. And

>

> $$ \langle 1, 1, 1 \rangle \cdot \langle -2, 1, 1 \rangle = 0 $$



Thus ending the proof.



\textbf{Skew lines}



Skew lines are lines that do not intersect and are not parallel with each other. For example, lines $L _1:\langle2,5,4\rangle+\langle1,1,1\rangle t$ and $L _2:\langle5,-1,0\rangle+\langle 3,-3,-3\rangle s$ are skew lines.



\textbf{Distance between two skew lines}



To solve this problem, we need to know more about planes in 3-dimensional space, which will be discussed below.



\paragraph{Planes in 3-dimensional space}



\paragraph{Introduction}



The same idea can be used to describe a plane in 3-dimensional space, which is uniquely defined by three points (which do not lie on a line) in space ($P _1, P _2, P _3$). Let $\mathbf{x} _i$ be the vectors from the origin to $P _i$. Then



$$\mathbf{x} = \mathbf{x} _1 + \lambda \mathbf{a} + \mu \mathbf{b}$$



with:



$$\mathbf{a} = \mathbf{x} _2 - \mathbf{x} _1 \,\, \text{and} \,\, \mathbf{b} = \mathbf{x} _3 - \mathbf{x} _1$$



Note that the starting point does not have to be $\mathbf{x} _1$, but can be any point in the plane. Similarly, the only requirement on the vectors $\mathbf{a}$ and $\mathbf{b}$ is that they have to be two non-collinear vectors in our plane.



Recall that in 2-dimensional vectors, if there are two vectors $\mathbf{i},\mathbf{j}$, any vector on the Cartesian plane can be expressed in terms of the vectors $\mathbf{i},\mathbf{j}$ by $a\mathbf{i}+b\mathbf{j} \quad (a,b\in\mathbb{R})$. Using the same method, we can deduce that



> the $\lambda \mathbf{a} + \mu \mathbf{b}$ part tells us that the graph

> should be a plane, and the $\mathbf{x} - \mathbf{x} _1$ part describes

> the "slope" and the axes intersections



However, there two other ways to express a plane in 3-dimensional space: the vector equation and the scalar equation.



\paragraph{Vector and scalar equations for planes}



The vector equation for planes requires us to understand the power of the dot product. We already knew that when the dot product of two vectors is zero, the two vectors should be perpendicular with each other. Now, imagine a vector $\mathbf{n}$ in 3-dimensional space. If we graph out all the vectors $\Delta\mathbf{r}$ that are perpendicular with $\mathbf{n}$, what will the result be?



The result should be a plane with a vector perpendicular with the plane. Thus, the vector equation for planes is simply $\mathbf{n}\,\cdot\,\Delta\mathbf{r}=0$.



> the vector $\mathbf{n}$ is the normal vector, which is perpendicular

> with the plane



> the vector $\Delta\mathbf{r}=\mathbf{r}-\mathbf{r _0}$ is the variable

> vector where $\mathbf{r}=\langle x,y,z \rangle$ (unknown points on the

> plane) and $\mathbf{r _0}=\langle x _0,y _0,z _0 \rangle$ (a given point

> on the plane). This expressions simply means all the vectors on the

> plane.



Of course, the vector equation for planes can be rewritten as $\mathbf{n}\,\cdot\,(\mathbf{r}-\mathbf{r _0})=0$ or $\mathbf{n}\,\cdot\,\mathbf{r}=\mathbf{n}\,\cdot\,\mathbf{r _0}$, depending on the writer.



To find out the scalar equation, we just need to calculate the dot product and do some simplifications. So, let's assume that



> $\mathbf{n}=\langle a,b,c \rangle$,

> $\mathbf{r}=\langle x,y,z \rangle$, and

> $\mathbf{r _0}=\langle x _0,y _0,z _0 \rangle$

>

> According to the vector equation

>

> > $\mathbf{n}\,\cdot\,(\mathbf{r}-\mathbf{r _0})=0$

>

> Thus, we have

>

> > $\langle a,b,c \rangle\,\cdot\,\langle x-x _0,y-y _0,z-z _0\rangle =0$

>

> After some algebraic manipulations, we can get

>

> > $ax+by+cz-(ax _0+by _0+cz _0)=0$

>

> Since the $ax _0+by _0+cz _0$ is a constant, we just let

> $ax _0+by _0+cz _0=-d$. As a result, the scalar equation for planes is

>

> $ax+by+cz+d=0$



Note that the constants $a,b,c$ are the same as the $\mathbf{i},\mathbf{j},\mathbf{k}$ components of the normal vector. This property will be extremely useful when discussing the relations between two planes. Also, $-d=\mathbf{n}\,\cdot\,\mathbf{r _0}$.



$\blacksquare$



To summarize, there are three ways to notate a plane in 3-dimensional space, with the latter two more commonly used:

$$\begin{array}{|c|c|c|} \hline

\textbf{Expansion from the equation for lines} & \textbf{Vector equation} & \textbf{Scalar equation} \\\\ \hline

\mathbf{x} = \mathbf{x} _1 + \lambda \mathbf{a} + \mu \mathbf{b} & \mathbf{n} \cdot (\mathbf{r} - \mathbf{r} _0) = 0 & ax + by + cz + d = 0 \\\\ \hline

\end{array}

$$



\paragraph{Relations between two planes}



\textbf{Parallel}



The normal vector is important because it determined the how the plane is shaped. So when we are discussing the relations between two planes, we are actually trying to find out the connections between the normal vectors of the two planes. In this case, assume there are two planes in the 3-dimensional space: $a _1x+b _1y+c _1z+d _1=0$ and $a _2x+b _2y+c _2z+d _2=0$. The normal vectors for the planes should be:



> $\mathbf{n _1}=\langle a _1,b _1,c _1 \rangle,\,\mathbf{n _2}=\langle a _2,b _2,c _2 \rangle$



Since normal vectors are perpendicular with their corresponding planes, if the normal vectors are parallel with each other, the respective planes should also be parallel with each other. Thus,



> If $\mathbf{n _1}=k\mathbf{n _2}\quad(k\in\mathbb{R})$, then

> $a _1x+b _1y+c _1z+d _1=0$ and $a _2x+b _2y+c _2z+d _2=0$ are parallel with

> each other.



\textbf{Perpendicular}



If the normal vectors are perpendicular with each other, the respective planes will be perpendicular. In other words, if $\mathbf{n _1}\perp\mathbf{n _2}$, then the planes are perpendicular with each other.



\textbf{Intersection}



To fully understand how to find the intersection between two lines, we should be familiar with the normal vector and its potential. If the two planes are not parallel and are not basically the same plane, then they must intersect with each other. The intersection should form a line. Imagine there are two planes $a _1x+b _1y+c _1z+d _1=0$ and $a _2x+b _2y+c _2z+d _2=0$, and their respective normal vectors $\mathbf{n _1}\ne k\mathbf{n _2}\quad(k\in\mathbb{R})$ (which means they are not parallel).



> Because normal vectors are perfectly perpendicular to all vectors on

> the plane, the opposite is also true: all vectors on the plane are

> perpendicular to their respective normal vectors. This is why

> $\mathbf{n}\,\cdot\,(\mathbf{r}-\mathbf{r _0})=0$ is the vector

> equation of the plane. Since the intersection of the two planes is a

> line, we can say that the direction vector $\mathbf{v}$ should be on

> both planes.



> Because $\mathbf{v}$ is on both planes, $\mathbf{v}$ should be perpendicular with both normal vectors

>

> > $\mathbf{v}\perp\mathbf{n _1},\,\mathbf{v}\perp\mathbf{n _2}$

>

> Recall that the cross product of two vectors will result in a new vector that is perpendicular with both the original vectors. We can calculate the cross product of $\mathbf{n _1}, \mathbf{n _2}$ to create $\mathbf{v}$.

>

> $\mathbf{n _1}\times\mathbf{n _2}= \begin{vmatrix} \mathbf{i}&\mathbf{j}&\mathbf{k} \\\\ a _1 & b _1 & c _1 \\\\ a _2 & b _2 & c _2 \end{vmatrix}= (b _1c _2-b _2c _1)\mathbf{i}-(a _1c _2-a _2c _1)\mathbf{j}+(a _1b _2-a _2b _1)\mathbf{k}$

>

> So the direction vector for the line is

> $\mathbf{v}=\mathbf{n _1}\times\mathbf{n _2}=\langle b _1c _2-b _2c _1,a _1c _2-a _2c _1,a _1b _2-a _2b _1\rangle$



> We still need to know a point on the line to finish the equation because $\mathbf{r} = \mathbf{r}_0 + \mathbf{v}t$. To find a point, simply let $z = 0$ and solve for $x, y$ in the following system of equations:

>

> $$ \begin{cases} a _1 x + b _1 y + d _1 = 0 \\\\ a _2 x + b _2 y + d _2 = 0 \end{cases} \Rightarrow \begin{cases} x = \frac{b _1 d _2 - b _2 d _1}{a _1 b _2 - a _2 b _1} \\\\ y = \frac{a _1 d _2 - a _2 d _1}{a _2 b _1 - a _1 b _2} \end{cases} $$

>

> Because the solution is too complicated to write down, we will let $x=\frac{b _1d _2-b _2d _1}{a _1b _2-a _2b _1}=x _0$ and $y=\frac{a _1d _2-a _2d _1}{a _2b _1-a _1b _2}=y _0$. Thus, the point $(x _0,y _0,0)$ is on both planes (recall that we let $z=0$) and $\mathbf{P _0}=\langle x _0,y _0,0 \rangle$.



> Now, we know a point on the line $(x _0,y _0,0)$ and the direction vector

> $\mathbf{v}=\mathbf{n _1}\times\mathbf{n _2}=\langle b _1c _2-b _2c _1,a _1c _2-a _2c _1,a _1b _2-a _2b _1\rangle$,

> the intersection between the two planes is:

>

> $L:\langle \frac{b _1d _2-b _2d _1}{a _1b _2-a _2b _1},\frac{a _1d _2-a _2d _1}{a _2b _1-a _1b _2},0 \rangle+\langle b _1c _2-b _2c _1,a _1c _2-a _2c _1,a _1b _2-a _2b _1\rangle t$



For those who prefer a cleaner expression:



> $L:\langle x _0,y _0,0 \rangle+\langle b _1c _2-b _2c _1,a _1c _2-a _2c _1,a _1b _2-a _2b _1\rangle t$



> $L:\mathbf{P _0}+(\mathbf{n _1}\times\mathbf{n _2})t$



And if we want to know the angle between the planes, similar to how we

find out the angle between two lines, we apply the dot product:



> $\theta=\arccos(\frac{\mathbf{n _1}\,\cdot\,\mathbf{n _2}}{\|\mathbf{n _1}\|\|\mathbf{n _2}\|})$



$\blacksquare$



\textbf{Distance between two parallel planes}



The distance between two non-parallel planes is zero because they intersect. So, we should focus on the distance between parallel planes. Before we do that, it will be more convenient if we know the distance from a point to a plane.



> Let the distance be $D$, the point be $P=(x _1,y _1,z _1)$, the plane be

> $ax+by+cz+d=0$.



> Knowing the equation for the plane can help us know the normal vector,

> since the normal vector is perpendicular to the plane, the exact

> direction we need. $\mathbf{n}=\langle a,b,c \rangle$.



Now, we start solving. First, assume there is point $P _0=(x _0,y _0,z _0)$ on the plane $ax+by+cz+d=0$. Then, we create a vector from $P _0$ to $P$: $\overrightarrow{P _0 P}=\langle x _1-x _0,y _1-y _0,z _1-z _0 \rangle$. We will also let the angle between vectors $\overrightarrow{P _0 P}$ and $\mathbf{n}$ be $\theta$. If we graph out what it looks like, we can easily understand how the distance is deduced.



> $D= | \|\overrightarrow{P _0 P}\|\ \cos\theta |$ (the absolute value is

> needed because the distance should always be larger than zero)



However, we do not know $\theta$. We can calculate $\theta$ by applying the dot product. But, there is a simpler way:



> Using some very interesting manipulation, we can get

>

> > $$D=\frac{|\|\mathbf{n}\|\|\overrightarrow{P _0 P}\|\cos\theta|}{\|\mathbf{n}\|}$$

>

> As we can see, the nominator is actually the other way of expressing

> the dot product of two vectors. So, we can stop worrying about not

> knowing the value of $\theta$.

>

> > $D=\frac{|\mathbf{n}\ \cdot\ \overrightarrow{P _0 P}|}{\|\mathbf{n}\|}$

>

> We can now substitute those vectors with their coordinates. After some

> algebra, we get

>

> > $$D=\frac{|\mathbf{n}\ \cdot\ \overrightarrow{P _0 P}|}{\|\mathbf{n}\|}  =\frac{|ax _1+by _1+cz _1-(ax _0+by _0+cz _0)|}{\sqrt{a^2+b^2+c^2}}$$

>

> Since $P _0=(x _0,y _0,z _0)$ is on the plane, $ax _0+by _0+cz _0=-d$. We can

> further simplify the formula into

> $D=\frac{|ax _1+by _1+cz _1+d|}{\sqrt{a^2+b^2+c^2}}$, thus finishing our

> deduction.



There are other ways to write down the formula. For people who prefer simpler notations, the following formulae are also ways to express the distance:



> $D=\frac{|\mathbf{n}\ \cdot\ \overrightarrow{P _0 P}|}{\|\mathbf{n}\|}$

> or $D=\mathrm{comp} _\mathbf{n} \overrightarrow{P _0 P}$



$\blacksquare$



After we do some further "investigation" on the distance formula, we can find that only the equation of the plane and one point are needed to calculate the distance between them, which is very convenient because the fewer things we need to solve a problem, the more convenient the problem is. When we try to find the distance between two parallel planes, we just need one point on one plane and the equation for the other plane to solve.



\textbf{Distance between two skew lines (continue)}



Assume that lines $L _1,L _2$ are skew lines, we can calculate the distance by assuming that those lines belong in two parallel planes $P _1,P _2$ respectively. Then, the problem changes from solving the distance between two skew lines to solving the distance between two parallel planes.



We still need to know the normal vector of both planes. We can simply apply the cross product on the two direction vectors of both lines: $\mathbf{n}=\mathbf{v _1}\times\mathbf{v _2}$. $\mathbf{n}\ne0$ because $\mathbf{v _1},\mathbf{v _2}$ are not pointing towards the same direction. Now we can apply the newly deduced distance formula on the two skew lines.



\paragraph{Cylinders and quadric surfaces}



This section of the chapter requires some understanding of conic sections.



\paragraph{Cylinders}



% [Image: An example of a parabolic cylinder in 3-dimensional space]



Figure: An example of a parabolic cylinder in 3-dimensional space.



A cylinder is a surface that consists of all lines that are parallel to a given line and pass through a given plane curve. There are several special cylinders such as the parabolic cylinder and the circular cylinder. For example, the image on the right is a parabolic cylinder. Parabolic cylinder typically has an equation of:



> $z=x^2$ etc.



If we want to move the cylinder around without rotating it, we can have the equation:



> $z-k=(x-h)^2$, where $h,k$ are constants.



It is just like the parabola we have discussed early but with a new dimension. Circular cylinders, similar to how we derive the parabolic cylinder, looks just like a cylinder with a circle as its "base", with an equation of:



> $x^2+y^2=1$ etc.



If we want to make the circular cylinder more "elliptic", just like how we derive the equation for an ellipse, the elliptic cylinder has an equation of:



> $\frac{x^2}{a^2}+\frac{y^2}{b^2}=1$, where $a,b$ are constants.



It is just like an ellipse on the $xy$-plane but infinitely stretched in the $z$ direction.



\paragraph{Quadric surfaces}



General equation for quadric surfaces: The general equation for quadric surfaces is:



$$Ax^2+By^2+Cz^2+Dxy+Eyz+Fxz+Gx+Hy+Iz+J=0$$ where $A,B,C,D,E,F,G,H,I,J$ are constants. It looks similar to the general equation for conic sections ($Ax^2+Bxy+Cy^2+Dx+Ey+F=0$), except it has one more variable $z$. After some translations and rotations, we can simplify the general equation into the standard equations:



The standard equations for quadric surfaces are:



$$Ax^2+By^2+Cz^2+J=0\quad\text{or}\quad Ax^2+By^2+Iz=0$$ where $A,B,C,I,J$ are constants.}} Of course, depending on the specific quadric surface, there are different forms of standard equations, which we will see in the following discussions.



\paragraph{Ellipsoid}



% [Image: An ellipsoid]



Figure: An ellipsoid with an equation $x^2+\frac{y^2}{2}+\frac{z^2}{3}=1$. See that if we intersect the ellipsoid with planes parallel to axis planes, the intersection will always be an ellipse.



Ellipsoids have the equation:

> $\frac{x^2}{a^2}+\frac{y^2}{b^2}+\frac{z^2}{c^2}=1$



It is difficult to sketch quadric surfaces because instead of two variables, there are three, making the process very complicated without the help of calculators. However, there is a way to make us easier to understand. We can analyze each plane to see what the shape looks like, and then combine what we analyzed from each plane to form a rather complete graph of the surface. Take this for example:



> $x^2+\frac{y^2}{9}+\frac{z^2}{4}=1$



Let us first examine the $xy$-plane. To examine the $xy$-plane, we need to imagine $z$ as a constant. In this case, imagine that $z=k$, so



> $x^2+\frac{y^2}{9}=1-\frac{k^2}{4}$, where $-2\leq k\leq2$.



We can see that in the $xy$-plane, the graph will be like an ellipse. We call that the horizontal trace in the plane $z=k$ is an ellipse.



Let us furthermore analyze the vertical traces in the planes $x=k,y=k$.



% [Image: An elliptic paraboloid]



Figure: An elliptic paraboloid with equation $z=x^2+y^2$. Note that the horizontal trace is an ellipse and both vertical traces are parabolas.



> $\frac{y^2}{9}+\frac{z^2}{4}=1-k^2\quad x=k$, where $-1\leq k\leq 1$.

>

> $x^2+\frac{z^2}{4}=1-\frac{k^2}{9}\quad y=k$, where $-3\leq k\leq3$.



We can see that on both vertical traces in the planes $x=k,y=k$ are ellipses. In other words, the graphs in both $xz$-plane and $yz$-plane look like ellipses. Since all traces are ellipses, the surface is an ellipsoid with vertices $(\pm1,0,0),(0,\pm3,0),(0,0,\pm2)$.



\paragraph{Elliptic paraboloid}



Elliptic paraboloids have the equation:



> $\frac{z}{c}=\frac{x^2}{a^2}+\frac{y^2}{b^2}$



% [Image: A hyperbolic paraboloid]



Figure: A hyperbolic paraboloid with equation $z=x^2-y^2$



There are variations of the equation depending on which trace is an ellipse. For the equation above, the horizontal trace is an ellipse as seen in the image on the right.



Take $z=x^2+y^2$ for example. the horizontal trace, which means $z=k$, is:



> $x^2+y^2=k$, where $k>0$.



For both vertical traces, we can see that those traces have the shape of a parabola.



> $z=x^2+k^2\quad y=k$

>

> $z=y^2+k^2\quad x=k$



The vertex for this elliptic paraboloid is $(0,0,0)$.



\paragraph{Hyperbolic paraboloid}



Hyperbolic paraboloids have the equation:



> $\frac{z}{c}=\frac{x^2}{a^2}-\frac{y^2}{b^2}$



% [Image: An image of a cone]



Figure: An image of a cone. Note that vertical traces in the planes $x=k$ and $y=k$ are hyperbolas if $k\ne0$ but are pairs of lines if $k=0$.



The respective traces are



> $\frac{x^2}{a^2}-\frac{y^2}{b^2}=\frac{k}{c}\quad z=k$, where

> $\frac{x^2}{a^2}-\frac{y^2}{b^2}=\frac{k}{c}\text{ when }k>0

> \quad\text{and}\quad

> \frac{y^2}{b^2}-\frac{x^2}{a^2}=\frac{k}{c}\text{ when }k

> $\frac{z}{c}=\frac{x^2}{a^2}-\frac{k^2}{b^2}\quad y=k$. The vertical trace is a parabola.

>

> $\frac{z}{c}=\frac{k^2}{a^2}-\frac{y^2}{b^2}\quad x=k$. The vertical trace is a parabola.



\paragraph{Cone, hyperboloid of one sheet, and hyperboloid of two sheets}



% [Image: An image of a hyperboloid of one sheet.]



Figure: An image of a hyperboloid of one sheet.



% [Image: An image of a hyperboloid of two sheets.]



Figure: An image of a hyperboloid of two sheets. Note that horizontal traces in $z=k$ are ellipses if $k>c$ or $k<-c$.



Cones have the equation:

> $\frac{z^2}{c^2}=\frac{x^2}{a^2}+\frac{y^2}{b^2}$



Note that $\frac{x^2}{a^2}+\frac{y^2}{b^2}=0$ when $z=0$.



Hyperboloids of one sheet have the equation:



> $\frac{x^2}{a^2}+\frac{y^2}{b^2}-\frac{z^2}{c^2}=1$



Note that $\frac{x^2}{a^2}+\frac{y^2}{b^2}=1$ when $z=0$.



Hyperboloids of two sheets have the equation:



> $-\frac{x^2}{a^2}-\frac{y^2}{b^2}+\frac{z^2}{c^2}=1$



Note that when $-c\leq z\leq c$, there is no real solution for $x,y$.



All of the three quadric surfaces have the same respective traces. The horizontal traces are all ellipses and the vertical traces are all hyperbolas.



\paragraph{Resources}



\begin{itemize}
    \item \href{https://math.libretexts.org/Bookshelves/Calculus/Book\%3A_Calculus_(OpenStax}{Equations of Lines and Planes in Space}/12\%3A_Vectors_in_Space/12.5\%3A_Equations_of_Lines_and_Planes_in_Space), Mathematics LibreTexts
    \item \href{https://openstax.org/books/calculus-volume-3/pages/2-6-quadric-surfaces}{Quadric Surfaces}, OpenStax
    \item \href{https://tutorial.math.lamar.edu/classes/calciii/eqnsoflines.aspx}{Equations of Lines}, Paul's Online Notes
    \item \href{https://tutorial.math.lamar.edu/classes/calciii/eqnsofplanes.aspx}{Equations of Planes}, Paul's Online Notes
    \item \href{https://tutorial.math.lamar.edu/classes/calciii/quadricsurfaces.aspx}{Quadric Surfaces}, Paul's Online Notes
\end{itemize}

\paragraph{Licensing}



Content obtained and/or adapted from:



\begin{itemize}
    \item \href{https://en.wikibooks.org/wiki/Calculus/Curves_and_Surfaces_in_Space}{Curves and Surfaces in Space, Wikibooks: Calculus} under a CC BY-SA license
\end{itemize}

\subsection{Learning Objectives}

\begin{enumerate}
    \item \textbf{Understanding and Applying Vector Equations for Lines in 3D Space:}
    \item \textbf{Outcome:} Students will be able to construct and interpret vector equations of lines in three-dimensional space.
    \item \textbf{Details:} This includes identifying direction vectors and points through which the line passes, as well as converting between vector, parametric, and symmetric forms of line equations.
    \item \textbf{Analyzing Relationships Between Lines in 3D Space:}
    \item \textbf{Outcome:} Students will be able to determine the relationships between two lines in three-dimensional space, such as intersection, parallelism, and perpendicularity.
    \item \textbf{Details:} This includes solving systems of equations to find points of intersection and applying the dot product to determine angles between direction vectors.
    \item \textbf{Graphing and Visualizing Lines in 3D Space:}
    \item \textbf{Outcome:} Students will be able to graph lines in three-dimensional space and interpret their geometric properties.
    \item \textbf{Details:} This includes plotting points and direction vectors, visualizing the spatial relationships of lines, and understanding the implications of the direction vectors on the orientation and position of lines in space.
\end{enumerate}

\subsection{Assessment Instrument}

\begin{enumerate}
    \item \textbf{Assessment Instrument: Equations of Lines, Planes, and Surfaces in Space Quiz} \textbf{Format:}
    \item \textbf{Type:} Mixed-format quiz (multiple-choice, short answer, and problem-solving questions)
    \item \textbf{Duration:} 45 minutes
    \item \textbf{Total Marks:} 30 points \textbf{Sections:}
    \item \textbf{Equations of Lines in 3D Space (10 points):}
    \item \textbf{Question 1:} (Multiple Choice) Which of the following is the vector equation of a line that passes through the point $(1, 2, 3)$ and is parallel to the vector $\langle 4, -1, 2 \rangle$?
    \item a. $\mathbf{r} = \langle 1, 2, 3 \rangle + t \langle 4, -1, 2 \rangle$
    \item b. $\mathbf{r} = \langle 4, -1, 2 \rangle + t \langle 1, 2, 3 \rangle$
    \item c. $\mathbf{r} = \langle 4, -1, 2 \rangle + t \langle -1, 2, 3 \rangle$
    \item d. $\mathbf{r} = \langle 1, 2, 3 \rangle + t \langle -1, 2, 4 \rangle$
    \item \textbf{Question 2:} (Short Answer) Write the parametric equations of the line given by the vector equation $\mathbf{r} = \langle -2, 3, 5 \rangle + t \langle 1, -1, 2 \rangle$.
    \item \textbf{Question 3:} (Problem Solving) Determine if the lines given by $\mathbf{r} = \langle 1, 2, 3 \rangle + t \langle 4, 1, 0 \rangle$ and $\mathbf{r} = \langle 5, -1, 2 \rangle + s \langle -2, 1, 1 \rangle$ intersect. If they do, find the point of intersection.
    \item \textbf{Equations of Planes in 3D Space (10 points):}
    \item \textbf{Question 4:} (Multiple Choice) Which of the following is the scalar equation of a plane passing through the point $(2, 1, -3)$ and normal to the vector $\langle 1, 4, -2 \rangle$?
    \item a. $x + 4y - 2z = 3$
    \item b. $x + 4y - 2z = -9$
    \item c. $x + 4y - 2z = -3$
    \item d. $x + 4y - 2z = 9$
    \item \textbf{Question 5:} (Short Answer) Write the vector equation of the plane passing through the points $(1, 2, 3)$, $(2, 0, 1)$, and $(0, -1, 2)$.
    \item \textbf{Question 6:} (Problem Solving) Find the distance between the planes $2x - y + z = 1$ and $4x - 2y + 2z = 5$.
    \item \textbf{Surfaces in 3D Space (10 points):}
    \item \textbf{Question 7:} (Multiple Choice) Which of the following equations represents an ellipsoid?
    \item a. $x^2 + y^2 - z^2 = 1$
    \item b. $\frac{x^2}{4} + \frac{y^2}{9} + \frac{z^2}{16} = 1$
    \item c. $\frac{x^2}{4} - \frac{y^2}{9} + \frac{z^2}{16} = 1$
    \item d. $x^2 + y^2 + z = 1$
    \item \textbf{Question 8:} (Short Answer) Write the equation of a cylinder with a circular base of radius 3 centered at $(0, 0, 0)$ and oriented along the z-axis.
    \item \textbf{Question 9:} (Problem Solving) Determine the traces of the hyperboloid $\frac{x^2}{4} + \frac{y^2}{9} - \frac{z^2}{16} = 1$ in the planes $z = 0$, $x = 0$, and $y = 0$. \textbf{Scoring Rubric:}
    \item \textbf{Equations of Lines in 3D Space:} 3 points for multiple choice, 3 points for short answer, 4 points for problem solving (total 10 points)
    \item \textbf{Equations of Planes in 3D Space:} 3 points for multiple choice, 3 points for short answer, 4 points for problem solving (total 10 points)
    \item \textbf{Surfaces in 3D Space:} 3 points for multiple choice, 3 points for short answer, 4 points for problem solving (total 10 points) This quiz will help assess students' understanding and application of the equations of lines, planes, and surfaces in three-dimensional space. It ensures they can construct and interpret equations, solve geometric problems, and visualize these objects in space.
\end{enumerate}

%%===============================================================
\section{Cylinders and Quadric Surfaces (incomplete)}
\label{sec:cylinders-and-quadric-surfaces-incomplete}
%%===============================================================

Cylinders and Quadric Surfaces

%%===============================================================
\section{Curves in Space and Vector Functions}
\label{sec:curves-in-space-and-vector-functions}
%%===============================================================

\paragraph{Introduction}



If we have a function $\vec{f}:\R\to\R^n$ , we say that $\vec{f}$'s

image (the set $\Big\\{\vec{f}(t)\Big|t\in\R\Big\\}$ - or some subset of

$\R$) is a \textit{curve} in $\R^n$ and $\vec{f}$ is its parametrization.



Parameterizations are not necessarily unique - for example, $\vec{f}(t)=(\cos(t),\sin(t))$ such that $t\in[0,2\pi)$ is one parametrization of the unit circle, and $\vec{g}(t)=(\cos(at),\sin(at))$ such that $t\in\left[0,\frac{2\pi}{a}\right)$ is a whole family of parameterizations of that circle.



\paragraph{Collision and intersection points}



Say we have two different curves. It may be important to consider



\begin{itemize}
    \item points the two curves share - where they \textit{intersect}
    \item intersections which occur for the same value of $t$ - where they \textit{collide}.
\end{itemize}

\paragraph{Intersection points}



Firstly, we have two parameterizations $\vec{f}(t)$ and $\vec{g}(t)$ , and we want to find out when they intersect, this means that we want to know when the function values of each parametrization are the same. This means that we need to solve



$$\vec{f}(t)=\vec{g}(s)$$ because we're seeking the function values independent of the times they intersect.



For example, if we have $\vec{f}(t)=(t,3t)$ and $\vec{g}(t)=(t,t^2)$ , and we want to find intersection points:



$$\vec{f}(t)=\vec{g}(s)$$



$$(t,3t)=(s,s^2)$$



$$t=s\ ,\ 3t=s^2$$ with solutions $(t,s)=(0,0)\ ,\ (3,9)$



So, the two curves intersect at the points $(0,0)\ ,\ (3,9)$ .



\paragraph{Collision points}



However, if we want to know when the points "collide", with $\vec{f}(t)$ and $\vec{g}(t)$ , we need to know when both the function values \textit{and} the times are the same, so we need to solve instead



$$\vec{f}(t)=\vec{g}(t)$$



For example, using the same functions as before, $\vec{f}(t)=(t, 3t)$ and $\vec{g}(t)=(t,t^2)$ , and we want to find collision points:



$$\vec{f}(t)=\vec{g}(t)$$



$$(t,3t)=(t,t^2)$$



$$t=t\ ,\ 3t=t^2$$ which gives solutions $t=0,3$ So the collision points are $(0,0)\ ,\ (3,9)$ .



We may want to do this to actually model physical problems, such as in ballistics.



\paragraph{Intersection vector functions}



We can also use vector functions to represent the curve of intersection of two surfaces. For example, we want to know the curve of intersection of the cylinder $x^2+y^2=1$ and the plane $y+z=2$.



> Vector functions rely on parameterizations, so we can rewrite the equation of the cylinder into: $\cos^2 t+\sin^2 t=1$, where $x=\cos t,y=\sin t$.

>

> From the equation of the plane, we know that $z=2-y=2-\sin t$. Thus

> the corresponding vector equation is:

>

> $\mathbf{r}(t)=\langle \cos t,\sin t,2-\sin t\rangle$



\paragraph{Limits and Continuity}



The limit of a vector function $\mathbf{r}$ is defined by taking the limits of its component functions.



There is a vector function $\mathbf{r}(t)=\langle f(t),g(t),h(t) \rangle$. If $\lim _{t \rightarrow a}f(t),\lim _{t \rightarrow a}g(t),\lim _{t \rightarrow a}h(t)$ exist, then



$$\lim _{t \rightarrow a}\mathbf{r}(t)=\langle \lim _{t \rightarrow a}f(t),\lim _{t \rightarrow a}g(t),\lim _{t \rightarrow a}h(t) \rangle$$



And the requirement for continuity is also simple:



> A vector function $\mathbf{r}$ is continuous at $a$ if $\lim _{t\rightarrow a}\mathbf{r}(t)=\mathbf{r}(a)$.



\paragraph{Derivatives and Integrals}



\paragraph{Differentiation}



Recall that the first derivative of a scalar function $f(x)$ is defined as:



$$f'(x)=\lim _{\Delta x\rightarrow 0}\frac{f(x+\Delta x)-f(x)}{\Delta x}$$



The first derivative of a vector function $\mathbf{r}(t)$ is defined in much the same way:



$$\mathbf{r}'(t)=\lim _{\Delta t\rightarrow 0}\frac{\mathbf{r}(t+\Delta t)-\mathbf{r}(t)}{\Delta t}$$



We can use this definition to prove that the derivative of a vector function can be presented as the derivative of its component functions.



>$$\begin{aligned} \mathbf{r}'(t) & = \lim _{\Delta t\rightarrow 0}\frac{1}{\Delta t}[\mathbf{r}(t+\Delta t)-\mathbf{r}(t)] \quad \text{definition} \\\\ & = \lim _{\Delta t\rightarrow 0}\frac{1}{\Delta t}[\langle f(t+\Delta t),g(t+\Delta t),h(t+\Delta t)\rangle-\langle f(t),g(t),h(t) \rangle] \quad \text{component form} \\\\ & = \lim _{\Delta t\rightarrow 0}\langle \frac{f(t+\Delta t)-f(t)}{\Delta t},\frac{g(t+\Delta t)-g(t)}{\Delta t},\frac{h(t+\Delta t)-h(t)}{\Delta t}\rangle \quad\text{vector addition and multiplication} \\\\ & = \langle \lim _{\Delta t\rightarrow 0}\frac{f(t+\Delta t)-f(t)}{\Delta t},\lim _{\Delta t\rightarrow 0}\frac{g(t+\Delta t)-g(t)}{\Delta t},\lim _{\Delta t\rightarrow 0}\frac{h(t+\Delta t)-h(t)}{\Delta t}\rangle\quad\text{definition} \\\\ & = \langle f'(t),g'(t),h'(t)\rangle \quad\text{definition}\end{aligned}$$



Thus, using the same method, we can derive the second derivative and so on.



\paragraph{Derivatives of a vector function}



There is a vector function

$\mathbf{r}(t)=\langle f(t),g(t),h(t) \rangle$. The first derivative of

this vector function is:



$$\mathbf{r}'(t)=\lim _{\Delta t\rightarrow 0}\frac{\mathbf{r}(t+\Delta t)-\mathbf{r}(t)}{\Delta t}= \langle f'(t),g'(t),h'(t)\rangle$$}}So

the $n$th-order derivative should look like this:



> $\mathbf{r}^{(n)}(t)=\langle f^{(n)}(t),g^{(n)}(t),h^{(n)}(t)\rangle$



\paragraph{Differentiation rules}



Just like real-valued functions, there are some differentiation rules in the world of vector functions. The factor that makes vector differentiation rules slight more complicated is the product rule because there are two kinds of multiplication in vectors: dot product and cross product.



Suppose $\mathbf{u},\mathbf{v}$ are differentiable vector functions and $f$ is a real-valued function. Then



1.  $\frac{d}{dt}[\mathbf{u}(t)+\mathbf{v}(t)]=\mathbf{u}'(t)+\mathbf{v}'(t)\quad$ (addition)

2.  $\frac{d}{dt}[f(t)\mathbf{u}(t)]=f'(t)\mathbf{u}(t)+f(t)\mathbf{u}'(t)\quad$ (scalar multiplication)

3.  $\frac{d}{dt}[\mathbf{u}(t)\cdot\mathbf{v}(t)]=\mathbf{u}'(t)\cdot\mathbf{v}(t)+\mathbf{u}(t)\cdot\mathbf{v}'(t)\quad$ (dot product)

4.  $\frac{d}{dt}[\mathbf{u}(t)\times\mathbf{v}(t)]=\mathbf{u}'(t)\times\mathbf{v}(t)+\mathbf{u}(t)\times\mathbf{v}'(t)\quad$ (cross product)

5.  $\frac{d}{dt}[\mathbf{u}(f(t))]=f'(t)\mathbf{u}'(f(t))\quad$ (chain rule)



Naturally, we will prove that those rules are correct. Let us assume that $\mathbf{u}(t)=\langle f _1(t),f _2(t),f _3(t)\rangle$ and $\mathbf{v}(t)=\langle g _1(t),g _2(t),g _3(t)\rangle$.



Rule 1: the addition rule



> $$ \begin{aligned} \frac{d}{dt}[\mathbf{u}(t)+\mathbf{v}(t)] & = \frac{d}{dt}[\langle f _1(t), f _2(t), f _3(t) \rangle + \langle g _1(t), g _2(t), g _3(t) \rangle] \quad \text{component form} \\\\ & = \frac{d}{dt} \langle f _1(t) + g _1(t), f _2(t) + g _2(t), f _3(t) + g _3(t) \rangle \quad \text{vector addition} \\\\ & = \langle \frac{d}{dt}(f _1(t) + g _1(t)), \frac{d}{dt}(f _2(t) + g _2(t)), \frac{d}{dt}(f _3(t) + g _3(t)) \rangle \quad \text{distribution} \\\\ & = \langle f _1'(t) + g _1'(t), f _2'(t) + g _2'(t), f _3'(t) + g _3'(t) \rangle \quad \text{addition rule for real-valued functions} \\\\ & = \langle f _1'(t), f _2'(t), f _3'(t) \rangle + \langle g _1'(t), g _2'(t), g _3'(t) \rangle \quad \text{vector addition} \\\\ & = \mathbf{u}'(t) + \mathbf{v}'(t) \quad \text{vector form} \end{aligned} $$



Rule 2: the scalar multiplication rule



> $$ \begin{aligned} \frac{d}{dt}[f(t)\mathbf{u}(t)] & = \frac{d}{dt}[f(t)\begin{pmatrix} f _1(t) \\\\ f _2(t) \\\\ f _3(t) \end{pmatrix}] \quad \text{component form} \\\\ & = \frac{d}{dt} \begin{pmatrix} f(t) f _1(t) \\\\ f(t) f _2(t) \\\\ f(t) f _3(t) \end{pmatrix} \quad \text{scalar multiplication} \\\\ & = \begin{pmatrix} \frac{d}{dt}(f(t) f _1(t)) \\\\ \frac{d}{dt}(f(t) f _2(t)) \\\\ \frac{d}{dt}(f(t) f _3(t)) \end{pmatrix} \quad \text{distribution} \\\\ & = \begin{pmatrix} f'(t) f _1(t) + f(t) f _1'(t) \\\\ f'(t) f _2(t) + f(t) f _2'(t) \\\\ f'(t) f _3(t) + f(t) f _3'(t) \end{pmatrix} \quad \text{multiplication rule for real-valued functions} \\\\ & = \begin{pmatrix} f'(t) f _1(t) \\\\ f'(t) f _2(t) \\\\ f'(t) f _3(t) \end{pmatrix} + \begin{pmatrix} f(t) f _1'(t) \\\\ f(t) f _2'(t) \\\\ f(t) f _3'(t) \end{pmatrix} \quad \text{vector addition} \\\\ & = f'(t) \begin{pmatrix} f _1(t) \\\\ f _2(t) \\\\ f _3(t) \end{pmatrix} + f(t) \begin{pmatrix} f _1'(t) \\\\ f _2'(t) \\\\ f _3'(t) \end{pmatrix} \quad \text{scalar multiplication} \\\\ & = f'(t)\mathbf{u}(t) + f(t)\mathbf{u}'(t) \quad \text{vector form} \end{aligned} $$



Rule 3: the dot product rule



> $$ \begin{aligned} \frac{d}{dt}[\mathbf{u}(t) \cdot \mathbf{v}(t)] & = \frac{d}{dt}[\langle f _1(t), f _2(t), f _3(t) \rangle \cdot \langle g _1(t), g _2(t), g _3(t) \rangle] \quad \text{component form} \\\\ & = \frac{d}{dt}(f _1(t) g _1(t) + f _2(t) g _2(t) + f _3(t) g _3(t)) \quad \text{dot product} \\\\ & = \frac{d}{dt} \sum _{i=1}^3 f _i(t) g _i(t) \quad \text{simplification} \\\\ & = \sum _{i=1}^3 \frac{d}{dt} [f _i(t) g _i(t)] \quad \text{distribution} \\\\ & = \sum _{i=1}^3 [f _i'(t) g _i(t) + f _i(t) g _i'(t)] \quad \text{multiplication rule for real-valued functions} \\\\ & = \sum _{i=1}^3 [f _i'(t) g _i(t) + f _i(t) g _i'(t)] \quad \text{vector addition} \\\\ & = \mathbf{u}'(t) \cdot \mathbf{v}(t) + \mathbf{u}(t) \cdot \mathbf{v}'(t) \quad \text{vector form} \end{aligned} $$



Rule 4: the cross product rule



> $$ \begin{aligned} \frac{d}{dt}[\mathbf{u}(t) \times \mathbf{v}(t)] & = \frac{d}{dt}[\langle f _1(t), f _2(t), f _3(t) \rangle \times \langle g _1(t), g _2(t), g _3(t) \rangle] \quad \text{component form} \\\\ & = \frac{d}{dt} \begin{vmatrix} \mathbf{\hat{i}} & \mathbf{\hat{j}} & \mathbf{\hat{k}} \\\\ f _1(t) & f _2(t) & f _3(t) \\\\ g _1(t) & g _2(t) & g _3(t) \\\\ \end{vmatrix} \quad \text{cross product} \\\\ & = \frac{d}{dt} \begin{pmatrix} f _2(t) g _3(t) - f _3(t) g _2(t) \\\\ f _3(t) g _1(t) - f _1(t) g _3(t) \\\\ f _1(t) g _2(t) - f _2(t) g _1(t) \\\\ \end{pmatrix} \quad \\\\ & = \begin{pmatrix} \frac{d}{dt}(f _2(t) g _3(t) - f _3(t) g _2(t)) \\\\ \frac{d}{dt}(f _3(t) g _1(t) - f _1(t) g _3(t)) \\\\ \frac{d}{dt}(f _1(t) g _2(t) - f _2(t) g _1(t)) \\\\ \end{pmatrix} \quad \text{distribution} \\\\ & = \begin{pmatrix} f _2'(t) g _3(t) + f _2(t) g _3'(t) - f _3'(t) g _2(t) - f _3(t) g _2'(t) \\\\ f _3'(t) g _1(t) + f _3(t) g _1'(t) - f _1'(t) g _3(t) - f _1(t) g _3'(t) \\\\ f _1'(t) g _2(t) + f _1(t) g _2'(t) - f _2'(t) g _1(t) - f _2(t) g _1'(t) \\\\ \end{pmatrix} \quad \text{multiplication rule for real-valued functions} \\\\ & = \begin{pmatrix} f _2'(t) g _3(t) - f _3'(t) g _2(t) + f _2(t) g _3'(t) - f _3(t) g _2'(t) \\\\ f _3'(t) g _1(t) - f _1'(t) g _3(t) + f _3(t) g _1'(t) - f _1(t) g _3'(t) \\\\ f _1'(t) g _2(t) - f _2'(t) g _1(t) + f _1(t) g _2'(t) - f _2(t) g _1'(t) \\\\ \end{pmatrix} \quad \text{rearrangement} \\\\ & = \begin{pmatrix} f _2'(t) g _3(t) - f _3'(t) g _2(t) \\\\ f _3'(t) g _1(t) - f _1'(t) g _3(t) \\\\ f _1'(t) g _2(t) - f _2'(t) g _1(t) \\\\ \end{pmatrix} + \begin{pmatrix} f _2(t) g _3'(t) - f _3(t) g _2'(t) \\\\ f _3(t) g _1'(t) - f _1(t) g _3'(t) \\\\ f _1(t) g _2'(t) - f _2(t) g _1'(t) \\\\ \end{pmatrix} \quad \text{vector addition} \\\\ & = \begin{vmatrix} \mathbf{\hat{i}} & \mathbf{\hat{j}} & \mathbf{\hat{k}} \\\\ f _1'(t) & f _2'(t) & f _3'(t) \\\\ g _1(t) & g _2(t) & g _3(t) \\\\ \end{vmatrix} + \begin{vmatrix} \mathbf{\hat{i}} & \mathbf{\hat{j}} & \mathbf{\hat{k}} \\\\ f _1(t) & f _2(t) & f _3(t) \\\\ g _1'(t) & g _2'(t) & g _3'(t) \\\\ \end{vmatrix} \quad \text{cross product} \\\\ & = \begin{pmatrix} f _1'(t) \\\\ f _2'(t) \\\\ f _3'(t) \\\\ \end{pmatrix} \times \begin{pmatrix} g _1(t) \\\\ g _2(t) \\\\ g _3(t) \\\\ \end{pmatrix} + \begin{pmatrix} f _1(t) \\\\ f _2(t) \\\\ f _3(t) \\\\ \end{pmatrix} \times \begin{pmatrix} g _1'(t) \\\\ g _2'(t) \\\\ g _3'(t) \\\\ \end{pmatrix} \quad \\\\ & = \mathbf{u}'(t) \times \mathbf{v}(t) + \mathbf{u}(t) \times \mathbf{v}'(t) \quad \text{vector form} \end{aligned} $$



Rule 5: the chain rule



> $$ \begin{aligned} \frac{d}{dt}[\mathbf{u}(f(t))] & = \frac{d}{dt} \begin{pmatrix} f _1(f(t)) \\\\ f _2(f(t)) \\\\ f _3(f(t)) \end{pmatrix} \quad \text{component form} \\\\ & = \begin{pmatrix} \frac{d}{dt} f _1(f(t)) \\\\ \frac{d}{dt} f _2(f(t)) \\\\ \frac{d}{dt} f _3(f(t)) \end{pmatrix} \quad \text{distribution} \\\\ & = \begin{pmatrix} f'(t) f _1'(f(t)) \\\\ f'(t) f _2'(f(t)) \\\\ f'(t) f _3'(f(t)) \end{pmatrix} \quad \text{chain rule for real-valued functions} \\\\ & = f'(t) \begin{pmatrix} f _1'(f(t)) \\\\ f _2'(f(t)) \\\\ f _3'(f(t)) \end{pmatrix} \quad \text{scalar multiplication} \\\\ & = f'(t) \mathbf{u}'(f(t)) \quad \text{vector form} \end{aligned} $$



$\blacksquare$



\paragraph{Angle between curves}



% [Image: ]



Figure: Inner product angle.



We can then formulate the concept of the \textit{angle} between two curves by considering the angle between the two tangent vectors. If two curves, parametrized by $\vec{f _1}$ and $\vec{f _2}$ intersect at some point, which means that



$$\vec{f _1}(s)=\vec{f _2}(t)=c$$ the angle between these two curves at $c$ is the angle between the tangent vectors $\vec{f _1'}(s)$ and $\vec{f _2'}(t)$, using the dot product, is given by



$$\arccos\left(\frac{\vec{f _1'}(s)\cdot\vec{f _2'}(t)}{\Big\|\vec{f _1'}(s)\Big\|\Big\|\vec{f _2'}(t)\Big\|}\right)$$



\paragraph{Integration}



Similar to real-valued functions, the definite integral of a vector function $\mathbf{r}(t)$ is defined as:



> $$ \begin{aligned} \int _a^b \mathbf{r}(t) \, dt & = \lim _{n \rightarrow \infty} \sum _{i=1}^n \mathbf{r}(t _i^\*) \Delta t \\\\ & = \lim _{n \rightarrow \infty} \left[ \left( \sum _{i=1}^n f(t _i^\*) \Delta t \right) \mathbf{i} + \left( \sum _{i=1}^n g(t _i^\*) \Delta t \right) \mathbf{j} + \left( \sum _{i=1}^n h(t _i^\*) \Delta t \right) \mathbf{k} \right] \\\\ & = \left( \int _a^b f(t) \, dt \right) \mathbf{i} + \left( \int _a^b g(t) \, dt \right) \mathbf{j} + \left( \int _a^b h(t) \, dt \right) \mathbf{k} \end{aligned} $$



We can extend the Fundamental Theorem of Calculus to continuous vector functions as follows:



$$\int _a^b\mathbf{r}(t)dt=\mathbf{R}(t)\biggr| _a^b=\mathbf{R}(b)-\mathbf{R}(a)$$

For indefinite integrals, the definition is:



> $\int\mathbf{r}(t)dt=\biggl(\int f(t)dt\biggr)\mathbf{i}+\biggl(\int g(t)dt\biggr)\mathbf{j}+\biggl(\int h(t)dt\biggr)\mathbf{k}+\mathbf{C}$



\paragraph{Arc length}



We can deduce the length of a curve with parametric equations $\begin{cases} x=f(t) \\ y=g(t) \end{cases}$, $a\le t\le b$ to be:



> $L=\int _a^b\sqrt{\biggl(\frac{dx}{dt}\biggr)^2+\biggl(\frac{dy}{dt}\biggr)^2}dt$



Since vector functions are fundamentally parametric equations with directions, we can utilize the formula above into the length of a space curve.



If the curve has the vector equation $\mathbf{r}(t)=\langle f(t),g(t),h(t)\rangle,a\le t\le b$, or, equivalently, the parametric equations $x=f(t),y=g(t),z=h(t)$, where $f',g',h'$ are continuous, then the length of the curve from $t=a$ to $t=b$ is:



$$L=\int _a^b\sqrt{[f'(t)]^2+[g'(t)]^2+[h'(t)]^2}dt=\int _a^b\sqrt{\biggl(\frac{dx}{dt}\biggr)^2+\biggl(\frac{dy}{dt}\biggr)^2+\biggl(\frac{dx}{dz}\biggr)^2}dt$$

For those who prefer simplicity, the formula can be rewritten into:



$$L=\int _a^b|\mathbf{r}'(t)|dt\quad \text{ or } \quad\frac{dL}{dt}=|\mathbf{r}'(t)|$$



\paragraph{Reparametrization}



Suppose that there is a vector function which describes the displacement

of a particle with respect to time and has the equation:



> $\mathbf{r}(t)=\cos t\ \mathbf{i}+\sin t\ \mathbf{j}+t\ \mathbf{k}$



However, for some reason, we do not want to know the displacement of this particle with respect to time. Instead, we want to know its displacement with respect to its traveled distance ($s$) from $(1,0,0)$ in the direction of increasing $t$. In order to do so, we need to find a way to describe time as a function of distance. In other words, we need to find $t=t(s)$. We can use the formula for the arc length to establish the relationship between time and distance because arc length in this case describes the distance traveled by the particle.



Before we start calculating, we need to introduce the arc length function.



Suppose that a piecewise-smooth curve with a vector function $\mathbf{r}(t)=\langle f(t),g(t),h(t)\rangle,a \le t \le b$, and the curve is traversed exactly once as $t$ increases from $a$ to $b$. The arc length function $s(t)$ is:



$$s(t)=\int _a^t\sqrt{\biggl(\frac{dx}{du}\biggr)^2+\biggl(\frac{dy}{du}\biggr)^2+\biggl(\frac{dz}{du}\biggr)^2}du$$



According to the definition, the arc length function for our curve $\mathbf{r}(t)$ should be



> $s(t)=\int _0^t|\mathbf{r}'(u)|du=\int _0^t\sqrt{(-\sin u)^2+(\cos u)^2+1^2}du=\sqrt{2}u\bigg| _0^t=\sqrt{2}t$

>

> Note that $a=0$ because the initial point $(1,0,0)$ corresponds to the parameter value $t=0$. Since it is in the direction of increasing $t$, the integration direction should be from $0$ to $t$.

>

> $t(s)=\frac{s}{\sqrt{2}}$



Then we substitute the original function and get the answer:



> $\mathbf{r}(s)=\cos \frac{s}{\sqrt{2}}\ \mathbf{i}+\sin \frac{s}{\sqrt{2}}\ \mathbf{j}+\frac{s}{\sqrt{2}}\ \mathbf{k}$



Reparametrization has important applications in real life because sometimes we want to know a value with respect to different variables. In this case, instead of describing the path of a particle with respect to time, we described its path with respect to its distance, which will be very useful in certain situations.



\paragraph{Curvature}



\paragraph{Terminology}



Before we start discussing curvature, there are some important vectors and concepts we need to be at least aware of.



\paragraph{The unit tangent vector}



In the differentiation section of this chapter, we discussed the derivatives of a vector function. We know that $\mathbf{v}(t)=\mathbf{r}'(t)$ at $t=t _0$ is tangent to the curve $\mathbf{r}(t)$ at $t=t _0$. $\mathbf{r}'(t)$ is called the tangent vector. The unit tangent vector, however, eliminates the aspect of magnitude because it is defined as:



$$\mathbf{T}(t)=\frac{\mathbf{r}'(t)}{|\mathbf{r}'(t)|}$$



As we can see, the magnitude of the unit tangent vector is always $1$. We can imagine that $\mathbf{r}(t)$ as the displacement of a particle with respect to time. So, the unit tangent vector can be perceived as the direction of the velocity of the particle with respect to time. It can also be perceived as the direction of the tangential acceleration of the particle with respect to time. We will discuss motion in space in the next section, but this is a useful method to intuitively understand some vectors.



\paragraph{The unit normal vector}



The unit normal vector is defined as



$$\mathbf{N}(t)=\frac{\mathbf{T}'(t)}{|\mathbf{T}'(t)|}$$



The unit normal is orthogonal to the unit tangent because since $|\mathbf{T}(t)|=1$, we can get that:



> $\frac{d}{dt}|\mathbf{T}(t)|^2=0=\frac{d}{dt}[\mathbf{T}(t)\cdot\mathbf{T}(t)]=2\mathbf{T}'(t)\cdot\mathbf{T}(t)$

>

> $\Leftrightarrow \mathbf{T}'(t)\cdot\mathbf{T}(t)=0$



This means that $\mathbf{T}'(t)$ is orthogonal to $\mathbf{T}(t)$. Therefore, $\mathbf{N}(t)$ is orthogonal to $\mathbf{T}(t)$. We can imagine that the unit normal vector is the direction of the normal acceleration of the particle with respect to time.



\paragraph{The binormal vector}



The binormal vector is defined as



> $\mathbf{B}(t)=\mathbf{T}(t)\times\mathbf{N}(t)$



The binormal vector is perpendicular to both the unit tangent and the unit normal because of the properties of the cross product. The magnitude of the binormal is always 1 because



> $|\mathbf{B}(t)|=|\mathbf{T}(t)||\mathbf{N}(t)|\sin\theta=1$



\paragraph{The normal plane, osculating plane, and the osculating circle}



\begin{itemize}
    \item The normal plane is the plane determined by the normal and binormal vectors $\mathbf{N}\text{ and }\mathbf{B}$. The normal plane consists of all lines that are orthogonal to the tangent vector $\mathbf{T}$.
    \item The osculating plane is the plane determined by the unit tangent and unit normal $\mathbf{T}\text{ and }\mathbf{N}$. It is the plane that comes closest to containing the part of the curve near a point where $t=t_0$.
    \item The osculating circle is the circle that lies in the osculating plane towards the direction of $\mathbf{N}$ with a radius $r=\frac{1}{\kappa}$ (the inverse of the curvature, which we will immediately discuss after this). It best describes how the curve behaves near the point where $t=t_0$ because it shares the same tangent, normal, and curvature at that point.
\end{itemize}

These concepts are very important in the branch of differential geometry

and in its applications to the motion of spacecraft.



\paragraph{Curvature}



The curvature of a curve at a given point is a measure of how quickly the curve changes direction at that point. We define it to be the magnitude of the rate of change of the unit tangent with respect to arc length. We use arc length so that the curvature will be independent of parametrization.



> Suppose that a space curve has vector function $\mathbf{r}$, unit

> tangent vector $\mathbf{T}$, and arc length $s$. The curvature of this

> curve is: $\kappa=\biggl|\frac{d\mathbf{T}(s)}{ds}\biggr|$.



There are two other ways to express the curvature. We can express curvature in terms of $t$ instead of $s$ by utilizing the chain rule (recall that $\frac{dL}{dt}=|\mathbf{r}'(t)|$):



$$\kappa=\biggl|\frac{d\mathbf{T}}{ds}\biggr|=\Biggl|\frac{\frac{d\mathbf{T}}{ds}\frac{ds}{dt}}{\frac{ds}{dt}}\Biggr|=\Biggl|\frac{\frac{d\mathbf{T}(t)}{dt}}{\frac{ds}{dt}}\Biggr|=\frac{|\mathbf{T}'(t)|}{|\mathbf{r}'(t)|}$$



The third way is more complicated to deduce, but it is often more convenient to apply because it only requires $\mathbf{r}(t)$ and its derivatives.



$$\kappa(t)=\frac{|\mathbf{r}'(t)\times\mathbf{r}''(t)|}{|\mathbf{r}'(t)|^3}$$



And now for the proof for this theorem:



> According to the definition of the unit tangent vector, we know that

> $\mathbf{r}'=|\mathbf{r}'|\mathbf{T}$. So the second derivative of

> $\mathbf{r}$ should be:

>

> $ \begin{aligned} \mathbf{r}'' & = (|\mathbf{r}'| \ \mathbf{T})' \\\\ & = |\mathbf{r}''| \ \mathbf{T} + |\mathbf{r}'| \ \mathbf{T}' \quad \text{the product rule} \end{aligned} $

>

> Now we calculate $\mathbf{r}'\times\mathbf{r}''$.

>

> $ \begin{aligned} \mathbf{r}' \times \mathbf{r}'' & = \left( |\mathbf{r}'| \mathbf{T} \right) \times \left( |\mathbf{r}''| \mathbf{T} + |\mathbf{r}'| \mathbf{T}' \right) \quad \text{substitution} \\\\ & = \left( |\mathbf{r}'| \mathbf{T} \times |\mathbf{r}''| \mathbf{T} \right) + \left( |\mathbf{r}'| \mathbf{T} \times |\mathbf{r}'| \mathbf{T}' \right) \quad \text{distribution} \\\\ & = |\mathbf{r}'||\mathbf{r}''| (\mathbf{T} \times \mathbf{T}) + |\mathbf{r}'|^2 (\mathbf{T} \times \mathbf{T}') \quad \text{rearrangement} \\\\ & = |\mathbf{r}'|^2 (\mathbf{T} \times \mathbf{T}') \quad \text{realizing that } \mathbf{T} \times \mathbf{T} = 0 \end{aligned} $

>

> And then we calculate $|\mathbf{r}'\times\mathbf{r}''|$.

>

> $ \begin{aligned} |\mathbf{r}' \times \mathbf{r}''| & = |\mathbf{r}'|^2 |\mathbf{T} \times \mathbf{T}'| \quad \text{substitution} \\\\ & = |\mathbf{r}'|^2 |\mathbf{T}| |\mathbf{T}'| \sin \theta \quad \text{the magnitude for the cross product} \\\\ & = |\mathbf{r}'|^2 |\mathbf{T}'| \quad \text{when you realize } |\mathbf{T}| = 1 \text{ and } \mathbf{T} \perp \mathbf{T}' \end{aligned} $

>

> We rearrange the equation into:

>

> $|\mathbf{T}'|=\frac{|\mathbf{r}'\times\mathbf{r}''|}{|\mathbf{r}'|^2}$

>

> Since $\kappa=\frac{|\mathbf{T}'|}{|\mathbf{r}'|}$, we can substitute

> $|\mathbf{T}'|$ with

> $\frac{|\mathbf{r}'\times\mathbf{r}''|}{|\mathbf{r}'|^2}$ and get:

> $\kappa=\frac{|\mathbf{r}'\times\mathbf{r}''|}{|\mathbf{r}'|^3}$



Here is a little summary on ways to calculate the curvature.





$$\begin{array}{|c|c|c|} \hline \textbf{Definition} & \textbf{With parametrization with respect to } \mathbf{t} & \textbf{In terms of } \mathbf{r}(t) \textbf{ and its derivatives} \\\\ \hline  \kappa=\left|\frac{d\mathbf{T}}{ds}\right| & \kappa=\frac{|\mathbf{T}'|}{|\mathbf{r}'|} & \kappa= \frac{|\mathbf{r}'\times\mathbf{r}'' |}{| \mathbf{r}' |^3} \\\\ \hline \end{array} $$



\paragraph{Motion in space}



\paragraph{Velocity and acceleration}



Remember in 2-dimensional calculus, we mentioned that a particle with displacement function $f(x)$ has velocity $f'(x)$ and acceleration $f''(x)$. In vector functions, the definition is basically the same. Suppose a particle moves through space so that its position vector at time $t$ is $\mathbf{r}(t)$, its velocity function and acceleration function are:



$$\mathbf{v}(t)=\lim _{h\rightarrow 0}\frac{\mathbf{r}(t+h)-\mathbf{r}(t)}{h}=\mathbf{r}'(t)\quad \text{ and }\quad\mathbf{a}(t)=\lim _{h\rightarrow 0}\frac{\mathbf{v}(t+h)-\mathbf{v}(t)}{h}=\mathbf{v}'(t)=\mathbf{r}''(t)$$



Simply put: $\mathbf{r}''(t)=\mathbf{v}'(t)=\mathbf{a}(t)$.



The speed of the particle ignores the direction. It is the magnitude of the velocity vector: $|\mathbf{v}(t)|$. The distance traveled by the particle from $t=a \text{ to } b$ is $\int _a^b|\mathbf{v}(t)|dt$, which is also the formula for the arc length.



With the help of the Fundamental Theorem of Calculus, we can deduce the velocity function and position function when we know the particle's acceleration.



$$\mathbf{v}(t)=\mathbf{v}(t _0)+\int _{t _0}^t\mathbf{a}(u)du\quad  \text{ and } \quad\mathbf{r}(t)=\mathbf{r}(t _0)+\int _{t _0}^t\mathbf{v}(u)du$$



\paragraph{Tangential and normal acceleration}



We can split the acceleration vector into two components: the tangential acceleration $a _T$ and the normal acceleration $a _N$. The tangential acceleration faces the same direction as the unit tangent vector ($\mathbf{T}$), and the normal acceleration faces the same direction as the unit normal vector ($\mathbf{N}$). Since both $\mathbf{T}$ and $\mathbf{N}$ are unit vectors, the acceleration vector can be written as the sum of two vectors:



> $\mathbf{a}=a _T\mathbf{T}+a _N\mathbf{N}$



Our goal is to figure out how to describe the two components.



> Recall that $\mathbf{T}=\frac{\mathbf{v}}{|\mathbf{v}|}$, thus

>

> > $\mathbf{v}=|\mathbf{v}|\mathbf{T}$

>

> Now we differentiate both sides of the equation,

>

> > $ \begin{aligned} \mathbf{v}' & = \left( |\mathbf{v}| \mathbf{T} \right)' \\\\ & = |\mathbf{v}|' \mathbf{T} + |\mathbf{v}| \mathbf{T}'  \end{aligned} $

>

> So we get

> $\mathbf{a}=|\mathbf{v}|'\mathbf{T}+|\mathbf{v}|\mathbf{T}'$.

>

> > Recall that

> > $\kappa=\frac{\big|\mathbf{T}'\big|}{\big|\mathbf{r}'\big|}=\frac{\big|\mathbf{T}'\big|}{\big|\mathbf{v}\big|}$,

> > thus $\big|\mathbf{T}'\big|=\kappa|\mathbf{v}|$.

>

> > Recall that $\mathbf{N}=\frac{\mathbf{T}'}{\big|\mathbf{T}'\big|}$,

> > thus

> > $\mathbf{T}'=\big|\mathbf{T}'\big|\mathbf{N}=\kappa|\mathbf{v}|\mathbf{N}$.

>

> We substitute $\mathbf{T}'$ for $\kappa|\mathbf{v}|\mathbf{N}$ to

> yield:

> $\mathbf{a}=|\mathbf{v}|'\mathbf{T}+\kappa|\mathbf{v}|^2\mathbf{N}$



That leaves us with:



> $a _T=|\mathbf{v}|'\quad$ and $\quad a _N=\kappa|\mathbf{v}|^2$



$\blacksquare$



Of course, it will be more convenient if those components can be written in terms if $\mathbf{r}(t)$ and its derivatives. Suppose that $\theta$ is the angle between $\mathbf{T}$ and $\mathbf{a}$, then we can write $a _T,a _N$ like this:



$$ \begin{array}{cc} \begin{aligned} a _T & = |\mathbf{a}(t)|\cos\theta \\\\ & = \frac{|\mathbf{v}(t)||\mathbf{a}(t)|\cos\theta}{|\mathbf{v}(t)|} \quad\text{algebraic manipulation} \\\\ & = \frac{\mathbf{v}(t)\cdot\mathbf{a}(t)}{|\mathbf{v}(t)|} \quad\text{dot product} \\\\ & = \frac{\mathbf{r}'(t)\cdot\mathbf{r}''(t)}{|\mathbf{r}'(t)|} \quad\text{in terms of }\mathbf{r}(t)\text{ and its derivatives} \end{aligned}  & \quad \begin{aligned} a _N & = |\mathbf{a}(t)|\sin\theta \\\\ & = \frac{|\mathbf{v}(t)||\mathbf{a}(t)|\sin\theta}{|\mathbf{v}(t)|} \quad\text{algebraic manipulation} \\\\ & = \frac{|\mathbf{v}(t)\times\mathbf{a}(t)|}{|\mathbf{v}(t)|} \quad\text{cross product} \\\\ & = \frac{|\mathbf{r}'(t)\times\mathbf{r}''(t)|}{|\mathbf{r}'(t)|} \quad\text{in terms of }\mathbf{r}(t)\text{ and its derivatives} \end{aligned} \end{array} $$



$\blacksquare$



To sum it up:



$$a _T=|\mathbf{v}|'=\frac{\mathbf{r}'(t)\cdot\mathbf{r}''(t)}{|\mathbf{r}'(t)|}\quad \text{ and } \quad a _N=\kappa|\mathbf{v}|^2= \frac{|\mathbf{r}'(t)\times\mathbf{r}''(t)|}{|\mathbf{r}'(t)|}$$



\paragraph{Expansion}



We only discussed vector functions with three variables $\mathbf{r}(t)=\langle f(t),g(t),h(t)\rangle$. How about expanding our understanding of vector functions into $n$ variables? Suppose we have a curve with vector function:



> $\mathbf{r}(t)= \begin{pmatrix} f _1(t) \\\\ f _2(t) \\\\ \vdots \\\\ f _n(t) \\\\ \end{pmatrix}$



\paragraph{Limits}



The limit of a vector function is defined as:



> $\lim _{t\rightarrow c} \mathbf{r}(t)= \begin{pmatrix} \lim _{t\rightarrow c} f _1(t) \\\\ \lim _{t\rightarrow c} f _2(t) \\\\ \vdots \\\\ \lim _{t\rightarrow c} f _n(t) \\\\ \end{pmatrix}$



\paragraph{Differentiation and integration}



The derivative of a vector function is defined as:



$$\mathbf{r}'(t)=\lim _{\Delta t\rightarrow 0}\frac{\mathbf{r}(t+\Delta t)-\mathbf{r}(t)}{\Delta t}, text{, thus } \quad\mathbf{r}'(t)= \begin{pmatrix} f _1'(t) \\\\ f _2'(t) \\\\ \vdots \\ f _n'(t) \\\\ \end{pmatrix}$$



All differentiation rules apply.



We can expand integration into:



$$\int\mathbf{r}(t)dt= \begin{pmatrix} \int f _1(t)dt \\\\ \int f _2(t)dt \\\\ \vdots \\\\ \int f _n(t)dt \\\\ \end{pmatrix} +\mathbf{C}$$



Then, the arc length will become:



$$L=\int _a^b|\mathbf{r}'(t)|dt\quad \text{ or } \quad L=\int _a^b \sqrt{[f _1'(t)]^2+[f _2'(t)]^2+\cdots+[f _n'(t)]^2}dt$$



\paragraph{Resources}



\begin{itemize}
    \item \href{https://math.libretexts.org/Bookshelves/Calculus/Book\%3A_Calculus_(OpenStax}{Vector-Valued Functions and Space Curves}/13\%3A_Vector-Valued_Functions/13.1\%3A_Vector-Valued_Functions_and_Space_Curves), Mathematics LibreTexts
\end{itemize}

\paragraph{Licensing}



Content obtained and/or adapted from:



\begin{itemize}
    \item \href{https://en.wikibooks.org/wiki/Calculus/Vector_Functions}{Vector Functions, WikiBooks: Calculus} under a CC BY-SA license
\end{itemize}

\subsection{Learning Objectives}

\begin{enumerate}
    \item Student Learning Objectives for Vector Functions and Parametrization
    \item \textbf{Understand the Concept of Parametrization:}
    \item \textbf{Outcome:} Students will be able to define and understand the concept of parametrization of curves in $\R^n$.
    \item \textbf{Details:} This includes recognizing that a vector function $\vec{f}(t) = (f_1(t), f_2(t), \ldots, f_n(t))$ can describe a curve and identifying different parametrizations of the same curve.
    \item \textbf{Identify Intersection and Collision Points:}
    \item \textbf{Outcome:} Students will be able to determine intersection and collision points of two parameterized curves.
    \item \textbf{Details:} This includes solving systems of equations $\vec{f}(t) = \vec{g}(s)$ for intersections and $\vec{f}(t) = \vec{g}(t)$ for collisions, with examples such as $\vec{f}(t) = (t, 3t)$ and $\vec{g}(t) = (t, t^2)$.
    \item \textbf{Differentiate and Integrate Vector Functions:}
    \item \textbf{Outcome:} Students will be able to differentiate and integrate vector functions.
    \item \textbf{Details:} This includes applying the definitions of the first derivative $\mathbf{r}'(t) = \lim_{\Delta t \to 0} \frac{\mathbf{r}(t + \Delta t) - \mathbf{r}(t)}{\Delta t}$ and definite integrals $\int_a^b \mathbf{r}(t) \, dt$, along with utilizing these concepts to solve problems involving vector functions. These objectives align with the detailed introduction and examples provided, ensuring students grasp the foundational concepts of parametrization, intersections, collisions, and calculus operations involving vector functions.
\end{enumerate}

\subsection{Assessment Instrument}

\begin{enumerate}
    \item \textbf{Assessment Instrument: Curves in Space and Vector-Valued Functions Quiz} \textbf{Format:}
    \item \textbf{Type:} Mixed-format quiz (multiple-choice, short answer, and problem-solving questions)
    \item \textbf{Duration:} 45 minutes
    \item \textbf{Total Marks:} 30 points \textbf{Sections:}
    \item \textbf{Concept of Parametrization (10 points):}
    \item \textbf{Question 1:} Which of the following is a correct parametrization of the unit circle?
    \item a. $\vec{f}(t) = (\cos(t), \sin(t))$ for $t \in [0, 2\pi)$
    \item b. $\vec{g}(t) = (t, t^2)$ for $t \in [0, 1]$
    \item c. $\vec{h}(t) = (e^t, e^{-t})$ for $t \in \mathbb{R}$
    \item d. $\vec{k}(t) = (t, \cos(t))$ for $t \in [0, 2\pi)$
    \item \textbf{Question 2:} Short answer: Explain why parametrizations of curves are not unique. Provide an example to support your explanation. (5 points)
    \item \textbf{Intersection and Collision Points (10 points):}
    \item \textbf{Question 3:} Determine whether the following parametrized curves intersect, and if so, find the intersection points:
    \item $\vec{f}(t) = (t, 3t)$
    \item $\vec{g}(s) = (s, s^2)$
    \item \textbf{Question 4:} Short answer: Explain the difference between intersection and collision points using the example from Question 3. (5 points)
    \item \textbf{Differentiation and Integration of Vector Functions (10 points):}
    \item \textbf{Question 5:} Find the derivative of the following vector function:
    \item $\mathbf{r}(t) = \langle t^2, \sin(t), e^t \rangle$
    \item \textbf{Question 6:} Evaluate the integral of the following vector function from $t = 0$ to $t = 1$:
    \item $\mathbf{r}(t) = \langle t, t^2, t^3 \rangle$ \textbf{Scoring Rubric:}
    \item \textbf{Concept of Parametrization:}
    \item Question 1: 5 points for the correct multiple-choice answer
    \item Question 2: 5 points for a correct and well-explained short answer
    \item \textbf{Intersection and Collision Points:}
    \item Question 3: 5 points for correctly identifying intersection points (partial credit for correct setup)
    \item Question 4: 5 points for a correct and clear explanation of the difference between intersection and collision points
    \item \textbf{Differentiation and Integration of Vector Functions:}
    \item Question 5: 5 points for correctly finding the derivative (partial credit for correct partial derivatives)
    \item Question 6: 5 points for correctly evaluating the integral (partial credit for correct integrals of individual components) The quiz will help assess the students' understanding of parametrization, their ability to identify intersection and collision points, and their proficiency in differentiating and integrating vector functions.
\end{enumerate}

%%===============================================================
\section{Integrals of Vector Functions (incomplete)}
\label{sec:integrals-of-vector-functions-incomplete}
%%===============================================================

Integrals of Vector Functions

%%===============================================================
\section{Arc Length}
\label{sec:arc-length}
%%===============================================================

We can deduce that the length of a curve with parametric equations

$\begin{cases} x=f(t) \\\\ y=g(t) \end{cases}$, $a\le t\le b$ should be:



$$L=\int _a^b\sqrt{\biggl(\frac{dx}{dt}\biggr)^2+\biggl(\frac{dy}{dt}\biggr)^2}dt$$



Since vector functions are fundamentally parametric equations with directions, we can utilize the formula above into the length of a space curve.



\paragraph{Arc length of a space curve}



If the curve has the vector equation $\mathbf{r}(t)=\langle f(t),g(t),h(t)\rangle,a\le t\le b$, or, equivalently, the parametric equations $x=f(t),y=g(t),z=h(t)$, where $f',g',h'$ are continuous, then the length of the curve from $t=a$ to $t=b$ is:



$$L=\int _a^b\sqrt{[f'(t)]^2+[g'(t)]^2+[h'(t)]^2}dt=\int _a^b\sqrt{\biggl(\frac{dx}{dt}\biggr)^2+\biggl(\frac{dy}{dt}\biggr)^2+\biggl(\frac{dx}{dz}\biggr)^2}dt$$}}

For those who prefer simplicity, the formula can be rewritten into:



$$L=\int _a^b|\mathbf{r}'(t)|dt\quad \text{ or } \quad\frac{dL}{dt}=|\mathbf{r}'(t)|$$



\paragraph{Example Problems}



1\. Find the circumference of the circle given by the parametric

equations $x(t)=R\cos(t),y(t)=R\sin(t)$ , with $t\in[0,2\pi]$.



$$ \begin{aligned} s &= \int\limits _0^{2\pi} \sqrt{\left(\tfrac{d}{dt}\big(R\cos(t)\big)\right)^2 + \left(\tfrac{d}{dt}\big(R\sin(t)\big)\right)^2} \, dt \hspace{15em} \\\\   &= \int\limits _0^{2\pi} \sqrt{\big(-R\sin(t)\big)^2 + \big(R\cos(t)\big)^2} \, dt \\\\   &= \int\limits _0^{2\pi} \sqrt{R^2\big(\sin^2(t) + \cos^2(t)\big)} \, dt \\\\   &= \int\limits _0^{2\pi} R \, dt \\\\   &= R \cdot t \Big| _0^{2\pi} \\\\   &= 2\pi R \end{aligned} $$



2\. Find the length of the curve $y=\frac{e^x+e^{-x}}{2}$ from $x=0$ to $x=1$.



>$$ \begin{aligned} s &= \int\limits _0^1 \sqrt{1 + \left(\frac{d}{dx}\left(\frac{e^{x} + e^{-x}}{2}\right)\right)^2} \, dx \hspace{17em} \\\\   &= \int\limits _0^1 \sqrt{1 + \left(\frac{e^{x} - e^{-x}}{2}\right)^2} \, dx \\\\   &= \int\limits _0^1 \sqrt{1 + \frac{e^{2x} - 2 + e^{-2x}}{4}} \, dx \\\\   &= \int\limits _0^1 \sqrt{\frac{e^{2x} + 2 + e^{-2x}}{4}} \, dx \\\\   &= \int\limits _0^1 \sqrt{\left(\frac{e^{x} + e^{-x}}{2}\right)^2} \, dx \\\\   &= \int\limits _0^1 \frac{e^{x} + e^{-x}}{2} \, dx \\\\   &= \frac{e^{x} - e^{-x}}{2} \bigg| _0^1 \\\\   &= \frac{e - \frac{1}{e}}{2} \end{aligned} $$



\paragraph{Resources}



\begin{itemize}
    \item \href{https://en.wikibooks.org/wiki/Calculus/Arc_length}{Arc Length}, WikiBooks: Calculus
    \item \href{https://openstax.org/books/calculus-volume-3/pages/3-3-arc-length-and-curvature}{Arc Length and Curvature}, OpenStax
\end{itemize}

\paragraph{Licensing}



Content obtained and/or adapted from:



\begin{itemize}
    \item \href{https://en.wikibooks.org/wiki/Calculus/Arc_length}{Arc Length, WikiBooks: Calculus} under a CC BY-SA license
\end{itemize}

\subsection{Learning Objectives}

\begin{enumerate}
    \item Student Learning Objectives (SLOs) for Arc Length of Curves
    \item \textbf{Calculate the Arc Length of a Parametric Curve:}
    \item \textbf{Outcome:} Students will be able to calculate the arc length of a curve given in parametric form.
    \item \textbf{Details:} This includes setting up and evaluating the integral $$L=\int_a^b \sqrt{\left(\frac{dx}{dt}\right)^2 + \left(\frac{dy}{dt}\right)^2} \, dt$$ for curves defined by the parametric equations $x=f(t)$ and $y=g(t)$ over the interval $[a,b]$. Students will practice finding the arc length for various parametric curves, ensuring they understand the derivation and application of the formula.
    \item \textbf{Determine the Arc Length of a Space Curve:}
    \item \textbf{Outcome:} Students will be able to determine the arc length of a space curve given in vector form.
    \item \textbf{Details:} This involves using the formula $$L=\int_a^b \sqrt{\left(f'(t)\right)^2 + \left(g'(t)\right)^2 + \left(h'(t)\right)^2} \, dt$$ for curves defined by the vector equation $\mathbf{r}(t)=\langle f(t), g(t), h(t) \rangle$. Students will learn to handle space curves and calculate their lengths over given intervals, emphasizing the transition from parametric to vector representations.
    \item \textbf{Solve Real-World Problems Involving Arc Length:}
    \item \textbf{Outcome:} Students will be able to apply their knowledge of arc length to solve real-world problems.
    \item \textbf{Details:} This includes interpreting and solving problems where calculating the length of a curve is necessary, such as determining the circumference of objects modeled by parametric equations or the length of trajectories in space. Students will develop skills to set up integrals for arc lengths in applied contexts and interpret the results in terms of the physical or geometric situations described.
\end{enumerate}

\subsection{Assessment Instrument}

\begin{enumerate}
    \item \textbf{Assessment Instrument: Arc Length Quiz} \textbf{Format:}
    \item \textbf{Type:} Mixed-format quiz (multiple-choice, short answer, and problem-solving questions)
    \item \textbf{Duration:} 45 minutes
    \item \textbf{Total Marks:} 30 points \textbf{Sections:}
    \item \textbf{Arc Length of Parametric Curves (10 points):}
    \item \textbf{Question 1 (Multiple Choice):} Given the parametric equations $x(t)=3\cos(t)$ and $y(t)=3\sin(t)$ for $t \in [0, \pi]$, what is the arc length of the curve?
    \item a. $3\pi$
    \item b. $6\pi$
    \item c. $9\pi$
    \item d. $12\pi$
    \item \textbf{Question 2 (Short Answer):} Find the arc length of the curve given by the parametric equations $x(t)=t^2$ and $y(t)=t^3$ for $t \in [0, 1]$.
    \item \textbf{Arc Length of Space Curves (10 points):}
    \item \textbf{Question 3 (Multiple Choice):} The vector equation of a space curve is given by $\mathbf{r}(t)=\langle t, t^2, t^3 \rangle$ for $t \in [0, 1]$. What is the arc length of this curve?
    \item a. $\frac{\sqrt{3}}{2}$
    \item b. $\frac{\sqrt{10}}{2}$
    \item c. $\sqrt{14}$
    \item d. $\sqrt{10}$
    \item \textbf{Question 4 (Short Answer):} Calculate the arc length of the space curve with the parametric equations $x(t)=\cos(t)$, $y(t)=\sin(t)$, and $z(t)=t$ for $t \in [0, 2\pi]$.
    \item \textbf{Applications of Arc Length (10 points):}
    \item \textbf{Question 5 (Problem Solving):} Find the circumference of the circle given by the parametric equations $x(t)=R\cos(t)$ and $y(t)=R\sin(t)$ for $t \in [0, 2\pi]$.
    \item \textbf{Question 6 (Problem Solving):} Find the length of the curve $y=\frac{e^x + e^{-x}}{2}$ from $x=0$ to $x=1$. \textbf{Scoring Rubric:}
    \item \textbf{Arc Length of Parametric Curves:}
    \item \textbf{Question 1:} 5 points for the correct answer.
    \item \textbf{Question 2:} 5 points for a correctly set up and solved integral.
    \item \textbf{Arc Length of Space Curves:}
    \item \textbf{Question 3:} 5 points for the correct answer.
    \item \textbf{Question 4:} 5 points for a correctly set up and solved integral.
    \item \textbf{Applications of Arc Length:}
    \item \textbf{Question 5:} 5 points for a correctly solved problem.
    \item \textbf{Question 6:} 5 points for a correctly solved problem. This quiz assesses students' understanding and application of calculating arc lengths for parametric and space curves, as well as solving related real-world problems. It ensures they can set up and evaluate the necessary integrals accurately.
\end{enumerate}

%%===============================================================
\section{Motion in Space}
\label{sec:motion-in-space}
%%===============================================================

\paragraph{Curvature}



\paragraph{Terminology}



Before we start discussing curvature, there are some important vectors and concepts we need to be at least aware of.



\paragraph{The unit tangent vector}



In the differentiation section of this chapter, we discussed the derivatives of a vector function. We know that $\mathbf{v}(t)=\mathbf{r}'(t)$ at $t=t _0$ is tangent to the curve $\mathbf{r}(t)$ at $t=t _0$. $\mathbf{r}'(t)$ is called the tangent vector. The unit tangent vector, however, eliminates the aspect of magnitude because it is defined as:



$$\mathbf{T}(t)=\frac{\mathbf{r}'(t)}{|\mathbf{r}'(t)|}$$



As we can see, the magnitude of the unit tangent vector is always $1$. We can imagine that $\mathbf{r}(t)$ as the displacement of a particle with respect to time. So, the unit tangent vector can be perceived as the direction of the velocity of the particle with respect to time. It can also be perceived as the direction of the tangential acceleration of the particle with respect to time. We will discuss motion in space in the next section, but this is a useful method to intuitively understand some vectors.



\paragraph{The unit normal vector}



The unit normal vector is defined as



$$\mathbf{N}(t)=\frac{\mathbf{T}'(t)}{|\mathbf{T}'(t)|}$$



The unit normal is orthogonal to the unit tangent because since $|\mathbf{T}(t)|=1$, we can get that:



$$\frac{d}{dt}|\mathbf{T}(t)|^2=0=\frac{d}{dt}[\mathbf{T}(t)\cdot\mathbf{T}(t)]=2\mathbf{T}'(t)\cdot\mathbf{T}(t)$$

$$\Leftrightarrow \mathbf{T}'(t)\cdot\mathbf{T}(t)=0$$



This means that $\mathbf{T}'(t)$ is orthogonal to $\mathbf{T}(t)$. Therefore, $\mathbf{N}(t)$ is orthogonal to $\mathbf{T}(t)$. We can imagine that the unit normal vector is the direction of the normal acceleration of the particle with respect to time.



\paragraph{The binormal vector}



The binormal vector is defined as



> $\mathbf{B}(t)=\mathbf{T}(t)\times\mathbf{N}(t)$



The binormal vector is perpendicular to both the unit tangent and the unit normal because of the properties of the cross product. The magnitude of the binormal is always 1 because



> $|\mathbf{B}(t)|=|\mathbf{T}(t)||\mathbf{N}(t)|\sin\theta=1$



\paragraph{The normal plane, osculating plane, and the osculating circle}



\begin{itemize}
    \item The normal plane is the plane determined by the normal and binormal vectors $\mathbf{N}$ and $\mathbf{B}$. The normal plane consists of all lines that are orthogonal to the tangent vector $\mathbf{T}$.
    \item The osculating plane is the plane determined by the unit tangent and unit normal $\mathbf{T}$ and $\mathbf{N}$. It is the plane that comes closest to containing the part of the curve near a point where $t=t _0$.
    \item The osculating circle is the circle that lies in the osculating plane towards the direction of $\mathbf{N}$ with a radius $r=\frac{1}{\kappa}$ (the inverse of the curvature, which we will immediately discuss after this). It best describes how the curve behaves near the point where $t=t _0$ because it shares the same tangent, normal, and curvature at that point.
\end{itemize}

These concepts are very important in the branch of differential geometry and in its applications to the motion of spacecraft.



\paragraph{Curvature}



The curvature of a curve at a given point is a measure of how quickly the curve changes direction at that point. We define it to be the magnitude of the rate of change of the unit tangent with respect to arc length. We use arc length so that the curvature will be independent of parametrization.



> Suppose that a space curve has vector function $\mathbf{r}$, unit

> tangent vector $\mathbf{T}$, and arc length $s$. The curvature of this

> curve is: $\kappa=\biggl|\frac{d\mathbf{T}(s)}{ds}\biggr|$.



There are two other ways to express the curvature. We can express curvature in terms of $t$ instead of $s$ by utilizing the chain rule (recall that $\frac{dL}{dt}=|\mathbf{r}'(t)|$):



$$\kappa=\biggl|\frac{d\mathbf{T}}{ds}\biggr|=\Biggl|\frac{\frac{d\mathbf{T}}{ds}\frac{ds}{dt}}{\frac{ds}{dt}}\Biggr|=\Biggl|\frac{\frac{d\mathbf{T}(t)}{dt}}{\frac{ds}{dt}}\Biggr|=\frac{|\mathbf{T}'(t)|}{|\mathbf{r}'(t)|}$$



The third way is more complicated to deduce, but it is often more

convenient to apply because it only requires $\mathbf{r}(t)$ and its

derivatives.



$$\kappa(t)=\frac{|\mathbf{r}'(t)\times\mathbf{r}''(t)|}{|\mathbf{r}'(t)|^3}$$



And now for the proof for this theorem:



> According to the definition of the unit tangent vector, we know that $\mathbf{r}'=|\mathbf{r}'|\mathbf{T}$. So the second derivative of $\mathbf{r}$ should be:

>

> > $ \begin{aligned} \mathbf{r}'' &= (|\mathbf{r}'|\ \mathbf{T})' \\\\ &= |\mathbf{r}''|\ \mathbf{T} + |\mathbf{r}'|\ \mathbf{T}' \quad\text{the product rule} \end{aligned} $

>

> Now we calculate $\mathbf{r}'\times\mathbf{r}''$.

>

> > $ \begin{aligned} \mathbf{r}' \times \mathbf{r}'' &= \bigl(|\mathbf{r}'| \mathbf{T}\bigr) \times \bigl(|\mathbf{r}''| \mathbf{T} + |\mathbf{r}'| \mathbf{T}'\bigr) \quad\text{substitution} \\\\ &= \bigl(|\mathbf{r}'| \mathbf{T} \times |\mathbf{r}''| \mathbf{T}\bigr) + \bigl(|\mathbf{r}'| \mathbf{T} \times |\mathbf{r}'| \mathbf{T}'\bigr) \quad\text{distribution} \\\\ &= |\mathbf{r}'| |\mathbf{r}''| (\mathbf{T} \times \mathbf{T}) + |\mathbf{r}'|^2 (\mathbf{T} \times \mathbf{T}') \quad\text{rearrangement} \\\\ &= |\mathbf{r}'|^2 (\mathbf{T} \times \mathbf{T}') \quad\text{realizing that } \mathbf{T} \times \mathbf{T} = 0 \end{aligned} $

>

> And then we calculate $|\mathbf{r}'\times\mathbf{r}''|$.

>

> > $ \begin{aligned} |\mathbf{r}' \times \mathbf{r}''| &= |\mathbf{r}'|^2 |\mathbf{T} \times \mathbf{T}'| \quad\text{substitution} \\\\ &= |\mathbf{r}'|^2 |\mathbf{T}| |\mathbf{T}'| \sin\theta \quad\text{the magnitude for the cross product} \\\\ &= |\mathbf{r}'|^2 |\mathbf{T}'| \quad\text{when you realize } |\mathbf{T}| = 1 \text{ and } \mathbf{T} \perp \mathbf{T}' \end{aligned} $

>

> We rearrange the equation into:

>

> > $|\mathbf{T}'|=\frac{|\mathbf{r}'\times\mathbf{r}''|}{|\mathbf{r}'|^2}$

>

> Since $\kappa=\frac{|\mathbf{T}'|}{|\mathbf{r}'|}$, we can substitute

> $|\mathbf{T}'|$ with

> $\frac{|\mathbf{r}'\times\mathbf{r}''|}{|\mathbf{r}'|^2}$ and get:

> $\kappa=\frac{|\mathbf{r}'\times\mathbf{r}''|}{|\mathbf{r}'|^3}$



Here is a little summary on ways to calculate the curvature.



$$\begin{array}{|c|c|c|} \hline

\textbf{Definition} & \textbf{With parametrization with respect to } \mathbf{t} & \textbf{In terms of } \mathbf{r}(t) \textbf{ and its derivatives} \\\\ \hline

\kappa=\biggl|\frac{d\mathbf{T}}{ds}\biggr| & \kappa=\frac{|\mathbf{T}'|}{|\mathbf{r}'|} & \kappa=\frac{|\mathbf{r}'\times\mathbf{r}''|}{|\mathbf{r}'|^3} \\\\ \hline

\end{array}

$$



\paragraph{Motion in space}



\paragraph{Velocity and acceleration}



Remember in 2-dimensional calculus, we mentioned that a particle with displacement function $f(x)$ has velocity $f'(x)$ and acceleration $f''(x)$. In vector functions, the definition is basically the same. Suppose a particle moves through space so that its position vector at time $t$ is $\mathbf{r}(t)$, its velocity function and acceleration function are:



$$\mathbf{v}(t)=\lim _{h\rightarrow 0}\frac{\mathbf{r}(t+h)-\mathbf{r}(t)}{h}=\mathbf{r}'(t)\quad \text{ and } \quad\mathbf{a}(t)=\lim _{h\rightarrow 0}\frac{\mathbf{v}(t+h)-\mathbf{v}(t)}{h}=\mathbf{v}'(t)=\mathbf{r}''(t)$$



Simply put: $\mathbf{r}''(t)=\mathbf{v}'(t)=\mathbf{a}(t)$.



The speed of the particle ignores the direction. It is the magnitude of the velocity vector: $|\mathbf{v}(t)|$. The distance traveled by the particle from $t=a \text{ to } b$ is $\int _a^b|\mathbf{v}(t)|dt$, which is also the formula for the arc length.



With the help of the Fundamental Theorem of Calculus, we can deduce the velocity function and position function when we know the particle's acceleration.



$$\mathbf{v}(t)=\mathbf{v}(t _0)+\int _{t _0}^t\mathbf{a}(u)du\quad \text{ and } \quad\mathbf{r}(t)=\mathbf{r}(t _0)+\int _{t _0}^t\mathbf{v}(u)du$$



\paragraph{Tangential and normal acceleration}



We can split the acceleration vector into two components: the tangential acceleration $a _T$ and the normal acceleration $a _N$. The tangential acceleration faces the same direction as the unit tangent vector ($\mathbf{T}$), and the normal acceleration faces the same direction as the unit normal vector ($\mathbf{N}$). Since both $\mathbf{T}$ and $\mathbf{N}$ are unit vectors, the acceleration vector can be written as the sum of two vectors:



> $\mathbf{a}=a _T\mathbf{T}+a _N\mathbf{N}$



Our goal is to figure out how to describe the two components.



> Recall that $\mathbf{T}=\frac{\mathbf{v}}{|\mathbf{v}|}$, thus

>

> > $\mathbf{v}=|\mathbf{v}|\mathbf{T}$

>

> Now we differentiate both sides of the equation,

>

> > $ \begin{aligned} \mathbf{v}' &= \bigl(|\mathbf{v}| \mathbf{T}\bigr)' \\\\ &= |\mathbf{v}|' \mathbf{T} + |\mathbf{v}| \mathbf{T}' \end{aligned} $

>

> So we get

> $\mathbf{a}=|\mathbf{v}|'\mathbf{T}+|\mathbf{v}|\mathbf{T}'$.

>

> > Recall that

> > $\kappa=\frac{\big|\mathbf{T}'\big|}{\big|\mathbf{r}'\big|}=\frac{\big|\mathbf{T}'\big|}{\big|\mathbf{v}\big|}$,

> > thus $\big|\mathbf{T}'\big|=\kappa|\mathbf{v}|$.

>

> > Recall that $\mathbf{N}=\frac{\mathbf{T}'}{\big|\mathbf{T}'\big|}$,

> > thus

> > $\mathbf{T}'=\big|\mathbf{T}'\big|\mathbf{N}=\kappa|\mathbf{v}|\mathbf{N}$.

>

> We substitute $\mathbf{T}'$ for $\kappa|\mathbf{v}|\mathbf{N}$ to

> yield:

> $\mathbf{a}=|\mathbf{v}|'\mathbf{T}+\kappa|\mathbf{v}|^2\mathbf{N}$



That leaves us with:



> $a _T=|\mathbf{v}|'\quad$ and $\quad a _N=\kappa|\mathbf{v}|^2$



$\blacksquare$



Of course, it will be more convenient if those components can be written in terms if $\mathbf{r}(t)$ and its derivatives. Suppose that $\theta$ is the angle between $\mathbf{T}$ and $\mathbf{a}$, then we can write $a _T,a _N$ like this:



$$ \begin{array}{cc} \begin{aligned} a _T & = |\mathbf{a}(t)|\cos\theta \\\\ & = \frac{|\mathbf{v}(t)||\mathbf{a}(t)|\cos\theta}{|\mathbf{v}(t)|} \quad\text{algebraic manipulation} \\\\ & = \frac{\mathbf{v}(t)\cdot\mathbf{a}(t)}{|\mathbf{v}(t)|} \quad\text{dot product} \\\\ & = \frac{\mathbf{r}'(t)\cdot\mathbf{r}''(t)}{|\mathbf{r}'(t)|} \quad\text{in terms of }\mathbf{r}(t)\text{ and its derivatives} \end{aligned}  & \quad \begin{aligned} a _N & = |\mathbf{a}(t)|\sin\theta \\\\ & = \frac{|\mathbf{v}(t)||\mathbf{a}(t)|\sin\theta}{|\mathbf{v}(t)|} \quad\text{algebraic manipulation} \\\\ & = \frac{|\mathbf{v}(t)\times\mathbf{a}(t)|}{|\mathbf{v}(t)|} \quad\text{cross product} \\\\ & = \frac{|\mathbf{r}'(t)\times\mathbf{r}''(t)|}{|\mathbf{r}'(t)|} \quad\text{in terms of }\mathbf{r}(t)\text{ and its derivatives} \end{aligned} \end{array} $$



$\blacksquare$



To sum it up:



$$a _T=|\mathbf{v}|'=\frac{\mathbf{r}'(t)\cdot\mathbf{r}''(t)}{|\mathbf{r}'(t)|}\quad \text{ and } \quad a _N=\kappa|\mathbf{v}|^2= \frac{|\mathbf{r}'(t)\times\mathbf{r}''(t)|}{|\mathbf{r}'(t)|}$$



\paragraph{Resources}



\begin{itemize}
    \item \href{https://en.wikibooks.org/wiki/Calculus/Vector_Functions}{Vector Functions}, Wikibooks: Calculus
    \item \href{https://openstax.org/books/calculus-volume-3/pages/3-4-motion-in-space}{Motion in Space}, OpenStax
\end{itemize}

\paragraph{Licensing}



Content obtained and/or adapted from:



\begin{itemize}
    \item \href{https://en.wikibooks.org/wiki/Calculus/Vector_Functions}{Vector Functions, Wikibooks: Calculus} under a CC BY-SA license
\end{itemize}

\subsection{Learning Objectives}

\begin{enumerate}
    \item Here are three Student Learning Objectives (SLOs) related to curvature, using the format provided:
    \item \textbf{Unit Tangent Vector:}
    \item \textbf{Outcome:} Students will be able to determine the unit tangent vector $\mathbf{T}(t)$ for a given vector function $\mathbf{r}(t)$.
    \item \textbf{Details:} This includes understanding the concept of the tangent vector, computing its magnitude, and normalizing the vector to find the unit tangent vector. Students will practice finding $\mathbf{T}(t)$ for different parametric curves.
    \item \textbf{Curvature Calculation:}
    \item \textbf{Outcome:} Students will be able to calculate the curvature $\kappa$ of a space curve using the formula $\kappa(t)=\frac{|\mathbf{r}'(t)\times\mathbf{r}''(t)|}{|\mathbf{r}'(t)|^3}$.
    \item \textbf{Details:} This includes understanding the geometric interpretation of curvature, the derivation of the curvature formula, and applying the formula to compute the curvature for specific curves. Students will solve problems involving different parametric representations of curves.
    \item \textbf{Tangential and Normal Acceleration:}
    \item \textbf{Outcome:} Students will be able to decompose the acceleration vector $\mathbf{a}(t)$ into its tangential and normal components $a_T$ and $a_N$.
    \item \textbf{Details:} This includes understanding the physical significance of tangential and normal acceleration, using the formulas $a_T=\frac{\mathbf{r}'(t)\cdot\mathbf{r}''(t)}{|\mathbf{r}'(t)|}$ and $a_N=\frac{|\mathbf{r}'(t)\times\mathbf{r}''(t)|}{|\mathbf{r}'(t)|}$, and applying these formulas to solve problems related to particle motion along a curve. Students will practice decomposing the acceleration vector for different parametric curves and interpreting the results.
\end{enumerate}

\subsection{Assessment Instrument}

\begin{enumerate}
    \item \textbf{Assessment Instrument: Motion in Space Quiz} \textbf{Format:}
    \item \textbf{Type:} Mixed-format quiz (multiple-choice, short answer, and problem-solving questions)
    \item \textbf{Duration:} 45 minutes
    \item \textbf{Total Marks:} 30 points \textbf{Sections:}
    \item \textbf{Unit Tangent Vector (10 points):}
    \item \textbf{Question 1 (Multiple Choice):} Given the vector function $\mathbf{r}(t) = \langle t^2, t^3, t^4 \rangle$, what is the unit tangent vector $\mathbf{T}(t)$ at $t=1$?
    \item a. $\langle 1, 2, 3 \rangle$
    \item b. $\langle \frac{1}{2}, \frac{3}{4}, \frac{1}{4} \rangle$
    \item c. $\langle \frac{1}{2\sqrt{29}}, \frac{3}{2\sqrt{29}}, \frac{4}{\sqrt{29}} \rangle$
    \item d. $\langle \frac{2}{3}, \frac{3}{2}, \frac{4}{3} \rangle$
    \item \textbf{Question 2 (Short Answer):} For the vector function $\mathbf{r}(t) = \langle \cos t, \sin t, t \rangle$, find the unit tangent vector $\mathbf{T}(t)$.
    \item \textbf{Curvature Calculation (10 points):}
    \item \textbf{Question 3 (Multiple Choice):} For the curve given by $\mathbf{r}(t) = \langle t, t^2, t^3 \rangle$, what is the curvature $\kappa$ at $t=1$?
    \item a. $\frac{1}{\sqrt{2}}$
    \item b. $\frac{\sqrt{14}}{9}$
    \item c. $\frac{2}{3}$
    \item d. $\frac{3}{\sqrt{7}}$
    \item \textbf{Question 4 (Problem-Solving):} Calculate the curvature $\kappa(t)$ for the vector function $\mathbf{r}(t) = \langle t, t^2, 0 \rangle$.
    \item \textbf{Tangential and Normal Acceleration (10 points):}
    \item \textbf{Question 5 (Multiple Choice):} If the position of a particle is given by $\mathbf{r}(t) = \langle t, t^2, t^3 \rangle$, what is the tangential acceleration $a_T$ at $t=1$?
    \item a. $2$
    \item b. $3$
    \item c. $4$
    \item d. $5$
    \item \textbf{Question 6 (Problem-Solving):} Given the vector function $\mathbf{r}(t) = \langle \cos t, \sin t, t \rangle$, find the tangential acceleration $a_T$ and the normal acceleration $a_N$ at $t=0$. \textbf{Scoring Rubric:}
    \item \textbf{Unit Tangent Vector:} 5 points for each correctly identified or calculated unit tangent vector (total 10 points)
    \item \textbf{Curvature Calculation:} 5 points for each correctly calculated curvature (total 10 points)
    \item \textbf{Tangential and Normal Acceleration:} 5 points for each correctly calculated tangential and normal acceleration (total 10 points) The quiz will assess the students' understanding and application of concepts related to motion in space, including the unit tangent vector, curvature, and the decomposition of acceleration into tangential and normal components. It ensures that they can perform the necessary calculations and interpret the physical significance of these quantities accurately.
\end{enumerate}

%%===============================================================
\section{Functions of Several Variables}
\label{sec:functions-of-several-variables}
%%===============================================================

This chapter serves as an introduction to multivariable calculus. Multivariable calculus is more complicated than when we were dealing with single-variable functions because more variables means more situations to be concerned about. In the following chapters, we will be discussing limits, differentiation, and integration of multivariable functions, using single-variable calculus as our basis.



\paragraph{Topology in $\mathbb{R} ^n$}



In your previous study of calculus, we have looked at functions and their behavior. Most of these functions we have examined have been all in the form



$$f(x):\mathbb{R}\to\mathbb{R}$$ and only occasional examination of functions of two variables. However, the study of functions of \textit{several} variables is quite rich in itself, and has applications in several fields.



We write functions of vectors - many variables - as follows:



$$f:\mathbb{R} ^m\to\mathbb{R} ^n$$ and $f(\vec{x})$ for the function that maps a vector in $\mathbb{R} ^m$ to a vector in $\mathbb{R} ^n$ .



Before we can do calculus in $\mathbb{R} ^n$ , we must familiarize ourselves with the structure of $\mathbb{R} ^n$ . We need to know which properties of $\mathbb{R}$ can be extended to $\mathbb{R} ^n$ . This page assumes at least some familiarity with basic linear algebra.



\paragraph{Lengths and distances}



If we have a vector in $\mathbb{R} ^2$ we can calculate its length using the Pythagorean theorem. For instance, the length of the vector $(2,3)$ is



$$\sqrt{3^2+2^2}=\sqrt{13}$$



We can generalize this to $\mathbb{R} ^n$ . We define a vector's length, written $\|\vec{x}\|$ , as the square root of the sum of the squares of each of its components. That is, if we have a vector $\vec{x}=(x _1,\ldots,x _n)$,



$$\|\vec{x}\|=\sqrt{x _1^2+\cdots+x _n^2}$$



Now that we have established some concept of length, we can establish the distance between two vectors. We define this distance to be the length of the two vectors' difference. We write this distance $d(\vec{x},\vec{y})$ , and it is



$$d(\vec{x},\vec{y})=\big\|\vec{x}-\vec{y}\big\|=\sqrt{\sum _{i=1}^n(x _i-y _i)^2}$$



This distance function is sometimes referred to as a \textit{metric}. Other metrics arise in different circumstances. The metric we have just defined is known as the \textit{Euclidean} metric.



\paragraph{Open and closed balls}



In $\mathbb{R}$ , we have the concept of an \textit{interval}, in that we choose a certain number of other points about some central point. For example, the interval $[-1,1]$ is centered about the point 0, and includes points to the left and right of 0.



In $\mathbb{R} ^2$ and up, the idea is a little more difficult to carry on. For $\mathbb{R} ^2$ , we need to consider points to the left, right, above, and below a certain point. This may be fine, but for $\mathbb{R} ^3$ we need to include points in more directions.



We generalize the idea of the interval by considering all the points that are a given, fixed distance from a certain point - now we know how to calculate distances in $\mathbb{R} ^n$ , we can make our generalization as follows, by introducing the concept of an \textit{open ball} and a \textit{closed ball} respectively, which are analogous to the open and closed interval respectively.



> An \textit{open ball} $B(\vec{a},r)$ is a set in the form $\left\\{ \vec{x} \in \mathbb{R} ^n \Big| d(\vec{x}, \vec{a}) < r \right\\}$

>

>

> A \textit{closed ball} $\overline{B}(\vec{a},r)$ is a set in the form $\left\\{\vec{x}\in \mathbb{R} ^n\Big| d(\vec{x},\vec{a}) \le r\right\\}$



In $\mathbb{R}$ , we have seen that the open ball is simply an open interval centered about the point $x=a$ . In $\mathbb{R} ^2$ this is a circle with no boundary, and in $\mathbb{R} ^3$ it is a sphere with no outer surface. (\textit{What would the closed ball be?})



\paragraph{Boundary points}



If we have some area, say a field, then the common sense notion of the \textit{boundary} is the points 'next to' both the inside and outside of the field. For a set, $S$ , we can define this rigorously by saying the boundary of the set contains all those points such that we can find points both inside and outside the set. We call the set of such points $\partial S$ .



Typically, when it exists the dimension of $\partial S$ is one lower than the dimension of $S$ . e.g. the boundary of a volume is a surface and the boundary of a surface is a curve.



This isn't always true; but it is true of all the sets we will be using.



A set $S$ is \textit{bounded} if there is some positive number such that we can encompass this set by a closed ball about $\vec0$ . $\rightarrow$ if every point in it is within a finite distance of the origin, i.e there exists some $r>0$ such that $\vec{x} \in S$ implies $\vec{x} $y=f(x)$



In order to ensure that $\lim _{x\rightarrow c}f(x)$ exists, we need to test it from two directions: one approaching $c$ from the left side ($x\rightarrow c^\-$) and the other approaching $c$ from the right side ($x\rightarrow c^\+$). Recall that



$$\lim _{x\rightarrow c}f(x) \text{ exists when } \lim _{x\rightarrow c^\-}f(x)=\lim _{x\rightarrow c^\+}f(x).$$



For example, $\lim _{x\rightarrow 0}\frac{1}{x}$ does not exist because $\lim _{x\rightarrow 0^\-}\frac{1}{x}=-\infty$ and $\lim _{x\rightarrow 0^\+}\frac{1}{x}=\infty$. Now, assume that there is a function with two variables:



> $z=f(x,y)$



If we want to take a limit, for example, $\lim _{(x,y)\rightarrow(a,b)}f(x,y)$, not only do we need to consider the limit from the direction of the $x$-axis, we also need to consider the limit from all directions, which includes the $y$-axis, lines, curves, etc. Generally speaking, if there is one direction where the calculated limit is different from others, the limit does not exist. We will be discussing this in detail \href{Calculus/Multivariable_Calculus/Limits_and_Continuity "wikilink"}{here}.



\paragraph{Differentiable functions}



% [Image: A two variable function]



Figure: A two variable function $f(x,y)=(x^2+3y^2)^{1-x^2-y^2}$



When we expand our scope into the 3-dimensional world, we have significantly more situations to consider. For example, derivatives. In previous chapters, derivatives only have one direction (the $x$-axis) because there is only one variable.



> $\frac{d(x^3+4x)}{dx}=3x^2+4$



When we have two or more variables, the rate of change can be calculated in different directions. For example, take a look at the image on the right. This is the graph of a two-variable function. Since there are two variables, the domain will be the whole $xy$-plane. We will graph the output $f(x,y)$ on the $z$-axis. The equation for the function on the right is:



> $f(x,y)=(x^2+3y^2)^{1-x^2-y^2}$



How can we calculate a derivative? The answer is to use partial derivatives. As the name suggests, it can only calculate a derivative "partially" because we can only calculate the rate of change of a graph in one direction.



\paragraph{Partial derivatives}



Notations are important for partial derivatives.



> $\frac{\partial}{\partial x}f(x,y)$ means the derivative of $f(x,y)$ in the $x$-axis direction, where we only view the $x$ as a variable while $y$ as a constant.

>

> $\frac{\partial}{\partial y}f(x,y)$ means the derivative of $f(x,y)$ in the $y$-axis direction, where we only view the $y$ as a variable while $x$ as a constant.



For simplicity, we will often use various standard abbreviations, so we can write most of the formulae on one line. This can make it easier to see the important details.



We can abbreviate partial differentials with a subscript, e.g.,



$$f _x(x,y)= \frac{\partial f}{\partial x} \quad f _y(x,y)= \frac{\partial f}{\partial y} \quad f _{xy}=\frac{\partial^2f}{\partial x \partial y}=\frac{\partial^2f}{\partial y \partial x}=f _{yx}$$



When we are using a subscript this way we will generally use the Heaviside $D$ (which stands for "directional") rather than $\partial$,



$$D _x h(x,y)= \frac{\partial h}{\partial x}  \quad D _x D _y h= D _y D _x h$$

$$D _{\mathbf{\hat{u}}}f(x,y) \text{ means the derivative of } f \text{ in the direction } \mathbf{\hat{u}}=\langle a,b\rangle$$



If we are using subscripts to label the axes, $x _1$, $x _2$ ..., then, rather than having two layers of subscripts, we will use the number as the subscript.



> $h _1 = D _1 h = \partial _1 h = \partial _{x _1}h = \frac{\partial h}{\partial x _1}$



We can also use subscripts for the components of a vector function, $\mathbf{u}=\langle u _x, u _y, u _z \rangle \text{ or } \mathbf{u}=(u _1,u _2,...,u _n)$



If we are using subscripts for both the components of a vector and for partial derivatives we will separate them with a comma.



$$u _{x,y}=\frac{\partial u _x}{\partial y}$$



The most widely used notation is $f _x$.



We will use whichever notation best suits the equation we are working with.



\paragraph{Directional derivatives}



Normally, a partial derivative of a function with respect to one of its variables, say, $x _j$, takes the derivative of that "slice" of that function parallel to the $x _j$th axis.



More precisely, we can think of cutting a function $f(x _1, ..., x _n)$ in space along the $x _j$th axis, with keeping everything but the $x _j$ variable constant.



From the definition, we have the partial derivative at a point $\mathbf{p}$ of the function along this slice as



$${\partial \mathbf{f} \over \partial x _j} = \lim _{t\rightarrow 0} {\mathbf{f}(\mathbf{p}+t\mathbf{e} _j) - \mathbf{f}(\mathbf{p}) \over t}$$



provided this limit exists.



Instead of the basis vector, which corresponds to taking the derivative along that axis, we can pick a vector in any direction (which we usually take as being a unit vector), and we take the \textit{directional derivative} of a function as



$${\partial \mathbf{f} \over \partial \mathbf{d}} = \lim _{t\rightarrow 0} {\mathbf{f}(\mathbf{p}+t\mathbf{d}) - \mathbf{f}(\mathbf{p}) \over t}$$

where $\mathbf{d}$ is the direction vector.



If we want to calculate directional derivatives, calculating them from the limit definition is rather painful, but, we have the following: if $\mathbb{f} : \mathbb{R} \rightarrow \mathbb{R}$ is differentiable at a point $\mathbf{p}$, $|\mathbf{p}| = 1$,



$${\partial \mathbf{f} \over \partial \mathbf{d}} = D _\mathbf{p} \mathbf{f}(\mathbf{d})$$



There is a closely related formulation which we'll look at in the next section.



\paragraph{Gradient vectors}



The partial derivatives of a scalar tell us how much it changes if we move along one of the axes. What if we move in a different direction?



We'll call the scalar $f$, and consider what happens if we move an infintesimal direction $\mathbf{dr}=(dx,dy,dz)$, using the chain rule.



$$\mathbf{df}=dx\frac{\partial f}{\partial x} + dy\frac{\partial f}{\partial y}+dz\frac{\partial f}{\partial z}$$



This is the dot product of $\mathbf{dr}$ with a vector whose components are the partial derivatives of $\mathbf{f}$, called the gradient of $\mathbf{f}$



$$\operatorname{grad} \mathbf{f} = \nabla \mathbf{f} = \left(\frac{\partial \mathbf{f}(\mathbf{p})}{\partial x _1},\cdots, \frac{\partial \mathbf{f}(\mathbf{p})}{\partial x _n}\right)$$



We can form directional derivatives at a point $\mathbf{p}$, in the direction $\mathbf{d}$ then by taking the dot product of the gradient with $\mathbf{d}$



$${\partial \mathbf{f}(\mathbf{p}) \over \partial \mathbf{d}} =\mathbf{d} \cdot \nabla \mathbf{f}(\mathbf{p})$$.



Notice that grad $f$ looks like a vector multiplied by a scalar. This particular combination of partial derivatives is commonplace, so we abbreviate it to



$$\nabla = \left( \frac{\partial }{\partial x}, \frac{\partial }{\partial y}, \frac{\partial }{\partial z}\right)$$



We can write the action of taking the gradient vector by writing this as an \textit{operator}. Recall that in the one-variable case we can write $d/dx$ for the action of taking the derivative with respect to $x$. This case is similar, but $\nabla$ acts like a vector.



We can also write the action of taking the gradient vector as:



$$\nabla  = \left( \frac{\partial }{\partial x _1}, \frac{\partial }{\partial x _2}, \cdots \frac{\partial }{\partial x _n}\right)$$



\paragraph{Properties of the gradient vector}



\paragraph{Geometry}



\begin{itemize}
    \item Grad $\mathbf{f}(\mathbf{p})$ is a vector pointing in the direction of steepest slope of $\mathbf{f}$. \|grad $\mathbf{f}(\mathbf{p})$\| is the rate of change of that slope at that point.
\end{itemize}

For example, if we consider $h(x,y)=x^2 + y^2$. The level sets of $h$ are concentric circles, centred on the origin, and



$$\nabla h = (h _x,h _y) = 2(x,y)= 2 \mathbf{r}$$ grad $h$ points directly away from the origin, at right angles to the contours.



\begin{itemize}
    \item Along a level set, $(\grad \mathbf{f})(\mathbf{p})$ is perpendicular to the level set $\\{\mathbf{x}\|\mathbf{f}(\mathbf{x})=\mathbf{f}(\mathbf{p}) \text{ at } \mathbf{x}=\mathbf{p}\\}$.
\end{itemize}

If $\mathbf{dr}$ points along the contours of $f$, where the function is constant, then $df$ will be zero. Since $df$ is a dot product, that means that the two vectors, $\mathbf{dr}$ and grad $f$, must be at right angles, i.e. the gradient is at right angles to the contours.



\paragraph{Algebraic properties}



Like $d/dx$, $\nabla$ is linear. For any pair of constants, $a$ and $b$, and any pair of scalar functions, $f$ and $g$



$$\frac{d}{dx} (af+bg)= a\frac{d}{dx}f + b\frac{d}{dx}g \quad \nabla (af+bg) = a \nabla f + b \nabla g$$



Since it's a vector, we can try taking its dot and cross product with other vectors, and with itself.



\paragraph{Product and chain rules}



Just as with ordinary differentiation, there are product rules for grad, div and curl.



\begin{itemize}
    \item If $g$ is a scalar and $\mathbf{v}$ is a vector, then
\end{itemize}

> the divergence of $g \mathbf{v}$ is

> $\newline \nabla \cdot (g\mathbf{v})=g \nabla \cdot \mathbf{v} + (\mathbf{v} \cdot \nabla) g$

> $\newline$the curl of $g \mathbf{v}$ is

> $\newline \nabla \times (g\mathbf{v}) = g(\nabla \times \mathbf{v})+ (\nabla g) \times \mathbf{v}$



\begin{itemize}
    \item If $u$ and $v$ are both vectors then the gradient of their dot product is
\end{itemize}

> $$\nabla ( \mathbf{u} \cdot \mathbf{v} ) =  \mathbf{u} \times (\nabla \times \mathbf{v} ) +  \mathbf{v} \times (\nabla \times \mathbf{u} ) +  (\mathbf{u} \cdot \nabla) \mathbf{v} + (\mathbf{v} \cdot \nabla) \mathbf{u}$$

> the divergence of their cross product is



> $$\nabla \cdot ( \mathbf{u} \times \mathbf{v} ) = \mathbf{v}\cdot ( \nabla \times \mathbf{u} ) - \mathbf{u}\cdot ( \nabla \times \mathbf{v} )$$

> the curl of their cross product is

> $$\nabla \times ( \mathbf{u} \times \mathbf{v} ) = (\mathbf{v} \cdot \nabla ) \mathbf{u} - (\mathbf{u} \cdot \nabla) \mathbf{v} + \mathbf{u}(\nabla \cdot \mathbf{v}) - \mathbf{v}(\nabla \cdot \mathbf{u})$$



We can also write chain rules. In the general case, when both functions

are vectors and the composition is defined, we can use the Jacobian

defined earlier.



$$\left. \nabla \mathbf{u}(\mathbf{v}) \right| _{\mathbf{r}}= \mathbf{J} _{\mathbf{v}} \left. \nabla \mathbf{v} \right| _{\mathbf{r}}$$

where $\mathbf{J} _{\mathbf{u}}$ is the Jacobian of $\mathbf{u}$ at the point $\mathbf{v}$.



Normally $\mathbf{J}$ is a matrix but if either the range or the domain of $\mathbf{u}$ is $\mathbb{R} ^1$ then it becomes a vector. In these special cases we can compactly write the chain rule using only vector notation.



\begin{itemize}
    \item If $g$ is a scalar function of a vector and $h$ is a scalar function of $g$ then
\end{itemize}

$$\nabla h(g) = \frac{dh}{dg} \nabla g$$



\begin{itemize}
    \item If $g$ is a scalar function of a vector then
\end{itemize}

$$\nabla = (\nabla g) \frac{d}{dg}$$ This substitution can be made in

any of the equations containing $\nabla$



\paragraph{Integration}



We have already considered differentiation of functions of more than one variable, which leads us to consider how we can meaningfully look at integration.



In the single variable case, we interpret the definite integral of a function to mean the area under the function. There is a similar interpretation in the multiple variable case: for example, if we have a paraboloid in $\mathbb{R} ^3$, we may want to look at the integral of that paraboloid over some region of the $xy$ plane, which will be the \textit{volume} under that curve and inside that region.



\paragraph{Riemann sums}



When looking at these forms of integrals, we look at the Riemann sum. Recall in the one-variable case we divide the interval we are integrating over into rectangles and summing the areas of these rectangles as their widths get smaller and smaller. For the multiple-variable case, we need to do something similar, but the problem arises how to split up $\mathbb{R} ^2$, or $\mathbb{R} ^3$, for instance.



To do this, we extend the concept of the interval, and consider what we call a $n$-interval. An $n$-interval is a set of points in some rectangular region with sides of some fixed width in each dimension, that is, a set in the form $$\\{\mathbf{x} \in \mathbb{R} ^n | a _i \leq x _i \leq b _i \text{ with } i = 0,...,n\\},$$ and its area/size/volume (which we simply call its \textit{measure} to avoid confusion) is the product of the lengths of all its sides.



So, an $n$-interval in $\mathbb{R} ^2$ could be some rectangular partition of the plane, such as $\\{(x,y) | x \in [0,1] \text{ and } y \in [0,2] \\}$. Its measure is 2.



If we are to consider the Riemann sum now in terms of sub-$n$-intervals of a region $\Omega$, it is



$$\sum _{i; S _i \subset \Omega} f(x^\* _i)m(S _i)$$



where $m(S _i)$ is the measure of the division of $\Omega$ into $k$ sub-$n$-intervals $S _i$, and $x^\* _i$ is a point in $S _i$. The index is important - we only perform the sum where $S _i$ falls completely within $\Omega$ - any $S _i$ that is not completely contained in $\Omega$ we ignore.



As we take the limit as $k$ goes to infinity, that is, we divide up $\Omega$ into finer and finer sub-$n$-intervals, and this sum is the same no matter how we divide up $\Omega$, we get the \textit{integral} of $f$ over $\Omega$ which we write



$$\int _\Omega f$$



For two dimensions, we may write



$$\int\int _\Omega f $$



and likewise for $n$ dimensions.



\paragraph{Iterated integrals}



Thankfully, we need not always work with Riemann sums every time we want to calculate an integral in more than one variable. There are some results that make life a bit easier for us.



For $\mathbb{R} ^2$, if we have some region bounded between two functions of the other variable (so two functions in the form $f(x)=y$, or $f(y)=x$), between a constant boundary (so, between $x=a$ and $x=b$ or $y=a$ and $y=b$), we have



$$\int _a^b\,\int _{f(x)}^{g(x)} h(x,y)\,dy dx$$



An important theorem (called \textit{Fubini's theorem}) assures us that this integral is the same as



$$\int\int _\Omega f, $$



if $f$ is continuous on the domain of integration.



\paragraph{Resources}



\begin{itemize}
    \item \href{https://tutorial.math.lamar.edu/classes/calciii/multivrblefcns.aspx}{Functions of Several Variables}, Paul's Online Notes
    \item \href{https://math.libretexts.org/Bookshelves/Calculus/Map\%3A_Calculus__Early_Transcendentals_(Stewart}{Functions of Several Variables}/14\%3A_Partial_Derivatives/14.01\%3A_Functions_of_Several_Variables), Mathematics LibreTexts
    \item \href{https://en.wikibooks.org/wiki/Calculus/Multivariable_calculus}{Multivariable Calculus}, WikiBooks: Calculus
\end{itemize}

\paragraph{Licensing}



Content obtained and/or adapted from:



\begin{itemize}
    \item \href{https://en.wikibooks.org/wiki/Calculus/Multivariable_calculus}{Multivariable Calculus, WikiBooks: Calculus} under a CC BY-SA license
\end{itemize}

\subsection{Learning Objectives}

\begin{enumerate}
    \item Student Learning Objectives (SLOs) for Multivariable Calculus
    \item \textbf{Understand and Apply Concepts of Length and Distance in $\mathbb{R}^n$:}
    \item \textbf{Outcome:} Students will be able to calculate the length of vectors and the distance between vectors in $\mathbb{R}^n$ using the Euclidean metric.
    \item \textbf{Details:} This includes the ability to generalize the Pythagorean theorem to higher dimensions and apply the formula $$\|\vec{x}\| = \sqrt{x_1^2 + x_2^2 + \cdots + x_n^2}$$ for vector length, and $$d(\vec{x}, \vec{y}) = \sqrt{\sum_{i=1}^n (x_i - y_i)^2}$$ for the distance between vectors.
    \item \textbf{Identify and Describe Open and Closed Balls in Multivariable Spaces:}
    \item \textbf{Outcome:} Students will be able to define and visualize open and closed balls in $\mathbb{R}^n$, and understand their significance in multivariable calculus.
    \item \textbf{Details:} This includes being able to write and interpret the mathematical definitions of open balls $B(\vec{a}, r) = \{\vec{x} \in \mathbb{R}^n \mid d(\vec{x}, \vec{a}) < r\}$ and closed balls $\overline{B}(\vec{a}, r) = \{\vec{x} \in \mathbb{R}^n \mid d(\vec{x}, \vec{a}) \le r\}$, and explain the geometric intuition behind these concepts.
    \item \textbf{Calculate and Interpret Partial Derivatives and Directional Derivatives:}
    \item \textbf{Outcome:} Students will be able to compute partial derivatives and directional derivatives of multivariable functions and interpret their geometric meaning.
    \item \textbf{Details:} This includes the ability to use the notation $$\frac{\partial}{\partial x_i} f(x_1, x_2, \ldots, x_n)$$ for partial derivatives, and $$D_{\vec{u}} f(\vec{x}) = \lim_{t \to 0} \frac{f(\vec{x} + t\vec{u}) - f(\vec{x})}{t}$$ for directional derivatives. Students should also be able to relate these derivatives to the gradient vector and understand how they describe the rate of change of the function in different directions.
\end{enumerate}

\subsection{Assessment Instrument}

\begin{enumerate}
    \item \textbf{Assessment Instrument: Functions of Several Variables Quiz} \textbf{Format:}
    \item \textbf{Type:} Mixed-format quiz (multiple-choice, short answer, and problem-solving questions)
    \item \textbf{Duration:} 45 minutes
    \item \textbf{Total Marks:} 30 points \textbf{Sections:}
    \item \textbf{Understanding Vector Lengths and Distances (10 points):}
    \item \textbf{Question 1:} Calculate the length of the vector $\vec{x} = (3, 4)$.
    \item \textbf{Question 2:} Determine the distance between the vectors $\vec{a} = (1, 2, 3)$ and $\vec{b} = (4, 5, 6)$.
    \item \textbf{Question 3:} Multiple choice: Which of the following represents the length of the vector $\vec{v} = (a, b, c)$?
    \item a. $\sqrt{a^2 + b^2}$
    \item b. $a + b + c$
    \item c. $\sqrt{a^2 + b^2 + c^2}$
    \item d. $|a| + |b| + |c|$
    \item \textbf{Question 4:} Short answer: Explain the concept of the Euclidean distance between two points in $\mathbb{R}^n$.
    \item \textbf{Question 5:} Problem-solving: Calculate the length of the vector $\vec{x} = (1, -2, 2)$ and the distance between $\vec{x}$ and $\vec{y} = (3, 0, -1)$.
    \item \textbf{Open and Closed Balls (10 points):}
    \item \textbf{Question 6:} Define an open ball in $\mathbb{R}^3$ centered at $\vec{a} = (1, 1, 1)$ with radius $2$.
    \item \textbf{Question 7:} Multiple choice: Which of the following describes a closed ball in $\mathbb{R}^2$?
    \item a. $\{\vec{x} \in \mathbb{R}^2 \mid d(\vec{x}, \vec{a}) < r\}$
    \item b. $\{\vec{x} \in \mathbb{R}^2 \mid d(\vec{x}, \vec{a}) \le r\}$
    \item c. $\{\vec{x} \in \mathbb{R}^2 \mid d(\vec{x}, \vec{a}) = r\}$
    \item d. $\{\vec{x} \in \mathbb{R}^2 \mid d(\vec{x}, \vec{a}) \ge r\}$
    \item \textbf{Question 8:} Short answer: Describe the difference between an open ball and a closed ball in $\mathbb{R}^n$.
    \item \textbf{Question 9:} Problem-solving: Find the set of points that form the closed ball $\overline{B}((0, 0), 3)$ in $\mathbb{R}^2$.
    \item \textbf{Question 10:} Multiple choice: In $\mathbb{R}$, the open ball $B(a, r)$ corresponds to which of the following intervals?
    \item a. $[a - r, a + r]$
    \item b. $(a - r, a + r)$
    \item c. $(a - r, a + r]$
    \item d. $[a - r, a + r)$
    \item \textbf{Partial Derivatives and Directional Derivatives (10 points):}
    \item \textbf{Question 11:} Calculate the partial derivative of the function $f(x, y) = x^2 y + y^3$ with respect to $x$.
    \item \textbf{Question 12:} Determine the partial derivative of $f(x, y) = e^{xy} + \sin(x) \cos(y)$ with respect to $y$.
    \item \textbf{Question 13:} Multiple choice: Which notation represents the partial derivative of $f(x, y)$ with respect to $y$?
    \item a. $f_x(x, y)$
    \item b. $f_{xy}(x, y)$
    \item c. $f_y(x, y)$
    \item d. $\partial f / \partial x$
    \item \textbf{Question 14:} Short answer: Explain the concept of a directional derivative and how it is different from a partial derivative.
    \item \textbf{Question 15:} Problem-solving: Given $f(x, y) = x^2 + y^2$, calculate the directional derivative at the point $(1, 1)$ in the direction of the vector $\vec{d} = (1, 1)$. \textbf{Scoring Rubric:}
    \item \textbf{Understanding Vector Lengths and Distances:} 2 points for each correct calculation or explanation (total 10 points)
    \item \textbf{Open and Closed Balls:} 2 points for each correct definition, description, or solution (total 10 points)
    \item \textbf{Partial Derivatives and Directional Derivatives:} 2 points for each correct derivative or explanation (total 10 points) The quiz will help assess students' understanding and application of vector lengths, distances, open and closed balls, partial derivatives, and directional derivatives in the context of functions of several variables. This ensures that they can perform necessary calculations and understand the geometric interpretations accurately.
\end{enumerate}

%%===============================================================
\section{Limits and Continuity in Higher Dimensions}
\label{sec:limits-and-continuity-in-higher-dimensions}
%%===============================================================

\paragraph{Limits and continuity}



Before we can look at derivatives of multivariate functions, we need to look at how limits work with functions of several variables first, just like in the single variable case.



If we have a function $f : \mathbb{R}^m \rightarrow \mathbb{R}^n$, we say that $f(\mathbf{x})$ approaches $\mathbf{b}$ (in $\mathbb{R}^n$) as $\mathbf{x}$ approaches $\mathbf{a}$ (in $\mathbb{R}^m$) if, for all positive $\varepsilon$, there is a corresponding positive number $\delta$, $\| f(\mathbf{x}) - \mathbf{b} \| < \varepsilon$ whenever $\| \mathbf{x} - \mathbf{a} \| < \delta$, with $\mathbf{x} \neq \mathbf{a}$.



This means that by making the difference between $\mathbf{x}$ and $\mathbf{a}$ smaller, we can make the difference between $f(\mathbf{x})$ and $\mathbf{b}$ as small as we want.



If the above is true, we say

\begin{itemize}
    \item $f(\mathbf{x})$ has \textit{limit} $\mathbf{b}$ at $\mathbf{a}$
    \item $\lim_{\mathbf{x} \to \mathbf{a}} f(\mathbf{x}) = \mathbf{b}$
    \item $f(\mathbf{x})$ approaches $\mathbf{b}$ as $\mathbf{x}$ approaches $\mathbf{a}$
    \item $f(\mathbf{x}) \to \mathbf{b}$ as $\mathbf{x} \to \mathbf{a}$
\end{itemize}

These four statements are all equivalent.



\paragraph{Rules}



Since this is an almost identical formulation of limits in the single variable case, many of the limit rules in the one variable case are the same as in the multivariate case.



For $\mathbf{f}$ and $\mathbf{g}$, mapping $\mathbb{R}^m$ to $\mathbb{R}^n$, and $h(\mathbf{x})$ a scalar function mapping $\mathbb{R}^m$ to $\mathbb{R}$, with

\begin{itemize}
    \item $\mathbf{f}(\mathbf{x}) \to \mathbf{b}$ as $\mathbf{x} \to \mathbf{a}$
    \item $\mathbf{g}(\mathbf{x}) \to \mathbf{c}$ as $\mathbf{x} \to \mathbf{a}$
    \item $h(\mathbf{x}) \to \mathbf{H}$ as $\mathbf{x} \to \mathbf{a}$
\end{itemize}

then:

\begin{itemize}
    \item $\lim _{\mathbf{x} \to \mathbf{a}} (\mathbf{f} + \mathbf{g}) = \mathbf{b} + \mathbf{c}$
    \item $\lim _{\mathbf{x} \to \mathbf{a}} (h\mathbf{f}) = H\mathbf{b}$
\end{itemize}

and consequently

\begin{itemize}
    \item $\lim _{\mathbf{x} \to \mathbf{a}} (\mathbf{f} \cdot \mathbf{g}) = \mathbf{b} \cdot \mathbf{c}$
    \item $\lim _{\mathbf{x} \to \mathbf{a}} (\mathbf{f} \times \mathbf{g}) = \mathbf{b} \times \mathbf{c}$
\end{itemize}

when $H \neq 0$

\begin{itemize}
    \item $\lim _{\mathbf{x} \to \mathbf{a}} \left( \frac{\mathbf{f}}{h} \right) = \frac{\mathbf{b}}{H}$
\end{itemize}

\paragraph{Continuity}

Again, we can use a similar definition to the one variable case to

formulate a definition of continuity for multiple variables.







Given that $\mathbf{f} : \mathbb{R}^m \rightarrow \mathbb{R}^n$, then we say that $\mathbf{f}$ is continuous at a point $\mathbf{a}$ in $\mathbb{R}^m$ if $\mathbf{f}(\mathbf{a})$ is defined and

$$\lim_{{\mathbf{x} \to \mathbf{a}}} \mathbf{f}(\mathbf{x}) = \mathbf{f}(\mathbf{a})$$



Just as for functions of one dimension, if $f$ and $g$ are both continuous at $\mathbf{p}$, then $f+g$, $\lambda f$ (for any scalar $\lambda$), $f \cdot g$, and $f \times g$ are all continuous as well. If $\varphi : \mathbb{R}^m \rightarrow \mathbb{R}$ is continuous at $\mathbf{p}$, then $q\varphi$ is continuous and, where $\varphi \neq 0$, $f/\varphi$ is also continuous.



From these facts, we also have that if $A$ is some matrix which is $n \times m$ in size, with $\mathbf{x} \in \mathbb{R}^m$, a function $\mathbf{f}(\mathbf{x})=A \mathbf{x}$ is continuous in that the function can be expanded in the form $x _1 a _1 + \ldots + x _m a _m$. This can be verified from the points above.



If $\mathbf{f} : \mathbb{R}^m \rightarrow \mathbb{R}^n$ which is in the form $\mathbf{f}(\mathbf{x}) = (f _1(\mathbf{x}), \ldots, f _n(\mathbf{x}))$ is continuous if and only if each of its component functions is a polynomial or rational function, whenever it is defined.



Finally, if $\mathbf{f}$ is continuous at $\mathbf{p}$ and $\mathbf{g}$ is continuous at $\mathbf{f}(\mathbf{p})$, then $\mathbf{g}(\mathbf{f}(\mathbf{x}))$ is continuous at $\mathbf{p}$.



\paragraph{Special note about limits}



It is important to note that we can approach a point \textit{in more than one direction}, and thus, the direction that we approach that point counts in our evaluation of the limit. It may be the case that a limit may exist moving in one direction, but not in another.



\paragraph{Resources}



\begin{itemize}
    \item \href{https://math.libretexts.org/Bookshelves/Calculus/Book\%3A_Calculus_(Apex}{Limits and Continuity of Multivariable Functions}/12\%3A_Functions_of_Several_Variables/12.02\%3A_Limits_and_Continuity_of_Multivariable_Functions), Mathematics LibreTexts
    \item \href{https://en.wikibooks.org/wiki/Calculus/Multivariable_Calculus/Limits_and_Continuity}{Limits and Continuity}, WikiBooks: Calculus/Multivariable Calculus
\end{itemize}

\paragraph{Licensing}



Content obtained and/or adapted from:



\begin{itemize}
    \item \href{https://en.wikibooks.org/wiki/Calculus/Multivariable_Calculus/Limits_and_Continuity}{Limits and Continuity, WikiBooks: Calculus/Multivariable Calculus} under a CC BY-SA license
\end{itemize}

\subsection{Learning Objectives}

\begin{enumerate}
    \item \textbf{Understand Multivariate Limits:}
    \item \textbf{Outcome:} Students will be able to evaluate limits of functions of several variables.
    \item \textbf{Details:} This includes understanding and applying the definition of a limit for multivariate functions, identifying the conditions under which limits exist, and evaluating limits using different approaches.
    \item \textbf{Apply Limit Laws to Multivariate Functions:}
    \item \textbf{Outcome:} Students will be able to use limit laws to simplify and evaluate limits of multivariate functions.
    \item \textbf{Details:} This involves applying the limit laws, such as sum, product, and quotient rules, to functions of several variables, and understanding the conditions under which these laws hold.
    \item \textbf{Evaluate Continuity of Multivariate Functions:}
    \item \textbf{Outcome:} Students will be able to determine the continuity of functions of several variables at a given point.
    \item \textbf{Details:} This includes understanding the definition of continuity for multivariate functions, using the limit definition to verify continuity, and applying the properties of continuous functions, such as the continuity of sums, products, and compositions of continuous functions.
\end{enumerate}

\subsection{Assessment Instrument}

\begin{enumerate}
    \item \textbf{Assessment Instrument: Limits and Continuity of Functions of Several Variables Quiz} \textbf{Format:}
    \item \textbf{Type:} Mixed-format quiz (multiple-choice, short answer, and problem-solving questions)
    \item \textbf{Duration:} 45 minutes
    \item \textbf{Total Marks:} 30 points \textbf{Sections:}
    \item \textbf{Understanding Multivariate Limits (10 points):}
    \item \textbf{Question 1:} Determine the limit if it exists:
    \item a. $\lim_{\mathbf{x} \to (1,1)} (x_1^2 + x_2^2)$ (2 points)
    \item b. $\lim_{\mathbf{x} \to (0,0)} \dfrac{x_1^2 - x_2^2}{x_1^2 + x_2^2}$ (4 points)
    \item c. $\lim_{\mathbf{x} \to (0,0)} \dfrac{x_1x_2}{x_1^2 + x_2^2}$ (4 points)
    \item \textbf{Applying Limit Laws (10 points):}
    \item \textbf{Question 2:} Evaluate the following limits using limit laws:
    \item a. $\lim_{\mathbf{x} \to (2,3)} (3x_1 - 2x_2)$ (2.5 points)
    \item b. $\lim_{\mathbf{x} \to (1,4)} (x_1^2 \cdot x_2)$ (2.5 points)
    \item c. $\lim_{\mathbf{x} \to (0,2)} \dfrac{x_1 + x_2}{x_1 - x_2}$ if it exists (2.5 points)
    \item d. $\lim_{\mathbf{x} \to (3,3)} \left( \dfrac{x_1^2}{x_2} + x_2^3 \right)$ (2.5 points)
    \item \textbf{Evaluating Continuity (10 points):}
    \item \textbf{Question 3:} Determine whether the following functions are continuous at the given point:
    \item a. $f(x_1, x_2) = x_1^2 + x_2^2$ at $(0,0)$ (2.5 points)
    \item b. $f(x_1, x_2) = \dfrac{x_1x_2}{x_1^2 + x_2^2}$ at $(1,1)$ (2.5 points)
    \item c. $f(x_1, x_2) = \sin(x_1) \cos(x_2)$ at $(\pi, \pi/2)$ (2.5 points)
    \item d. $f(x_1, x_2) = \dfrac{x_1^3 - x_2^3}{x_1 - x_2}$ at $(1,1)$ (2.5 points) \textbf{Scoring Rubric:}
    \item \textbf{Understanding Multivariate Limits:} Points are awarded based on the correct evaluation of the limits and appropriate justification of the steps taken (total 10 points).
    \item \textbf{Applying Limit Laws:} Points are awarded for the correct application of limit laws and accurate simplification of the limits (total 10 points).
    \item \textbf{Evaluating Continuity:} Points are awarded for correctly determining the continuity of the functions at the given points and providing appropriate justification (total 10 points). The quiz will help assess the students' understanding and application of limits and continuity in multivariate functions. It ensures that they can correctly evaluate limits, apply limit laws, and determine the continuity of functions of several variables.
\end{enumerate}

%%===============================================================
\section{Partial Derivatives}
\label{sec:partial-derivatives}
%%===============================================================

In mathematics, a \textbf{partial derivative} of a function of several variables is its derivative with respect to one of those variables, with the others held constant (as opposed to the total derivative, in which all variables are allowed to vary). Partial derivatives are used in vector calculus and differential geometry.



The partial derivative of a function $f(x, y, \dots)$ with respect to the variable $x$ is variously denoted by



> $f' _x$, $f _x$, $\partial _x f$, $\ D _xf$, $D _1f$, $\frac{\partial}{\partial x}f$, or $\frac{\partial f}{\partial x}$.



Sometimes, for $z=f(x, y, \ldots)$, the partial derivative of $z$ with respect to $x$ is denoted as $\tfrac{\partial z}{\partial x}.$ Since a partial derivative generally has the same arguments as the original function, its functional dependence is sometimes explicitly signified by the notation, such as in:



$$f _x(x, y, \ldots), \frac{\partial f}{\partial x} (x, y, \ldots).$$



The symbol used to denote partial derivatives is $\partial$. One of the first known uses of this symbol in mathematics is by Marquis de Condorcet from 1770, who used it for partial differences. The modern partial derivative notation was created by Adrien-Marie Legendre (1786) (although he later abandoned it, Carl Gustav Jacob Jacobi reintroduced the symbol in 1841).



\paragraph{Definition}



Like ordinary derivatives, the partial derivative is defined as a limit. Let \textit{U} be an open subset of $\R^n$ and $f:U\to\R$ a function. The partial derivative of $f$ at the point $\mathbf{a}=(a _1, \ldots, a _n) \in U$ with respect to the $i$-th variable $x _i$ is defined as



$$\begin{aligned}

\frac{\partial }{\partial x _i }f(\mathbf{a}) & = \lim _{h \to 0} \frac{f(a _1, \ldots , a _{i-1}, a _i+h, a _{i+1}, \ldots ,a _n) -

f(a _1, \ldots, a _i, \dots ,a _n)}{h} \\\& = \lim _{h \to 0} \frac{f(\mathbf{a}+h\mathbf{e _i}) -

f(\mathbf{a})}{h}

\end{aligned}$$



Even if all partial derivatives $\frac{\partial f}{\partial x _i} (a)$ exist at a given point $a$, the function need not be continuous there. However, if all partial derivatives exist in a neighborhood of $a$ and are continuous there, then $f$ is totally differentiable in that neighborhood and the total derivative is continuous. In this case, it is said that $f$ is a $C^1$ function. This can be used to generalize for vector valued functions, $f:U \to \mathbb{R} ^m$, by carefully using a componentwise argument.



The partial derivative $\frac{\partial f}{\partial x}$ can be seen as another function defined on \textit{U} and can again be partially differentiated. If all mixed second order partial derivatives are continuous at a point (or on a set), $f$ is termed a $C^2$ function at that point (or on that set); in this case, the partial derivatives can be exchanged by Clairaut's theorem:



$$\frac{\partial^2f}{\partial x _i \partial x _j} = \frac{\partial^2f} {\partial x _j \partial x _i}$$



\paragraph{Notation}



For the following examples, let $f$ be a function in $x, y$ and $z$.



First-order partial derivatives:



$$\frac{ \partial f}{ \partial x} = f _x = \partial _x f$$



Second-order partial derivatives:



$$\frac{ \partial^2 f}{ \partial x^2} = f _{xx} = \partial _{xx} f = \partial _x^2 f$$



Second-order mixed derivatives:



$$\frac{\partial^2 f}{\partial y \,\partial x} = \frac{\partial}{\partial y} \left( \frac{\partial f}{\partial x} \right) = (f _{x}) _{y} = f _{xy} = \partial _{yx} f = \partial _y \partial _x f$$



Higher-order partial and mixed derivatives:



$$\frac{\partial^{i+j+k} f}{\partial x^i \partial y^j \partial z^k } = f^{(i, j, k)} = \partial _x^i \partial _y^j \partial _z^k f$$



When dealing with functions of multiple variables, some of these variables may be related to each other, thus it may be necessary to specify explicitly which variables are being held constant to avoid ambiguity. In fields such as statistical mechanics, the partial derivative of $f$ with respect to $x$, holding $y$ and $z$ constant, is often expressed as



$$\left( \frac{\partial f}{\partial x} \right) _{y,z}$$



Conventionally, for clarity and simplicity of notation, the partial derivative \textit{function} and the \textit{value} of the function at a specific point are conflated by including the function arguments when the partial derivative symbol (Leibniz notation) is used. Thus, an expression like



$$\frac{\partial f(x,y,z)}{\partial x}$$



is used for the function, while



$$\frac{\partial f(u,v,w)}{\partial u}$$



might be used for the value of the function at the point $(x,y,z)=(u,v,w)$. However, this convention breaks down when we want to evaluate the partial derivative at a point like $(x,y,z)=(17, u+v, v^2)$. In such a case, evaluation of the function must be expressed in an unwieldy manner as



$$\frac{\partial f(x,y,z)}{\partial x}(17, u+v, v^2)$$



or



$$\left. \frac{\partial f(x,y,z)}{\partial x}\right | _{(x,y,z)=(17, u+v, v^2)}$$



in order to use the Leibniz notation. Thus, in these cases, it may be preferable to use the Euler differential operator notation with $D _i$ as the partial derivative symbol with respect to the \textit{i}th variable. For instance, one would write $D _1 f(17, u+v, v^2)$ for the example described above, while the expression $D _1 f$ represents the partial derivative \textit{function} with respect to the 1st variable.



For higher order partial derivatives, the partial derivative (function) of $D _i f$ with respect to the \textit{j}th variable is denoted $D _j(D _i f)=D _{i,j} f$. That is, $D _j\circ D _i =D _{i,j}$, so that the variables are listed in the order in which the derivatives are taken, and thus, in reverse order of how the composition of operators is usually notated. Of course, Clairaut's theorem implies that $D _{i,j}=D _{j,i}$ as long as comparatively mild regularity conditions on $f$ are satisfied.



\paragraph{Gradient}



An important example of a function of several variables is the case of a scalar-valued function $f(x _1, ..., x _n)$ on a domain in Euclidean space $\mathbb{R} ^n$ (e.g., on $\mathbb{R} ^2$ or $\mathbb{R} ^3$). In this case $f$ has a partial derivative $\frac{\partial f}{\partial x _j}$ with respect to each variable $x _j$. At the point $a$, these partial derivatives define the vector



$$\nabla f(a) = \left(\frac{\partial f}{\partial x _1}(a), \ldots, \frac{\partial f}{\partial x _n}(a)\right).$$



This vector is called the \textit{gradient} of $f$ at $a$. If $f$ is differentiable at every point in some domain, then the gradient is a vector-valued function $\nabla f$ which takes the point $a$ to the vector $\nabla f(a)$. Consequently, the gradient produces a vector field.



A common abuse of notation is to define the del operator ($\nabla$) as follows in three-dimensional Euclidean space $\R^3$ with unit vectors $\hat{\mathbf{i}}, \hat{\mathbf{j}}, \hat{\mathbf{k}}$:



$$\nabla = \left[{\frac{\partial}{\partial x}} \right] \hat{\mathbf{i}} + \left[{\frac{\partial}{\partial y}} \right] \hat{\mathbf{j}} + \left[{\frac{\partial}{\partial z}}\right] \hat{\mathbf{k}}$$



Or, more generally, for \textit{n}-dimensional Euclidean space $\R^n$ with coordinates $x _1, \ldots, x _n$ and unit vectors $\hat{\mathbf{e}} _1, \ldots, \hat{\mathbf{e}} _n$:



$$\nabla = \sum _{j=1}^n \left[\frac{\partial}{\partial x _j} \right] \hat{\mathbf{e}} _j = \left[\frac{\partial}{\partial x _1} \right] \hat{\mathbf{e}} _1 + \left[\frac{\partial}{\partial x _2} \right] \hat{\mathbf{e}} _2 + \dots + \left[\frac{\partial}{\partial x _n} \right] \hat{\mathbf{e}} _n$$



\paragraph{Directional derivative}



The \textit{directional derivative} of a scalar function



$$f(\mathbf{x}) = f(x _1, x _2, \ldots, x _n)$$ along a vector



$$\mathbf{v} = (v _1, \ldots, v _n)$$ is the function $\nabla _{\mathbf{v}}{f}$ defined by the limit



$$\nabla _{\mathbf{v}}{f}(\mathbf{x}) = \lim _{h \to 0}{\frac{f(\mathbf{x} + h\mathbf{v}) - f(\mathbf{x})}{h}}.$$



This definition is valid in a broad range of contexts, for example where the norm of a vector (and hence a unit vector) is undefined.



\paragraph{Example}



Suppose that $f$ is a function of more than one variable. For instance,



$$z = f(x,y) = x^2 + xy + y^2$$.



% [Image: A graph]



Figure: A graph of $z = x^2 + xy + y^2$. For the partial derivative at (1, 1) that leaves $y$ constant, the corresponding tangent line is parallel to the \textit{xz}-plane.



% [Image: A slice of the graph above]



Figure: A slice of the graph above showing the function in the \textit{xz}-plane at $y=1$. Note that the two axes are shown here with different scales. The slope of the tangent line is 3.



The graph of this function defines a surface in Euclidean space. To every point on this surface, there are an infinite number of tangent lines. Partial differentiation is the act of choosing one of these lines and finding its slope. Usually, the lines of most interest are those that are parallel to the $xz$-plane, and those that are parallel to the $yz$-plane (which result from holding either $y$ or $x$ constant, respectively).



To find the slope of the line tangent to the function at $P(1, 1)$ and parallel to the $xz$-plane, we treat $y$ as a constant. The graph and this plane are shown on the right. Below, we see how the function looks on the plane $y=1$ . By finding the derivative of the equation while assuming that $y$ is a constant, we find that the slope of $f$ at the point $(x, y)$ is:



> $\frac{\partial z}{\partial x} = 2x+y$.



So at $(1, 1)$, by substitution, the slope is 3. Therefore,



> $\frac{\partial z}{\partial x} = 3$



at the point $(1, 1)$. That is, the partial derivative of $z$ with respect to $x$ at $(1, 1)$ is 3, as shown in the graph.



The function $f$ can be reinterpreted as a family of functions of one variable indexed by the other variables:



> $f(x,y) = f _y(x) = x^2 + xy + y^2.$



In other words, every value of $y$ defines a function, denoted $f _y$, which is a function of one variable $x$. That is,



> $f _y(x) = x^2 + xy + y^2.$



In this section the subscript notation $f _y$ denotes a function contingent on a fixed value of $y$, and not a partial derivative.



Once a value of $y$ is chosen, say $a$, then $f(x,y)$ determines a function $f _a$ which traces a curve $x^2 + ax + a^2$ on the $xz$-plane:



> $f _a(x) = x^2 + ax + a^2$.



In this expression, $a$ is a \textit{constant}, not a \textit{variable}, so $f _a$ is a function of only one real variable, that being $x$. Consequently, the definition of the derivative for a function of one variable applies:



> $f_a'(x) = 2x + a$.



The above procedure can be performed for any choice of $a$. Assembling the derivatives together into a function gives a function which describes the variation of $f$ in the $x$ direction:



> $\frac{\partial f}{\partial x}(x,y) = 2x + y.$



This is the partial derivative of $f$ with respect to $x$. Here $\partial$ is a rounded \textit{d} called the \textit{partial derivative symbol}; to distinguish it from the letter \textit{d}, $\partial$ is sometimes pronounced "partial".



\paragraph{Higher order partial derivatives}



Second and higher order partial derivatives are defined analogously to the higher order derivatives of univariate functions. For the function $f(x, y, ...)$ the "own" second partial derivative with respect to $x$ is simply the partial derivative of the partial derivative (both with respect to $x$):



$$\frac{\partial ^2 f}{\partial x^2} \equiv \partial \frac{{\partial f / \partial x}}{{\partial x}} \equiv \frac{{\partial f _x }}{{\partial x }} \equiv f _{xx}.$$



The cross partial derivative with respect to $x$ and $y$ is obtained by taking the partial derivative of $f$ with respect to $x$, and then taking the partial derivative of the result with respect to $y$, to obtain



$$\frac{{\partial ^2 f}}{{\partial y\, \partial x}} \equiv \partial \frac{{\partial f / \partial x}}{{\partial y}} \equiv \frac{{\partial f _x}}{{\partial y}} \equiv f _{{xy}}.$$



Schwarz's theorem states that if the second derivatives are continuous the expression for the cross partial derivative is unaffected by which variable the partial derivative is taken with respect to first and which is taken second. That is,



$$\frac {{\partial ^2 f}}{{\partial x\, \partial y}} = \frac{{\partial ^2 f}}{{\partial y\, \partial x}}$$



or equivalently $f _{{yx}}=f _{{xy}}.$



Own and cross partial derivatives appear in the Hessian matrix which is used in the second order conditions in optimization problems.



\paragraph{Antiderivative analogue}



There is a concept for partial derivatives that is analogous to antiderivatives for regular derivatives. Given a partial derivative, it allows for the partial recovery of the original function.



Consider the example of



$$\frac{\partial z}{\partial x} = 2x+y.$$



The "partial" integral can be taken with respect to $x$ (treating $y$ as constant, in a similar manner to partial differentiation):



$$z = \int \frac{\partial z}{\partial x} \,dx = x^2 + xy + g(y)$$



Here, the "constant" of integration is no longer a constant, but instead a function of all the variables of the original function except $x$. The reason for this is that all the other variables are treated as constant when taking the partial derivative, so any function which does not involve $x$ will disappear when taking the partial derivative, and we have to account for this when we take the antiderivative. The most general way to represent this is to have the "constant" represent an unknown function of all the other variables.



Thus the set of functions $x^2 + xy + g(y)$, where \textit{g} is any one-argument function, represents the entire set of functions in variables $x$, $y$ that could have produced the $x$-partial derivative $2x + y$.



If all the partial derivatives of a function are known (for example, with the gradient), then the antiderivatives can be matched via the above process to reconstruct the original function up to a constant. Unlike in the single-variable case, however, not every set of functions can be the set of all (first) partial derivatives of a single function. In other words, not every vector field is conservative.



\paragraph{Applications}



\paragraph{Geometry}



% [Image: The volume of a cone depends on height and radius]



Figure: The volume of a cone depends on height and radius.



The volume $V$ of a cone depends on the cone's height \textit{h} and its radius \textit{r} according to the formula



$$V(r, h) = \frac{\pi r^2 h}{3}.$$



The partial derivative of $V$ with respect to $r$ is



$$\frac{ \partial V}{\partial r} = \frac{ 2 \pi r h}{3},$$



which represents the rate with which a cone's volume changes if its radius is varied and its height is kept constant. The partial derivative with respect to $h$ equals $\frac{\pi r^2}{3},$ which represents the rate with which the volume changes if its height is varied and its radius is kept constant.



By contrast, the \textit{total} derivative of \textit{V} with respect to $r$ and $h$ are respectively



$$\frac{dV}{dr} = \overbrace{\frac{2 \pi r h}{3}}^\frac{ \partial V}{\partial r} + \overbrace{\frac{\pi r^2}{3}}^\frac{ \partial V}{\partial h}\frac{dh}{dr}$$



and



$$\frac{dV}{dh} = \overbrace{\frac{\pi r^2}{3}}^\frac{\partial V}{\partial h} + \overbrace{\frac{2 \pi r h}{3}}^\frac{ \partial V}{\partial r}\frac{dr}{dh}$$



The difference between the total and partial derivative is the elimination of indirect dependencies between variables in partial derivatives.



If (for some arbitrary reason) the cone's proportions have to stay the same, and the height and radius are in a fixed ratio $k$,



$$k = \frac{h}{r} = \frac{dh}{dr}.$$



This gives the total derivative with respect to $r$:



$$\frac{dV}{dr} = \frac{2 \pi r h}{3} + \frac{\pi r^2}{3}k$$



which simplifies to:



$$\frac{dV}{dr} = k \pi r^2$$



Similarly, the total derivative with respect to $h$ is:



$$\frac{dV}{dh} = \pi r^2$$



The total derivative with respect to \textit{both} $r$ and $h$ of the volume intended as scalar function of these two variables is given by the gradient vector



$$\nabla V = \left(\frac{\partial V}{\partial r},\frac{\partial V}{\partial h}\right) = \left(\frac{2}{3}\pi rh, \frac{1}{3}\pi r^2\right)$$.



\paragraph{Optimization}



Partial derivatives appear in any calculus-based optimization problem with more than one choice variable. For example, in economics a firm may wish to maximize profit $\pi(x,y)$ with respect to the choice of the quantities $x$ and $y$ of two different types of output. The first order conditions for this optimization are $\pi _x = 0 = \pi _y$. Since both partial derivatives $\pi _x$ and $\pi _y$ will generally themselves be functions of both arguments $x$ and $y$, these two first order conditions form a system of two equations in two unknowns.



\paragraph{Thermodynamics, quantum mechanics and mathematical physics}



Partial derivatives appear in thermodynamic equations like Gibbs-Duhem equation, in quantum mechanics as Schrodinger wave equation as well in other equations from mathematical physics. Here the variables being held constant in partial derivatives can be ratio of simple variables like mole fractions \textit{x\textasciitilde{}i\textasciitilde{}} in the following example involving the Gibbs energies in a ternary mixture system:



$$\bar{G _2}= G + (1-x _2) \left(\frac{{\partial G}}{{\partial x _2}}\right) _{\frac{x _1}{x _3}}$$



Express mole fractions of a component as functions of other components' mole fraction and binary mole ratios:



$$x _1 = \frac{1-x _2}{1+\frac{x _3}{x _1}}$$



$$x _3 = \frac{1-x _2}{1+\frac{x _1}{x _3}}$$



Differential quotients can be formed at constant ratios like those above:



$$\left(\frac{\partial x _1}{\partial x _2}\right) _{\frac{x _1}{x _3}} = - \frac{x _1}{1-x _2}$$



$$\left(\frac{\partial x _3}{\partial x _2}\right) _{\frac{x _1}{x _3}} = - \frac{x _3}{1-x _2}$$



Ratios X, Y, Z of mole fractions can be written for ternary and multicomponent systems:



$$X = \frac{x _3}{x _1 + x _3}$$



$$Y = \frac{x _3}{x _2 + x _3}$$



$$Z = \frac{x _2}{x _1 + x _2}$$



which can be used for solving partial differential equations like:



$$\left(\frac{\partial \mu _2}{\partial n _1}\right) _{n _2, n _3} = \left(\frac{\partial \mu _1}{\partial n _2}\right) _{n _1, n _3}$$



This equality can be rearranged to have differential quotient of mole fractions on one side.



\paragraph{Image resizing}



Partial derivatives are key to target-aware image resizing algorithms. Widely known as seam carving, these algorithms require each pixel in an image to be assigned a numerical 'energy' to describe their dissimilarity against orthogonal adjacent pixels. The algorithm then progressively removes rows or columns with the lowest energy. The formula established to determine a pixel's energy (magnitude of gradient at a pixel) depends heavily on the constructs of partial derivatives.



\paragraph{Economics}



Partial derivatives play a prominent role in economics, in which most functions describing economic behaviour posit that the behaviour depends on more than one variable. For example, a societal consumption function may describe the amount spent on consumer goods as depending on both income and wealth; the marginal propensity to consume is then the partial derivative of the consumption function with respect to income.



\paragraph{Licensing}



Content obtained and/or adapted from:



\begin{itemize}
    \item \href{https://en.wikipedia.org/wiki/Partial_derivative}{Partial derivative, Wikipedia} under a CC BY-SA license
\end{itemize}

\subsection{Learning Objectives}

\begin{enumerate}
    \item \textbf{Compute Partial Derivatives:}
    \item \textbf{Outcome:} Students will be able to compute the partial derivatives of functions with multiple variables.
    \item \textbf{Details:} This includes applying the definition of partial derivatives to various functions, such as polynomials, exponentials, and trigonometric functions, and interpreting the results.
    \item \textbf{Understand Notation and Symbolism:}
    \item \textbf{Outcome:} Students will understand and correctly use the different notations for partial derivatives, including $\frac{\partial}{\partial x}$, $f_x$, and $\partial_x f$.
    \item \textbf{Details:} This includes recognizing and translating between different notations used in textbooks and research papers, and understanding the context in which each notation is used.
    \item \textbf{Apply Partial Derivatives in Optimization Problems:}
    \item \textbf{Outcome:} Students will be able to apply partial derivatives to solve optimization problems involving functions of multiple variables.
    \item \textbf{Details:} This includes finding critical points by setting partial derivatives to zero, using the second partial derivative test to classify these points, and applying these techniques to practical problems in physics, economics, and engineering.
\end{enumerate}

\subsection{Assessment Instrument}

\begin{enumerate}
    \item \textbf{Assessment Instrument: Partial Derivatives Quiz} \textbf{Format:}
    \item \textbf{Type:} Mixed-format quiz (multiple-choice, short answer, and problem-solving questions)
    \item \textbf{Duration:} 45 minutes
    \item \textbf{Total Marks:} 30 points \textbf{Sections:}
    \item \textbf{Compute Partial Derivatives (10 points):}
    \item \textbf{Question 1:} Compute the partial derivative of the following functions with respect to the specified variable:
    \item a. $f(x, y) = x^2y + y^3$, with respect to $x$
    \item b. $g(x, y, z) = \sin(xy) + z^2x$, with respect to $z$
    \item c. $h(x, y) = e^{xy} + \cos(y)$, with respect to $y$
    \item d. $k(x, y, z) = x^2 + y^2 + z^2$, with respect to $x$
    \item \textbf{Understand Notation and Symbolism (10 points):}
    \item \textbf{Question 2:} Match the partial derivative notations to their equivalent expressions:
    \item a. $\frac{\partial f}{\partial x}$
    \item b. $f_x$
    \item c. $\partial_x f$
    \item d. $\frac{\partial^2 f}{\partial x \partial y}$
    \item e. $D_1 f$
    \item \textbf{Apply Partial Derivatives in Optimization Problems (10 points):}
    \item \textbf{Question 3:} Solve the following optimization problems using partial derivatives:
    \item a. Find the critical points of $f(x, y) = x^3 + y^3 - 3xy$.
    \item b. Determine whether the critical points found in part (a) are local minima, local maxima, or saddle points. \textbf{Scoring Rubric:}
    \item \textbf{Compute Partial Derivatives:} 2.5 points for each correctly computed partial derivative (total 10 points)
    \item \textbf{Understand Notation and Symbolism:} 2 points for each correctly matched notation (total 10 points)
    \item \textbf{Apply Partial Derivatives in Optimization Problems:} 5 points for correctly finding the critical points, and 5 points for correctly classifying them (total 10 points) The quiz assesses students' understanding and application of partial derivatives, ensuring they can compute derivatives, understand various notations, and solve optimization problems involving functions of multiple variables.
\end{enumerate}

%%===============================================================
\section{Directional Derivatives and Gradient Vectors}
\label{sec:directional-derivatives-and-gradient-vectors}
%%===============================================================

\paragraph{Directional derivatives}



Normally, a partial derivative of a function with respect to one of its variables, say, $x _j$, takes the derivative of that "slice" of that function parallel to the $x _j$th axis.



More precisely, we can think of cutting a function $\mathbf{f}(x _1, ..., x _n)$ in space along the $x _j$th axis, with keeping everything but the $x _j$ variable constant.



From the definition, we have the partial derivative at a point \textbf{\textit{p}} of the function along this slice as



$${\partial \mathbf{f} \over \partial x _j} = \lim _{t\rightarrow 0} {\mathbf{f}(\mathbf{p}+t\mathbf{e} _j) - \mathbf{f}(\mathbf{p}) \over t}$$



provided this limit exists.



Instead of the basis vector, which corresponds to taking the derivative along that axis, we can pick a vector in any direction (which we usually take as being a unit vector), and we take the \textit{directional derivative} of a function as



$${\partial \mathbf{f} \over \partial \mathbf{d}} = \lim _{t\rightarrow 0} {\mathbf{f}(\mathbf{p}+t\mathbf{d}) - \mathbf{f}(\mathbf{p}) \over t}$$

where \textbf{\textit{d}} is the direction vector.



If we want to calculate directional derivatives, calculating them from the limit definition is rather painful, but, we have the following: if $f : \mathbb{R} ^n \rightarrow \mathbb{R}$ is differentiable at a point $\mathbf{p}$, $|\mathbf{p}|=1$,



> ${\partial \mathbf{f} \over \partial \mathbf{d}} = D _\mathbf{p} \mathbf{f}(\mathbf{d})$



There is a closely related formulation which we'll look at in the next section.



\paragraph{Gradient vectors}



The partial derivatives of a scalar tell us how much it changes if we move along one of the axes. What if we move in a different direction?



We'll call the scalar $f$, and consider what happens if we move an infintesimal direction $\mathbf{dr}=(dx,dy,dz)$, using the chain rule.



$$\mathbf{df}=dx\frac{\partial f}{\partial x} + dy\frac{\partial f}{\partial y}+dz\frac{\partial f}{\partial z}$$



This is the dot product of $\mathbf{dr}$ with a vector whose components are the partial derivatives of $\mathbf{f}$, called the gradient of $\mathbf{f}$



$\operatorname{grad} \mathbf{f} = \nabla \mathbf{f} = \left(\frac{\partial \mathbf{f}(\mathbf{p})}{\partial x _1},\cdots, \frac{\partial \mathbf{f}(\mathbf{p})}{\partial x _n}\right)$



We can form directional derivatives at a point $\mathbf{p}$, in the direction $\mathbf{d}$ by taking the dot product of the gradient with $\mathbf{d}$.



$${\partial \mathbf{f}(\mathbf{p}) \over \partial \mathbf{d}} =\mathbf{d} \cdot \nabla \mathbf{f}(\mathbf{p})$$.



Notice that grad $\mathbf{f}$ looks like a vector multiplied by a scalar. This particular combination of partial derivatives is commonplace, so we abbreviate it to



$$\nabla = \left( \frac{\partial }{\partial x},  \frac{\partial }{\partial y}, \frac{\partial }{\partial z}\right)$$



We can write the action of taking the gradient vector by writing this as an \textit{operator}. Recall that in the one-variable case we can write $d/dx$ for the action of taking the derivative with respect to $x$. This case is similar, but $\nabla$ acts like a vector.



We can also write the action of taking the gradient vector as:



$$\nabla  = \left( \frac{\partial }{\partial x _1}, \frac{\partial }{\partial x _2}, \cdots \frac{\partial }{\partial x _n}\right)$$



\paragraph{Properties of the gradient vector}



\paragraph{Geometry}



\begin{itemize}
    \item Grad $\mathbf{f}(\mathbf{p})$ is a vector pointing in the direction of steepest slope of $\mathbf{f}$. \|grad $\mathbf{f}(\mathbf{p})$ \| is the rate of change of that slope at that point.
\end{itemize}

For example, we consider $h(x,y)=x^2+y^2$. The level sets of $h$ are concentric circles, centred on the origin, and



$$\nabla h = (h _x,h _y) = 2(x,y)= 2 \mathbf{r}$$ grad \textit{h} points directly away from the origin, at right angles to the contours.



\begin{itemize}
    \item Along a level set, $(\nabla \mathbf{f})(\mathbf{p})$ is perpendicular to the level set $\\{ \mathbf{x} | \mathbf{f} (\mathbf{x}) = \mathbf{f}(\mathbf{p}) \text{ at } \mathbf{x}=\mathbf{p} \\} $.
\end{itemize}

If $\mathbf{dr}$ points along the contours of $f$, where the function is constant, then $df$ will be zero. Since $df$ is a dot product, that means that the two vectors, $\mathbf{dr}$ and grad $f$, must be at right angles, i.e. the gradient is at right angles to the contours.



\paragraph{Algebraic properties}



Like $d/dx$, $\nabla$ is linear. For any pair of constants, $a$ and $b$, and any pair of scalar functions, $f$ and $g$



$$\frac{d}{dx} (af+bg)= a\frac{d}{dx}f + b\frac{d}{dx}g \quad \nabla (af+bg) = a \nabla f + b \nabla g$$



Since it's a vector, we can try taking its dot and cross product with other vectors, and with itself.



\paragraph{Resources}



\begin{itemize}
    \item \href{https://www.youtube.com/watch?v=CnVes9TdnPo}{How To Find The Directional Derivative and The Gradient Vector} Video by The Organic Chemistry Tutor 2019
    \item \href{https://en.wikibooks.org/wiki/Calculus/Multivariable_Calculus/Directional_Derivative_and_the_Gradient_Vector}{Directional Derivative and the Gradient Vector}, WikiBooks: Calculus/Multivariable Calculus
    \item \href{https://www.youtube.com/watch?v=OBELQIPH5xY}{Implicit Differentiation With Partial Derivatives Using The Implicit Function Theorem} Video by The Organic Chemistry Tutor 2019
\end{itemize}

\paragraph{Licensing}



Content obtained and/or adapted from:



\begin{itemize}
    \item \href{https://en.wikibooks.org/wiki/Calculus/Multivariable_Calculus/Directional_Derivative_and_the_Gradient_Vector}{Directional Derivative and the Gradient Vector, WikiBooks: Calculus/Multivariable Calculus} under a CC BY-SA license
\end{itemize}

\subsection{Learning Objectives}

\begin{enumerate}
    \item Directional Derivatives
    \item \textbf{Understand the Definition and Calculation of Directional Derivatives:}
    \item \textbf{Outcome:} Students will be able to understand the definition of directional derivatives and calculate them using the limit definition.
    \item \textbf{Details:} This includes understanding the limit definition of directional derivatives and being able to apply it to find the directional derivative of a given function in a specified direction.
    \item \textbf{Compute Directional Derivatives Using Gradient Vectors:}
    \item \textbf{Outcome:} Students will be able to compute directional derivatives using the gradient vector and the dot product.
    \item \textbf{Details:} This involves understanding the relationship between directional derivatives and gradient vectors, and being able to calculate the directional derivative by taking the dot product of the gradient vector and the direction vector.
    \item \textbf{Interpret the Geometric Meaning of Directional Derivatives:}
    \item \textbf{Outcome:} Students will be able to interpret the geometric meaning of directional derivatives and their relationship to the gradient vector.
    \item \textbf{Details:} This includes understanding how the directional derivative represents the rate of change of the function in a specified direction and how the gradient vector points in the direction of the steepest ascent of the function.
\end{enumerate}

\subsection{Assessment Instrument}

\begin{enumerate}
    \item \textbf{Assessment Instrument: Directional Derivatives and Gradient Vectors Quiz} \textbf{Format:}
    \item \textbf{Type:} Mixed-format quiz (multiple-choice, short answer, and problem-solving questions)
    \item \textbf{Duration:} 45 minutes
    \item \textbf{Total Marks:} 30 points \textbf{Sections:}
    \item \textbf{Understanding Directional Derivatives (10 points):}
    \item \textbf{Question 1:} Explain in your own words what a directional derivative is and how it differs from a partial derivative. (Short Answer)
    \item \textbf{Answer:} (3 points)
    \item \textbf{Question 2:} Given a function $f(x, y) = x^2 + y^2$, compute the directional derivative at the point (1, 1) in the direction of the vector $\mathbf{d} = (3, 4)$. (Problem-Solving)
    \item \textbf{Answer:} (7 points)
    \item \textbf{Computing Directional Derivatives Using Gradient Vectors (10 points):}
    \item \textbf{Question 3:} For the function $f(x, y, z) = xyz$, find the gradient vector $\nabla f$ at the point (1, -2, 3). (Short Answer)
    \item \textbf{Answer:} (5 points)
    \item \textbf{Question 4:} Using the gradient vector found in Question 3, compute the directional derivative of $f$ at the point (1, -2, 3) in the direction of the unit vector $\mathbf{u} = \frac{1}{\sqrt{14}}(1, 2, 3)$. (Problem-Solving)
    \item \textbf{Answer:} (5 points)
    \item \textbf{Geometric Interpretation of Directional Derivatives and Gradient Vectors (10 points):}
    \item \textbf{Question 5:} Describe the geometric meaning of the gradient vector at a point on the function surface. What does the magnitude and direction of the gradient vector signify? (Short Answer)
    \item \textbf{Answer:} (5 points)
    \item \textbf{Question 6:} Consider the function $h(x, y) = x^2 + y^2$. If you are standing at the point (3, 4) on the surface, in which direction should you move to increase $h$ most rapidly? Justify your answer using the concept of gradient vectors. (Problem-Solving)
    \item \textbf{Answer:} (5 points) \textbf{Scoring Rubric:}
    \item \textbf{Understanding Directional Derivatives:} 3 points for correct explanation; 7 points for correct calculation and interpretation.
    \item \textbf{Computing Directional Derivatives Using Gradient Vectors:} 5 points for correct gradient vector calculation; 5 points for correct directional derivative calculation using the gradient vector.
    \item \textbf{Geometric Interpretation of Directional Derivatives and Gradient Vectors:} 5 points for correct description and interpretation of the gradient vector's geometric meaning; 5 points for correctly identifying the direction of the steepest ascent and justification using the gradient vector. The quiz assesses students' understanding of directional derivatives, their ability to compute these derivatives using the gradient vector, and their comprehension of the geometric interpretation of gradient vectors. It ensures that students can correctly perform calculations and apply theoretical concepts to practical scenarios.
\end{enumerate}

%%===============================================================
\section{Tangent Plane}
\label{sec:tangent-plane}
%%===============================================================

\paragraph{Tangent Plane}



Intuitively, it seems clear that, in a plane, only one line can be tangent to a curve at a point. However, in three-dimensional space, many lines can be tangent to a given point. If these lines lie in the same plane, they determine the tangent plane at that point. A tangent plane at a regular point contains all of the lines tangent to that point. A more intuitive way to think of a tangent plane is to assume the surface is smooth at that point (no corners). Then, a tangent line to the surface at that point in any direction does not have any abrupt changes in slope because the direction changes smoothly.



\textbf{Definition of Tangent Plane}



% [Image: Tangent plane]



Figure: Tangent plane



> Let $P _0 = (x _0,y _0,z _0)$ be a point on a surface $S$, and let $C$ be any curve passing through $P _0$ and lying entirely in $S$. If the tangent lines to all such curves $C$ at $P _0$ lie in the same plane, then this plane is called the tangent plane to $S$ at $P _0$.



For a tangent plane to a surface to exist at a point on that surface, it

is sufficient for the function that defines the surface to be

differentiable at that point, defined later in this section. We define

the term tangent plane here and then explore the idea intuitively.



\textbf{Equation of Tangent Plane}



> Let $S$ be a surface defined by a differentiable function $z=f(x,y)$, and let $P _0 = (x _0,y _0)$ be a point in the domain of $f$. Then, the equation of the tangent plane to $S$ at $P _0$ is given by



> $z = f(x _0,y _0) + f _x(x _0,y _0)(x - x _0) + f _y(x _0,y _0)(y - y _0)$



To see why this formula is correct, let's first find two tangent lines to the surface S. The equation of the tangent line to the curve that is represented by the intersection of S with the vertical trace given by $x = x _0$ is $z=f(x _0,y _0)+f _y(x _0,y _0)(y-y _0)$. Similarly, the equation of the tangent line to the curve that is represented by the intersection of $S$ with the vertical trace given by $y=y _0$ is $z = f(x _0,y _0) + f _x(x _0,y _0)(x - x _0)$. A parallel vector to the first tangent line is $a = \mathbf{j} + f _y(x _0,y _0)\mathbf{k}$; a parallel vector to the second tangent line is $b = \mathbf{i} + f _x(x _0,y _0)\mathbf{k}$. We can take the cross product of these two vectors:



$$a\times b = (\mathbf{j} + f _y(x _0,y _0)\mathbf{k}) \times (\mathbf{i}+f _x(x _0,y _0)\mathbf{k})$$



> $= f _x(x _0,y _0)\mathbf{i}+f _y(x _0,y _0)\mathbf{j} - \mathbf{k}$



This vector is perpendicular to both lines and is therefore perpendicular to the tangent plane. We can use this vector as a normal vector to the tangent plane, along with the point $P _0=(x _0,y _0,f(x _0,y _0))$ in the equation for a plane:



> $n\cdot ((x-x _0)\mathbf{i}+(y-y _0)\mathbf{j}+(z-f(x _0,y _0))\mathbf{k}) = 0$



> $\implies (f _x(x _0,y _0)\mathbf{i}+f _y(x _0,y _0)\mathbf{j}-\mathbf{k})\cdot ((x-x _0)i+(y-y _0)\mathbf{j}+(z-f(x _0,y _0))\mathbf{k}) = 0$



> $\implies f _x(x _0,y _0)(x-x _0)+f _y(x _0,y _0)(y-y _0)-(z-f(x _0,y _0)) = 0$



> $\implies z = f(x _0,y _0) + f _x(x _0,y _0)(x - x _0) + f _y(x _0,y _0)(y - y _0)$



\paragraph{Linear Approximations}



When working with a function of two variables, the tangent line is

replaced by a tangent plane, but the approximation idea is much the

same.



\textbf{Definition of Linear Approximation} Given a function $z=f(x,y)$ with

continuous partial derivatives that exist at the point $(x _0,y _0)$, the

linear approximation of $f$ at the point $(x _0,y _0)$ is given by the

equation



> $L(x,y) = f(x _0,y _0) + f _x(x _0,y _0)(x - x _0) + f _y(x _0,y _0)(y-y _0)$



Notice that this equation also represents the tangent plane to the

surface defined by $z = f(x,y)$ at the point $(x _0,y _0)$. The idea

behind using a linear approximation is that, if there is a point

$(x _0,y _0)$ at which the precise value of $f(x,y)$ is known, then for

values of $(x,y)$ reasonably close to $(x _0,y _0)$, the linear

approximation (i.e., tangent plane) yields a value that is also

reasonably close to the exact value of $f(x,y)$. Furthermore the plane

that is used to find the linear approximation is also the tangent plane

to the surface at the point $(x _0,y _0)$.



\paragraph{Resources}



\begin{itemize}
    \item \href{https://openstax.org/books/calculus-volume-3/pages/4-4-tangent-planes-and-linear-approximations}{Tangent planes and linear approximations}, OpenStax
    \item \href{https://math.libretexts.org/Bookshelves/Calculus/Book\%3A_Vector_Calculus_(Corral}{Tangent Plane to a Surface}/02\%3A_Functions_of_Several_Variables/2.03\%3A_Tangent_Plane_to_a_Surface), Mathematics LibreTexts
\end{itemize}

\paragraph{Videos}



\begin{itemize}
    \item \href{https://www.youtube.com/watch?v=6paZkmBMZwQ}{Tangent planes \| MIT 18.02SC Multivariable Calculus, Fall 2010}
\end{itemize}

\begin{itemize}
    \item \href{https://www.youtube.com/watch?v=UfnpcOcZAO0}{Determining the Equation of a Tangent Plane} Video by Mathispower4u 2011
\end{itemize}

\begin{itemize}
    \item \href{https://www.youtube.com/watch?v=pxmW8_Cpd7U}{Ex 1: Find the Equation of a Tangent Plane to a Surface} video by Mathispower4u 2014
\end{itemize}

\begin{itemize}
    \item \href{https://www.youtube.com/watch?v=lEF2mmR3CWU}{How To Find The Equation of the Normal Line-The Organic Chemistry Tutor 2014} Video by The Organic Chemistry Tutor 2014
\end{itemize}

\begin{itemize}
    \item \href{https://www.youtube.com/watch?v=nDOVQeNHCJw}{quation of the normal line at a point (KristaKingMath)} Video by KristaKingMath 2012
\end{itemize}

\begin{itemize}
    \item \href{https://www.youtube.com/watch?v=e1Kp-fUIJCU}{Tangent Plane Approximations} Video by -patrickJMT
\end{itemize}

\paragraph{Licensing}



Content obtained and/or adapted from:



\begin{itemize}
    \item \href{https://openstax.org/books/calculus-volume-3/pages/4-4-tangent-planes-and-linear-approximations}{Tangent planes and linear approximations, OpenStax} under a CC BY-SA license
    \item \href{https://math.libretexts.org/Bookshelves/Calculus/Book\%3A_Vector_Calculus_(Corral}{Tangent Plane to a Surface, Mathematics LibreTexts}/02\%3A_Functions_of_Several_Variables/2.03\%3A_Tangent_Plane_to_a_Surface) under a CC BY-SA license
\end{itemize}

\subsection{Learning Objectives}

\begin{enumerate}
    \item Tangent Plane
    \item \textbf{Understand Tangent Planes:}
    \item \textbf{Outcome:} Students will be able to explain the concept of a tangent plane to a surface at a given point.
    \item \textbf{Details:} This includes describing how the tangent plane relates to the surface and the conditions under which it exists.
    \item \textbf{Derive the Equation of a Tangent Plane:}
    \item \textbf{Outcome:} Students will be able to derive the equation of a tangent plane to a surface at a given point using partial derivatives.
    \item \textbf{Details:} This includes using the formula $z = f(x_0, y_0) + f_x(x_0, y_0)(x - x_0) + f_y(x_0, y_0)(y - y_0)$ and understanding the geometric interpretation.
    \item \textbf{Apply Tangent Plane Concepts:}
    \item \textbf{Outcome:} Students will be able to apply the concept of the tangent plane to solve problems involving the approximation of functions near a given point.
    \item \textbf{Details:} This includes using the tangent plane to approximate values of a function and analyzing the accuracy of these approximations in various scenarios.
\end{enumerate}

\subsection{Assessment Instrument}

\begin{enumerate}
    \item \textbf{Assessment Instrument: Tangent Planes Quiz} \textbf{Format:}
    \item \textbf{Type:} Mixed-format quiz (multiple-choice, short answer, and problem-solving questions)
    \item \textbf{Duration:} 45 minutes
    \item \textbf{Total Marks:} 30 points \textbf{Sections:}
    \item \textbf{Conceptual Understanding of Tangent Planes (10 points):}
    \item \textbf{Question 1:} (Multiple Choice) What is the definition of a tangent plane to a surface at a given point?
    \item a. A plane that touches the surface at exactly one point and lies completely within the surface.
    \item b. A plane that intersects the surface at multiple points.
    \item c. A plane that touches the surface at exactly one point and contains all tangent lines to the surface at that point.
    \item d. A plane that lies parallel to the surface at a given point.
    \item \textbf{Derivation of Tangent Plane Equation (10 points):}
    \item \textbf{Question 2:} (Short Answer) Derive the equation of the tangent plane to the surface $z = f(x, y)$ at the point $(x_0, y_0)$ using partial derivatives. Show all steps.
    \item \textbf{Application of Tangent Plane (10 points):}
    \item \textbf{Question 3:} (Problem-Solving) Given the surface $z = x^2 + y^2$ and the point $(1, 1)$:
    \item a. Find the partial derivatives $f_x$ and $f_y$ at the point $(1, 1)$. (2 points)
    \item b. Write the equation of the tangent plane to the surface at the point $(1, 1)$. (4 points)
    \item c. Use the tangent plane equation to approximate the value of the surface at the point $(1.1, 1.1)$. (4 points) \textbf{Scoring Rubric:}
    \item \textbf{Conceptual Understanding of Tangent Planes:} 10 points for the correct choice with an explanation (total 10 points)
    \item \textbf{Derivation of Tangent Plane Equation:} 10 points for a complete and correct derivation (total 10 points)
    \item \textbf{Application of Tangent Plane:}
    \item a. 2 points for correct partial derivatives
    \item b. 4 points for the correct tangent plane equation
    \item c. 4 points for correct approximation using the tangent plane equation (total 10 points) This quiz will assess the students' understanding of the concept of tangent planes, their ability to derive the equation of a tangent plane using partial derivatives, and their skill in applying this concept to approximate values on a surface.
\end{enumerate}

%%===============================================================
\section{Differentiability}
\label{sec:differentiability}
%%===============================================================

\paragraph{Differentiable functions}



We will start from the one-variable definition of the derivative at a point $p$, namely



$$\lim _{x\rightarrow p} {f(x)-f(p) \over x-p} = f'(p)$$



Let's change above to equivalent form of



$$\lim _{x\rightarrow p} {f(x)-f(p)-f'(p)(x-p) \over x-p} = 0$$



which achieved after pulling f'(p) inside and putting it over a common denominator.



We can't divide by vectors, so this definition can't be immediately extended to the multiple variable case. Nonetheless, we don't have to: the thing we took interest in was the quotient of two small distances (magnitudes), not their other properties (like sign). It's worth noting that 'other' property of vector neglected is its direction. Now we can divide by the absolute value of a vector, so lets rewrite this definition in terms of absolute values



$$\lim _{x\rightarrow p} \frac{\left|f(x)-f(p)-f'(p)(x-p)\right|}{\left|  x-p\right|} = 0$$



Another form of formula above is obtained by letting $h=x-p$ we have $x=p+h$ and if $x \rightarrow p$, the $h = x - p \rightarrow 0$, so



$$\lim _{h\rightarrow 0} \frac{\left|f(p+h)-f(p)-f'(p)h\right|}{\left| h\right|} = 0,$$



where $h$ can be thought of as a 'small change'.



So, how can we use this for the several-variable case?



If we switch all the variables over to vectors and replace the constant (which performs a linear map in one dimension) with a matrix (which denotes also a linear map), we have



$$\lim _{\mathbf{x}\rightarrow\mathbf{p}} {|\mathbf{f}(\mathbf{x})-\mathbf{f}(\mathbf{p})-\mathbf{A}(\mathbf{x}-\mathbf{p})| \over |\mathbf{x}-\mathbf{p}|} = 0$$



or



$$\lim _{\mathbf{h}\rightarrow\mathbf{0}} {|\mathbf{f}(\mathbf{p}+\mathbf{h})-\mathbf{f}(\mathbf{p})-\mathbf{A}\mathbf{h}| \over |\mathbf{h}|} = 0$$



If this limit exists for some $f: \mathbb{R} ^m \rightarrow \mathbb{R} ^n$, and there is a linear map $A: \mathbb{R} ^m \rightarrow \mathbb{R} ^n$ (denoted by matrix $A$ which is $m \times n$), we refer to this map as being the derivative and we write it as $D _{\mathbf{p}} \mathbf{f}$.



A point on terminology, in referring to the action of taking the derivative (giving the linear map $A$), we write $D \_{\mathbf{p}} \mathbf{f}$, but in referring to the matrix $A$ itself, it is known as the \textit{Jacobian matrix} and is also written $J_p f$. More on the Jacobian later.



=== Properties ===

There are a number of important properties of this formulation of the derivative.



==== Affine approximations ====

If $\mathbf{f}$ is differentiable at $\mathbf{p}$ for $\mathbf{x}$ close to $\mathbf{p}$, $|\mathbf{f}(\mathbf{x})-(\mathbf{f}(\mathbf{p})+\mathbf{A}(\mathbf{x}-\mathbf{p}))|$ is small compared to $|\mathbf{x}-\mathbf{p}|$, which means that $\mathbf{f}(\mathbf{x})$ is approximately equal to $\mathbf{f}(\mathbf{p})+\mathbf{A}(\mathbf{x}-\mathbf{p})$.



We call an expression of the form $\mathbf{g}(\mathbf{x})+c$ affine, when $\mathbf{g}(\mathbf{x})$ is linear and $c$ is a constant. $\mathbf{f}(\mathbf{p})+\mathbf{A}(\mathbf{x}-\mathbf{p})$ is an affine approximation to $\mathbf{f}(\mathbf{x})$.



\paragraph{Jacobian matrix and partial derivatives}



The Jacobian matrix of a function is in the form



$$(J _{\mathbf{p}} \mathbf{f}) _{ij} = \frac{\partial f _i}{\partial x _j} \Bigg| _{\mathbf{p}}$$



for a $\mathbf{f} : \mathbb{R}^m \rightarrow \mathbb{R}^n$, $J _{\mathbf{p}} \mathbf{f}$ is a $n \times m$ matrix.



The consequence of this is that if $\mathbf{f}$ is differentiable at $\mathbf{p}$, \textit{all} the partial derivatives of $\mathbf{f}$ exist at $\mathbf{p}$.



However, it is possible that all the partial derivatives of a function exist at some point yet that function is not differentiable there, so it is very important not to mix derivative (linear map) with the Jacobian (matrix) especially in situations akin to the one cited.



\textbf{Rules of taking Jacobians}



If $\mathbf{f} : \mathbb{R}^m \rightarrow \mathbb{R}^n$, and $h(x) : \mathbb{R}^m \rightarrow \mathbb{R}$ are differentiable at $\mathbf{p}$:





\begin{itemize}
    \item $J _{\mathbf{p}} (\mathbf{f} + \mathbf{g}) = J _{\mathbf{p}} \mathbf{f} + J _{\mathbf{p}} \mathbf{g}$
    \item $J _{\mathbf{p}} (h \mathbf{f}) = h J _{\mathbf{p}} \mathbf{f} + \mathbf{f}(\mathbf{p}) J _{\mathbf{p}} h$
    \item $J _{\mathbf{p}} (\mathbf{f} \cdot \mathbf{g}) = \mathbf{g}^T J _{\mathbf{p}} \mathbf{f} + \mathbf{f}^T J _{\mathbf{p}} \mathbf{g}$
\end{itemize}



\textit{Important: make sure the order is right - matrix multiplication is not commutative!}



\paragraph{Chain rule}



The chain rule for functions of several variables is as follows. For $\mathbf{f} : \mathbb{R}^ m \rightarrow \mathbb{R}^ n$ and $\mathbf{g} : \mathbb{R}^ n \rightarrow \mathbb{R}^ p$, and $\mathbf{g} \circ \mathbf{f}$ differentiable at $\mathbf{p}$, then the Jacobian is given by



$$(J _{\mathbf{f}(\mathbf{p})} \mathbf{g}) (J _{\mathbf{p}} \mathbf{f})$$



Again, we have matrix multiplication, so one must preserve this exact order. Compositions in one order may be defined, but not necessarily in the other way.



\paragraph{Continuity and differentiability}



Furthermore, if all the partial derivatives exist, and are continuous in some neighbourhood of a point $\mathbf{p}$, then $\mathbf{f}$ is differentiable at $\mathbf{p}$. This has the consequence that for a function $\mathbf{f}$ which has its component functions built from continuous functions (such as rational functions, differentiable functions or otherwise), $\mathbf{f}$ is differentiable everywhere $\mathbf{f}$ is defined.



We use the terminology \textit{continuously differentiable} for a function differentiable at $\mathbf{p}$ which has all its partial derivatives existing and are continuous in some neighbourhood at $\mathbf{p}$.



\paragraph{Resources}



\begin{itemize}
    \item \href{https://en.wikibooks.org/wiki/Calculus/Multivariable_Calculus/Partial_Derivatives}{Partial Derivatives}, WikiBooks: Calculus/Multivariable Calculus
    \item \href{https://en.wikibooks.org/wiki/Calculus/Multivariable_Calculus/Chain_Rule}{Chain Rule}, WikiBooks: Calculus/Multivariable Calculus
    \item \href{https://math.libretexts.org/Courses/University_of_California_Davis/UCD_Mat_21C\%3A_Multivariate_Calculus/13\%3A_Partial_Derivatives/13.6\%3A_Tangent_Planes_and_Differentials#:\textasciitilde{}:text=1\%3A\%20The\%20tangent\%20plane\%20to,be\%20differentiable\%20at\%20that\%20point.}{Tangent Planes and Differentials}, Mathematics LibreTexts
\end{itemize}

\paragraph{Licensing}



Content obtained and/or adapted from:



\begin{itemize}
    \item \href{https://en.wikibooks.org/wiki/Calculus/Multivariable_Calculus/Partial_Derivatives}{Partial Derivatives, WikiBooks: Calculus/Multivariable Calculus} under a CC BY-SA license
    \item \href{https://en.wikibooks.org/wiki/Calculus/Multivariable_Calculus/Chain_Rule}{Chain Rule, WikiBooks: Calculus/Multivariable Calculus} under a CC BY-SA license
\end{itemize}

\subsection{Learning Objectives}

\begin{enumerate}
    \item Certainly! Here are three Student Learning Objectives (SLOs) for the topic of differentiable functions, based on the provided format. ---
    \item \textbf{Understanding Differentiability in One Variable:}
    \item \textbf{Outcome:} Students will be able to define and explain the concept of a derivative at a point for functions of one variable.
    \item \textbf{Details:} This includes understanding and applying the limit definition of the derivative at a point \( p \): $$\lim _{x\rightarrow p} \frac{f(x)-f(p)}{x-p} = f'(p)$$ and its equivalent form: $$\lim _{x\rightarrow p} \frac{f(x)-f(p)-f'(p)(x-p)}{x-p} = 0.$$
    \item \textbf{Differentiability in Multiple Variables:}
    \item \textbf{Outcome:} Students will be able to extend the concept of differentiability to functions of several variables and interpret the definition involving limits.
    \item \textbf{Details:} This involves understanding the limit definition for differentiability in multiple variables: $$\lim _{\mathbf{x}\rightarrow\mathbf{p}} \frac{|\mathbf{f}(\mathbf{x})-\mathbf{f}(\mathbf{p})-\mathbf{A}(\mathbf{x}-\mathbf{p})|}{|\mathbf{x}-\mathbf{p}|} = 0,$$ where \( \mathbf{A} \) is a linear map (Jacobian matrix).
    \item \textbf{Jacobian Matrix and Partial Derivatives:}
    \item \textbf{Outcome:} Students will be able to compute the Jacobian matrix of a multivariable function and understand its significance.
    \item \textbf{Details:} This includes recognizing the relationship between partial derivatives and the Jacobian matrix, expressed as: $$(J _{\mathbf{p}} \mathbf{f})_{ij} = \frac{\partial f_i}{\partial x_j} \Bigg|_{\mathbf{p}}.$$ Students will learn to apply this to check differentiability by ensuring the existence and continuity of all partial derivatives. --- These SLOs ensure that students grasp the foundational concepts of differentiable functions, both in one variable and in multiple variables, along with the practical skills needed to compute and apply the Jacobian matrix.
\end{enumerate}

\subsection{Assessment Instrument}

\begin{enumerate}
    \item \textbf{Assessment Instrument: Differentiability Quiz} \textbf{Format:}
    \item \textbf{Type:} Mixed-format quiz (multiple-choice, short answer, and problem-solving questions)
    \item \textbf{Duration:} 45 minutes
    \item \textbf{Total Marks:} 30 points \textbf{Sections:}
    \item \textbf{Understanding Differentiability in One Variable (10 points):}
    \item \textbf{Question 1:} Define the derivative of a function at a point $p$ using the limit definition. (2 points)
    \item \textbf{Question 2:} Explain the equivalent form of the limit definition for the derivative at a point $p$. (3 points)
    \item \textbf{Question 3:} For the function $f(x) = x^2$, calculate the derivative at the point $p = 3$ using the limit definition. Show your work. (5 points)
    \item \textbf{Differentiability in Multiple Variables (10 points):}
    \item \textbf{Question 4:} Describe the limit definition of differentiability for a function of several variables. (2 points)
    \item \textbf{Question 5:} Explain the role of the linear map $\mathbf{A}$ in the definition of differentiability for multiple variables. (3 points)
    \item \textbf{Question 6:} Given $\mathbf{f}(\mathbf{x}) = (x_1^2 + x_2, x_1 - x_2^2)$, check if $\mathbf{f}$ is differentiable at $\mathbf{p} = (1, 1)$ by verifying the limit definition. (5 points)
    \item \textbf{Jacobian Matrix and Partial Derivatives (10 points):}
    \item \textbf{Question 7:} Define the Jacobian matrix of a function $\mathbf{f} : \mathbb{R}^m \rightarrow \mathbb{R}^n$. (2 points)
    \item \textbf{Question 8:} Explain the relationship between partial derivatives and the Jacobian matrix. (3 points)
    \item \textbf{Question 9:} Compute the Jacobian matrix for the function $\mathbf{f}(\mathbf{x}) = (x_1^2 + x_2, x_1 - x_2^2)$ at the point $\mathbf{p} = (1, 1)$. (5 points) \textbf{Scoring Rubric:}
    \item \textbf{Understanding Differentiability in One Variable:} 2 points for defining the derivative, 3 points for explaining the equivalent form, 5 points for correct calculation and work shown.
    \item \textbf{Differentiability in Multiple Variables:} 2 points for describing the limit definition, 3 points for explaining the role of $\mathbf{A}$, 5 points for verifying differentiability using the limit definition.
    \item \textbf{Jacobian Matrix and Partial Derivatives:} 2 points for defining the Jacobian matrix, 3 points for explaining the relationship with partial derivatives, 5 points for correct computation of the Jacobian matrix. This quiz will help assess the students' understanding and application of differentiability concepts, including the limit definitions, role of linear maps, and computation of the Jacobian matrix. It ensures that they can correctly define, explain, and compute the necessary elements involved in differentiability for both one-variable and multi-variable functions.
\end{enumerate}

%%===============================================================
\section{The Chain Rule}
\label{sec:the-chain-rule}
%%===============================================================

The generalization of the chain rule to multi-variable functions is rather technical. However, it is simpler to write in the case of functions of the form



$$f(g _1(x), \dots, g _k(x)).$$



As this case occurs often in the study of functions of a single variable, it is worth describing it separately.



\paragraph{Case of $f(g _1(x), \dots, g _k(x))$}



For writing the chain rule for a function of the form



$f(g _1(x), \dots, g _k(x))$



one needs the partial derivatives of $f$ with respect to its $k$ arguments. The usual notations for partial derivatives involve names for the arguments of the function. As these arguments are not named in the above formula, it is simpler and clearer to denote by $D _{i} f$, the derivative of $f$ with respect to its $i$th argument, and by $D _i f(z)$, the value of this derivative at $z$.



With this notation, the chain rule is



$$\frac{d}{dx}f(g _1(x), \dots, g _k (x))=\sum _{i=1}^k  \left(\frac{d}{dx}{g _i}(x)\right) D _i f(g _1(x), \dots, g _k (x)).$$



\paragraph{Example: arithmetic operations}



If the function $f$ is addition, that is, if



$$f(u,v)=u+v,$$ then $D _1 f = \frac{\partial f}{\partial u} = 1$ and $D _2 f = \frac{\partial f}{\partial v} = 1$. Thus, the chain rule gives



$$\frac{d}{dx}(g(x)+h(x)) = \left( \frac{d}{dx}g(x) \right) D _1 f+\left( \frac{d}{dx}h(x)\right) D _2 f=\frac{d}{dx}g(x) +\frac{d}{dx}h(x).$$



For multiplication



$$f(u,v)=uv,$$ the partials are $D _1 f = v$ and $D _2 f = u$. Thus,



$$\frac{d}{dx}(g(x)h(x)) = h(x) \frac{d}{dx} g(x) + g(x) \frac{d}{dx} h(x).$$



The case of exponentiation



$$f(u,v)=u^v$$ is slightly more complicated, as



$$D _1 f = vu^{v-1},$$ and, as $u^v=e^{v\ln u},$



$$D _2 f = u^v\ln u.$$ It follows that



$$\frac{d}{dx}\left(g(x)^{h(x)}\right) = h(x)g(x)^{h(x)-1} \frac{d}{dx}g(x) + g(x)^{h(x)} \ln g(x) \frac{d}{dx}h(x).$$



\paragraph{General rule}



The simplest way for writing the chain rule in the general case is to use the total derivative, which is a linear transformation that captures all directional derivatives in a single formula. Consider differentiable functions $f : \mathbb{R} ^m \to \mathbb{R} ^K$ and $g : \mathbb{R} ^n \to \mathbb{R} ^m$, and a point $a$ in $\mathbb{R} ^n$. Let $D _{a} g$ denote the total derivative of $g$ at $a$ and $D _{g} (a) f$ denote the total derivative of $f$ at $g(a)$. These two derivatives are linear transformations $\mathbb{R} ^n \to \mathbb{R} ^m$ and $\mathbb{R} ^m \to \mathbb{R} ^K$, respectively, so they can be composed. The chain rule for total derivatives is that their composite is the total derivative of $f \circ g$ at $a$:



$$D _a (f \circ g) = D _{g(a)} f \circ D _a g,$$



or for short,



$$D(f \circ g) = Df \circ Dg.$$



The higher-dimensional chain rule can be proved using a technique similar to the second proof given above.



Because the total derivative is a linear transformation, the functions appearing in the formula can be rewritten as matrices. The matrix corresponding to a total derivative is called a Jacobian matrix, and the composite of two derivatives corresponds to the product of their Jacobian matrices. From this perspective the chain rule therefore says:



$$J _{f \circ g}(\mathbf{a}) = J _{f}(g(\mathbf{a})) J _{g}(\mathbf{a}),$$ or for short,



$$J _{f \circ g} = (J _f \circ g)J _g.$$



That is, the Jacobian of a composite function is the product of the Jacobians of the composed functions (evaluated at the appropriate points).



The higher-dimensional chain rule is a generalization of the one-dimensional chain rule. If $k, m,$ and $n$ are 1, so that $f : \mathbb{R} \to \mathbb{R}$ and $g : \mathbb{R} \to \mathbb{R}$, then the Jacobian matrices of $f$ and $g$ are $1 \times 1$. Specifically, they are:



$$\begin{aligned}

J _g(a) & = (g'(a)), \\\J _f(g(a)) & = (f'(g(a))).

\end{aligned}$$



The Jacobian of $f \circ g$ is the product of these $1 \times 1$ matrices, so it is $f'(g(a)) \cdot g'(a)$, as expected from the one-dimensional chain rule. In the language of linear transformations, $D _{a}(g)$ is the function which scales a vector by a factor of $g'(a)$ and $D _{g(a)}(f)$ is the function which scales a vector by a factor of $f'(g(a))$. The chain rule says that the composite of these two linear transformations is the linear transformation $D _a(f \circ g)$, and therefore it is the function that scales a vector by $f'(g(a)) \cdot g'(a)$.



Another way of writing the chain rule is used when $f$ and $g$ are expressed in terms of their components as $\mathbf{y} = f(\mathbf{u}) = (f _1(\mathbf{u}), \ldots, f _k(\mathbf{u}))$ and $\mathbf{u} = g(\mathbf{x}) = (g _1(\mathbf{x}), \ldots, g _m(\mathbf{x}))$. In this case, the above rule for Jacobian matrices is usually written as:



$$\frac{\partial(y _1, \ldots, y _k)}{\partial(x _1, \ldots, x _n)} = \frac{\partial(y _1, \ldots, y _k)}{\partial(u _1, \ldots, u _m)} \frac{\partial(u _1, \ldots, u _m)}{\partial(x _1, \ldots, x _n)}.$$



The chain rule for total derivatives implies a chain rule for partial derivatives. Recall that when the total derivative exists, the partial derivative in the $i$th coordinate direction is found by multiplying the Jacobian matrix by the $i$th basis vector. By doing this to the formula above, we find:



$$\frac{\partial(y _1, \ldots, y _k)}{\partial x _i} = \frac{\partial(y _1, \ldots, y _k)}{\partial(u _1, \ldots, u _m)} \frac{\partial(u _1, \ldots, u _m)}{\partial x _i}.$$



Since the entries of the Jacobian matrix are partial derivatives, we may simplify the above formula to get:



$$\frac{\partial(y _1, \ldots, y _k)}{\partial x _i} = \sum _{\ell = 1}^m \frac{\partial(y _1, \ldots, y _k)}{\partial u _\ell} \frac{\partial u _\ell}{\partial x _i}.$$ More conceptually, this rule expresses the fact that a change in the $x _i$ direction may change all of $g _1$ through $g _m$, and any of these changes may affect \textit{f}.



In the special case where $k=1$, so that \textit{f} is a real-valued function, then this formula simplifies even further:



$$\frac{\partial y}{\partial x _i} = \sum _{\ell = 1}^m \frac{\partial y}{\partial u _\ell} \frac{\partial u _\ell}{\partial x _i}.$$



This can be rewritten as a dot product. Recalling that $\mathbf{u} = (g _1, \ldots, g _m)$, the partial derivative $\frac{\partial \mathbf{u}}{\partial x _i}$ is also a vector, and the chain rule says that:



$$\frac{\partial y}{\partial x _i} = \nabla y \cdot \frac{\partial \mathbf{u}}{\partial x _i}.$$



\paragraph{Example}



Given $u(x, y) = x^2 + 2y$ where $x(r, t) = r \sin(t)$ and $y(r, t) = \sin^2(t)$, determine the value of $\frac{\partial u}{\partial r}$ and $\frac{\partial u}{\partial t}$ using the chain rule.



$$\frac{\partial u}{\partial r}=\frac{\partial u}{\partial x} \frac{\partial x}{\partial r}+\frac{\partial u}{\partial y} \frac{\partial y}{\partial r} = (2x)(\sin(t)) + (2)(0) = 2r \sin^2(t),$$ and



$$ \begin{aligned}

\frac{\partial u}{\partial t} &= \frac{\partial u}{\partial x} \frac{\partial x}{\partial t} + \frac{\partial u}{\partial y} \frac{\partial y}{\partial t} \\\&= (2x)(r \cos(t)) + (2)(2 \sin(t) \cos(t)) \\\&= (2r \sin(t))(r \cos(t)) + 4 \sin(t) \cos(t) \\\&= 2(r^2 \sin(t) \cos(t)) + 4 \sin(t) \cos(t) \\\&= 2(r^2 + 2) \sin(t) \cos(t) \\\&= (r^2 + 2) \sin(2t).

\end{aligned} $$



\paragraph{Higher derivatives of multivariable functions}



Faà di Bruno's formula for higher-order derivatives of single-variable functions generalizes to the multivariable case. If $y = f(\mathbf{u})$ is a function of $\mathbf{u} = g(\mathbf{x})$ as above, then the second derivative of $f \circ g$ is:



$$

\frac{\partial^2 y}{\partial x _i \partial x _j} = \sum _k \left( \frac{\partial y}{\partial u _k} \frac{\partial^2 u _k}{\partial x _i \partial x _j} \right) + \sum _{k, \ell} \left( \frac{\partial^2 y}{\partial u _k \partial u _\ell} \frac{\partial u _k}{\partial x _i} \frac{\partial u _\ell}{\partial x _j} \right).

$$



\paragraph{Resources}



\begin{itemize}
    \item \href{https://en.wikibooks.org/wiki/Calculus/Multivariable_Calculus/Chain_Rule}{Chain Rule}, WikiBooks: Multivariable Calculus
    \item \href{https://math.hmc.edu/calculus/hmc-mathematics-calculus-online-tutorials/multivariable-calculus/multi-variable-chain-rule/}{Multivariable Chain Rule}, Harvey Mudd College
    \item \href{https://www.youtube.com/watch?v=XipB_uEexF0}{Chain Rule With Partial Derivatives - Multivariable Calculus} Video by The Organic Chemistry Tutor 2019
\end{itemize}

\paragraph{Licensing}



Content obtained and/or adapted from:



\begin{itemize}
    \item \href{https://en.wikipedia.org/wiki/Chain_rule}{Chain rule, Wikipedia} under a CC BY-SA license
\end{itemize}

\subsection{Learning Objectives}

\begin{enumerate}
    \item \textbf{Understand the Multi-Variable Chain Rule:}
    \item \textbf{Outcome:} Students will be able to apply the chain rule to multi-variable functions of the form \( f(g_1(x), \dots, g_k(x)) \).
    \item \textbf{Details:} This includes identifying and calculating the partial derivatives \( D_i f \) and using them to find the derivative of the composite function. For example, given \( f(u, v) = uv \), students should be able to derive the chain rule result \( \frac{d}{dx}(g(x)h(x)) = h(x) \frac{d}{dx} g(x) + g(x) \frac{d}{dx} h(x) \).
    \item \textbf{Compute Partial Derivatives in Composite Functions:}
    \item \textbf{Outcome:} Students will be able to compute partial derivatives for composite functions and apply these in the chain rule.
    \item \textbf{Details:} This involves calculating \( \frac{d}{dx} f(g_1(x), \dots, g_k(x)) \) using the sum of the derivatives of each \( g_i(x) \) multiplied by the respective partial derivative of \( f \). For instance, students should be able to derive that for \( f(u,v) = u+v \), the chain rule gives \( \frac{d}{dx}(g(x) + h(x)) = \frac{d}{dx} g(x) + \frac{d}{dx} h(x) \).
    \item \textbf{Apply the Chain Rule to Higher Derivatives:}
    \item \textbf{Outcome:} Students will be able to generalize the chain rule to higher-order derivatives in multi-variable functions.
    \item \textbf{Details:} This includes understanding and applying Faà di Bruno's formula for higher-order derivatives of composite functions. For example, students should be able to derive that the second derivative of \( f \circ g \) is given by \( \frac{\partial^2 y}{\partial x_i \partial x_j} = \sum_k \left( \frac{\partial y}{\partial u_k} \frac{\partial^2 u_k}{\partial x_i \partial x_j} \right) + \sum_{k, \ell} \left( \frac{\partial^2 y}{\partial u_k \partial u_\ell} \frac{\partial u_k}{\partial x_i} \frac{\partial u_\ell}{\partial x_j} \right) \).
\end{enumerate}

\subsection{Assessment Instrument}

\begin{enumerate}
    \item \textbf{Assessment Instrument: Chain Rule for Functions of More than One Variable Quiz} \textbf{Format:}
    \item \textbf{Type:} Mixed-format quiz (multiple-choice, short answer, and problem-solving questions)
    \item \textbf{Duration:} 45 minutes
    \item \textbf{Total Marks:} 30 points \textbf{Sections:}
    \item \textbf{Understanding the Multi-Variable Chain Rule (10 points):}
    \item \textbf{Question 1 (Multiple Choice):} Which of the following is the correct expression of the chain rule for the function $f(g_1(x), \dots, g_k(x))$?
    \item a) $\frac{d}{dx} f(g_1(x), \dots, g_k(x)) = \sum_{i=1}^k \left(\frac{d}{dx} g_i(x)\right) D_i f$
    \item b) $\frac{d}{dx} f(g_1(x), \dots, g_k(x)) = \sum_{i=1}^k g_i(x) \frac{d}{dx} D_i f$
    \item c) $\frac{d}{dx} f(g_1(x), \dots, g_k(x)) = \prod_{i=1}^k \left(\frac{d}{dx} g_i(x)\right) D_i f$
    \item d) $\frac{d}{dx} f(g_1(x), \dots, g_k(x)) = \sum_{i=1}^k g_i(x) D_i f$
    \item \textbf{Question 2 (Short Answer):} Given $f(u, v) = u + v$ and $g(x) = x^2, h(x) = x^3$, use the chain rule to find $\frac{d}{dx} f(g(x), h(x))$.
    \item \textbf{Question 3 (Problem-Solving):} For $f(u, v) = u^v$ and $g(x) = x^2, h(x) = \ln(x)$, compute $\frac{d}{dx} f(g(x), h(x))$.
    \item \textbf{Computing Partial Derivatives in Composite Functions (10 points):}
    \item \textbf{Question 4 (Multiple Choice):} Which of the following represents the partial derivative $D_1 f$ for $f(u, v) = u^v$?
    \item a) $v u^{v-1}$
    \item b) $u^v \ln u$
    \item c) $v u^{v+1}$
    \item d) $u^v$
    \item \textbf{Question 5 (Short Answer):} Calculate the partial derivatives $D_1 f$ and $D_2 f$ for $f(u, v) = u \cdot v$.
    \item \textbf{Question 6 (Problem-Solving):} Using $f(u, v) = u v$, and $g(x) = \sin(x), h(x) = \cos(x)$, compute $\frac{d}{dx} f(g(x), h(x))$.
    \item \textbf{Applying the Chain Rule to Higher Derivatives (10 points):}
    \item \textbf{Question 7 (Multiple Choice):} What is the general formula for the second derivative $\frac{\partial^2 y}{\partial x_i \partial x_j}$ of a composite function?
    \item a) $\sum_k \left( \frac{\partial y}{\partial u_k} \frac{\partial^2 u_k}{\partial x_i \partial x_j} \right)$
    \item b) $\sum_k \left( \frac{\partial^2 y}{\partial u_k \partial u_\ell} \frac{\partial u_k}{\partial x_i} \frac{\partial u_\ell}{\partial x_j} \right)$
    \item c) $\sum_k \left( \frac{\partial y}{\partial u_k} \frac{\partial^2 u_k}{\partial x_i \partial x_j} \right) + \sum_{k, \ell} \left( \frac{\partial^2 y}{\partial u_k \partial u_\ell} \frac{\partial u_k}{\partial x_i} \frac{\partial u_\ell}{\partial x_j} \right)$
    \item d) $\sum_{k, \ell} \left( \frac{\partial^2 y}{\partial u_k \partial u_\ell} \frac{\partial u_k}{\partial x_i} \frac{\partial u_\ell}{\partial x_j} \right)$
    \item \textbf{Question 8 (Short Answer):} Given $y = f(u, v)$ where $u = x^2 + y^2$, $v = x - y$, find $\frac{\partial^2 y}{\partial x \partial y}$ using the chain rule.
    \item \textbf{Question 9 (Problem-Solving):} If $u(x, y) = x^2 + 2y$ and $x(r, t) = r \sin(t)$, $y(r, t) = \sin^2(t)$, use the chain rule to determine $\frac{\partial^2 u}{\partial r^2}$ and $\frac{\partial^2 u}{\partial t^2}$. \textbf{Scoring Rubric:}
    \item \textbf{Understanding the Multi-Variable Chain Rule:} 3 points for correct multiple-choice answers, 3.5 points for correct short answers, 3.5 points for correct problem-solving (total 10 points)
    \item \textbf{Computing Partial Derivatives in Composite Functions:} 3 points for correct multiple-choice answers, 3.5 points for correct short answers, 3.5 points for correct problem-solving (total 10 points)
    \item \textbf{Applying the Chain Rule to Higher Derivatives:} 3 points for correct multiple-choice answers, 3.5 points for correct short answers, 3.5 points for correct problem-solving (total 10 points) This quiz assesses students' understanding and application of the chain rule in multi-variable functions, including the calculation of partial derivatives and higher-order derivatives.
\end{enumerate}

%%===============================================================
\section{Maxima and minima problems}
\label{sec:maxima-and-minima-problems}
%%===============================================================

% [Image: The four types of extrema.]



 Figure: The four types of extrema.



 \textbf{Maxima} and \textbf{minima} are points where a function reaches a highest or lowest value, respectively. There are two kinds of \textbf{extrema} (a word meaning maximum \textit{or} minimum): \textbf{global} and \textbf{local}, sometimes referred to as "absolute" and "relative", respectively. A global maximum is a point that takes the largest value on the entire range of the function, while a global minimum is the point that takes the smallest value on the range of the function. On the other hand, local extrema are the largest or smallest values of the function in the immediate vicinity.



In many cases, extrema look like the crest of a hill or the bottom of a bowl on a graph of the function. A global extremum is always a local extremum too, because it is the largest or smallest value on the entire range of the function, and therefore also its vicinity. It is also possible to have a function with no extrema, global or local: $y=x$ is a simple example.



At any extremum, the slope of the graph is necessarily 0 (or is undefined, as in the case of $-|x|$), as the graph must stop rising or falling at an extremum, and begin to head in the opposite direction. Because of this, extrema are also commonly called \textbf{stationary points} or \textbf{turning points}. Therefore, the first derivative of a function is equal to 0 at extrema. If the graph has one or more of these stationary points, these may be found by setting the first derivative equal to 0 and finding the roots of the resulting equation.



% [Image: The function $f(x)=x^3$ , which contains a saddle point at the point $(0,0)$.]



Figure: The function $f(x)=x^3$ , which contains a saddle point at the point $(0,0)$.



However, a slope of zero does not guarantee a maximum or minimum: there is a third class of stationary point called a saddle point. Consider the function



$$f(x)=x^3$$ The derivative is



$$f'(x)=3x^2$$ The slope at $x=0$ is 0. We have a slope of 0, but while this makes it a stationary point, this doesn't mean that it is a maximum or minimum. Looking at the graph of the function you will see that $x=0$ is neither, it's just a spot at which the function flattens out. True extrema require a sign change in the first derivative. This makes sense - you have to rise (positive slope) to and fall (negative slope) from a maximum. In between rising and falling, on a smooth curve, there will be a point of zero slope - the maximum. A minimum would exhibit similar properties, just in reverse.



% [Image: Good and bad points to check to classify extremum]



Figure: Good (\textit{B} and \textit{C}, green) and bad (\textit{D} and \textit{E}, blue) points to check in order to classify the extremum ($a$, black). The bad points lead to an incorrect classification of $a$ as a minimum.



This leads to a simple method to classify a stationary point - plug x values slightly left and right into the derivative of the function. If the results have opposite signs then it is a true maximum/minimum. You can also use these slopes to figure out if it is a maximum or a minimum: the left side slope will be positive for a maximum and negative for a minimum. However, you must exercise caution with this method, as, if you pick a point too far from the extremum, you could take it on the far side of \textit{another} extremum and incorrectly classify the point.



\paragraph{The Extremum Test}



A more rigorous method to classify a stationary point is called the \textbf{extremum test}, or \textbf{2nd Derivative Test}. As we mentioned before, the sign of the first derivative must change for a stationary point to be a true extremum. Now, the \textit{second} derivative of the function tells us the rate of change of the first derivative. It therefore follows that if the second derivative is positive at the stationary point, then the gradient is increasing. The fact that it is a stationary point in the first place means that this can only be a minimum. Conversely, if the second derivative is negative at that point, then it is a maximum.



Now, if the second derivative is 0, we have a problem. It could be a point of inflexion, or it could still be an extremum. Examples of each of these cases are below - all have a second derivative equal to 0 at the stationary point in question:



\begin{itemize}
    \item $y=x^3$ has a point of inflexion at $x=0$
    \item $y=x^4$ has a minimum at $x=0$
    \item $y=-x^4$ has a maximum at $x=0$
\end{itemize}

However, this is not an insoluble problem. What we must do is continue to differentiate until we get, at the $(n+1)$th derivative, a non-zero result at the stationary point:



$$f'(x)=0\ ,\ f''(x)=0\ ,\ \ldots\ ,\ f^{(n)}(x)=0\ ,\ f^{(n+1)}(x)\ne0$$ If $n$ is odd, then the stationary point is a true extremum. If the $(n+1)$th derivative is positive, it is a minimum; if the $(n+1)$th derivative is negative, it is a maximum. If $n$ is even, then the stationary point is a point of inflexion.



As an example, let us consider the function



$$f(x)=-x^4$$ We now differentiate until we get a non-zero result at the stationary point at $x=0$ (assume we have already found this point as usual):



$$f'(x)=-4x^3$$



$$f''(x)=-12x^2$$



$$f'''(x)=-24x$$



$$f^{(4)}(x)=-24$$ Therefore, $n+1$ is 4, so $n$ is 3. This is odd, and the fourth derivative is negative, so we have a maximum. Note that none of the methods given can tell you if this is a global extremum or just a local one. To do this, you would have to set the function equal to the height of the extremum and look for other roots.



\paragraph{Critical Points}



\textbf{Critical points} are the points where a function's derivative is 0 or not defined. Suppose we are interested in finding the maximum or minimum on given closed interval of a function that is continuous on that interval. The extreme values of the function on that interval will be at one or more of the critical points and/or at one or both of the endpoints. We can prove this by contradiction. Suppose that the function $f(x)$ has maximum at a point $x=c$ in the interval $(a,b)$ where the derivative of the function is defined and not $0$ . If the derivative is positive, then $x$ values slightly greater than $c$ will cause the function to increase. Since $c$ is not an endpoint, at least some of these values are in $[a,b]$ . But this contradicts the assumption that $f(c)$ is the maximum of $f(x)$ for $x$ in $[a,b]$ . Similarly, if the derivative is negative, then $x$ values slightly less than $c$ will cause the function to increase. Since $c$ is not an endpoint, at least some of these values are in $[a,b]$ . This contradicts the assumption that $f(c)$ is the maximum of $f(x)$ for $x$ in $[a,b]$ . A similar argument could be made for the minimum.



\paragraph{Example 1}



Consider the function $f(x)=x$ on the interval $[-1,1]$ . The unrestricted function $f(x)=x$ has no maximum or minimum. On the interval $[-1,1]$ , however, it is obvious that the minimum will be $-1$ , which occurs at $=-1$ and the maximum will be $1$, which occurs at $x=1$ . Since there are no critical points ($f'(x)$ exists and equals $1$ everywhere), the extreme values must occur at the endpoints.



\paragraph{Example 2}



Find the maximum and minimum of the function $f(x)=x^3-2x^2-5x+6$ on the interval $[-3,3]$ .First start by finding the roots of the function derivative:



$$f'(x)=3x^2-4x-5=0$$



$$x=\frac{2\pm\sqrt{19}}{3}$$



:   Now evaluate the function at all critical points and endpoints to find the extreme values.

:   $f(-3)=-24$

:   $f(\frac{2-\sqrt{19}}{3})=\frac{56+38\sqrt{19}}{27}\approx 8.2088$

:   $f(\frac{2+\sqrt{19}}{3})=\frac{56-38\sqrt{19}}{27}\approx -4.0607$

:   $f(3)=0$

:   From this we can see that the minimum on the interval is -24 when $x=-3$ and the maximum on the interval is $\frac{56+38\sqrt{19}}{27}$ when $x=\frac{2-\sqrt{19}}{3}$



\paragraph{Saddle Point}



% [Image: A saddle point (in red) on the graph of z=x^2^-y^2^ (hyperbolic paraboloid)]



Figure: A saddle point (in red) on the graph of z=x^2^-y^2^ (hyperbolic paraboloid).



% [Image: Saddle point between two hills (the intersection of the figure-eight $z$-contour)]



Figure: Saddle point between two hills (the intersection of the figure-eight $z$-contour.



In mathematics, a \textbf{saddle point} or \textbf{minimax point} is a point on the surface of the graph of a function where the slopes (derivatives) in orthogonal directions are all zero (a critical point), but which is not a local extremum of the function. An example of a saddle point is when there is a critical point with a relative minimum along one axial direction (between peaks) and at a relative maximum along the crossing axis. However, a saddle point need not be in this form. For example, the function $f(x,y) = x^2 + y^3$ has a critical point at $(0, 0)$ that is a saddle point since it is neither a relative maximum nor relative minimum, but it does not have a relative maximum or relative minimum in the $y$-direction.



The name derives from the fact that the prototypical example in two dimensions is a surface that \textit{curves up} in one direction, and \textit{curves down} in a different direction, resembling a riding saddle or a mountain pass between two peaks forming a landform saddle. In terms of contour lines, a saddle point in two dimensions gives rise to a contour graph or trace in which the contour corresponding to the saddle point's value appears to intersect itself.



% [Image: Saddle point on the countour plot is the point where level curves cross ]



Figure: Saddle point on the countour plot is the point where level curves cross.



\paragraph{Mathematical discussion}



A simple criterion for checking if a given stationary point of a real-valued function \textit{F}(\textit{x},\textit{y}) of two real variables is a saddle point is to compute the function's Hessian matrix at that point: if the Hessian is indefinite, then that point is a saddle point. For example, the Hessian matrix of the function $z=x^2-y^2$ at the stationary point $(x, y, z)=(0, 0, 0)$ is the matrix



> $\begin{bmatrix}

    2 \& 0 \\\    0 \& -2 \\\    \end{bmatrix}$



which is indefinite. Therefore, this point is a saddle point. This criterion gives only a sufficient condition. For example, the point $(0, 0, 0)$ is a saddle point for the function $z=x^4-y^4,$ but the Hessian matrix of this function at the origin is the null matrix, which is not indefinite.



In the most general terms, a \textbf{saddle point} for a smooth function (whose graph is a curve, surface or hypersurface) is a stationary point such that the curve/surface/etc. in the neighborhood of that point is not entirely on any side of the tangent space at that point.



% [Image: The plot of \textit{y} = \textit{x}^3^ with a saddle point at 0]



Figure: The plot of $y= x^3$ with a saddle point at 0



In a domain of one dimension, a saddle point is a point which is both a stationary point and a point of inflection. Since it is a point of inflection, it is not a local extremum.



\paragraph{Saddle surface}



% [Image: Hyperbolic paraboloid]



Figure: Hyperbolic paraboloid



% [Image: A model of an elliptic hyperboloid of one sheet]



Figure: A model of an elliptic hyperboloid of one sheet.



% [Image: A monkey saddle]



Figure: A monkey saddle.



A \textbf{saddle surface} is a smooth surface containing one or more saddle points.



Classical examples of two-dimensional saddle surfaces in the Euclidean space are second order surfaces, the hyperbolic paraboloid $z=x^2-y^2$ (which is often referred to as "\textit{the} saddle surface" or "the standard saddle surface") and the hyperboloid of one sheet. The Pringles potato chip or crisp is an everyday example of a hyperbolic paraboloid shape.



Saddle surfaces have negative Gaussian curvature which distinguish them from convex/elliptical surfaces which have positive Gaussian curvature. A classical third-order saddle surface is the monkey saddle.



\paragraph{Examples}



In a two-player zero sum game defined on a continuous space, the equilibrium point is a saddle point.



For a second-order linear autonomous system, a critical point is a saddle point if the characteristic equation has one positive and one negative real eigenvalue.



In optimization subject to equality constraints, the first-order conditions describe a saddle point of the Lagrangian.



\paragraph{Other uses}



In dynamical systems, if the dynamic is given by a differentiable map $f$ then a point is hyperbolic if and only if the differential of $f ^n$ (where $n$ is the period of the point) has no eigenvalue on the (complex) unit circle when computed at the point. Then a \textit{saddle point} is a hyperbolic periodic point whose stable and unstable manifolds have a dimension that is not zero.



A saddle point of a matrix is an element which is both the largest element in its column and the smallest element in its row.



\paragraph{Extrema for functions of two variables}



Suppose that $f(x, y)$ is a differentiable real function of two variables whose second partial derivatives exist and are continuous. The Hessian matrix $H$ of $f$ is the $2 \times 2$ matrix of partial derivatives of $f$:



$$

H(x,y) = \begin{bmatrix}

f _{xx}(x,y) & f _{xy}(x,y) \\\f _{yx}(x,y) & f _{yy}(x,y)

\end{bmatrix}.

$$



Define $D(x,y)$ to be the determinant



$$

D(x,y) = \det(H(x,y)) = f _{xx}(x,y) f _{yy}(x,y) - \left( f _{xy}(x,y) \right)^2

$$



of $H$. Finally, suppose that $(a, b)$ is a critical point of $f$ (that is, $f _x(a, b) = f _y(a, b) = 0$). Then the second partial derivative test asserts the following:



1.  If $D(a, b) > 0$ and $f _{xx}(a, b) > 0$ then $(a, b)$ is a local minimum of $f$.

2.  If $D(a, b) > 0$ and $f _{xx}(a, b) < 0$ then $(a, b)$ is a local maximum of $f$.

3.  If $D(a, b) < 0$ then $(a, b)$ is a saddle point of $f$.

4.  If $D(a, b) = 0$ then the second derivative test is inconclusive, and the point $(a, b)$ could be any of a minimum, maximum, or saddle point.



Sometimes other equivalent versions of the test are used. Note that in cases 1 and 2, the requirement that $f _{xx} f _{yy} - f _{xy}^2$ is positive at $(x, y)$ implies that $f _{xx}$ and $f _{yy}$ have the same sign there. Therefore the second condition, that $f _{xx}$ be greater (or less) than zero, could equivalently be that $f _{yy}$ or $\text{tr}(H) = f _{xx} + f _{yy}$ be greater (or less) than zero at that point.



\paragraph{Functions of many variables}



For a function $f$ of three or more variables, there is a generalization of the rule above. In this context, instead of examining the determinant of the Hessian matrix, one must look at the eigenvalues of the Hessian matrix at the critical point. The following test can be applied at any critical point $a$ for which the Hessian matrix is invertible:



1.  If the Hessian is positive definite (equivalently, has all eigenvalues positive) at $a$, then \textit{f} attains a local minimum at $a$.

2.  If the Hessian is negative definite (equivalently, has all eigenvalues negative) at $a$, then \textit{f} attains a local maximum at $a$.

3.  If the Hessian has both positive and negative eigenvalues then $a$ is a saddle point for \textit{f} (and in fact this is true even if $a$ is degenerate).



In those cases not listed above, the test is inconclusive.



For functions of three or more variables, the \textit{determinant} of the Hessian does not provide enough information to classify the critical point, because the number of jointly sufficient second-order conditions is equal to the number of variables, and the sign condition on the determinant of the Hessian is only one of the conditions. Note that in the one-variable case, the Hessian condition simply gives the usual second derivative test.



In the two variable case, $D(a, b)$ and $f _{xx}(a,b)$ are the principal minors of the Hessian. The first two conditions listed above on the signs of these minors are the conditions for the positive or negative definiteness of the Hessian. For the general case of an arbitrary number \textit{n} of variables, there are \textit{n} sign conditions on the \textit{n} principal minors of the Hessian matrix that together are equivalent to positive or negative definiteness of the Hessian (Sylvester's criterion): for a local minimum, all the principal minors need to be positive, while for a local maximum, the minors with an odd number of rows and columns need to be negative and the minors with an even number of rows and columns need to be positive.



\paragraph{Resources}



\begin{itemize}
    \item \href{https://www.youtube.com/watch?v=UID893EosM8}{Local Extrema, Critical Points, \& Saddle Points of Multivariable Functions - Calculus 3-} Video by he Organic Chemistry Tutor 2019
\end{itemize}

<!-- -->



\begin{itemize}
    \item \href{https://www.youtube.com/watch?v=Hg38kfK5w4E}{Absolute Maximum and Minimum Values of Multivariable Functions - Calculus 3} Video by The Organic Chemistry Tutor 2019
\end{itemize}

\paragraph{Licensing}



Content obtained and/or adapted from:



\begin{itemize}
    \item \href{https://en.wikibooks.org/wiki/Calculus/Extrema_and_Points_of_Inflection}{Extrema and Points of Inflection, Wikibooks: Calculus} under a CC BY-SA license
    \item \href{https://en.wikipedia.org/wiki/Second_partial_derivative_test}{Second partial derivative test, Wikipedia} under a CC BY-SA license
    \item \href{https://en.wikipedia.org/wiki/Saddle_point}{Saddle point, Wikipedia} under a CC BY-SA license
\end{itemize}

\subsection{Learning Objectives}

\begin{enumerate}
    \item \textbf{Identify Critical Points and Classify Extrema:}
    \item \textbf{Outcome:} Students will be able to find the critical points of a function and classify them as local maxima, local minima, or saddle points using the first derivative test.
    \item \textbf{Details:} This includes understanding how to find the derivative of a function, set it to zero to find critical points, and analyze the sign changes of the derivative around these points to classify the type of extrema.
    \item \textbf{Apply the Second Derivative Test:}
    \item \textbf{Outcome:} Students will be able to apply the second derivative test to classify critical points as local maxima, local minima, or saddle points.
    \item \textbf{Details:} This includes calculating the second derivative of a function at critical points and using its sign to determine whether each critical point is a maximum, minimum, or point of inflection. Students will also learn how to handle cases where the second derivative is zero.
    \item \textbf{Determine Global Extrema on a Closed Interval:}
    \item \textbf{Outcome:} Students will be able to determine the global maxima and minima of a function on a closed interval.
    \item \textbf{Details:} This involves evaluating the function at critical points and endpoints of the interval, and comparing these values to find the highest and lowest values. Students will understand how to apply the extreme value theorem and use critical points and boundary points in their analysis.
\end{enumerate}

\subsection{Assessment Instrument}

\begin{enumerate}
    \item \textbf{Assessment Instrument: Maxima and Minima Quiz} \textbf{Format:}
    \item \textbf{Type:} Mixed-format quiz (multiple-choice, short answer, and problem-solving questions)
    \item \textbf{Duration:} 45 minutes
    \item \textbf{Total Marks:} 30 points \textbf{Sections:}
    \item \textbf{Identifying Critical Points (10 points):}
    \item \textbf{Question 1:} Find the critical points of the following functions by setting the first derivative equal to zero:
    \item a. \( f(x) = x^3 - 3x^2 + 2 \)
    \item b. \( g(x) = x^4 - 4x^2 \)
    \item \textbf{Classifying Extrema Using the First Derivative Test (10 points):}
    \item \textbf{Question 2:} Use the first derivative test to classify the critical points found in Question 1 as local maxima, local minima, or saddle points:
    \item a. \( f(x) = x^3 - 3x^2 + 2 \)
    \item b. \( g(x) = x^4 - 4x^2 \)
    \item \textbf{Applying the Second Derivative Test (10 points):}
    \item \textbf{Question 3:} Use the second derivative test to classify the critical points found in Question 1:
    \item a. \( f(x) = x^3 - 3x^2 + 2 \)
    \item b. \( g(x) = x^4 - 4x^2 \) \textbf{Scoring Rubric:}
    \item \textbf{Identifying Critical Points:} 5 points for each correctly identified critical point (total 10 points)
    \item \textbf{Classifying Extrema Using the First Derivative Test:} 5 points for each correctly classified critical point (total 10 points)
    \item \textbf{Applying the Second Derivative Test:} 5 points for each correctly classified critical point using the second derivative test (total 10 points) The quiz will assess the students' understanding and application of finding and classifying critical points, local extrema, and using both the first and second derivative tests. It ensures that students can correctly identify critical points and apply the appropriate methods to classify these points accurately.
\end{enumerate}

%%===============================================================
\section{Lagrange Multipliers}
\label{sec:lagrange-multipliers}
%%===============================================================

The method of Lagrange multipliers solves the constrained optimization problem by transforming it into a non-constrained optimization problem of the form:



\begin{itemize}
    \item $\operatorname{\mathcal{L}}(x _1,x _2,\ldots, x _n,\lambda)= \operatorname{f}(x _1,x _2,\ldots, x _n)+\operatorname{\lambda}(k-g(x _1,x _2,\ldots, x _n))$
\end{itemize}

Then finding the gradient and Hessian as was done above will determine any optimum values of $\operatorname{\mathcal{L}}(x _1,x _2,\ldots, x _n,\lambda)$.



Suppose we now want to find optimum values for $f(x,y)=2x^2+y^2$ subject to $x+y=1$ from \[2\].



Then the Lagrangian method will result in a non-constrained function.



\begin{itemize}
    \item $\operatorname{\mathcal{L}}(x,y,\lambda)= 2x^2+y^2+\lambda (1-x-y)$
\end{itemize}

The gradient for this new function is



\begin{itemize}
    \item $\frac{\partial \mathcal{L}}{\partial x}(x,y,\lambda)= 4x+\lambda (-1)=0$
    \item $\frac{\partial \mathcal{L}}{\partial y}(x,y,\lambda)= 2y+\lambda (-1)=0$
    \item $\frac{\partial \mathcal{L}}{\partial \lambda}(x,y,\lambda)=1-x-y=0$
\end{itemize}

Finding the stationary points of the above equations can be obtained from their matrix from.



$$\begin{bmatrix} 4 & 0 & -1 \\\\ 0 & 2 & -1 \\\\ -1 & -1 & 0 \end{bmatrix} \begin{bmatrix} x \\\\ y \\\\ \lambda \end{bmatrix} = \begin{bmatrix} 0 \\\\ 0 \\\\ -1 \end{bmatrix}$$



This results in $x=1/3, y=2/3, \lambda=4/3$.



Next we can use the Hessian as before to determine the type of this stationary point.



> $H(\mathcal{L})=

     \begin{bmatrix}

    4 \& 0 \& -1 \\\    0\& 2 \& -1 \\\    -1\&-1\&0

    \end{bmatrix}$



Since $H(\mathcal{L}) >0$ then the solution $(1/3,2/3,4/3)$ minimizes $f(x,y)=2x^2+y^2$ subject to $x+y=1$ with $f(x,y)=2/3$.



\paragraph{Resources}



\begin{itemize}
    \item \href{https://en.wikibooks.org/wiki/Calculus_Optimization_Methods/Lagrange_Multipliers}{Lagrange Multipliers}, WikiBooks: Calculus Optimization Methods
\end{itemize}

\paragraph{Videos}



\begin{itemize}
    \item \href{https://www.youtube.com/watch?v=ry9cgNx1QV8\&list=RDCMUCFe6jenM1Bc54qtBsIJGRZQ\&index=5}{LaGrange Multipliers - Finding Maximum or Minimum Values}Video by patrickJMT
\end{itemize}

\begin{itemize}
    \item \href{https://www.youtube.com/watch?v=x6j6yFzTUgU}{Lagrange Multipliers Practice Problems} Video by ames Hamblin 2017
\end{itemize}

\begin{itemize}
    \item \href{https://www.youtube.com/watch?v=HyqBcD_e_Uw}{Lagrange multipliers \| MIT 18.02SC Multivariable Calculus, Fall 2010} Video by MIT OpenCourseWare
\end{itemize}

\begin{itemize}
    \item \href{https://www.youtube.com/watch?v=qXhcpqslNUU}{Lagrange Multipliers - Two Constraints -patrickJMT 2009} Video by patrickJMT 2009
\end{itemize}

\begin{itemize}
    \item \href{https://www.youtube.com/watch?v=nDuS5uQ7-lo}{Lagrange multipliers (3 variables) \| MIT 18.02SC Multivariable Calculus, Fall 2010} Video by MIT OpenCourseWare
\end{itemize}

\paragraph{Licensing}



Content obtained and/or adapted from:



\begin{itemize}
    \item \href{https://en.wikibooks.org/wiki/Calculus_Optimization_Methods/Lagrange_Multipliers}{Lagrange Multipliers, WikiBooks: Calculus Optimization Methods} under a CC BY-SA license
\end{itemize}

\subsection{Learning Objectives}

\begin{enumerate}
    \item Here are three Student Learning Objectives (SLOs) based on the method of Lagrange multipliers for solving constrained optimization problems:
    \item \textbf{Formulate Lagrangian Functions:}
    \item \textbf{Outcome:} Students will be able to formulate the Lagrangian function for a given constrained optimization problem.
    \item \textbf{Details:} This includes transforming a constrained optimization problem into a Lagrangian function by incorporating the constraint(s) with a Lagrange multiplier. For example, given $f(x, y) = 2x^2 + y^2$ subject to $x + y = 1$, students will formulate the Lagrangian function as $\operatorname{\mathcal{L}}(x, y, \lambda) = 2x^2 + y^2 + \lambda (1 - x - y)$.
    \item \textbf{Compute Partial Derivatives of the Lagrangian:}
    \item \textbf{Outcome:} Students will be able to compute the partial derivatives of the Lagrangian function with respect to all variables including the Lagrange multiplier.
    \item \textbf{Details:} This includes taking partial derivatives with respect to each variable to set up the system of equations. For example, for the Lagrangian function $\operatorname{\mathcal{L}}(x, y, \lambda) = 2x^2 + y^2 + \lambda (1 - x - y)$, students will compute $\frac{\partial \mathcal{L}}{\partial x} = 4x - \lambda = 0$, $\frac{\partial \mathcal{L}}{\partial y} = 2y - \lambda = 0$, and $\frac{\partial \mathcal{L}}{\partial \lambda} = 1 - x - y = 0$.
    \item \textbf{Solve for Stationary Points and Determine Optima:}
    \item \textbf{Outcome:} Students will be able to solve the system of equations obtained from the partial derivatives to find stationary points and determine if these points are maxima, minima, or saddle points using the Hessian matrix.
    \item \textbf{Details:} This includes solving the system of equations and analyzing the Hessian matrix to classify the stationary points. For example, solving $\frac{\partial \mathcal{L}}{\partial x} = 4x - \lambda = 0$, $\frac{\partial \mathcal{L}}{\partial y} = 2y - \lambda = 0$, and $\frac{\partial \mathcal{L}}{\partial \lambda} = 1 - x - y = 0$ results in $x = \frac{1}{3}$, $y = \frac{2}{3}$, $\lambda = \frac{4}{3}$. Using the Hessian matrix $H(\mathcal{L}) = \begin{bmatrix} 4 & 0 & -1 \\ 0 & 2 & -1 \\ -1 & -1 & 0 \end{bmatrix}$, students will analyze its properties to determine the nature of the stationary point $(\frac{1}{3}, \frac{2}{3}, \frac{4}{3})$ and confirm that it minimizes the function $f(x, y) = 2x^2 + y^2$ subject to the constraint $x + y = 1$ with $f(\frac{1}{3}, \frac{2}{3}) = \frac{2}{3}$.
\end{enumerate}

\subsection{Assessment Instrument}

\begin{enumerate}
    \item \textbf{Assessment Instrument: Lagrange Multipliers Quiz} \textbf{Format:}
    \item \textbf{Type:} Mixed-format quiz (multiple-choice, short answer, and problem-solving questions)
    \item \textbf{Duration:} 45 minutes
    \item \textbf{Total Marks:} 30 points \textbf{Sections:}
    \item \textbf{Formulating the Lagrangian Function (10 points):}
    \item \textbf{Question 1:} Given the objective function $f(x, y) = 3x^2 + 4y^2$ subject to the constraint $2x + y = 5$, formulate the Lagrangian function.
    \item \textbf{Answer:} $\operatorname{\mathcal{L}}(x, y, \lambda) = 3x^2 + 4y^2 + \lambda (5 - 2x - y)$
    \item \textbf{Question 2:} Multiple Choice: Which of the following correctly represents the Lagrangian for the function $f(x, y) = x^2 + y^2$ with the constraint $x + 2y = 3$?
    \item a. $\operatorname{\mathcal{L}}(x, y, \lambda) = x^2 + y^2 + \lambda (3 - x - 2y)$
    \item b. $\operatorname{\mathcal{L}}(x, y, \lambda) = x^2 + y^2 - \lambda (3 - x - 2y)$
    \item c. $\operatorname{\mathcal{L}}(x, y, \lambda) = x^2 + y^2 + \lambda (x + 2y - 3)$
    \item d. $\operatorname{\mathcal{L}}(x, y, \lambda) = x^2 + y^2 + \lambda (3 - x + 2y)$
    \item \textbf{Computing Partial Derivatives (10 points):}
    \item \textbf{Question 3:} Compute the partial derivatives of the Lagrangian function $\operatorname{\mathcal{L}}(x, y, \lambda) = 2x^2 + y^2 + \lambda (1 - x - y)$ with respect to $x$, $y$, and $\lambda$.
    \item \textbf{Answer:}
    \item $\frac{\partial \mathcal{L}}{\partial x} = 4x - \lambda$
    \item $\frac{\partial \mathcal{L}}{\partial y} = 2y - \lambda$
    \item $\frac{\partial \mathcal{L}}{\partial \lambda} = 1 - x - y$
    \item \textbf{Question 4:} Short Answer: For the Lagrangian function $\operatorname{\mathcal{L}}(x, y, \lambda) = x^2 + y^2 + \lambda (4 - x - y)$, write down the partial derivative equations.
    \item \textbf{Solving for Stationary Points and Determining Optima (10 points):}
    \item \textbf{Question 5:} Solve the system of equations obtained from the partial derivatives of the Lagrangian function $\operatorname{\mathcal{L}}(x, y, \lambda) = x^2 + y^2 + \lambda (1 - x - y)$.
    \item \textbf{Answer:}
    \item $\frac{\partial \mathcal{L}}{\partial x} = 2x - \lambda = 0$
    \item $\frac{\partial \mathcal{L}}{\partial y} = 2y - \lambda = 0$
    \item $\frac{\partial \mathcal{L}}{\partial \lambda} = 1 - x - y = 0$
    \item Solving these: $x = \frac{\lambda}{2}$, $y = \frac{\lambda}{2}$, $\frac{\lambda}{2} + \frac{\lambda}{2} = 1 \Rightarrow \lambda = 2 \Rightarrow x = 1, y = 1$
    \item \textbf{Question 6:} Problem Solving: For the Lagrangian $\operatorname{\mathcal{L}}(x, y, \lambda) = x^2 + 2y^2 + \lambda (3 - x - 2y)$, find the stationary points and use the Hessian matrix to determine if it is a maximum, minimum, or saddle point. \textbf{Scoring Rubric:}
    \item \textbf{Formulating the Lagrangian Function:} 5 points for each correctly formulated Lagrangian (total 10 points)
    \item \textbf{Computing Partial Derivatives:} 5 points for each set of correct partial derivatives (total 10 points)
    \item \textbf{Solving for Stationary Points and Determining Optima:} 5 points for each correctly solved stationary point and determination using the Hessian matrix (total 10 points) This quiz will assess the students' understanding and application of the Lagrange multipliers method for constrained optimization problems, ensuring they can correctly formulate the Lagrangian function, compute partial derivatives, and determine stationary points and their nature.
\end{enumerate}

%%===============================================================
\section{Multiple Integrals (too much for one lesson)}
\label{sec:multiple-integrals-too-much-for-one-lesson}
%%===============================================================

% [Image: Integral as area between two curves.]



Figure: Integral as area between two curves.



% [Image: Double integral as volume under a surface]



Figure: Double integral as volume under a surface $z = 10 - \frac{x^2 - y^2}{8}$. The rectangular region at the bottom of the body is the domain of integration, while the surface is the graph of the two-variable function to be integrated.



In mathematics (specifically multivariable calculus), a \textbf{multiple integral} is a definite integral of a function of several real variables, for instance, $f(x,y)$ or $f(x,y,z)$. Integrals of a function of two variables over a region in $\mathbb{R} ^2$ (the real-number plane) are called \textbf{double integrals}, and integrals of a function of three variables over a region in $\mathbb{R} ^3$ (real-number 3D space) are called \textbf{triple integrals}. For multiple integrals of a single-variable function, see the Cauchy formula for repeated integration.



\paragraph{Introduction}



Just as the definite integral of a positive function of one variable represents the area of the region between the graph of the function and the $x$-axis, the \textbf{double integral} of a positive function of two variables represents the volume of the region between the surface defined by the function (on the three-dimensional Cartesian plane where $z = f(x,y)$) and the plane which contains its domain. If there are more variables, a multiple integral will yield hypervolumes of multidimensional functions.



Multiple integration of a function in $n$ variables: $f(x _1, x _2, ..., x _n)$ over a domain $D$ is most commonly represented by nested integral signs in the reverse order of execution (the leftmost integral sign is computed last), followed by the function and integrand arguments in proper order (the integral with respect to the rightmost argument is computed last). The domain of integration is either represented symbolically for every argument over each integral sign, or is abbreviated by a variable at the rightmost integral sign:



$$\int \cdots \int _\mathbf{D}\, f(x _1,x _2,\ldots,x _n) \,dx _1 \cdots dx _n$$



Since the concept of an antiderivative is only defined for functions of a single real variable, the usual definition of the indefinite integral does not immediately extend to the multiple integral.



\paragraph{Mathematical definition}



For $n > 1$, consider a so-called "half-open" $n$-dimensional hyperrectangular domain $T$, defined as:



$$T= [ a _1, b _1) \times [ a _2, b _2) \times \cdots \times [ a _n, b _n) \subseteq \mathbb{R} ^n.$$



Partition each interval $[a _j, b _j)$ into a finite family $I _j$ of non-overlapping subintervals $i _{j _{\alpha}}$, with each subinterval closed at the left end, and open at the right end.



Then the finite family of subrectangles $C$ given by



$$C=I _1\times I _2\times \cdots \times I _n$$



is a partition of $T$; that is, the subrectangles $C _k$ are non-overlapping and their union is $T$.



Let $f : T \to \mathbb{R}$ be a function defined on $T$. Consider a partition $C$ of $T$ as defined above, such that $C$ is a family of $m$ subrectangles $C _m$ and



$$T=C _1\cup C _2\cup \cdots \cup C _m$$



We can approximate the total $(n + 1)$th-dimensional volume bounded below by the $n$-dimensional hyperrectangle $T$ and above by the $n$-dimensional graph of $f$ with the following Riemann sum:



$$\sum _{k=1}^m f(P _k)\, \operatorname{m}(C _k)$$



where $P _k$ is a point in $C _k$ and $m(C _k)$ is the product of the lengths of the intervals whose Cartesian product is $C _k$, also known as the measure of $C _k$.



The \textbf{diameter} of a subrectangle $C _k$ is the largest of the lengths of the intervals whose Cartesian product is $C _k$. The diameter of a given partition of $T$ is defined as the largest of the diameters of the subrectangles in the partition. Intuitively, as the diameter of the partition $C$ is restricted smaller and smaller, the number of subrectangles $m$ gets larger, and the measure $m(C _k)$ of each subrectangle grows smaller. The function $f$ is said to be \textbf{Riemann integrable} if the limit



$$S=\lim _{\delta \to 0} \sum _{k=1}^m f(P _k)\, \operatorname{m} (C _k)$$



exists, where the limit is taken over all possible partitions of $T$ of diameter at most $\delta$.



If $f$ is Riemann integrable, $S$ is called the \textbf{Riemann integral} of $f$ over $T$ and is denoted



$$\int \cdots \int _T\, f(x _1,x _2,\ldots,x _n) \,dx _1 \cdots dx _n$$



Frequently this notation is abbreviated as



> $\int _T f(\mathbf{x})\,d^n\mathbf{x}.$



where $\mathbf{x}$ represents the $n$-tuple $(x _1, \ldots, x _n)$ and $d^n\mathbf{x}$ is the $n$-dimensional volume differential.



The Riemann integral of a function defined over an arbitrary bounded $n$-dimensional set can be defined by extending that function to a function defined over a half-open rectangle whose values are zero outside the domain of the original function. Then the integral of the original function over the original domain is defined to be the integral of the extended function over its rectangular domain, if it exists.



In what follows the Riemann integral in $n$ dimensions will be called the \textbf{multiple integral}.



\paragraph{Properties}



Multiple integrals have many properties common to those of integrals of functions of one variable (linearity, commutativity, monotonicity, and so on). One important property of multiple integrals is that the value of an integral is independent of the order of integrands under certain conditions. This property is popularly known as Fubini's theorem.



\paragraph{Particular cases}



In the case of $T \subseteq \R^2$, the integral



$$l = \iint _T f(x,y)\, dx\, dy$$



is the \textbf{double integral} of $f$ on $T$, and if $T \subseteq \mathbb{R} ^3$ the integral



$$l = \iiint _T f(x,y,z)\, dx\, dy\, dz$$



is the \textbf{triple integral} of $f$ on $T$.



Notice that, by convention, the double integral has two integral signs, and the triple integral has three; this is a notational convention which is convenient when computing a multiple integral as an iterated integral, as shown later in this article.



\paragraph{Methods of integration}



The resolution of problems with multiple integrals consists, in most cases, of finding a way to reduce the multiple integral to an iterated integral, a series of integrals of one variable, each being directly solvable. For continuous functions, this is justified by Fubini's theorem. Sometimes, it is possible to obtain the result of the integration by direct examination without any calculations.



The following are some simple methods of integration:



\paragraph{Integrating constant functions}



When the integrand is a constant function $c$, the integral is equal to the product of $c$ and the measure of the domain of integration. If $c = 1$ and the domain is a subregion of $\mathbb{R} ^2$, the integral gives the area of the region, while if the domain is a subregion of $\mathbb{R} ^3$, the integral gives the volume of the region.



> \textbf{Example.} Let $f(x, y) = 2$ and

>

> $$D = \left\\{ (x,y) \in \mathbb{R} ^2 \ : \ 2 \le x \le 4 \ ; \ 3 \le y \le 6 \right\\}$$

>

> in which case

>

> $$\int _3^6 \int _2^4 \ 2 \, dx\, dy = 2 \int _3^6 \int _2^4 \ 1 \, dx\, dy = 2 \cdot \operatorname{area}(D) = 2 \cdot (2 \cdot 3) = 12,$$

>

> since by definition we have:

>

> $$\int _3^6 \int _2^4 \ 1 \, dx\, dy = \operatorname{area}(D).$$



\paragraph{Use of symmetry}



When the domain of integration is symmetric about the origin with respect to at least one of the variables of integration and the integrand is odd with respect to this variable, the integral is equal to zero, as the integrals over the two halves of the domain have the same absolute value but opposite signs. When the integrand is even with respect to this variable, the integral is equal to twice the integral over one half of the domain, as the integrals over the two halves of the domain are equal.



> \textbf{Example 1.} Consider the function $f(x,y) = 2 \sin(x) - 3y^3 + 5$ integrated over the domain

>

> $$T = \left\\{ (x,y) \in \mathbb{R} ^2 \ : \ x^2 + y^2 \le 1 \right\\},$$

>

> a disc with radius 1 centered at the origin with the boundary included.

>

> Using the linearity property, the integral can be decomposed into three pieces:

>

> $$\iint _T \left(2 \sin(x) - 3y^3 + 5\right) \, dx \, dy = \iint _T 2 \sin(x) \, dx \, dy - \iint _T 3y^3 \, dx \, dy + \iint _T 5 \, dx \, dy$$

>

> The function $2 \sin(x)$ is an odd function in the variable $x$ and the disc $T$ is symmetric with respect to the $y$-axis, so the value of the first integral is 0. Similarly, the function $3y^3$ is an odd function of $y$, and $T$ is symmetric with respect to the $x$-axis, and so the only contribution to the final result is that of the third integral. Therefore, the original integral is equal to the area of the disk times 5, or $5\pi$.





> \textbf{Example 2.} Consider the function $f(x, y, z) = x \exp(y^2 + z^2)$ and as the integration region the ball with radius 2 centered at the origin,

>

> $$T = \left \\{ ( x,y, z) \in \R^3 \ : \ x^2+y^2+z^2 \le 4 \right \\}.$$ The "ball" is symmetric about all three axes, but it is sufficient to integrate with respect to $x$-axis to show that the integral is 0, because the function is an odd function of that variable.



\paragraph{Normal domains on $\mathbb{R} ^2$}



This method is applicable to any domain $D$ for which:



\begin{itemize}
    \item the projection of $D$ onto either the $x$-axis or the $y$-axis is bounded by the two values, $A$ and $B$
    \item any line perpendicular to this axis that passes between these two values intersects the domain in an interval whose endpoints are given by the graphs of two functions, $\alpha$ and $\beta$.
\end{itemize}

Such a domain will be here called a \textit{normal domain}. Elsewhere in the literature, normal domains are sometimes called type I or type II domains, depending on which axis the domain is fibred over. In all cases, the function to be integrated must be Riemann integrable on the domain, which is true (for instance) if the function is continuous.



\paragraph{$x$-axis}



If the domain $D$ is normal with respect to the $x$-axis, and $f : D \to \mathbb{R}$ is a continuous function; then $\alpha(x)$ and $\beta(x)$ (both of which are defined on the interval $[a, b]$) are the two functions that determine $D$. Then, by Fubini's theorem:



$$\iint _D f(x,y)\, dx\, dy = \int _a^b dx \int _{ \alpha (x)}^{ \beta (x)} f(x,y)\, dy.$$



\paragraph{$y$-axis}



If $D$ is normal with respect to the $y$-axis and $f : D \to \mathbb{R}$ is a continuous function; then $\alpha(y)$ and $\beta(y)$ (both of which are defined on the interval $[a, b]$) are the two functions that determine $D$. Again, by Fubini's theorem:



$$\iint _D f(x,y)\, dx\, dy = \int _a^b dy \int _{\alpha (y)}^{ \beta (y)} f(x,y)\, dx.$$



\paragraph{Normal domains on $\mathbb{R} ^3$}



If $T$ is a domain that is normal with respect to the $xy$-plane and determined by the functions $\alpha(x, y)$ and $\beta(x, y)$, then



$$\iiint _T f(x,y,z) \, dx\, dy\, dz = \iint _D \int _{\alpha (x,y)}^{\beta (x,y)} f(x,y,z) \, dz\, dx\, dy$$



This definition is the same for the other five normality cases on $\mathbb{R} ^3$. It can be generalized in a straightforward way to domains in $\mathbb{R} ^n$.



\paragraph{Change of variables}



The limits of integration are often not easily interchangeable (without normality or with complex formulae to integrate). One makes a change of variables to rewrite the integral in a more "comfortable" region, which can be described in simpler formulae. To do so, the function must be adapted to the new coordinates.



> \textbf{Example 1a.} The function is $f(x, y) = (x - 1)^2 + \sqrt{y}$; if one adopts the substitution $\newline$ $u = x - 1$, $v = y$ therefore $x = u + 1$, $y = v$ one obtains the new function $f _2(u, v) = u^2 + \sqrt{v}$.



\begin{itemize}
    \item Similarly for the domain because it is delimited by the original variables that were transformed before ($x$ and $y$ in example).
    \item the differentials $dx$ and $dy$ transform via the absolute value of the determinant of the Jacobian matrix containing the partial derivatives of the transformations regarding the new variable (consider, as an example, the differential transformation in polar coordinates).
\end{itemize}

There exist three main "kinds" of changes of variable (one in $\mathbb{R} ^2$, two in $\mathbb{R} ^3$); however, more general substitutions can be made using the same principle.



\paragraph{Polar coordinates}



% [Image: Transformation from cartesian to polar coordinates.]



 Figure: Transformation from cartesian to polar coordinates.



In $\mathbb{R} ^2$ if the domain has a circular symmetry and the function has some particular characteristics one can apply the \textit{transformation to polar coordinates} (see the example in the picture) which means that the generic points $P(x,y)$ in Cartesian coordinates switch to their respective points in polar coordinates. That allows one to change the shape of the domain and simplify the operations.



The fundamental relation to make the transformation is the following:



$$f(x,y) \rightarrow f(\rho \cos \varphi,\rho \sin \varphi ).$$



> \textbf{Example 2a.} The function is $f(x,y)=x+y$ and applying the transformation one obtains

>

> $$f(\rho, \varphi) = \rho \cos \varphi + \rho \sin \varphi = \rho(\cos \varphi + \sin \varphi ).$$



> \textbf{Example 2b.} The function is $f(x,y)=x^2+y^2$, in this case one has:

>

> $$f(\rho, \varphi) = \rho^2 \left(\cos^2 \varphi + \sin^2 \varphi\right) = \rho^2$$ using the Pythagorean trigonometric identity (very useful to simplify this operation).



The transformation of the domain is made by defining the radius' crown length and the amplitude of the described angle to define the $\rho, \varphi$ intervals starting from $x, y$.



% [Image: Example of a domain transformation from cartesian to polar.]



Figure: Example of a domain transformation from cartesian to polar.



> \textbf{Example 2c.} The domain is $D = \\{x^2 + y^2 \le 4\\}$, that is, a circumference of radius 2; it's evident that the covered angle is the circle angle, so $\varphi$ varies from 0 to $2\pi$, while the crown radius varies from 0 to 2 (the crown with the inside radius null is just a circle).



> \textbf{Example 2d.} The domain is $D = \\{x^2 + y^2 \le 9, \ x^2 + y^2 \ge 4, \ y \ge 0\\}$, that is, the circular crown in the positive $y$ half-plane (please see the picture in the example); $\varphi$ describes a plane angle while $\rho$ varies from 2 to 3. Therefore, the transformed domain will be the following rectangle:



> $$T = \\{ 2 \le \rho \le 3, \ 0 \le \varphi \le \pi \\}.$$

>

> The Jacobian determinant of that transformation is the following:

>

> $$\frac{\partial (x,y)}{\partial (\rho, \varphi)} = \begin{vmatrix} \cos \varphi & - \rho \sin \varphi \\\\ \sin \varphi & \rho \cos \varphi \end{vmatrix} = \rho$$

>

> which has been obtained by inserting the partial derivatives of $x = \rho \cos(\varphi)$, $y = \rho \sin(\varphi)$ in the first column with respect to $\rho$ and in the second with respect to $\varphi$, so the $dx \, dy$ differentials in this transformation become $\rho \, d\rho \, d\varphi$.

>

> Once the function is transformed and the domain evaluated, it is possible to define the formula for the change of variables in polar coordinates:

>

> $$\iint _D f(x,y) \, dx\, dy = \iint _T f(\rho \cos \varphi, \rho \sin \varphi) \rho \, d \rho\, d \varphi.$$

>

> $\varphi$ is valid in the $[0, 2\pi]$ interval while $\rho$, which is a measure of a length, can only have positive values.



> \textbf{Example 2e.} The function is $f(x,y)=x$ and the domain is the same as in Example 2d. From the previous analysis of $D$ we know the intervals of $\rho$ (from 2 to 3) and of $\varphi$ (from 0 to $\pi$). Now we change the function:

>

> $$f(x,y) = x \longrightarrow f(\rho,\varphi) = \rho \cos \varphi.$$

>

> finally let's apply the integration formula:

>

> $$\iint _D x \, dx\, dy = \iint _T \rho \cos \varphi \rho \, d\rho\, d\varphi.$$

>

> Once the intervals are known, you have

>

> $$\int _0^\pi \int _2^3 \rho^2 \cos \varphi \, d \rho \, d \varphi = \int _0^\pi \cos \varphi \ d \varphi \left[ \frac{\rho^3}{3} \right] _2^3 = \Big[ \sin \varphi \Big] _0^\pi \ \left(9 - \frac{8}{3} \right) = 0.$$



\paragraph{Cylindrical coordinates}



% [Image: Cylindrical coordinates.]



Figure: Cylindrical coordinates.



In $\mathbb{R} ^3$ the integration on domains with a circular base can be made by the \textit{passage to cylindrical coordinates}; the transformation of the function is made by the following relation:



$$f(x,y,z) \rightarrow f(\rho \cos \varphi, \rho \sin \varphi, z)$$



The domain transformation can be graphically attained, because only the shape of the base varies, while the height follows the shape of the starting region.



> \textbf{Example 3a.} The region is $D = \\{x^2 + y^2 \le 9, \ x^2 + y^2 \ge 4, \ 0 \le z \le 5\\}$ (that is, the "tube" whose base is the circular crown of Example 2d and whose height is 5); if the transformation is applied, this region is obtained:

>

> $$T = \\{ 2 \le \rho \le 3, \ 0 \le \varphi \le 2 \pi, \ 0 \le z \le 5 \\}$$ (that is, the parallelepiped whose base is similar to the rectangle in Example 2d and whose height is 5).

>

> Because the $z$ component is unvaried during the transformation, the $dx \, dy \, dz$ differentials vary as in the passage to polar coordinates: therefore, they become $\rho \, d\rho \, d\varphi \, dz$.

>

> Finally, it is possible to apply the final formula to cylindrical coordinates:

>

> $$\iiint _D f(x,y,z) \, dx\, dy\, dz = \iiint _T f(\rho \cos \varphi, \rho \sin \varphi, z) \rho \, d\rho\, d\varphi\, dz.$$

>

> This method is convenient in case of cylindrical or conical domains or in regions where it is easy to individuate the $z$ interval and even transform the circular base and the function.



> \textbf{Example 3b.} The function is $f(x, y, z) = x^2 + y^2 + z$ and as integration domain this cylinder: $D = \\{x^2 + y^2 \le 9, \ -5 \le z \le 5\\}$. The transformation of $D$ in cylindrical coordinates is the following:

>

> $$T = \\{ 0 \le \rho \le 3, \ 0 \le \varphi \le 2 \pi, \ -5 \le z \le 5 \\}.$$

>

> while the function becomes

>

> $$f(\rho \cos \varphi, \rho \sin \varphi, z) = \rho^2 + z$$

>

> Finally one can apply the integration formula:

>

> $$\iiint _D \left(x^2 + y^2 +z\right) \, dx\, dy\, dz = \iiint _T \left( \rho^2 + z\right) \rho \, d\rho\, d\varphi\, dz;$$

>

> developing the formula you have

>

> $$ \begin{aligned}

  \int _{-5}^5 dz \int _0^{2\pi} d\varphi \int _0^3 \left( \rho^3 + \rho z \right)\, d\rho &= 2 \pi \int _{-5}^5 \left[ \frac{\rho^4}{4} + \frac{\rho^2 z}{2} \right] _0^3 \, dz \\\  &= 2 \pi \int _{-5}^5 \left( \frac{81}{4} + \frac{9}{2} z \right) \, dz \\\  &= 2 \pi \left[ \frac{81}{4} z + \frac{9}{4} z^2 \right] _{-5}^5 \\\  &= 2 \pi \left( \left[ \frac{81}{4} (5) + \frac{9}{4} (25) \right] - \left[ \frac{81}{4} (-5) + \frac{9}{4} (25) \right] \right) \\\  &= ... \\\  &= 405 \pi

  \end{aligned} $$



\paragraph{Spherical coordinates}



% [Image: Spherical coordinates.]



Figure: Spherical coordinates.



In $\mathbb{R} ^3$ some domains have a spherical symmetry, so it's possible to specify the coordinates of every point of the integration region by two angles and one distance. It's possible to use therefore the \textit{passage to spherical coordinates}; the function is transformed by this relation:



$$f(x,y,z) \longrightarrow f(\rho \cos \theta \sin \varphi, \rho \sin \theta \sin \varphi, \rho \cos \varphi)$$



Points on the $z$-axis do not have a precise characterization in spherical coordinates, so $\theta$ can vary between 0 and $2 \pi$.



The better integration domain for this passage is the sphere.



> \textbf{Example 4a.} The domain is $D = \\{x^2 + y^2 + z^2 \le 16\\}$ (a sphere with radius 4 and center at the origin); applying the transformation you get the region

>

> $$T = \\{ 0 \le \rho \le 4, \ 0 \le \varphi \le \pi, \ 0 \le \theta \le 2 \pi \\}.$$

>

> The Jacobian determinant of this transformation is the following:

>

> $$\frac{\partial (x,y,z)}{\partial (\rho, \theta, \varphi)} = \begin{vmatrix} \cos \theta \sin \varphi & - \rho \sin \theta \sin \varphi & \rho \cos \theta \cos \varphi \\\\ \sin \theta \sin \varphi & \rho \cos \theta \sin \varphi & \rho \sin \theta \cos \varphi \\\\ \cos \varphi & 0 & - \rho \sin \varphi \end{vmatrix} = \rho^2 \sin \varphi$$

>

> The $dx \, dy \, dz$ differentials therefore are transformed to $\rho^2 \sin(\varphi) \, d\rho \, d\theta \, d\varphi$.

>

> This yields the final integration formula:

>

> $$ \begin{aligned} \iiint _D f(x,y,z) \, dx\, dy\, dz &= \iiint _T f(\rho \sin \varphi \cos \theta, \rho \sin \varphi \sin \theta, \rho \cos \varphi) \rho^2 \sin \varphi \, d\rho\, d\theta\, d\varphi \end{aligned}  $$



It is better to use this method in case of spherical domains \textbf{and} in case of functions that can be easily simplified by the first fundamental relation of trigonometry extended to $\mathbb{R} ^3$ (see Example 4b); in other cases it can be better to use cylindrical coordinates (see Example 4c).



$$\iiint _T f(a,b,c) \rho^2 \sin \varphi \, d\rho\, d\theta\, d\varphi.$$



The extra $\rho^2$ and $\sin \varphi$ come from the Jacobian.



In the following examples, the roles of $\varphi$ and $\theta$ have been reversed.



> \textbf{Example 4b.} $D$ is the same region as in Example 4a and $f(x, y, z) = x^2 + y^2 + z^2$ is the function to integrate. Its transformation is very easy:

>

> $$f(\rho \sin \varphi \cos \theta, \rho \sin \varphi \sin \theta, \rho \cos \varphi) = \rho^2,$$

>

> while we know the intervals of the transformed region $T$ from $D$:

>

> $$T=\\{0 \le \rho \le 4, \ 0 \le \varphi \le \pi, \ 0 \le \theta \le 2 \pi\\}.$$

>

> We therefore apply the integration formula:

>

> $$\iiint _D \left(x^2 + y^2 + z^2\right) \, dx\, dy\, dz = \iiint _T \rho^2 \, \rho^2 \sin \theta \, d\rho\, d\theta\, d\varphi,$$

>

> and, developing, we get

>

> $$ \begin{aligned} \iiint _T \rho^4 \sin \varphi \, d\rho\, d\theta\, d\varphi &= \int _0^{\pi} \sin \varphi \, d\varphi \int _0^4 \rho^4 \, d\rho \int _0^{2 \pi} d\theta \\\\ &= 2 \pi \int _0^{\pi} \sin \varphi \left[ \frac{\rho^5}{5} \right] _0^4 \, d\varphi \\\\ &= 2 \pi \int _0^{\pi} \left( \frac{4^5}{5} \sin \varphi \right) \, d\varphi \\\\ &= 2 \pi \left( \frac{1024}{5} \right) \int _0^{\pi} \sin \varphi \, d\varphi \\\\ &= 2 \pi \left( \frac{1024}{5} \right) \left[ -\cos \varphi \right] _0^{\pi} \\\\ &= 2 \pi \left( \frac{1024}{5} \right) \left[ -\cos (\pi) + \cos (0) \right] \\\\ &= 2 \pi \left( \frac{1024}{5} \right) \left[ -(-1) + 1 \right] \\\\ &= 2 \pi \left( \frac{1024}{5} \right) \left[ 2 \right] \\\\ &= \frac{4096 \pi}{5} \end{aligned} $$



> \textbf{Example 4c.} The domain $D$ is the ball with center at the origin and radius $3a$,

>

> $$D = \left \\{ x^2 + y^2 + z^2 \le 9a^2 \right \\}$$

>

> and $f(x,y,z)=x^2+y^2$ is the function to integrate.

>

> Looking at the domain, it seems convenient to adopt the passage to spherical coordinates, in fact, the intervals of the variables that delimit the new $T$ region are obviously:

>

> $$T=\\{0 \le \rho \le 3a, \ 0 \le \varphi \le 2 \pi, \ 0 \le \theta \le \pi\\}.$$

>

> However, applying the transformation, we get

>

> $$f(x,y,z) = x^2 + y^2 \longrightarrow \rho^2 \sin^2 \theta \cos^2 \varphi + \rho^2 \sin^2 \theta \sin^2 \varphi = \rho^2 \sin^2 \theta.$$

>

> Applying the formula for integration we obtain:

>

> $$\iiint _T \rho^2 \sin^2 \theta \rho^2 \sin \theta \, d\rho\, d\theta\, d\varphi = \iiint _T \rho^4 \sin^3 \theta \, d\rho\, d\theta\, d\varphi$$

>

> which can be solved by turning it into an iterated integral.

>

> $\iiint _T \rho^4 \sin^3 \theta \, d\rho\, d\theta\, d\varphi = \underbrace{\int _0^{3a}\rho^4 d\rho} _{I} \,\underbrace{\int _0^\pi \sin^3\theta\,d\theta} _{II}\, \underbrace{\int _0^{2\pi} d \varphi} _{III}$.

>

> $I = \left.\int _0^{3a}\rho^4 d\rho = \frac{\rho^5}{5}\right\vert _0^{3a} = \frac{243}{5}a^5$,

>

> $$\begin{aligned} II &= \int _0^\pi \sin^3\theta \, d\theta \hspace{30em} \\\\ &= -\int _0^\pi \sin^2\theta \, d(\cos \theta) \\\\ &= \int _0^\pi (\cos^2\theta - 1) \, d(\cos \theta) \\\\ &= \left. \frac{\cos^3\theta}{3} \right| _0^\pi - \left. \cos\theta \right| _0^\pi \\\\ &= \left( \frac{\cos^3(\pi)}{3} - \frac{\cos^3(0)}{3} \right) - (\cos(\pi) - \cos(0)) \\\\ &= \left( \frac{(-1)^3}{3} - \frac{(1)^3}{3} \right) - (-1 - 1) \\\\ &= \left( -\frac{1}{3} - \frac{1}{3} \right) - (-2) \\\\ &= -\frac{2}{3} + 2 \\\\ &= -\frac{2}{3} + \frac{6}{3} \\\\ &= \frac{4}{3} \end{aligned},$$

>

> $III = \int _0^{2\pi} d \varphi = 2\pi$.

>

> Collecting all parts,

>

> $\iiint _T \rho^4 \sin^3 \theta \, d\rho\, d\theta\, d\varphi = I\cdot II\cdot III = \frac{243}{5}a^5\cdot \frac{4}{3}\cdot 2\pi = \frac{648}{5}\pi a^5$.

>

> Alternatively, this problem can be solved by using the passage to cylindrical coordinates. The new $T$ intervals are

>

> $$T=\left\\{0 \le \rho \le 3a, \ 0 \le \varphi \le 2 \pi, \ - \sqrt{9a^2 - \rho^2} \le z \le \sqrt{9a^2 - \rho^2}\right\\};$$

>

> the $z$ interval has been obtained by dividing the ball into two hemispheres simply by solving the inequality from the formula of $D$ (and then directly transforming $x^2 + y^2$ into $\rho^2$). The new function is simply $\rho^2$. Applying the integration formula:

>

> $$\iiint _T \rho^2 \rho \, d \rho \, d \varphi \, dz.$$

>

> Then we get

>

> $$\begin{aligned} \int _0^{2\pi} d\varphi \int _0^{3a} \rho^3 d\rho \int _{-\sqrt{9a^2 - \rho^2}}^{\sqrt{9 a^2 - \rho^2}}\, dz &= 2 \pi \int _0^{3a} 2 \rho^3 \sqrt{9 a^2 - \rho^2} \, d\rho \\\\ &= -2 \pi \int _{9 a^2}^0 (9 a^2 - t) \sqrt{t}\, dt \quad t = 9 a^2 - \rho^2 \\\\ &= 2 \pi \int _0^{9 a^2} \left( 9 a^2 \sqrt{t} - t \sqrt{t} \right) \, dt \\\\ &= 2 \pi \left( \int _0^{9 a^2} 9 a^2 \sqrt{t} \, dt - \int _0^{9 a^2} t \sqrt{t} \, dt \right) \\\\ &= 2 \pi \left[ 9 a^2 \frac{2}{3} t^{ \frac{3}{2} } - \frac{2}{5} t^{ \frac{5}{2} } \right] _0^{9 a^2} \\\\ &= 2 \cdot 27 \pi a^5 \left( 6 - \frac{18}{5} \right) \\\\ &= \frac{648 \pi}{5} a^5. \end{aligned}$$

>

> Thanks to the passage to cylindrical coordinates it was possible to reduce the triple integral to an easier one-variable integral.



\paragraph{Examples}



\paragraph{Double integral over a rectangle}



Let us assume that we wish to integrate a multivariable function $f$ over a region $A$:



$$A = \left \\{ (x,y) \in \mathbb{R} ^2 \ : \ 11 \le x \le 14 \ ; \ 7 \le y \le 10 \right \\} \text{ and } f(x,y) = x^2 + 4y\,$$



From this we formulate the iterated integral



$$\int _7^{10} \int _{11}^{14} (x^2 + 4y) \, dx\, dy$$



The inner integral is performed first, integrating with respect to $x$ and taking $y$ as a constant, as it is not the variable of integration. The result of this integral, which is a function depending only on $y$, is then integrated with respect to $y$.



$$\begin{aligned}

\int _{11}^{14} \left(x^2 + 4y\right) \, dx &= \left [\frac{1}{3} x^3 + 4yx \right] _{x=11}^{x=14} \\\&= \frac{1}{3}(14)^3 + 4y(14) - \frac{1}{3}(11)^3 - 4y(11) \\\&= 471 + 12y.

\end{aligned}$$





We then integrate the result with respect to $y$.



$$\begin{aligned}

\int _7^{10} (471 + 12y) \ dy & = \Big[471y + 6y^2\Big] _{y=7}^{y=10} \\\&= 471(10)+ 6(10)^2 - 471(7) - 6(7)^2 \\\&= 1719 \end{aligned}$$



In cases where the double integral of the absolute value of the function is finite, the order of integration is interchangeable, that is, integrating with respect to $x$ first and integrating with respect to $y$ first produce the same result. That is Fubini's theorem. For example, doing the previous calculation with order reversed gives the same result:



$$\begin{aligned}

\int _{11}^{14} \int _{7}^{10} \, \left(x^2 + 4y\right) \, dy\, dx & = \int _{11}^{14} \Big[x^2 y + 2y^2 \Big] _{y=7}^{y=10} \, dx \\\&= \int _{11}^{14} \, (3x^2 + 102) \, dx \\\&= \Big[x^3 + 102x \Big] _{x=11}^{x=14} \\\&= 1719. \end{aligned}$$



\paragraph{Double integral over a normal domain}



% [Image: A double integral over the normal region D]



Figure: A double integral over the normal region \textit{D}



Consider the region (please see the graphic in the example):



$$D = \\{ (x,y) \in \mathbb{R} ^2 \ : \ x \ge 0, y \le 1, y \ge x^2 \\}$$ Calculate



$$\iint _D (x+y) \, dx \, dy.$$



This domain is normal with respect to both the $x$- and $y$-axes. To apply the formulae it is required to find the functions that determine \textit{D} and the intervals over which these functions are defined. In this case the two functions are:



$$\alpha (x) = x^2\text{ and }\beta (x) = 1$$



while the interval is given by the intersections of the functions with $x$ = 0, so the interval is $[a, b] = [0, 1]$ (normality has been chosen with respect to the $x$-axis for a better visual understanding).



It is now possible to apply the formula:



$$\iint _D (x+y) \, dx \, dy = \int _0^1 dx \int _{x^2}^1 (x+y) \, dy = \int _0^1 dx \ \left[xy + \frac{y^2}{2} \right]^1 _{x^2}$$



(at first the second integral is calculated considering $x$ as a constant). The remaining operations consist of applying the basic techniques of integration:



$$\int _0^1 \left[xy + \frac{y^2}{2}\right]^1 _{x^2} \, dx = \int _0^1 \left(x + \frac{1}{2} - x^3 - \frac{x^4}{2} \right) dx = \cdots = \frac{13}{20}.$$



If we choose normality with respect to the $y$-axis we could calculate



$$\int _0^1 dy \int _0^{\sqrt{y}} (x+y) \, dx.$$



and obtain the same value.



% [Image: Example of domain in R3 that is normal with respect to the xy-plane.]



Figure: Example of domain in $\mathbb{R} ^3$ that is normal with respect to the $xy$-plane.



\paragraph{Calculating volume}



Using the methods previously described, it is possible to calculate the volumes of some common solids.



\begin{itemize}
    \item \textbf{Cylinder}: The volume of a cylinder with height $h$ and circular base of radius $R$ can be calculated by integrating the constant function $h$ over the circular base, using polar coordinates.
\end{itemize}

$$\mathrm{Volume} = \int _0^{2\pi} d \varphi\, \int _0^R h \rho \, d \rho = 2 \pi h \left[\frac{\rho^2}{2}\right] _0^R = \pi R^2 h$$



This is in agreement with the formula for the volume of a prism



$$\mathrm{Volume} = \text{base area} \times \text{height}.$$



\begin{itemize}
    \item \textbf{Sphere}: The volume of a sphere with radius $R$ can be calculated by integrating the constant function 1 over the sphere, using spherical coordinates.
\end{itemize}

$$\begin{aligned} \text{Volume} &= \iiint _D f(x,y,z) \, dx\, dy\, dz \\\&= \iiint _D 1 \, dV \\\&= \iiint _S \rho^2 \sin \varphi \, d\rho\, d\theta\, d\varphi \\\&= \int _0^{2\pi} \, d \theta \int _0^{ \pi } \sin \varphi\, d \varphi \int _0^R \rho^2\, d \rho \\\&= 2 \pi \int _0^\pi \sin \varphi\, d \varphi \int _0^R \rho^2\, d \rho \\\&= 2 \pi \int _0^\pi \sin \varphi \frac{R^3}{3 }\, d \varphi \\\&= \frac23 \pi R^3 \Big[-\cos \varphi\Big] _0^\pi = \frac43 \pi R^3.

\end{aligned}$$



\begin{itemize}
    \item \textbf{Tetrahedron} (triangular pyramid or 3-simplex): The volume of a tetrahedron with its apex at the origin and edges of length $\ell$ along the $x$-, $y$-, and $z$-axes can be calculated by integrating the constant function 1 over the tetrahedron.
\end{itemize}

$$\begin{aligned} \text{Volume} &= \int _0^\ell dx \int _0^{\ell-x}\, dy \int _0^{\ell-x-y }\, dz \\\&= \int _0^\ell dx \int _0^{\ell-x } (\ell - x - y)\, dy \\\&= \int _0^\ell \left( l^2 - 2 \ell x + x^2 - \frac{(\ell-x)^2 }{2}\right)\, dx \\\&= \ell^3 - \ell \ell^2 + \frac{\ell^3}{3 } - \left[\frac{\ell^2 x}{2} - \frac{ \ell x^2}{2} + \frac{x^3}{6 }\right] _0^ \ell \\\&= \frac{\ell^3}{3} - \frac{\ell^3}{6} = \frac{ \ell^3}{6}\end{aligned}$$



> This is in agreement with the formula for the volume of a pyramid



$$\mathrm{Volume} = \frac13 \times \text{base area} \times \text{height} = \frac13 \times \frac{\ell^2}{2} \times \ell = \frac{ \ell^3}{6}.$$



% [Image: Example of an improper domain.]



Figure: Example of an improper domain.



\paragraph{Multiple improper integral}



In case of unbounded domains or functions not bounded near the boundary of the domain, we have to introduce the \textbf{double improper integral} or the \textbf{triple improper integral}.



\paragraph{Multiple integrals and iterated integrals}



Fubini's theorem states that if



$$\iint _{A\times B} \left|f(x,y)\right|\,d(x,y)<\infty,$$ that is, if the integral is absolutely convergent, then the multiple integral will give the same result as either of the two iterated integrals:



$$\iint _{A\times B} f(x,y)\,d(x,y)=\int _A\left(\int _B f(x,y)\,dy\right)\,dx=\int _B\left(\int _A f(x,y)\,dx\right)\,dy.$$ In particular this will occur if $|f(x,y)|$ is a bounded function and $A$ and $B$ are bounded sets.



If the integral is not absolutely convergent, care is needed not to confuse the concepts of \textit{multiple integral} and \textit{iterated integral}, especially since the same notation is often used for either concept. The notation



$$\int _0^1\int _0^1 f(x,y)\,dy\,dx$$



means, in some cases, an iterated integral rather than a true double integral. In an iterated integral, the outer integral



$$\int _0^1 \cdots \, dx$$



is the integral with respect to $x$ of the following function of $x$:



$$g(x)=\int _0^1 f(x,y)\,dy.$$



A double integral, on the other hand, is defined with respect to area in the $xy$-plane. If the double integral exists, then it is equal to each of the two iterated integrals (either "$dy \, dx$" or "$dx \, dy$") and one often computes it by computing either of the iterated integrals. But sometimes the two iterated integrals exist when the double integral does not, and in some such cases the two iterated integrals are different numbers, i.e., one has



$$\int _0^1\int _0^1 f(x,y)\,dy\,dx \neq \int _0^1\int _0^1 f(x,y)\,dx\,dy.$$



This is an instance of rearrangement of a conditionally convergent integral.



On the other hand, some conditions ensure that the two iterated integrals are equal even though the double integral need not exist. By the Fichtenholz–Lichtenstein theorem: If $f$ is bounded on $[0, 1] \times [0, 1]$ and both iterated integrals exist, then they are equal. Moreover, existence of the inner integrals ensures existence of the outer integrals.



The notation



$$\int _{[0,1]\times[0,1]} f(x,y)\,dx\,dy$$



may be used if one wishes to be emphatic about intending a double integral rather than an iterated integral.



\paragraph{Some practical applications}



Quite generally, just as in one variable, one can use the multiple integral to find the average of a function over a given set. Given a set $D \subseteq \mathbb{R}^n$ and an integrable function $f$ over $D$, the average value of $f$ over its domain is given by



$$\bar{f} = \frac{1}{m(D)} \int _D f(x)\, dx,$$



where $m(D)$ is the measure of $D$.



Additionally, multiple integrals are used in many applications in physics. The examples below also show some variations in the notation.



In mechanics, the moment of inertia is calculated as the volume integral (triple integral) of the density weighed with the square of the distance from the axis:



$$I _z = \iiint _V \rho r^2\, dV.$$



The gravitational potential associated with a mass distribution given by a mass measure $dm$ on three-dimensional Euclidean space $\mathbb{R}^3$ is



$$V(\mathbf{x}) = -\iiint _{\mathbb{R} ^3} \frac{G}{|\mathbf{x} - \mathbf{y}|}\,dm(\mathbf{y}).$$



If there is a continuous function $\rho(\mathbf{x})$ representing the density of the distribution at $\mathbf{x}$, so that $dm(\mathbf{x}) = \rho (\mathbf{x})d^3\mathbf{x}$, where $d^3\mathbf{x}$ is the Euclidean volume element, then the gravitational potential is



$$V(\mathbf{x}) = -\iiint _{\mathbb{R} ^3} \frac{G}{|\mathbf{x}-\mathbf{y}|}\,\rho(\mathbf{y})\,d^3\mathbf{y}.$$



In electromagnetism, Maxwell's equations can be written using multiple integrals to calculate the total magnetic and electric fields. In the following example, the electric field produced by a distribution of charges given by the volume charge density $\rho\vec{r}$ is obtained by a \textit{triple integral} of a vector function:



$$\vec E = \frac {1}{4 \pi \varepsilon _0} \iiint \frac {\vec r - \vec r'}{\left \| \left \| \vec r - \vec r' \right \| \right \| ^3} \rho (\vec r')\, d^3 r'.$$



This can also be written as an integral with respect to a signed measure representing the charge distribution.



\paragraph{Licensing}



Content obtained and/or adapted from:



\begin{itemize}
    \item \href{https://en.wikipedia.org/wiki/Multiple_integral}{Multiple integral, Wikipedia} under a CC BY-SA license
\end{itemize}

\subsection{Learning Objectives}

\begin{enumerate}
    \item \textbf{Calculate Area Between Two Curves:}
    \item \textbf{Outcome:} Students will be able to calculate the area between two curves using definite integrals.
    \item \textbf{Details:} This includes setting up and evaluating the definite integral of the difference between two functions over a specified interval, understanding the geometric interpretation, and handling cases where the curves intersect.
    \item \textbf{Evaluate Double Integrals:}
    \item \textbf{Outcome:} Students will be able to evaluate double integrals to find the volume under a surface.
    \item \textbf{Details:} This includes setting up double integrals in Cartesian coordinates, interpreting the integral as a volume, and using iterated integration to compute the integral over rectangular and more general domains.
    \item \textbf{Transform Double Integrals to Polar Coordinates:}
    \item \textbf{Outcome:} Students will be able to convert and evaluate double integrals in polar coordinates.
    \item \textbf{Details:} This includes understanding the relationship between Cartesian and polar coordinates, converting the integral limits and integrand, and computing the integral for regions with circular symmetry.
\end{enumerate}

\subsection{Assessment Instrument}

\begin{enumerate}
    \item \textbf{Assessment Instrument: Double and Triple Integrals over General Regions Quiz} \textbf{Format:}
    \item \textbf{Type:} Mixed-format quiz (multiple-choice, short answer, and problem-solving questions)
    \item \textbf{Duration:} 45 minutes
    \item \textbf{Total Marks:} 30 points \textbf{Sections:}
    \item \textbf{Understanding Double Integrals Over General Regions (10 points):}
    \item \textbf{Question 1 (Multiple Choice):} Which of the following statements is true about double integrals over general regions?
    \item a. Double integrals over general regions can only be computed using rectangular coordinates.
    \item b. The order of integration in double integrals can be interchanged if the function is continuous over the region.
    \item c. Double integrals over general regions are always easier to compute than over rectangular regions.
    \item d. The double integral over a general region is always zero. (2 points)
    \item \textbf{Question 2 (Short Answer):} Explain how to set up a double integral over a region bounded by two functions $y = g_1(x)$ and $y = g_2(x)$ from $x = a$ to $x = b$. (4 points)
    \item \textbf{Question 3 (Problem Solving):} Set up the double integral to find the area between the curves $y = x^2$ and $y = x + 2$ over the interval $x \in [0,2]$. Do not evaluate the integral. (4 points)
    \item \textbf{Evaluating Triple Integrals (10 points):}
    \item \textbf{Question 4 (Multiple Choice):} In which coordinate system is it most appropriate to evaluate the triple integral for a region with spherical symmetry?
    \item a. Cartesian coordinates
    \item b. Polar coordinates
    \item c. Cylindrical coordinates
    \item d. Spherical coordinates (2 points)
    \item \textbf{Question 5 (Short Answer):} Describe the process of converting a triple integral from Cartesian to cylindrical coordinates. (4 points)
    \item \textbf{Question 6 (Problem Solving):} Evaluate the triple integral $\iiint_V z \, dV$ where $V$ is the region bounded by the paraboloid $z = x^2 + y^2$ and the plane $z = 4$. (4 points)
    \item \textbf{Applications of Double and Triple Integrals (10 points):}
    \item \textbf{Question 7 (Multiple Choice):} Which of the following represents the volume of a region bounded by the cylinder $x^2 + y^2 = 4$ and the planes $z = 0$ and $z = 5$?
    \item a. $\iiint_V 1 \, dV$ in Cartesian coordinates
    \item b. $\iiint_V 1 \, dV$ in cylindrical coordinates
    \item c. $\iint_R 5 \, dA$ in cylindrical coordinates
    \item d. $\iiint_V 5 \, dV$ in spherical coordinates (2 points)
    \item \textbf{Question 8 (Short Answer):} Provide an example of a physical problem that can be solved using triple integrals and explain the setup. (4 points)
    \item \textbf{Question 9 (Problem Solving):} Given the region $R$ bounded by $y = 0$, $y = \sqrt{x}$, and $x = 1$, set up and evaluate the double integral $\iint_R (x+y) \, dA$. Show all steps. (4 points) \textbf{Scoring Rubric:}
    \item \textbf{Understanding Double Integrals Over General Regions:} 2 points for correct multiple-choice answer, 4 points for clear and accurate short answer, and 4 points for correctly set up problem-solving question (total 10 points)
    \item \textbf{Evaluating Triple Integrals:} 2 points for correct multiple-choice answer, 4 points for clear and accurate short answer, and 4 points for correct and fully solved problem-solving question (total 10 points)
    \item \textbf{Applications of Double and Triple Integrals:} 2 points for correct multiple-choice answer, 4 points for clear and relevant short answer, and 4 points for correct and fully solved problem-solving question (total 10 points) The quiz will assess the students' understanding and application of double and triple integrals over general regions, ensuring they can correctly interpret, set up, and solve relevant integrals in various coordinate systems.
\end{enumerate}

%%===============================================================
\section{Vector Fields}
\label{sec:vector-fields}
%%===============================================================

In vector calculus and physics, a \textbf{vector field} is an assignment of a vector to each point in a subset of space. For instance, a vector field in the plane can be visualised as a collection of arrows with a given magnitude and direction, each attached to a point in the plane. Vector fields are often used to model, for example, the speed and direction of a moving fluid throughout space, or the strength and direction of some force, such as the magnetic or gravitational force, as it changes from one point to another point.



The elements of differential and integral calculus extend naturally to vector fields. When a vector field represents force, the line integral of a vector field represents the work done by a force moving along a path, and under this interpretation conservation of energy is exhibited as a special case of the fundamental theorem of calculus. Vector fields can usefully be thought of as representing the velocity of a moving flow in space, and this physical intuition leads to notions such as the divergence (which represents the rate of change of volume of a flow) and curl (which represents the rotation of a flow).



In coordinates, a vector field on a domain in n-dimensional Euclidean space can be represented as a vector-valued function that associates an n-tuple of real numbers to each point of the domain. This representation of a vector field depends on the coordinate system, and there is a well-defined transformation law in passing from one coordinate system to the other. Vector fields are often discussed on open subsets of Euclidean space, but also make sense on other subsets such as surfaces, where they associate an arrow tangent to the surface at each point (a tangent vector).



More generally, vector fields are defined on differentiable manifolds, which are spaces that look like Euclidean space on small scales, but may have more complicated structure on larger scales. In this setting, a vector field gives a tangent vector at each point of the manifold (that is, a section of the tangent bundle to the manifold). Vector fields are one kind of tensor field.



\paragraph{Fields in vector calculus}



% [Image: A depiction of xyz Cartesian coordinates with the ijk elementary basis vectors.]



Figure: A depiction of xyz Cartesian coordinates with the $ijk$ elementary basis vectors.



\paragraph{Scalar fields}



A scalar field is a function $f: \mathbb{R} ^3 \to \mathbb{R}$ that assigns a real number to each point in space. Scalar fields typically denote densities or potentials at each specific point. For the sake of simplicity, all scalar fields considered by this chapter will be assumed to be defined at all points and differentiable at all points.



\paragraph{Vector fields}



A vector field is a function $\mathbf{F}: \mathbb{R} ^3 \to \mathbb{R} ^3$ that assigns a vector to each point in space. Vector fields typically denote flow densities or potential gradients at each specific point. For the sake of simplicity, all vector fields considered by this chapter will be assumed to be defined at all points and differentiable at all points.



% [Image: A depiction of cylindrical coordinates and the accompanying orthonormal basis vectors.]



Figure: A depiction of cylindrical coordinates and the accompanying orthonormal basis vectors.



\paragraph{Vector fields in cylindrical coordinates}



The cylindrical coordinate system used here has the three parameters: $(\rho,\phi,z)$. The Cartesian coordinate equivalent to the point $(\rho,\phi,z)$ is



$x = \rho\cos\phi$



$y = \rho\sin\phi$



$z = z$



Any vector field in cylindrical coordinates is a linear combination of the following 3 mutually orthogonal unit length basis vectors:



$\hat{\mathbf{\rho}} = \cos\phi\mathbf{i} + \sin\phi\mathbf{j}$



$\hat{\mathbf{\phi}} = -\sin\phi\mathbf{i} + \cos\phi\mathbf{j}$



$\hat{\mathbf{z}} = \mathbf{k}$



Note that these basis vectors are not constant with respect to position. The fact that the basis vectors change from position to position should always be considered. The cylindrical basis vectors change according to the following rates:



$$\begin{array}{c|c|c}

 & \frac{\partial}{\partial \rho} & \frac{\partial}{\partial \phi} & \frac{\partial}{\partial z} \\\\ 

\hline

\hat{\mathbf{\rho}} & \frac{\partial \hat{\mathbf{\rho}}}{\partial \rho} = \mathbf{0} & \frac{\partial \hat{\mathbf{\rho}}}{\partial \phi} = \hat{\mathbf{\phi}} & \frac{\partial \hat{\mathbf{\rho}}}{\partial z} = \mathbf{0} \\\\ \hline

\hat{\mathbf{\phi}} & \frac{\partial \hat{\mathbf{\phi}}}{\partial \rho} = \mathbf{0} & \frac{\partial \hat{\mathbf{\phi}}}{\partial \phi} = -\hat{\mathbf{r}} & \frac{\partial \hat{\mathbf{\phi}}}{\partial z} = \mathbf{0} \\\\ \hline

\hat{\mathbf{z}} & \frac{\partial \hat{\mathbf{z}}}{\partial \rho} = \mathbf{0} & \frac{\partial \hat{\mathbf{z}}}{\partial \phi} = \mathbf{0} & \frac{\partial \hat{\mathbf{z}}}{\partial z} = \mathbf{0} \\\\ 

\end{array}$$



Any vector field $\mathbf{F}$ expressed in cylindrical coordinates has the form:



$\mathbf{F}(\mathbf{q}) = F _\rho(\mathbf{q})\hat{\mathbf{\rho}} + F _\phi(\mathbf{q})\hat{\mathbf{\phi}} + F _z(\mathbf{q})\hat{\mathbf{z}}$



Given an arbitrary position $\mathbf{q} = (\rho,\phi,z)$ that changes with time, the velocity of the position is:



$$ \frac{d \mathbf{q}}{dt} = \frac{d \rho}{dt}\hat{\mathbf{\rho}} + \rho\frac{d \phi}{dt}\hat{\mathbf{\phi}} + \frac{dz}{dt}\hat{\mathbf{z}} $$



The coefficient of $\rho$ for the term $\rho\frac{d \phi}{dt}\hat{\mathbf{\phi}}$ originates from the fact that as the azimuth angle $\phi$ increases, the position $\mathbf{q}$ swings around at a speed of $\rho$.



% [Image: A depiction of spherical coordinates and the accompanying orthonormal basis vectors.]



Figure: A depiction of spherical coordinates and the accompanying orthonormal basis vectors.



\paragraph{Vector fields in spherical coordinates}



The spherical coordinate system used here has the three parameters: $(r,\theta,\phi)$. The Cartesian coordinate equivalent to the point $(r,\theta,\phi)$ is



$x = r\sin\theta\cos\phi$



$y = r\sin\theta\sin\phi$



$z = r\cos\theta$



Any vector field in spherical coordinates is a linear combination of the following 3 mutually orthogonal unit length basis vectors:



$\hat{\mathbf{r}} = \sin\theta\cos\phi\mathbf{i} + \sin\theta\sin\phi\mathbf{j} + \cos\theta\mathbf{k}$



$\hat{\mathbf{\theta}} = \cos\theta\cos\phi\mathbf{i} + \cos\theta\sin\phi\mathbf{j} - \sin\theta\mathbf{k}$



$\hat{\mathbf{\phi}} = -\sin\phi\mathbf{i} + \cos\phi\mathbf{j}$



Note that these basis vectors are not constant with respect to position. The fact that the basis vectors change from position to position should always be considered. The spherical basis vectors change according to the following rates:



$$\begin{array}{c|c|c|c}

 & \frac{\partial}{\partial r} & \frac{\partial}{\partial \theta} & \frac{\partial}{\partial \phi} \\\\ 

\hline

\hat{\mathbf{r}} & \frac{\partial \hat{\mathbf{r}}}{\partial r} = \mathbf{0} & \frac{\partial \hat{\mathbf{r}}}{\partial \theta} = \hat{\mathbf{\theta}} & \frac{\partial \hat{\mathbf{r}}}{\partial \phi} = \sin\theta \hat{\mathbf{\phi}} \\\\ \hline

\hat{\mathbf{\theta}} & \frac{\partial \hat{\mathbf{\theta}}}{\partial r} = \mathbf{0} & \frac{\partial \hat{\mathbf{\theta}}}{\partial \theta} = -\hat{\mathbf{r}} & \frac{\partial \hat{\mathbf{\theta}}}{\partial \phi} = \cos\theta \hat{\mathbf{\phi}} \\\\ \hline

\hat{\mathbf{\phi}} & \frac{\partial \hat{\mathbf{\phi}}}{\partial r} = \mathbf{0} & \frac{\partial \hat{\mathbf{\phi}}}{\partial \theta} = \mathbf{0} & \frac{\partial \hat{\mathbf{\phi}}}{\partial \phi} = -(\sin\theta \hat{\mathbf{r}} + \cos\theta \hat{\mathbf{\theta}}) \\\\ 

\end{array}$$



Any vector field $\mathbf{F}$ expressed in spherical coordinates has the form: $\mathbf{F}(\mathbf{q}) = F _r(\mathbf{q})\hat{\mathbf{r}} + F _{\theta}(\mathbf{q})\hat{\mathbf{\theta}} + F _{\phi}(\mathbf{q})\hat{\phi}$



Given an arbitrary position $\mathbf{q} = (r,\theta,\phi)$ that changes with time, the velocity of this position is:



$$\frac{d\mathbf{q}}{dt} = \frac{dr}{dt}\hat{\mathbf{r}} + r\frac{d\theta}{dt}\hat{\mathbf{\theta}} + r\sin\theta\frac{d\phi}{dt}\hat{\mathbf{\phi}}$$



The coefficient of $r$ for the term $r\frac{d\theta}{dt}\hat{\mathbf{\theta}}$ arises from the fact that as the latitudinal angle $\theta$ changes, the position $\mathbf{q}$ traverses a great circle at a speed of $r$.



The coefficient of $r\sin\theta$ for the term $r\sin\theta\frac{d\phi}{dt}\hat{\mathbf{\phi}}$ arises from the fact that as the longitudinal angle $\phi$ changes, the position $\mathbf{q}$ traverses a latitude circle at a speed of $r\sin\theta$.



\paragraph{Operations on vector fields}



\paragraph{Line integral}



A common technique in physics is to integrate a vector field along a curve, also called determining its line integral. Intuitively this is summing up all vector components in line with the tangents to the curve, expressed as their scalar products. For example, given a particle in a force field (e.g. gravitation), where each vector at some point in space represents the force acting there on the particle, the line integral along a certain path is the work done on the particle, when it travels along this path. Intuitively, it is the sum of the scalar products of the force vector and the small tangent vector in each point along the curve.



The line integral is constructed analogously to the Riemann integral and it exists if the curve is rectifiable (has finite length) and the vector field is continuous.



Given a vector field $\mathbf{V}$ and a curve $\gamma$, parametrized by $t$ in $[a, b]$ (where $a$ and $b$ are real numbers), the line integral is defined as



$$\int _{\gamma}  V(\boldsymbol {x})\cdot \mathrm{d}\boldsymbol {x}  = \int _a^b  V(\gamma(t))\cdot\dot \gamma(t)\; \mathrm{d}t.$$



\paragraph{Divergence}



The divergence of a vector field on Euclidean space is a function (or scalar field). In three-dimensions, the divergence is defined by



$$\operatorname{div} \mathbf{F} = \nabla \cdot \mathbf{F} = \frac{\partial F _1}{\partial x} + \frac{\partial F _2}{\partial y}+\frac{\partial F _3}{\partial z},$$



with the obvious generalization to arbitrary dimensions. The divergence at a point represents the degree to which a small volume around the point is a source or a sink for the vector flow, a result which is made precise by the divergence theorem.



The divergence can also be defined on a Riemannian manifold, that is, a manifold with a Riemannian metric that measures the length of vectors.



\paragraph{Curl in three dimensions}



The curl is an operation which takes a vector field and produces another vector field. The curl is defined only in three dimensions, but some properties of the curl can be captured in higher dimensions with the exterior derivative. In three dimensions, it is defined by



$$\operatorname{curl}\,\mathbf{F} = \nabla \times \mathbf{F} = \left(\frac{\partial F _3}{\partial y}- \frac{\partial F _2}{\partial z}\right)\mathbf{e} _1 - \left(\frac{\partial F _3}{\partial x}- \frac{\partial F _1}{\partial z}\right)\mathbf{e} _2 + \left(\frac{\partial F _2}{\partial x}- \frac{\partial F _1}{\partial y}\right)\mathbf{e} _3.$$



The curl measures the density of the angular momentum of the vector flow at a point, that is, the amount to which the flow circulates around a fixed axis. This intuitive description is made precise by Stokes\' theorem.



\paragraph{Index of a vector field}



The index of a vector field is an integer that helps to describe the behaviour of a vector field around an isolated zero (i.e., an isolated singularity of the field). In the plane, the index takes the value -1 at a saddle singularity but +1 at a source or sink singularity.



Let the dimension of the manifold on which the vector field is defined be $n$. Take a small sphere $S$ around the zero so that no other zeros lie in the interior of $S$. A map from this sphere to a unit sphere of dimension $n - 1$ can be constructed by dividing each vector on this sphere by its length to form a unit length vector, which is a point on the unit sphere $S^{n-1}$. This defines a continuous map from $S$ to $S^{n-1}$. The index of the vector field at the point is the degree of this map. It can be shown that this integer does not depend on the choice of $S$, and therefore depends only on the vector field itself.



\textbf{The index of the vector field} as a whole is defined when it has just a finite number of zeroes. In this case, all zeroes are isolated, and the index of the vector field is defined to be the sum of the indices at all zeroes.



The index is not defined at any non-singular point (i.e., a point where the vector is non-zero). It is equal to $+1$ around a source, and more generally equal to $(-1)^k$ around a saddle that has $k$ contracting dimensions and $n-k$ expanding dimensions. For an ordinary (2-dimensional) sphere in three-dimensional space, it can be shown that the index of any vector field on the sphere must be 2. This shows that every such vector field must have a zero. This implies the hairy ball theorem, which states that if a vector in $\mathbb{R}^3$ is assigned to each point of the unit sphere $S^2$ in a continuous manner, then it is impossible to "comb the hairs flat," i.e., to choose the vectors in a continuous way such that they are all non-zero and tangent to $S^2$.



For a vector field on a compact manifold with a finite number of zeroes, the Poincaré-Hopf theorem states that the index of the vector field is equal to the Euler characteristic of the manifold.



\paragraph{Resources}



\begin{itemize}
    \item \href{https://en.wikibooks.org/wiki/Calculus/Vector_calculus}{Vector Calculus}, Wikibooks: Calculus
    \item \href{https://math.libretexts.org/Bookshelves/Calculus/Book\%3A_Calculus_(OpenStax}{Vector Fields}/16\%3A_Vector_Calculus/16.1\%3A_Vector_Fields), Mathematics LibreTexts
    \item \href{https://www.khanacademy.org/math/multivariable-calculus/thinking-about-multivariable-function/visualizing-vector-valued-functions/v/vector-fields-introduction}{Introduction to Vector Fields}, Khan Academy
    \item \href{https://www.khanacademy.org/math/multivariable-calculus/integrating-multivariable-functions/line-integrals-vectors/v/line-integrals-and-vector-fields}{Line Integrals and Vector Fields}, Khan Academy
\end{itemize}

\paragraph{Licensing}



Content obtained and/or adapted from:



\begin{itemize}
    \item \href{https://en.wikipedia.org/wiki/Vector_field}{Vector field, Wikipedia} under a CC BY-SA license
    \item \href{https://en.wikibooks.org/wiki/Calculus/Vector_calculus}{Vector calculus, Wikibooks: Calculus} under a CC BY-SA license
\end{itemize}

\subsection{Learning Objectives}

\begin{enumerate}
    \item Here are three Student Learning Objectives (SLOs) related to vector fields in vector calculus and physics:
    \item \textbf{Understand and Describe Vector Fields:}
    \item \textbf{Outcome:} Students will be able to define a vector field and describe how it assigns a vector to each point in a subset of space.
    \item \textbf{Details:} This includes recognizing examples of vector fields in both two-dimensional and three-dimensional spaces and visualizing them using diagrams.
    \item \textbf{Calculate Line Integrals in Vector Fields:}
    \item \textbf{Outcome:} Students will be able to compute the line integral of a vector field along a given curve in space.
    \item \textbf{Details:} This includes applying the concept of work done by a force along a path and using the parametric representation of curves to evaluate line integrals.
    \item \textbf{Apply Divergence and Curl to Analyze Vector Fields:}
    \item \textbf{Outcome:} Students will be able to calculate the divergence and curl of a vector field and interpret their physical significance.
    \item \textbf{Details:} This includes understanding divergence as a measure of a source or sink at a point and curl as a measure of rotation or angular momentum in a vector field.
\end{enumerate}

\subsection{Assessment Instrument}

\begin{enumerate}
    \item \textbf{Assessment Instrument: Vector Fields Quiz} \textbf{Format:}
    \item \textbf{Type:} Mixed-format quiz (multiple-choice, short answer, and problem-solving questions)
    \item \textbf{Duration:} 45 minutes
    \item \textbf{Total Marks:} 30 points \textbf{Sections:}
    \item \textbf{Understanding Vector Fields (10 points):}
    \item \textbf{Question 1:} Define a vector field and provide a physical example where vector fields are used.
    \item \textbf{(5 points)}
    \item \textbf{Question 2:} Consider a vector field in the plane represented by arrows. Which of the following best describes the nature of these arrows?
    \item a. They point in arbitrary directions.
    \item b. They all point in the same direction.
    \item c. Each arrow points towards a common center.
    \item d. Each arrow represents both magnitude and direction at its point of origin.
    \item \textbf{(5 points)}
    \item \textbf{Line Integrals in Vector Fields (10 points):}
    \item \textbf{Question 3:} Compute the line integral of the vector field $\mathbf{F}(x, y) = x\mathbf{i} + y\mathbf{j}$ along the curve $C$ given by the path $\gamma(t) = (t, t^2)$, where $t$ ranges from 0 to 1.
    \item \textbf{(10 points)}
    \item \textbf{Divergence and Curl in Vector Fields (10 points):}
    \item \textbf{Question 4:} Calculate the divergence of the vector field $\mathbf{F}(x, y, z) = x^2\mathbf{i} + yz\mathbf{j} + z^2\mathbf{k}$.
    \item \textbf{(5 points)}
    \item \textbf{Question 5:} Determine the curl of the vector field $\mathbf{F}(x, y, z) = -y\mathbf{i} + x\mathbf{j} + 0\mathbf{k}$.
    \item \textbf{(5 points)} \textbf{Scoring Rubric:}
    \item \textbf{Understanding Vector Fields:}
    \item Definition and physical example: 5 points
    \item Multiple-choice question on the nature of arrows: 5 points
    \item \textbf{Line Integrals in Vector Fields:}
    \item Correct computation of the line integral: 10 points
    \item \textbf{Divergence and Curl in Vector Fields:}
    \item Correct calculation of divergence: 5 points
    \item Correct determination of curl: 5 points The quiz is designed to assess students' understanding of vector fields, their ability to compute line integrals, and their grasp of key operations such as divergence and curl.
\end{enumerate}

%%===============================================================
\section{Line Integrals}
\label{sec:line-integrals}
%%===============================================================

In mathematics, a \textbf{line integral} is an integral where the function to be integrated is evaluated along a curve. The terms \textit{path integral}, \textit{curve integral}, and \textit{curvilinear integral} are also used; \textit{contour integral} is used as well, although that is typically reserved for line integrals in the complex plane.



The function to be integrated may be a scalar field or a vector field. The value of the line integral is the sum of values of the field at all points on the curve, weighted by some scalar function on the curve (commonly arc length or, for a vector field, the scalar product of the vector field with a differential vector in the curve). This weighting distinguishes the line integral from simpler integrals defined on intervals. Many simple formulae in physics, such as the definition of work as $W=\mathbf{F}\cdot\mathbf{s}$, have natural continuous analogues in terms of line integrals, in this case $\textstyle W = \int _L \mathbf{F}(\mathbf{s})\cdot d\mathbf{s}$, which computes the work done on an object moving through an electric or gravitational field \textbf{F} along a path $L$.



\paragraph{Volume, path (line), and surface integrals}



\paragraph{Volume Integrals}



Given a scalar field $\rho: \mathbb{R} ^3 \to \mathbb{R}$ that denotes a density at each specific point, and an arbitrary volume $\Omega \subseteq \mathbb{R} ^3$, the total "mass" $M$ inside of $\Omega$ can be determined by partitioning $\Omega$ into infinitesimal volumes. At each position $\mathbf{q} \in \Omega$, the volume of the infinitesimal volume is denoted by the infinitesimal $dV$. This gives rise to the following integral:



$$M = \iiint_{\mathbf{q} \in \Omega} \rho(\mathbf{q})dV$$



\paragraph{Path Integrals}



Given any oriented path $C$ (oriented means that there is a preferred direction), the differential $d\mathbf{q} = dx\mathbf{i} + dy\mathbf{j} + dz\mathbf{k}$ denotes an infinitesimal displacement along $C$ in the preferred direction. This differential can be used in various path integrals. Letting $f: \mathbb{R} ^3 \to \mathbb{R}$ denote an arbitrary scalar field, and $\mathbf{F}: \mathbb{R} ^3 \to \mathbb{R} ^3$ denote an arbitrary vector field, various path integrals include:



$\int _{\mathbf{q} \in C} f(\mathbf{q})d\mathbf{q}$, $\int _{\mathbf{q} \in C} f(\mathbf{q})|d\mathbf{q}|$, $\int _{\mathbf{q} \in C} \mathbf{F}(\mathbf{q}) \cdot d\mathbf{q}$, $\int _{\mathbf{q} \in C} \mathbf{F}(\mathbf{q})|d\mathbf{q}|$, and many more.



$\int _{\mathbf{q} \in C} d\mathbf{q}$ denotes the total displacement along $C$, and $\int _{\mathbf{q} \in C} |d\mathbf{q}|$ denotes the total length of $C$.



\paragraph{Calculating Path Integrals}



To compute a path integral, the continuous oriented curve $C$ must be parameterized. $\mathbf{q} _C(t)$ will denote the point along $C$ indexed by $t$ from the range $[t _0, t _1]$. $\mathbf{q} _C(t _0) = \mathbf{q} _0$ must be the starting point of $C$ and $\mathbf{q} _C(t _1) = \mathbf{q} _1$ must be the ending point of $C$. As $t$ increases, $\mathbf{q} _C(t)$ must proceed along $C$ in the preferred direction. An infinitesimal change in $t$, $dt$, results in the infinitesimal displacement $d\mathbf{q} = \frac{d\mathbf{q} _C}{dt}dt$ along $C$. In the path integral $\int _{\mathbf{q} \in C} \mathbf{F}(\mathbf{q}) \cdot d\mathbf{q}$, the differential $d\mathbf{q}$ can be replaced with $\frac{d\mathbf{q} _C}{dt}dt$ to get

$$\int _{t = t _0}^{t _1} \mathbf{F}(\mathbf{q} _C(t)) \cdot \frac{d\mathbf{q} _C}{dt}dt$$



> Example 1: As an example, consider the vector field $$\mathbf{F}(x,y,z) = 3\mathbf{i} - x\mathbf{j} + 5y\mathbf{k}$$ and the straight line curve $C$ that starts at $(1,1,1)$ and ends at $(7,-1,-2)$. $C$ can be parameterized by

> $$\mathbf{q} _C(t) = (1+6t, 1-2t, 1-3t) \text{ where } t \in [0, 1]$$

> We note that the partial derivative is $$\frac{d\mathbf{q} _C}{dt} = 6\mathbf{i} - 2\mathbf{j} - 3\mathbf{k}$$

> $\newline$ We can then evaluate the path integral $\int _{\mathbf{q} \in C} \mathbf{F}(\mathbf{q}) \cdot d\mathbf{q}$ as follows:

>

> $\begin{aligned} \int _{\mathbf{q} \in C} \mathbf{F}(\mathbf{q}) \cdot d\mathbf{q} &= \int _{\mathbf{q} \in C} (3\mathbf{i} - x\mathbf{j} + 5y\mathbf{k}) \cdot d\mathbf{q} \\\\ &= \int _{t=0}^{1} \left(3\mathbf{i} - (1+6t)\mathbf{j} + 5(1-2t)\mathbf{k}\right) \cdot (6\mathbf{i} - 2\mathbf{j} - 3\mathbf{k}) \, dt \\\\ &= \int _{t=0}^{1} \left(18 + (2+12t) + (-15+30t)\right) \, dt \\\\ &= \int _{t=0}^{1} \left(5 + 42t\right) \, dt \\\\ &= \left(5t + 21t^2\right) \bigg| _{t=0}^1 \\\\ &= 26 \end{aligned}$



If a vector field $\mathbf{F}$ denotes a "force field", which returns the force on an object as a function of position, the work performed on a point mass that traverses the oriented curve $C$ is $W = \int _{\mathbf{q} \in C} \mathbf{F}(\mathbf{q}) \cdot d\mathbf{q}$



> Example 2: Consider the gravitational field that surrounds a point mass of $M$ located at the origin: $$\mathbf{g}(\mathbf{q}) = -\frac{GM}{|\mathbf{q}|^2}\frac{\mathbf{q}}{|\mathbf{q}|}$$ using Newton's inverse square law. The force acting on a point mass of $m$ at position $\mathbf{q}$ is $$\mathbf{F}(\mathbf{q}) = m\mathbf{g}(\mathbf{q}) = -\frac{GMm}{|\mathbf{q}|^2}\frac{\mathbf{q}}{|\mathbf{q}|}$$. In spherical coordinates the force is $$\mathbf{F}(r,\theta,\phi) = -\frac{GMm}{r^2}\hat{\mathbf{r}}$$ (note that $\hat{\mathbf{r}},

>  \hat{\mathbf{\theta}}, \hat{\mathbf{\phi}}$ are the unit length mutually orthogonal basis vectors for spherical coordinates).

>

> Consider an arbitrary path $C$ that $m$ traverses that starts at an altitude of $r = r _1$ and ends at an altitude of $r = r _2$. The work done by the gravitational field is:

>

> $W = \int _{\mathbf{q} \in C} \mathbf{F}(\mathbf{q}) \cdot d\mathbf{q}$

>

> The infinitesimal displacement $d\mathbf{q}$ is equivalent to the displacement expressed in spherical coordinates:

$$d\mathbf{q} = dr \cdot \hat{\mathbf{r}} + r \cdot d\theta \cdot \hat{\mathbf{\theta}} + r\sin\theta \cdot d\phi \cdot \hat{\mathbf{\phi}}$$

>

> We use the previous formulas to expand and compute $W$

> $$\begin{aligned}

  W &= \int _{\mathbf{q} \in C} -\frac{GMm}{r^2}\hat{\mathbf{r}} \cdot \left(dr \cdot \hat{\mathbf{r}} + r \cdot d\theta \cdot \hat{\mathbf{\theta}} + r\sin\theta \cdot d\phi \cdot \hat{\mathbf{\phi}}\right) \\\  &= \int _{r=r _1}^{r _2} -\frac{GMm}{r^2} \, dr \\\  &= \left[ \frac{GMm}{r} \right] _{r=r _1}^{r _2} \\\  &= GMm \left(\frac{1}{r _2} - \frac{1}{r _1}\right)

  \end{aligned}$$

>

> The work is equal to the amount of gravitational potential energy lost, so one possible function for the gravitational potential energy is $\phi(r,\theta,\phi) = -\frac{GMm}{r}$ or equivalently, $\phi(\mathbf{q}) = -\frac{GMm}{|\mathbf{q}|}$.



> Example 3: Consider the spiral $C$ parameterized with respect to $t$ in cylindrical coordinates by $\mathbf{q} _C(t) = (\rho = R, \phi = t, z = k\frac{t}{2\pi})$. Consider the problem of determining the spiral's length with $t$ restricted to the range $[0, 2\pi]$. An infinitesimal change of $dt$ in $t$ results in the infinitesimal displacement:

>

> $$d\mathbf{q} = \frac{d\mathbf{q} _C}{dt}dt = \left(\frac{d\rho}{dt}\hat{\mathbf{\rho}} + \rho\frac{d\phi}{dt}\hat{\mathbf{\phi}} + \frac{dz}{dt}\hat{\mathbf{z}}\right)dt = \left(R\hat{\mathbf{\phi}} + \frac{k}{2\pi}\hat{\mathbf{z}}\right)dt$$

>

> The length of the infinitesimal displacement is $|d\mathbf{q}| = \sqrt{R^2 + \left(\frac{k}{2\pi}\right)^2} \cdot dt$.

>

> The length of the spiral is therefore: $$\begin{aligned} \int _{\mathbf{q} \in C} |d\mathbf{q}| &= \int _{t=0}^{2\pi} \sqrt{R^2 + \left(\frac{k}{2\pi}\right)^2} \, dt \\\\ &= 2\pi \sqrt{R^2 + \left(\frac{k}{2\pi}\right)^2} \\\\ &= \sqrt{(2\pi R)^2 + k^2} \end{aligned}$$



\paragraph{Path Independence}



If a vector field \textbf{F} is the gradient of a scalar field \textit{G} (i.e. if \textbf{F} is conservative), that is,



$$\mathbf{F} = \nabla G ,$$



then by the multivariable chain rule, the derivative of the composition of \textit{G} and \textbf{r}(\textit{t}) is



$$\frac{dG(\mathbf{r}(t))}{dt} = \nabla G(\mathbf{r}(t)) \cdot \mathbf{r}'(t) = \mathbf{F}(\mathbf{r}(t)) \cdot \mathbf{r}'(t)$$



which happens to be the integrand for the line integral of \textbf{F} on \textbf{r}(\textit{t}). It follows, given a path \textit{C}, that



$$\int _C \mathbf{F}(\mathbf{r})\cdot\,d\mathbf{r} = \int _a^b \mathbf{F}(\mathbf{r}(t))\cdot\mathbf{r}'(t)\,dt = \int _a^b \frac{dG(\mathbf{r}(t))}{dt}\,dt = G(\mathbf{r}(b)) - G(\mathbf{r}(a)).$$



In other words, the integral of \textbf{F} over \textit{C} depends solely on the values of \textit{G} at the points \textbf{r}(\textit{b}) and \textbf{r}(\textit{a}), and is thus independent of the path between them. For this reason, a line integral of a conservative vector field is called \textit{path independent}.



\paragraph{Surface Integrals}



Given any oriented surface $\sigma$ (oriented means that the there is a preferred direction to pass through the surface), an infinitesimal portion of the surface is defined by an infinitesimal area $|dS|$, and a unit length outwards oriented normal vector $\mathbf{n}$. $\mathbf{n}$ has a length of 1 and is perpendicular to the surface of $\sigma$, while penetrating $\sigma$ in the preferred direction. The infinitesimal portion of the surface is denoted by the infinitesimal "surface vector": $\mathbf{dS} = |dS|\mathbf{n}$. If a vector field $\mathbf{F}: \mathbb{R} ^3 \to \mathbb{R} ^3$ denotes a flow density, then the flow through the infinitesimal surface portion in the preferred direction is $\mathbf{F}(\mathbf{q}) \cdot \mathbf{dS}$.



The infinitesimal "surface vector" $\mathbf{dS} = \mathbf{n}|dS|$ describes the infinitesimal surface element in a manner similar to how the infinitesimal displacement $d\mathbf{q}$ describes an infinitesimal portion of a path. More specifically, similar to how the interior points on a path do not affect the total displacement, the interior points on a surface to not affect the total surface vector.



% [Image: The displacement between two points is independent of the path that connects them.]



Figure: The displacement between two points is independent of the path that connects them.



Consider for instance two paths $C _1$ and $C _2$ that both start at point $A$, and end at point $B$. The total displacements, $\int _{\mathbf{q} \in C _1} d\mathbf{q}$ and $\int _{\mathbf{q} \in C _2} d\mathbf{q}$, are both equivalent and equal to the displacement between $A$ and $B$. Note however that the total lengths $\int _{\mathbf{q} \in C _1} |d\mathbf{q}|$ and $\int _{\mathbf{q} \in C _2} |d\mathbf{q}|$ are not necessarily equivalent.



Similarly, given two surfaces $S _1$ and $S _2$ that both share the same counter-clockwise oriented boundary $C$, the total surface vectors $\int _{\mathbf{q} \in S _1} \mathbf{dS}$ and $\int _{\mathbf{q} \in S _2} \mathbf{dS}$ are both equivalent and are a function of the boundary $C$. This implies that a surface can be freely deformed within its boundaries without changing the total surface vector. Note however that the surface areas $\int _{\mathbf{q} \in S _1} |\mathbf{dS}|$ and $\int _{\mathbf{q} \in S _2} |\mathbf{dS}|$ are not necessarily equivalent.



The fact that the total surface vectors of $S _1$ and $S _2$ are equivalent is not immediately obvious. To prove this fact, let $\mathbf{F}$ be a constant vector field. $S _1$ and $S _2$ share the same boundary, so the flux/flow of $\mathbf{F}$ through $S _1$ and $S _2$ is equivalent. The flux through $S _1$ is $\Phi _1 = \int _{\mathbf{q} \in S _1} \mathbf{F} \cdot \mathbf{dS} = \mathbf{F} \cdot \int _{\mathbf{q} \in S _1} \mathbf{dS}$, and similarly for $S _2$ is $\Phi _2 = \int _{\mathbf{q} \in S _2} \mathbf{F} \cdot \mathbf{dS} = \mathbf{F} \cdot \int _{\mathbf{q} \in S _2} \mathbf{dS}$. Since $\mathbf{F} \cdot \int _{\mathbf{q} \in S _1} \mathbf{dS} = \mathbf{F} \cdot \int _{\mathbf{q} \in S _2} \mathbf{dS}$ for every choice of $\mathbf{F}$, it follows that $\int _{\mathbf{q} \in S _1} \mathbf{dS} = \int _{\mathbf{q} \in S _2} \mathbf{dS}$.



The geometric significance of the total surface vector is that each component measures the area of the projection of the surface onto the plane formed by the other two dimensions. Let $\sigma$ be a surface with surface vector $\mathbf{S} = S _x\mathbf{i} + S  _y\mathbf{j} + S _z\mathbf{k}$. It is then the case that: $S _x$ is the area of the projection of $\sigma$ onto the yz-plane; $S _y$ is the area of the projection of $\sigma$ onto the xz-plane; and $S _z$ is the area of the projection of $\sigma$ onto the xy-plane.



% [Image: The boundary d Sigma of Sigma is counter-clockwise oriented.]



Figure: The boundary $\partial\Sigma$ of $\Sigma$ is counter-clockwise oriented.



Given an oriented surface $\Sigma$, another important concept is the oriented boundary. The boundary of $\Sigma$ is an oriented curve $\partial\Sigma$ but how is the orientation chosen? If the boundary is "counter-clockwise" oriented, then the boundary must follow a counter-clockwise direction when the oriented surface normal vectors point towards the viewer. The counter-clockwise boundary also obeys the "right-hand rule": If you hold your right hand with your thumb in the direction of the surface normals (penetrating the surface in the "preferred" direction), then your fingers will wrap around in the direction of the counter-clockwise oriented boundary.



> Example 1: Consider the Cartesian points $(0,0,0)$; $(1,0,0)$; $(0,1,0)$; and $(0,0,1)$.

>

> Let $\sigma _1$ be the surface formed by the triangular planes $\{(0,1,0), (0,0,0), (1,0,0)\}$; $\{(0,0,1), (0,0,0), (0,1,0)\}$; and $\{(1,0,0), (0,0,0), (0,0,1)\}$ where the vertices are listed in a counterclockwise direction relative to the surface normal directions. The surface vectors of each plane are respectively $\frac{1}{2}\mathbf{k}$; $\frac{1}{2}\mathbf{i}$; and $\frac{1}{2}\mathbf{j}$ respectively which add to a total surface vector of $\mathbf{S} _1 = \frac{1}{2}(\mathbf{i}+\mathbf{j}+\mathbf{k})$.

>

> Let $\sigma _2$ be the surface formed by the single triangular plane $\{(1,0,0), (0,1,0), (0,0,1)\}$ where the vertices are listed in a counterclockwise direction relative to the normal direction. It can be seen that $\sigma _1$ and $\sigma _2$ share a the common counter clockwise boundary $(1,0,0) \to (0,1,0) \to (0,0,1)$The surface vector is $\mathbf{S} _2 = \frac{1}{2}(\mathbf{j} - \mathbf{i}) \times (\mathbf{k} - \mathbf{i}) = \frac{1}{2}(\mathbf{i} + \mathbf{j} + \mathbf{k})$ which is equivalent to $\mathbf{S} _1$.



> Example 2: This example will show how moving a point that is in the interior of a "triangular mesh" does not affect the total surface vector. Consider the points $P _0, P _1, P _2, \dots, P _n$ where $n \geq 3$. Let the closed path $C$ be defined by the cycle $P _1 \to P _2 \to \dots \to P _n \to P _1$. For simplicity, $P _{n+1} = P _1$. For each $i = 1, 2, \dots, n$, $\mathbf{v} _i$ will denote the displacement of $P _i$ relative to $P _0$. Like with $P _{n+1}$, $\mathbf{v} _{n+1} = \mathbf{v} _1$.

>

> Let $\sigma$ denote a surface that is a "triangular mesh" comprised of the closed fan of triangles: $\{P _2,P _0,P _1\}$; $\{P _3,P _0,P _2\}$; \...; $\{P _n,P _0,P _{n-1}\}$; $\{P _1,P _0,P _n\}$ where the vertices of each triangle are listed in a counterclockwise direction. It can be seen that the counterclockwise boundary of $\sigma$ is $C$ and does not depend on the location of $P _0$. The total surface vector for $\sigma$ is:

>

> $$\begin{aligned}

  \mathbf{S} &= \frac{1}{2}(\mathbf{v} _1 \times \mathbf{v} _2) + \frac{1}{2}(\mathbf{v} _2 \times \mathbf{v} _3) + \dots + \frac{1}{2}(\mathbf{v} _{n-1} \times \mathbf{v} _n) + \frac{1}{2}(\mathbf{v} _n \times \mathbf{v} _{n+1}) \\\  &= \frac{1}{2} \sum _{i=1}^n (\mathbf{v} _i \times \mathbf{v} _{i+1})

  \end{aligned}$$

>

> Now displace $P _0$ by $\mathbf{w}$ to get $P' _0$. The displacement vector of $P _i$ relative to $P _0'$ becomes $\mathbf{v}' _i = \mathbf{v} _i - \mathbf{w}$. The counterclockwise boundary is unaffected. The total surface vector is:

>

> $$\begin{aligned}

  \mathbf{S}' &= \frac{1}{2}\sum _{i=1}^n (\mathbf{v}' _i \times \mathbf{v}' _{i+1}) \\\  &= \frac{1}{2}\sum _{i=1}^n \left((\mathbf{v} _i - \mathbf{w}) \times (\mathbf{v} _{i+1} - \mathbf{w})\right) \\\  &= \frac{1}{2}\sum _{i=1}^n \left((\mathbf{v} _i \times \mathbf{v} _{i+1}) - (\mathbf{v} _i \times \mathbf{w}) - (\mathbf{w} \times \mathbf{v} _{i+1}) + (\mathbf{w} \times \mathbf{w})\right) \\\  &= \frac{1}{2}\sum _{i=1}^n \left((\mathbf{v} _i \times \mathbf{v} _{i+1}) - \mathbf{w} \times (\mathbf{v} _{i+1} - \mathbf{v} _i)\right) \\\  &= \mathbf{S} - \frac{1}{2}\mathbf{w} \times (\mathbf{v} _{n+1} - \mathbf{v} _1) \\\  &= \mathbf{S} - \frac{1}{2}\mathbf{w} \times \mathbf{0} \\\  &= \mathbf{S}

  \end{aligned}$$

>

> Therefore moving the interior point $P _0$ neither affects the boundary, nor the total surface vector.



\paragraph{Calculating Surface Integrals}



To calculate a surface integral, the oriented surface $\sigma$ must be parameterized. Let $\mathbf{q} _{\sigma}(u, v)$ be a continuous function that maps each point $(u,v)$ from a two-dimensional domain $D _{u,v}$ to a point in $\sigma$. $\mathbf{q} _{\sigma}(u, v)$ must be continuous and onto. While $\mathbf{q} _{\sigma}(u, v)$ does not necessarily have to be one to one, the parameterization should never "fold back" on itself. The infinitesimal increases in $u$ and $v$ are respectively $du$ and $dv$. These respectively give rise to the displacements $\frac{\partial \mathbf{q} _{\sigma}}{\partial u}du$ and $\frac{\partial \mathbf{q} _{\sigma}}{\partial v}dv$. Assuming that the surface's orientation follows the right hand rule with respect to the displacements $\frac{\partial \mathbf{q} _{\sigma}}{\partial u}du$ and $\frac{\partial \mathbf{q} _{\sigma}}{\partial v}dv$, the surface vector that arises is $\mathbf{dS} = (\frac{\partial \mathbf{q} _{\sigma}}{\partial u} \times \frac{\partial \mathbf{q} _{\sigma}}{\partial v})dudv$.



In the surface integral $\iint _{\mathbf{q} \in \sigma} \mathbf{F}(\mathbf{q}) \cdot \mathbf{dS}$, the differential $\mathbf{dS}$ can be replaced with $(\frac{\partial \mathbf{q} _{\sigma}}{\partial u} \times \frac{\partial \mathbf{q} _{\sigma}}{\partial v})dudv$ to get

$$\iint _{(u,v) \in D _{u,v}} \mathbf{F}(\mathbf{q} _{\sigma}(u,v)) \cdot (\frac{\partial \mathbf{q} _{\sigma}}{\partial u} \times \frac{\partial \mathbf{q} _{\sigma}}{\partial v})dudv$$



> Example 3: Consider the problem of computing the surface area of a sphere of radius $R$.

>

> Center the sphere $\sigma$ on the origin, and using $u$ and $v$ as the parameter variables, the sphere can be parameterized in spherical coordinates via $\mathbf{q} _{\sigma}(u,v) = (r = R, \theta = u, \phi = v)$ where $u \in [0, \pi]$ and $v \in [-\pi, +\pi]$. The infinitesimal displacements from small changes in the parameters are:

>

> $du$ causes $$\frac{\partial\mathbf{q} _{\sigma}}{\partial u}du = (\frac{\partial r}{\partial u}\hat{\mathbf{r}} + r\frac{\partial \theta}{\partial u}\hat{\mathbf{\theta}} + r\sin\theta\frac{\partial \phi}{\partial u}\hat{\mathbf{\phi}})du = (R\hat{\mathbf{\theta}})du$$

>

> $dv$ causes $$\frac{\partial\mathbf{q} _{\sigma}}{\partial v}dv = (\frac{\partial r}{\partial v}\hat{\mathbf{r}} + r\frac{\partial \theta}{\partial v}\hat{\mathbf{\theta}} + r\sin\theta\frac{\partial \phi}{\partial v}\hat{\mathbf{\phi}})dv = (R\sin(u)\hat{\mathbf{\phi}})dv$$

>

> The infinitesimal surface vector is hence $$\mathbf{dS} = (R\hat{\mathbf{\theta}})du \times (R\sin(u)\hat{\mathbf{\phi}})dv = (R^2\sin(u)\hat{\mathbf{r}})dudv$$. While not important to this example, note how the parameterization was chosen so that the surface vector points outwards. The area is $|\mathbf{dS}| = R^2\sin(u)dudv$.

>

> The total surface area is hence:

>

> $$\begin{aligned} \iint _{\mathbf{q} \in \sigma} |\mathbf{dS}| &= \int _{u = 0}^{\pi}\int _{v = -\pi}^{+\pi} R^2\sin(u) \, dv \, du \\\\ &= 2\pi R^2 \int _{u = 0}^{\pi} \sin(u) \, du \\\\ &= 2\pi R^2 \left(-\cos(u)\bigg| _{u=0}^{\pi}\right) \\\\ &= 2\pi R^2 \left(-\cos(\pi) + \cos(0)\right) \\\\ &= 2\pi R^2 \left(-(-1) + 1\right) \\\\ &= 2\pi R^2 (2) \\\\ &= 4\pi R^2 \end{aligned}$$ 



\href{Surface_Integrals "wikilink"}{More on Surface Integrals and Area}



\paragraph{Resources}



\begin{itemize}
    \item \href{https://math.libretexts.org/Bookshelves/Calculus/Supplemental_Modules_(Calculus}{Line Integrals}/Vector_Calculus/4\%3A_Integration_in_Vector_Fields/4.3\%3A_Line_Integrals), Mathematics LibreTexts
    \item \href{https://en.wikibooks.org/wiki/Calculus/Vector_calculus}{Vector Calculus}, Wikibooks: Calculus
    \item \href{https://en.wikipedia.org/wiki/Line_integral}{Line integral}, Wikipedia
    \item \href{https://www.khanacademy.org/math/multivariable-calculus/integrating-multivariable-functions/line-integrals/v/introduction-to-the-line-integral}{Introduction to The Line Integral}, Khan Academy
\end{itemize}

\paragraph{Licensing}



Content obtained and/or adapted from:



\begin{itemize}
    \item \href{https://en.wikibooks.org/wiki/Calculus/Vector_calculus#Volume,_path,_and_surface_integrals}{Vector calculus, Wikibooks: Calculus} under a CC BY-SA license
\end{itemize}

\subsection{Learning Objectives}

\begin{enumerate}
    \item Here are three Student Learning Objectives (SLOs) related to line integrals in the context of vector calculus:
    \item \textbf{Evaluate Line Integrals in Scalar Fields:}
    \item \textbf{Outcome:} Students will be able to evaluate line integrals of scalar fields along a given curve.
    \item \textbf{Details:} This includes parameterizing the curve, computing the differential element along the curve, and integrating the scalar field along the curve.
    \item \textbf{Compute Work Using Line Integrals in Vector Fields:}
    \item \textbf{Outcome:} Students will be able to compute the work done by a vector field along a given path using line integrals.
    \item \textbf{Details:} This includes understanding the concept of work in the context of vector fields, parameterizing the curve, and calculating the line integral of the vector field along the path.
    \item \textbf{Apply Path Independence in Conservative Fields:}
    \item \textbf{Outcome:} Students will be able to identify when a vector field is conservative and apply the concept of path independence to evaluate line integrals.
    \item \textbf{Details:} This includes recognizing conservative vector fields, using potential functions, and simplifying the calculation of line integrals by evaluating the difference in potential function values at the endpoints of the path.
\end{enumerate}

\subsection{Assessment Instrument}

\begin{enumerate}
    \item \textbf{Assessment Instrument: Line Integrals Quiz} \textbf{Format:}
    \item \textbf{Type:} Mixed-format quiz (multiple-choice, short answer, and problem-solving questions)
    \item \textbf{Duration:} 45 minutes
    \item \textbf{Total Marks:} 30 points \textbf{Sections:}
    \item \textbf{Conceptual Understanding of Line Integrals (10 points):}
    \item \textbf{Question 1 (Multiple Choice):} Which of the following best describes a line integral in a vector field?
    \item a. The sum of values of the vector field at all points on a curve.
    \item b. The total area under a curve in a 2D plane.
    \item c. The total volume enclosed by a 3D surface.
    \item d. The total displacement between two points in a plane.
    \item \textbf{Question 2 (Short Answer):} Explain the significance of path independence in conservative vector fields.
    \item \textbf{Evaluating Line Integrals in Scalar Fields (10 points):}
    \item \textbf{Question 3 (Problem-Solving):} Evaluate the line integral of the scalar field \( f(x,y) = x^2 + y^2 \) along the straight line segment from the point \( (0,0) \) to \( (1,1) \).
    \item \textit{Hint:} Parameterize the curve and compute the integral.
    \item \textbf{Question 4 (Multiple Choice):} If a scalar field \( f(x,y) = xy \) is integrated along the path \( C \) from \( (0,0) \) to \( (1,1) \) parameterized by \( \mathbf{r}(t) = (t, t) \), for \( t \) in \( [0, 1] \), what is the value of the line integral?
    \item a. 0
    \item b. \( \frac{1}{2} \)
    \item c. 1
    \item d. \( \frac{3}{2} \)
    \item \textbf{Computing Work Using Line Integrals in Vector Fields (10 points):}
    \item \textbf{Question 5 (Problem-Solving):} Consider the vector field \( \mathbf{F}(x,y) = (-y, x) \). Compute the work done by the field along the quarter-circle path \( C \) from \( (1, 0) \) to \( (0, 1) \).
    \item \textit{Hint:} Parameterize the curve using \( \mathbf{r}(t) = (\cos t, \sin t) \) for \( t \) in \( [0, \frac{\pi}{2}] \) and compute the line integral \( \int_C \mathbf{F} \cdot d\mathbf{r} \).
    \item \textbf{Question 6 (Short Answer):} What is the physical interpretation of the line integral of a force field along a given path? \textbf{Scoring Rubric:}
    \item \textbf{Conceptual Understanding of Line Integrals:} 5 points for each correct answer (total 10 points)
    \item \textbf{Evaluating Line Integrals in Scalar Fields:} 5 points for each correct evaluation (total 10 points)
    \item \textbf{Computing Work Using Line Integrals in Vector Fields:} 5 points for each correct solution (total 10 points) This quiz will assess students' understanding of line integrals, their ability to evaluate these integrals in scalar and vector fields, and their application to compute work in a physical context.
\end{enumerate}

%%===============================================================
\section{Path Independence, Conservative Fields, and Potential Functions}
\label{sec:path-independence-conservative-fields-and-potential-functions}
%%===============================================================

\paragraph{Path Independence}



If a vector field $\mathbf{F}$ is the gradient of a scalar field $G$, i.e.,



$$\mathbf{F} = \nabla G,$$



then by the multivariable chain rule, the derivative of the composition of $G$ and $\mathbf{r}(t)$ is



$$\begin{aligned}

\frac{dG(\mathbf{r}(t))}{dt} 

&= \nabla G(\mathbf{r}(t)) \cdot \mathbf{r}'(t) \\\&= \mathbf{F}(\mathbf{r}(t)) \cdot \mathbf{r}'(t) 

\end{aligned}$$





which happens to be the integrand for the line integral of $\mathbf{F}$ on $\mathbf{r}(t)$. It follows, given a path $C$, that



$$\begin{aligned}

\int _C \mathbf{F}(\mathbf{r})\cdot\,d\mathbf{r} 

&= \int _a^b \mathbf{F}(\mathbf{r}(t))\cdot\mathbf{r}'(t)\,dt \\\&= \int _a^b \frac{dG(\mathbf{r}(t))}{dt}\,dt \\\&= G(\mathbf{r}(b)) - G(\mathbf{r}(a)) 

\end{aligned}$$





In other words, the integral of $\mathbf{F}$ over $C$ depends solely on the values of $G$ at the points $\mathbf{r}(b)$ and $\mathbf{r}(a)$, and is thus independent of the path between them. For this reason, a line integral of a conservative vector field is called \textit{path independent}.



\paragraph{Stokes' Theorem}



Stokes' Theorem is effectively a generalization of Green's theorem to 3 dimensions, and the "curl" is a generalization of the quantity $\frac{\partial F _y}{\partial x} - \frac{\partial F _x}{\partial y}$ to 3 dimensions. An arbitrary oriented surface $\sigma$ can be articulated into a family of infinitesimal surfaces, some parallel to the $xy$-plane, others parallel to the $zx$-plane, and the remainder parallel to the $yz$-plane. Let $\mathbf{F}$ denote an arbitrary vector field.



Let $\sigma$ be a surface that is parallel to the $yz$-plane with counter-clockwise oriented boundary $C$. Green's theorem gives:



$$\begin{aligned}

\int _{\mathbf{q} \in C} \mathbf{F}(\mathbf{q}) \cdot d\mathbf{q} 

&= \iint _{\mathbf{q} \in \sigma}\left(\frac{\partial F _z}{\partial y} - \frac{\partial F _y}{\partial z}\right)(\mathbf{dS} \cdot \mathbf{i}) \\\&= \iint _{\mathbf{q} \in \sigma}\left(\left(\frac{\partial F _z}{\partial y} - \frac{\partial F _y}{\partial z}\right)\mathbf{i}\right) \cdot \mathbf{dS} 

\end{aligned}$$





$\mathbf{dS} \cdot \mathbf{i}$ is positive if the normal direction to $\sigma$ points in the positive $x$ direction and is negative if otherwise. If the normal direction to $\sigma$ points in the negative $x$ direction, then $C$ is oriented clockwise instead of counter-clockwise in the $yz$-plane.



% [Image: Decomposing a 3D loop into an ensemble of infinitesimal loops that are parallel to the yz, zx, or xy planes.]



Figure: Decomposing a 3D loop into an ensemble of infinitesimal loops that are parallel to the $yz$, $zx$, or $xy$ planes.



Repeating this argument for $\sigma$ being parallel to the $zx$-plane and $xy$-plane respectively gives:



$$\int _{\mathbf{q} \in C} \mathbf{F}(\mathbf{q}) \cdot d\mathbf{q} =

\iint _{\mathbf{q} \in \sigma}\left(\frac{\partial F _x}{\partial z} - \frac{\partial F _z}{\partial x}\right)(\mathbf{dS} \cdot \mathbf{j}) =

\iint _{\mathbf{q} \in \sigma}\left(\left(\frac{\partial F _x}{\partial z} - \frac{\partial F _z}{\partial x}\right)\mathbf{j}\right) \cdot \mathbf{dS}$$



and



$$\int _{\mathbf{q} \in C} \mathbf{F}(\mathbf{q}) \cdot d\mathbf{q} =

\iint _{\mathbf{q} \in \sigma}\left(\frac{\partial F _y}{\partial x} - \frac{\partial F _x}{\partial y}\right)(\mathbf{dS} \cdot \mathbf{k}) =

\iint _{\mathbf{q} \in \sigma}\left(\left(\frac{\partial F _y}{\partial x} - \frac{\partial F _x}{\partial y}\right)\mathbf{k}\right) \cdot \mathbf{dS}$$



Treating $\sigma$ as an ensemble of infinitesimal surfaces parallel to the $yz$-plane, $zx$-plane, or $xy$-plane gives:



$$\int _{\mathbf{q} \in C} \mathbf{F}(\mathbf{q}) \cdot d\mathbf{q} =

\iint _{\mathbf{q} \in \sigma}\left(

\left(\frac{\partial F _z}{\partial y} - \frac{\partial F _y}{\partial z}\right)\mathbf{i} +

\left(\frac{\partial F _x}{\partial z} - \frac{\partial F _z}{\partial x}\right)\mathbf{j} +

\left(\frac{\partial F _y}{\partial x} - \frac{\partial F _x}{\partial y}\right)\mathbf{k}

\right) \cdot \mathbf{dS}$$



This is Stokes' theorem, and 

$$\nabla \times \mathbf{F} =

\left(\frac{\partial F _z}{\partial y} - \frac{\partial F _y}{\partial z}\right)\mathbf{i} +

\left(\frac{\partial F _x}{\partial z} - \frac{\partial F _z}{\partial x}\right)\mathbf{j} +

\left(\frac{\partial F _y}{\partial x} - \frac{\partial F _x}{\partial y}\right)\mathbf{k}$$ 

is the "curl" of $\mathbf{F}$ which generalizes the "circulation density" to 3 dimensions.



The direction of $\nabla \times \mathbf{F}$ at $\mathbf{q}$ is effectively an "axis of rotation" around which the counterclockwise circulation density in a plane whose normal is parallel to $\nabla \times \mathbf{F}$ is $|\nabla \times \mathbf{F}|$. Out of all planes that pass through $\mathbf{q}$, the plane whose normal is parallel to $\nabla \times \mathbf{F}$ has the largest counterclockwise circulation density at $\mathbf{q}$ which is $|\nabla \times \mathbf{F}|$.



An arbitrary vector field $\mathbf{F}$ that is differentiable everywhere is considered to be "irrotational" or "conservative" if $\nabla \times \mathbf{F} = \mathbf{0}$ everywhere, or equivalently that $\int _{\mathbf{q} \in C} \mathbf{F}(\mathbf{q}) \cdot d\mathbf{q} = 0$ for all continuous closed curves $C$



\paragraph{Conservative vector fields}



A vector field $\mathbf{F}$ for which $\nabla \times \mathbf{F} = \mathbf{0}$ at all points is an "\textbf{conservative}" vector field. $\mathbf{F}$ can also be referred to as being "\textbf{irrotational}" since the gain around any closed curve is always 0.



A key property of a conservative vector field $\mathbf{F}$ is that the gain of $\mathbf{F}$ along a continuous curve is purely a function of the curve's end points. If $C _1$ and $C _2$ are two continuous curves which share the same starting point $\mathbf{q} _0$ and end point $\mathbf{q} _1$, then $\int _{\mathbf{q} \in C _1} \mathbf{F}(\mathbf{q}) \cdot d\mathbf{q} = \int _{\mathbf{q} \in C _2} \mathbf{F}(\mathbf{q}) \cdot d\mathbf{q}$. In other words, the gain is purely a function of $\mathbf{q} _0$ and $\mathbf{q} _1$. This property can be derived from Stokes' theorem as follows:



Invert the orientation of $C _2$ to get $-C _2$ and combine $C _1$ and $-C _2$ to get a continuous closed curve $C _3 = C _1 - C _2$, linking the curves together at the endpoints $\mathbf{q} _0$ and $\mathbf{q} _1$. Let $\sigma$ denote a surface for which $C _3$ is the counterclockwise oriented boundary.



Stokes' theorem states that

$$\int _{\mathbf{q} \in C _3} \mathbf{F}(\mathbf{q}) \cdot d\mathbf{q} = \iint _{\mathbf{q} \in \sigma} (\nabla \times \mathbf{F})| _\mathbf{q} \cdot \mathbf{dS} = 0.$$ The gain around $C _3$ is the gain along $C _1$ minus the gain along $C _2$:

$$\begin{aligned} \int _{\mathbf{q} \in C _3} \mathbf{F}(\mathbf{q}) \cdot d\mathbf{q} &= \int _{\mathbf{q} \in C _1} \mathbf{F}(\mathbf{q}) \cdot d\mathbf{q} + \int _{\mathbf{q} \in -C _2} \mathbf{F}(\mathbf{q}) \cdot d\mathbf{q} \\\\ &= \int _{\mathbf{q} \in C _1} \mathbf{F}(\mathbf{q}) \cdot d\mathbf{q} - \int _{\mathbf{q} \in C _2} \mathbf{F}(\mathbf{q}) \cdot d\mathbf{q} \end{aligned}$$

Therefore:



$$\int _{\mathbf{q} \in C _3} \mathbf{F}(\mathbf{q}) \cdot d\mathbf{q} = 0 \implies \int _{\mathbf{q} \in C _1} \mathbf{F}(\mathbf{q}) \cdot d\mathbf{q} = \int _{\mathbf{q} \in C _2} \mathbf{F}(\mathbf{q}) \cdot d\mathbf{q}$$



\paragraph{Scalar Potential}



If $\mathbf{F}$ is a conservative vector field (also called \textit{irrotational}, \textit{curl-free}, or \textit{potential}), and its components have continuous partial derivatives, the \textbf{potential} of $\mathbf{F}$ with respect to a reference point $\mathbf r _0$ is defined in terms of the line integral:



$$V(\mathbf r) = -\int _C \mathbf{F}(\mathbf{r})\cdot\,d\mathbf{r} = -\int _a^b \mathbf{F}(\mathbf{r}(t))\cdot\mathbf{r}'(t)\,dt,$$



where $C$ is a parametrized path from $\mathbf r _0$ to $\mathbf r,$



$$\mathbf{r}(t), a\leq t\leq b, \mathbf{r}(a)=\mathbf{r _0}, \mathbf{r}(b)=\mathbf{r}.$$



The fact that the line integral depends on the path $C$ only through its terminal points $\mathbf r _0$ and $\mathbf r$ is, in essence, the \textbf{path independence property} of a conservative vector field. The fundamental theorem of line integrals implies that if $V$ is defined in this way, then $\mathbf{F}= -\nabla V,$ so that $V$ is a scalar potential of the conservative vector field $\mathbf{F}$. Scalar potential is not determined by the vector field alone: indeed, the gradient of a function is unaffected if a constant is added to it. If $V$ is defined in terms of the line integral, the ambiguity of $V$ reflects the freedom in the choice of the reference point $\mathbf r _0.$



\paragraph{Resources}



\begin{itemize}
    \item \href{https://en.wikibooks.org/wiki/Calculus/Vector_calculus}{Vector Calculus}, Wikibooks: Calculus
\end{itemize}

$\textbf{Conservative Vector Fields}$



\begin{itemize}
    \item \href{https://www.youtube.com/watch?v=cVTcsHcispQ}{Conservative Vector Fields} Video by James Sousa, Math is Power 4U
    \item \href{https://www.youtube.com/watch?v=VLTUku45Qs4}{Conservative Vector Fields - The Definition and a Few Remarks} Video by Patrick JMT
    \item \href{https://www.youtube.com/watch?v=gAb1ZTD41wo}{Showing a Vector Field on R\^2 is Conservative} Video by Patrick JMT
\end{itemize}

$\textbf{Finding a Potential Function of a Conservative Vector Field}$



\begin{itemize}
    \item \href{https://www.youtube.com/watch?v=nQkHh2psLck}{Determining the Potential Function of a Conservative Vector Field} Video by James Sousa, Math is Power 4U
    \item \href{https://www.youtube.com/watch?v=iLAK2IsQ_Uo}{Finding a Potential for a Conservative Vector Field} Video by Patrick JMT
    \item \href{https://www.youtube.com/watch?v=tslJEOnt9aY}{Finding a Potential for a Conservative Vector Field Ex 2} Video by Patrick JMT
    \item \href{https://www.youtube.com/watch?v=nY4mW_R-T40}{Potential Function of a Conservative Vector Field} Video by Krista King
    \item \href{https://www.youtube.com/watch?v=Rb12YouJXyA}{Potential Function of a Conservative Vector Field in 3D} Video by Krista King
\end{itemize}

$\textbf{The Fundamental Theorem of Line Integrals}$



\begin{itemize}
    \item \href{https://www.youtube.com/watch?v=Td6g2bvJh18}{The Fundamental Theorem of Line Integrals Part 1} Video by James Sousa, Math is Power 4U
    \item \href{https://www.youtube.com/watch?v=I8BtXV-xgxM}{The Fundamental Theorem of Line Integrals Part 2} Video by James Sousa, Math is Power 4U
    \item \href{https://www.youtube.com/watch?v=00RN-hw9BXc}{The Fundamental Theorem of Line Integrals on a Closed Path} Video by James Sousa, Math is Power 4U
    \item \href{https://www.youtube.com/watch?v=62oBGKSjYiY}{Ex 1: Fundamental Theorem of Line Integrals in the Plane} Video by James Sousa, Math is Power 4U
    \item \href{https://www.youtube.com/watch?v=Eh6OinLiGWg}{Ex 2: Fundamental Theorem of Line Integrals in the Plane} Video by James Sousa, Math is Power 4U
    \item \href{https://www.youtube.com/watch?v=pbHlZR2TD7M}{Ex 3: Fundamental Theorem of Line Integrals in the Plane} Video by James Sousa, Math is Power 4U
    \item \href{https://www.youtube.com/watch?v=e2g7321ne48}{Ex 4: Fundamental Theorem of Line Integrals in Space} Video by James Sousa, Math is Power 4U
    \item \href{https://www.youtube.com/watch?v=xVhHow_usMQ}{The Fundamental Theorem for Line Integrals} Video by Patrick JMT
    \item \href{https://www.youtube.com/watch?v=hT3YSXQaUHQ}{Potential Function of a Conservative Vector Field to Evaluate a Line Integral} Video by Krista King
\end{itemize}

\paragraph{Licensing}



Content obtained and/or adapted from:



\begin{itemize}
    \item \href{https://en.wikibooks.org/wiki/Calculus/Vector_calculus#Volume,_path,_and_surface_integrals}{Vector calculus, Wikibooks: Calculus} under a CC BY-SA license
    \item \href{https://en.wikipedia.org/wiki/Scalar_potential}{Scalar potential, Wikipedia} under a CC BY-SA license
\end{itemize}

\subsection{Learning Objectives}

\begin{enumerate}
    \item Here are three Student Learning Objectives (SLOs) for the concepts related to path independence, Stokes' Theorem, and conservative vector fields:
    \item \textbf{Understand and Apply Path Independence in Conservative Vector Fields:}
    \item \textbf{Outcome:} Students will be able to determine if a given vector field is conservative by verifying path independence.
    \item \textbf{Details:} This includes understanding that for a conservative vector field, the line integral between two points is independent of the path taken. Students will solve problems where they compute line integrals along different paths and confirm that the results are the same, reinforcing the concept of path independence.
    \item \textbf{Apply Stokes' Theorem in the Evaluation of Surface Integrals:}
    \item \textbf{Outcome:} Students will be able to apply Stokes' Theorem to convert a surface integral over a surface $\sigma$ into a line integral over the boundary curve $C$ of $\sigma$.
    \item \textbf{Details:} This involves understanding the geometric interpretation of Stokes' Theorem, recognizing when it can be applied, and correctly setting up and evaluating both sides of the theorem. Students will practice by transforming surface integrals into line integrals using Stokes' Theorem and solving related problems.
    \item \textbf{Characterize and Verify Conservative Vector Fields:}
    \item \textbf{Outcome:} Students will be able to verify if a vector field is conservative by calculating its curl and determining if it equals zero.
    \item \textbf{Details:} This includes the ability to compute the curl of a vector field and interpret its physical meaning. Students will also relate the concept of a conservative vector field to the existence of a scalar potential function and verify the conservative nature of vector fields by checking for a zero curl and connecting it to path independence.
\end{enumerate}

\subsection{Assessment Instrument}

\begin{enumerate}
    \item \textbf{Assessment Instrument: Path Independence, Conservative Fields, and Potential Functions Quiz} \textbf{Format:}
    \item \textbf{Type:} Mixed-format quiz (multiple-choice, short answer, and problem-solving questions)
    \item \textbf{Duration:} 45 minutes
    \item \textbf{Total Marks:} 30 points \textbf{Sections:}
    \item \textbf{Path Independence (10 points):}
    \item \textbf{Question 1 (Multiple Choice):} Which of the following statements is true about a conservative vector field $\mathbf{F}$?
    \item a. The line integral of $\mathbf{F}$ depends on the path taken between two points.
    \item b. The line integral of $\mathbf{F}$ is zero for all paths.
    \item c. The line integral of $\mathbf{F}$ depends only on the endpoints of the path.
    \item d. The curl of $\mathbf{F}$ is non-zero everywhere.
    \item \textbf{(2 points)}
    \item \textbf{Question 2 (Short Answer):} Given the vector field $\mathbf{F} = \nabla G$, where $G$ is a scalar field, explain why the line integral of $\mathbf{F}$ is independent of the path taken between two points.
    \item \textbf{(3 points)}
    \item \textbf{Question 3 (Problem-Solving):} Calculate the line integral $\int_C \mathbf{F} \cdot d\mathbf{r}$ for the vector field $\mathbf{F}(x,y) = \langle 2x, 3y \rangle$ along two different paths from $(0,0)$ to $(2,2)$. Show that the results are identical, thus verifying path independence.
    \item \textbf{(5 points)}
    \item \textbf{Conservative Fields (10 points):}
    \item \textbf{Question 4 (Multiple Choice):} Which of the following conditions must hold for a vector field $\mathbf{F}$ to be considered conservative?
    \item a. $\nabla \times \mathbf{F} \neq 0$
    \item b. $\nabla \cdot \mathbf{F} = 0$
    \item c. $\nabla \times \mathbf{F} = 0$
    \item d. $\mathbf{F}$ is constant.
    \item \textbf{(2 points)}
    \item \textbf{Question 5 (Short Answer):} Explain why a vector field with a zero curl is called "irrotational" and how this relates to it being conservative.
    \item \textbf{(3 points)}
    \item \textbf{Question 6 (Problem-Solving):} Given the vector field $\mathbf{F}(x,y,z) = \langle y, -x, 0 \rangle$, determine whether it is conservative. If it is, find a potential function $G$ such that $\mathbf{F} = \nabla G$.
    \item \textbf{(5 points)}
    \item \textbf{Potential Functions (10 points):}
    \item \textbf{Question 7 (Multiple Choice):} Which of the following is a correct interpretation of a potential function $G$ for a conservative vector field $\mathbf{F}$?
    \item a. $G$ is the magnitude of $\mathbf{F}$.
    \item b. $G$ represents the work done by $\mathbf{F}$ in moving a particle between two points.
    \item c. $G$ is a constant scalar field.
    \item d. $G$ is the curl of $\mathbf{F}$.
    \item \textbf{(2 points)}
    \item \textbf{Question 8 (Short Answer):} Describe how the potential function $G$ is related to the concept of path independence.
    \item \textbf{(3 points)}
    \item \textbf{Question 9 (Problem-Solving):} For the conservative vector field $\mathbf{F}(x,y) = \langle 3x^2, 4y^3 \rangle$, find the potential function $G(x,y)$, assuming that $G(0,0) = 0$. Then, compute the line integral $\int_C \mathbf{F} \cdot d\mathbf{r}$ from $(0,0)$ to $(1,1)$ using $G$.
    \item \textbf{(5 points)} \textbf{Scoring Rubric:}
    \item \textbf{Path Independence:}
    \item Question 1: 2 points for the correct answer
    \item Question 2: 3 points for a clear and correct explanation
    \item Question 3: 5 points for correct calculations and demonstrating path independence
    \item \textbf{Conservative Fields:}
    \item Question 4: 2 points for the correct answer
    \item Question 5: 3 points for a clear and correct explanation
    \item Question 6: 5 points for correct determination and finding the potential function if applicable
    \item \textbf{Potential Functions:}
    \item Question 7: 2 points for the correct answer
    \item Question 8: 3 points for a clear and correct explanation
    \item Question 9: 5 points for correctly finding the potential function and calculating the line integral using it The quiz will assess students' understanding of path independence, the criteria for conservative vector fields, and the relationship between conservative fields and potential fu
\end{enumerate}

%%===============================================================
\section{Green´s Theorem}
\label{sec:green-s-theorem}
%%===============================================================

In vector calculus, \textbf{Green's theorem} relates a line integral around a simple closed curve $C$ to a double integral over the plane region $D$ bounded by $C$. It is the two-dimensional special case of Stokes' theorem.



\paragraph{Theorem}



Let $C$ be a positively oriented, piecewise smooth, simple closed curve in a plane, and let $D$ be the region bounded by $C$. If $L$ and $M$ are functions of $(x, y)$ defined on an open region containing $D$ and having continuous partial derivatives there, then



$$\oint _{\scriptstyle C} (L\, dx + M\, dy) = \iint _{D} \left(\frac{\partial M}{\partial x} - \frac{\partial L}{\partial y}\right) dx\, dy$$



where the path of integration along $C$ is anticlockwise.



In physics, Green's theorem finds many applications. One is solving two-dimensional flow integrals, stating that the sum of fluid outflowing from a volume is equal to the total outflow summed about an enclosing area. In plane geometry, and in particular, area surveying, Green's theorem can be used to determine the area and centroid of plane figures solely by integrating over the perimeter.



\paragraph{Proof when \textit{D} is a simple region}



% [Image: A region with its boundary consisting of curves C1, C2, C3, C4]



Figure: If $D$ is a simple type of region with its boundary consisting of the curves $C _1$, $C _2$, $C _3$, $C _4$, half of Green's theorem can be demonstrated.



The following is a proof of half of the theorem for the simplified area $D$, a type I region where $C _1$ and $C _3$ are curves connected by vertical lines (possibly of zero length). A similar proof exists for the other half of the theorem when $D$ is a type II region where $C _2$ and $C _4$ are curves connected by horizontal lines (again, possibly of zero length). Putting these two parts together, the theorem is thus proven for regions of type III (defined as regions which are both type I and type II). The general case can then be deduced from this special case by decomposing $D$ into a set of type III regions.



If it can be shown that if



$$\begin{align}

\oint _C L\, dx = \iint _{D} \left(- \frac{\partial L}{\partial y}\right) dA

\end{align}$$



and



$$\begin{align}

\oint _C M\, dy = \iint _{D} \left(\frac{\partial M}{\partial x}\right) dA

\end{align}$$



are true, then Green's theorem follows immediately for the region D. We can prove the first equation easily for regions of type I, and the second equation for regions of type II. Green's theorem then follows for regions of type III.



Assume region \textit{D} is a type I region and can thus be characterized, as pictured on the right, by



$$D = \\{(x,y)\mid a\le x\le b, \quad g _1(x) \le y \le g _2(x)\\}$$



where $g _1$ and $g _2$ are continuous functions on $[a, b]. Compute the double integral in (1):



$$\begin{align}\begin{aligned}

\iint _D \frac{\partial L}{\partial y}\, dA

& = \int _a^b\,\int _{g _1(x)}^{g _2(x)} \frac{\partial L}{\partial y} (x,y)\,dy\,dx \\\& = \int _a^b \left[L(x,g _2(x)) - L(x,g _1(x)) \right] \, dx.

\end{aligned}\end{align}$$



Now compute the line integral in (1). $C$ can be rewritten as the union of four curves: $C _1$, $C _2$, $C _3$, $C _4$.



With $C _1$, use the parametric equations: $x = x$, $y = g _1(x)$, $a \leq x \leq b$. Then



$$\int _{C _1} L(x,y)\, dx = \int _a^b L(x,g _1(x))\, dx.$$



With $C _3$, use the parametric equations: $x = x$, $y = g _2(x)$, $a \leq x \leq b$. Then



$$\int _{C _3} L(x,y)\, dx = -\int _{-C _3} L(x,y)\, dx = - \int _a^b L(x,g _2(x))\, dx.$$



The integral over $C _3$ is negated because it goes in the negative direction from $b$ to $a$, as $C$ is oriented positively (anticlockwise). On $C _2$ and $C _4$, $x$ remains constant, meaning



$$\int _{C _4} L(x,y)\, dx = \int _{C _2} L(x,y)\, dx = 0.$$



Therefore,



$$\begin{align}\begin{aligned}

\int _C L\, dx & = \int _{C _1} L(x,y)\, dx + \int _{C _2} L(x,y)\, dx + \int _{C _3} L(x,y)\, dx + \int _{C _4} L(x,y)\, dx \\\& = \int _a^b L(x,g _1(x))\, dx - \int _a^b L(x,g _2(x))\, dx .

\end{aligned}\end{align}$$



Combining (3) with (4), we get (1) for regions of type I. A similar treatment yields (2) for regions of type II. Putting the two together, we get the result for regions of type III.



\paragraph{Proof for rectifiable Jordan curves}



We are going to prove the following



\textbf{Theorem.} Let $\Gamma$ be a rectifiable, positively oriented Jordan curve in $\R^2$ and let $R$ denote its inner region. Suppose that $A, B:\overline{R} \to \R$ are continuous functions with the property that $A$ has second partial derivative at every point of $R$, $B$ has first partial derivative at every point of $R$ and that the functions $D _1 B, D _2 A : R \to \R$ are Riemann-integrable over $R$. Then



$$\int _{\Gamma}A\,dx+B\,dy = \int _R \left(D _1 B(x,y)-D _2 A(x,y)\right)\,d(x,y).$$



We need the following lemmas whose proofs can be found in:



\textbf{Lemma 1 (Decomposition Lemma).} Assume $\Gamma$ is a rectifiable, positively oriented Jordan curve in the plane and let $R$ be its inner region. For every positive real $\delta$, let $\mathcal{F}(\delta)$ denote the collection of squares in the plane bounded by the lines $x=m\delta, y=m\delta$, where $m$ runs through the set of integers. Then, for this $\delta$, there exists a decomposition of $\overline{R}$ into a finite number of non-overlapping subregions in such a manner that



1.  Each one of the subregions contained in $R$, say $R _1, R _2,\ldots, R _k$, is a square from $\mathcal{F}(\delta)$.

2.  Each one of the remaining subregions, say $R _{k+1},\ldots,R _s$, has as boundary a rectifiable Jordan curve formed by a finite number of arcs of $\Gamma$ and parts of the sides of some square from $\mathcal{F}(\delta)$.

3.  Each one of the border regions $R _{k+1},\ldots,R _s$ can be enclosed in a square of edge-length $2\delta$.

4.  If $\Gamma _i$ is the positively oriented boundary curve of $R _i$, then $\Gamma=\Gamma _1+\Gamma _2+\cdots+\Gamma _s.$

5.  The number $s-k$ of border regions is no greater than $4 \left(\frac{\Lambda}{\delta} + 1\right)$, where $\Lambda$ is the length of $\Gamma$.



\textbf{Lemma 2.} Let $\Gamma$ be a rectifiable curve in the plane and let $\Delta _\Gamma(h)$ be the set of points in the plane whose distance from (the range of) $\Gamma$ is at most $h$. The outer Jordan content of this set satisfies $\overline{c}\,\,\Delta _\Gamma(h)\le2h\Lambda +\pi h^2$.



\textbf{Lemma 3.} Let $\Gamma$ be a rectifiable curve in $\R^2$ and let $f:\text{range of }\Gamma \to \R$ be a continuous function. Then



$$\left\vert\int _\Gamma f(x,y)\,dy\right\vert \text{ and } \left\vert\int _\Gamma f(x,y)\,dx\right\vert \text{ are } \le\frac{1}{2} \Lambda\Omega _f,$$

where $\Omega _f$ is the oscillation of $f$ on the range of $\Gamma$.



Now we are in position to prove the Theorem:



\textbf{Proof of Theorem.} Let $\varepsilon$ be an arbitrary positive real number. By continuity of $A$, $B$ and compactness of $\overline{R}$, given $\varepsilon>0$, there exists $0<\delta 0.$$



We may as well choose $\delta$ so that the RHS of the last inequality is $<\varepsilon.$



The remark in the beginning of this proof implies that the oscillations of $A$ and $B$ on every border region is at most $\varepsilon$. We have



$$\left\vert\sum _{i=k+1}^s \int _{\Gamma _i}A\,dx+B\,dy\right\vert\le\frac{1}{2} \varepsilon\sum _{i=k+1}^s \Lambda _i.$$



By Lemma 1(iii),



$$\sum _{i=k+1}^s \Lambda _i \le\Lambda + (4\delta)\,4 \left(\frac{\Lambda}{\delta}+1\right)\le17\Lambda+16.$$



Combining these, we finally get



$$\left\vert\int _\Gamma A\,dx+B\,dy\quad-\int _R \varphi\right\vert< C \varepsilon,$$



for some $C > 0$. Since this is true for every $\varepsilon > 0$, we are done.



Validity under different hypotheses



The hypothesis of the last theorem are not the only ones under which Green's formula is true. Another common set of conditions is the following:



The functions $A, B:\overline{R} \to \R$ are still assumed to be continuous. However, we now require them to be Fréchet-differentiable at every point of $R$. This implies the existence of all directional derivatives, in particular $D _{e _i}A=:D _i A, D _{e _i}B=:D _i B, \,i=1,2$, where, as usual, $(e _1,e _2)$ is the canonical ordered basis of $\R^2$. In addition, we require the function $D _1 B-D _2 A$ to be Riemann-integrable over $R$.



As a corollary of this, we get the Cauchy Integral Theorem for rectifiable Jordan curves:



**Theorem (Cauchy).** If $\Gamma$ is a rectifiable Jordan curve in $\Complex$ and if $f:\text{closure of inner region of } \Gamma \to \Complex$ is a continuous mapping holomorphic throughout the inner region of $\Gamma$, then



> $\int _\Gamma f=0,$



the integral being a complex contour integral.



**Proof.** We regard the complex plane as $\R^2$. Now, define $u,v:\overline{R} \to \R$ to be such that $f(x+iy)=u(x,y)+iv(x,y).$ These functions are clearly continuous. It is well known that $u$ and $v$ are Fréchet-differentiable and that they satisfy the Cauchy-Riemann equations: $D _1 v+D _2 u=D _1 u-D _2 v =\text{zero function}$.



Now, analyzing the sums used to define the complex contour integral in question, it is easy to realize that



$$\int _\Gamma f=\int _\Gamma u\,dx-v\,dy\quad+i\int _\Gamma v\,dx+u\,dy,$$



the integrals on the RHS being usual line integrals. These remarks allow us to apply Green's Theorem to each one of these line integrals, finishing the proof.



Multiply-connected regions



**Theorem.** Let $\Gamma _0,\Gamma _1,\ldots,\Gamma _n$ be positively oriented rectifiable Jordan curves in $\R^{2}$ satisfying



$$\begin{aligned}

    \Gamma _i \subset R _0,\&\&\text{if } 1\le i\le n \\\    \Gamma _i \subset \R^2 \setminus \overline{R} _j,\&\&\text{if }1\le i,j \le n\text{ and }i\ne j,

    \end{aligned}$$



where $R _i$ is the inner region of $\Gamma _i$. Let



$$D = R _0 \setminus (\overline{R} _1 \cup \overline{R} _2 \cup \cdots \cup \overline{R} _n).$$



Suppose $p: \overline{D} \to \R$ and $q: \overline{D} \to \R$ are continuous functions whose restriction to $D$ is Fréchet-differentiable. If the function



$$(x,y)\longmapsto\frac{\partial q}{\partial e _1}(x,y)-\frac{\partial p}{\partial e _2}(x,y)$$



is Riemann-integrable over $D$, then



$\begin{aligned}

    & \int _{\Gamma _0} p(x,y)\,dx+q(x,y)\,dy-\sum _{i=1}^n \int _{\Gamma _i} p(x,y)\,dx + q(x,y)\,dy \\\\[5pt]

    = {} & \int _D\left\\{\frac{\partial q}{\partial e _1}(x,y)-\frac{\partial p}{\partial e _2}(x,y)\right\\} \, d(x,y).

    \end{aligned}$



Relationship to Stokes' theorem



Green's theorem is a special case of the Kelvin--Stokes theorem, when applied to a region in the $xy$-plane.



We can augment the two-dimensional field into a three-dimensional field with a *z* component that is always 0. Write **F** for the vector-valued function $\mathbf{F}=(L,M,0)$. Start with the left side of Green's theorem:



$$\oint _C (L\, dx + M\, dy) = \oint _C (L, M, 0) \cdot (dx, dy, dz) = \oint _C \mathbf{F} \cdot d\mathbf{r}.$$



The Kelvin--Stokes theorem:



$$\oint _C \mathbf{F} \cdot d\mathbf{r} = \iint _S \nabla \times \mathbf{F} \cdot \mathbf{\hat n} \, dS.$$



The surface $S$ is just the region in the plane $D$, with the unit normal $\mathbf{\hat n}$ defined (by convention) to have a positive z component in order to match the "positive orientation" definitions for both theorems.



The expression inside the integral becomes



$$\nabla \times \mathbf{F} \cdot \mathbf{\hat n} = \left[ \left(\frac{\partial 0}{\partial y}  - \frac{\partial M}{\partial z}\right) \mathbf{i} + \left(\frac{\partial L}{\partial z} - \frac{\partial 0}{\partial x}\right) \mathbf{j} + \left(\frac{\partial M}{\partial x} - \frac{\partial L}{\partial y}\right) \mathbf{k} \right] \cdot \mathbf{k} = \left(\frac{\partial M}{\partial x} - \frac{\partial L}{\partial y}\right).$$



Thus we get the right side of Green's theorem



$$\iint _S \nabla \times \mathbf{F} \cdot \mathbf{\hat n} \, dS = \iint _D \left(\frac{\partial M}{\partial x} - \frac{\partial L}{\partial y}\right) \, dA.$$



Green's theorem is also a straightforward result of the general Stokes' theorem using differential forms and exterior derivatives:



$$\begin{aligned}

    \& \oint _C L \,dx + M \,dy = \oint _{\partial D} \omega = \int _D \,d\omega \\\\[5pt]

    = {} \& \int _D \frac{\partial L}{\partial y} \,dy \wedge \,dx + \frac{\partial M}{\partial x} \,dx \wedge \,dy = \iint _D \left(\frac{\partial M}{\partial x} - \frac{\partial L}{\partial y} \right) \,dx \,dy.

\end{aligned}$$





Relationship to the divergence theorem



Considering only two-dimensional vector fields, Green's theorem is equivalent to the two-dimensional version of the divergence theorem:



$$\iiint _V\left(\mathbf{\nabla}\cdot\mathbf{F}\right)\,dV= \oiint _{\partial \scriptstyle V} (\mathbf{F}\cdot\mathbf{\hat n})\,dS .$$



where $\nabla\cdot\mathbf{F}$ is the divergence on the two-dimensional vector field $\mathbf{F}$, and $\mathbf{\hat n}$ is the outward-pointing unit normal vector on the boundary.



To see this, consider the unit normal $\mathbf{\hat n}$ in the right side of the equation. Since in Green's theorem $d\mathbf{r} = (dx, dy)$ is a vector pointing tangential along the curve, and the curve *C* is the positively oriented (i.e. anticlockwise) curve along the boundary, an outward normal would be a vector which points 90° to the right of this; one choice would be $(dy, -dx)$. The length of this vector is $\sqrt{dx^2 + dy^2} = ds.$ So $(dy, -dx) = \mathbf{\hat n}\,ds.$



Start with the left side of Green's theorem:



$$\oint _C (L\, dx + M\, dy) = \oint _C (M, -L) \cdot (dy, -dx) = \oint _C (M, -L) \cdot \mathbf{\hat n}\,ds.$$ Applying the two-dimensional divergence theorem with $\mathbf{F} = (M, -L)$, we get the right side of Green's theorem:



$$\oint _C (M, -L) \cdot \mathbf{\hat n}\,ds = \iint _D\left(\nabla \cdot (M, -L) \right) \, dA = \iint _D \left(\frac{\partial M}{\partial x} - \frac{\partial L}{\partial y}\right) \, dA.$$



Area calculation



Green's theorem can be used to compute area by line integral. The area of a planar region $D$ is given by



$$A = \iint _D dA.$$



Choose $L$ and $M$ such that $\frac{\partial M}{\partial x} - \frac{\partial L}{\partial y} = 1$, the area is given by



$$A = \oint _{C} (L\, dx + M\, dy).$$



Possible formulas for the area of $D$ include



$$A=\oint _C x\, dy = -\oint _C y\, dx = \tfrac 12 \oint _C (-y\, dx + x\, dy).$$



\paragraph{Resources}



\begin{itemize}
    \item \href{https://www.youtube.com/watch?v=sDby9VzPoj4}{Green's Theorem Part 1} by James Sousa, Math is Power 4U
    \item \href{https://www.youtube.com/watch?v=oXISOEqd0X4}{Green's Theorem Part 2} by James Sousa, Math is Power 4U
    \item \href{https://www.youtube.com/watch?v=eoS8mNyIJYo}{Evaluate a Line Integral Using Green's Theorem} by James Sousa, Math is Power 4U
    \item \href{https://www.youtube.com/watch?v=UhCXd204PL4}{Use Green's Theorem to Evaluate a Line Integral on a Rectangle} by James Sousa, Math is Power 4U
    \item \href{https://www.youtube.com/watch?v=Ak_09_hZ6GQ}{Use Green's Theorem to Evaluate a Line Integral of a Vector Field on a Circle} by James Sousa, Math is Power 4U
    \item \href{https://www.youtube.com/watch?v=SwMelGrLBUU}{Use Green's Theorem to Evaluate a Line Integral Using Polar Coordinates} by James Sousa, Math is Power 4U
    \item \href{https://www.youtube.com/watch?v=bdxFPZbF53E}{Use Green's Theorem to Evaluate a Line Integral with Negative Orientation} by James Sousa, Math is Power 4U
    \item \href{https://www.youtube.com/watch?v=R5-egVQ9-X0}{Use Green's Theorem to Determine the Area of a Region Enclosed by a Curve} by James Sousa, Math is Power 4U
    \item \href{https://www.youtube.com/watch?v=WmGa-OpDP80}{Determining Area Using Line Integrals} by James Sousa, Math is Power 4U
    \item \href{https://www.youtube.com/watch?v=0jKqk4pvEk8}{Flux Form of Green's Theorem} by James Sousa, Math is Power 4U
    \item \href{https://www.youtube.com/watch?v=ccW3F8M1BLY}{Determine the Flux of a 2D Vector Field Across a Rectangle Using Green's Theorem} by James Sousa, Math is Power 4U
    \item \href{https://www.youtube.com/watch?v=FB-JmM_dbBQ}{Determine the Flux of a 2D Vector Field Across a Parabolic Region Using Green's Theorem} by James Sousa, Math is Power 4U
    \item \href{https://www.youtube.com/watch?v=vpfz1-whGW4}{Determine the Flux of a 2D Vector Field Using Green's Theorem (Hole)} by James Sousa, Math is Power 4U
    \item \href{https://www.youtube.com/watch?v=a_zdFvYXX_c}{Green's Theorem} by Patrick JMT
    \item \href{https://www.youtube.com/watch?v=FJvW3dGkVVk}{Green's Theorem (One Region)} by Krista King
    \item \href{https://www.youtube.com/watch?v=8TilBgxQe80}{Green's Theorem (Two Regions)} by Krista King
    \item \href{https://www.youtube.com/watch?v=gGXnILbrhsM}{Green's Theorem Example 1} by Khan Academy
    \item \href{https://www.youtube.com/watch?v=sSyPAAyL8nQ}{Green's Theorem Example 2} by Khan Academy
\end{itemize}

\paragraph{Licensing}



Content obtained and/or adapted from:



\begin{itemize}
    \item \href{https://en.wikipedia.org/wiki/Green\%27s_theorem}{Green's Theorem, Wikipedia} under a CC BY-SA license
\end{itemize}

\subsection{Learning Objectives}

\begin{enumerate}
    \item Here are three Student Learning Objectives (SLOs) related to Green's theorem in vector calculus:
    \item \textbf{Apply Green's Theorem to Compute Line Integrals:}
    \item \textbf{Outcome:} Students will be able to apply Green's theorem to convert a line integral around a simple closed curve $C$ into a double integral over the region $D$ bounded by $C$.
    \item \textbf{Details:} This includes selecting appropriate functions $L(x, y)$ and $M(x, y)$ for a given vector field and performing the necessary calculations to evaluate the double integral, thereby determining the value of the line integral.
    \item \textbf{Demonstrate the Proof of Green's Theorem for Simple Regions:}
    \item \textbf{Outcome:} Students will be able to demonstrate the proof of Green's theorem for regions of type I and type II, and explain how the theorem generalizes to type III regions.
    \item \textbf{Details:} This includes understanding and reconstructing the steps involved in proving the theorem for specific regions, such as using parametric equations for the curve $C$ and computing the necessary line and double integrals.
    \item \textbf{Interpret Physical Applications of Green's Theorem:}
    \item \textbf{Outcome:} Students will be able to interpret and apply Green's theorem in physical contexts, such as calculating fluid flow across a boundary or determining the area of a planar region.
    \item \textbf{Details:} This includes analyzing real-world problems that can be modeled using Green's theorem, setting up the appropriate integrals, and interpreting the results in the context of the physical situation, such as determining the total outflow in a fluid dynamics scenario.
\end{enumerate}

\subsection{Assessment Instrument}

\begin{enumerate}
    \item \textbf{Assessment Instrument: Green's Theorem Quiz} \textbf{Format:}
    \item \textbf{Type:} Mixed-format quiz (multiple-choice, short answer, and problem-solving questions)
    \item \textbf{Duration:} 45 minutes
    \item \textbf{Total Marks:} 30 points \textbf{Sections:}
    \item \textbf{Understanding Green's Theorem (10 points):}
    \item \textbf{Question 1:} (Multiple-Choice) Which of the following correctly represents Green's theorem?
    \item a. $$\oint _{\scriptstyle C} (L\, dx + M\, dy) = \iint _{D} \left(\frac{\partial L}{\partial x} + \frac{\partial M}{\partial y}\right) dx\, dy$$
    \item b. $$\oint _{\scriptstyle C} (L\, dx + M\, dy) = \iint _{D} \left(\frac{\partial M}{\partial x} - \frac{\partial L}{\partial y}\right) dx\, dy$$
    \item c. $$\oint _{\scriptstyle C} (L\, dx + M\, dy) = \iint _{D} \left(\frac{\partial M}{\partial y} - \frac{\partial L}{\partial x}\right) dx\, dy$$
    \item d. $$\oint _{\scriptstyle C} (L\, dx + M\, dy) = \iint _{D} \left(\frac{\partial L}{\partial y} - \frac{\partial M}{\partial x}\right) dx\, dy$$
    \item \textbf{Question 2:} (Short Answer) Briefly explain the relationship between the line integral around a curve $C$ and the double integral over the region $D$ as stated in Green's theorem.
    \item \textbf{Application of Green's Theorem (10 points):}
    \item \textbf{Question 3:} (Problem-Solving) Given the vector field $$\mathbf{F} = (y^2, x^2),$$ use Green's theorem to evaluate the line integral $$\oint _{C} y^2 \, dx + x^2 \, dy,$$ where $C$ is the positively oriented boundary of the region enclosed by the curve $x^2 + y^2 = 1$.
    \item \textbf{Question 4:} (Problem-Solving) Determine the area of the region bounded by the ellipse $$\frac{x^2}{4} + \frac{y^2}{9} = 1$$ using Green's theorem. Use an appropriate vector field to simplify the calculation.
    \item \textbf{Proof and Physical Interpretation (10 points):}
    \item \textbf{Question 5:} (Short Answer) Prove Green's theorem for a type I region where the curve $C$ is composed of vertical lines and horizontal curves. Show the steps involved in converting the line integral into a double integral.
    \item \textbf{Question 6:} (Multiple-Choice) In which of the following physical contexts is Green's theorem typically applied?
    \item a. Calculating the mass of a three-dimensional object
    \item b. Determining the flux of a fluid across a surface
    \item c. Evaluating the circulation of a vector field along a closed curve in two dimensions
    \item d. Finding the potential energy in a conservative vector field \textbf{Scoring Rubric:}
    \item \textbf{Understanding Green's Theorem:} 5 points for correctly answering each question (total 10 points)
    \item \textbf{Application of Green's Theorem:} 5 points for each correctly solved problem (total 10 points)
    \item \textbf{Proof and Physical Interpretation:} 5 points for each correct and well-explained response (total 10 points) This quiz will assess students' comprehension of Green's theorem, their ability to apply the theorem to solve problems, and their understanding of its proof and physical applications.
\end{enumerate}

%%===============================================================
\section{Stokes´ Theorem (old)}
\label{sec:stokes-theorem-old}
%%===============================================================

% [Image: An illustration of Stokes' theorem]



Figure: An illustration of Stokes' theorem, with surface $\Sigma$, its boundary $\partial \Sigma$, and the normal vector $\mathbf{n}$.



\textbf{Stokes' theorem}, also known as \textbf{Kelvin–Stokes theorem} after Lord Kelvin and George Stokes, the \textbf{fundamental theorem for curls} or simply the \textbf{curl theorem}, is a theorem in vector calculus on $\mathbb{R}^3$. Given a vector field, the theorem relates the integral of the curl of the vector field over some surface to the line integral of the vector field around the boundary of the surface. The classical Stokes' theorem can be stated in one sentence: The line integral of a vector field over a loop is equal to the \textit{flux of its curl} through the enclosed surface.



Stokes' theorem is a special case of the generalized Stokes' theorem. In particular, a vector field on $\mathbb{R}^3$ can be considered as a 1-form, in which case its curl is its exterior derivative, a 2-form.



\paragraph{Theorem}



Let $\Sigma$ be a smooth oriented surface in $\mathbb{R}^3$ with boundary $\partial \Sigma$. If a vector field

$$\mathbf{A} = (P(x, y, z), Q(x, y, z), R(x, y, z))$$

is defined and has continuous first-order partial derivatives in a region containing $\Sigma$, then



$$\iint _\Sigma (\nabla \times \mathbf{A}) \cdot \mathrm{d}\mathbf{a}  =  \oint _{\partial\Sigma} \mathbf{A} \cdot

\mathrm{d}\mathbf{l}.$$



More explicitly, the equality says that

$$\begin{aligned} &\iint _\Sigma \left(\left(\frac{\partial R}{\partial y}-\frac{\partial Q}{\partial z} \right)\,\mathrm{d}y\, \mathrm{d}z +\left(\frac{\partial P}{\partial z}-\frac{\partial R}{\partial x}\right)\, \mathrm{d}z\, \mathrm{d}x  +\left (\frac{\partial Q}{\partial x}-\frac{\partial P}{\partial y}\right)\, \mathrm{d}x\, \mathrm{d}y\right) \\\\ &= \oint _{\partial\Sigma} \Bigl(P\, \mathrm{d}x+Q\, \mathrm{d}y+R\, \mathrm{d}z\Bigr). \end{aligned}$$



The main challenge in a precise statement of Stokes' theorem is in defining the notion of a boundary. Surfaces such as the Koch snowflake, for example, are well-known not to exhibit a Riemann-integrable boundary, and the notion of surface measure in Lebesgue theory cannot be defined for a non-Lipschitz surface. One (advanced) technique is to pass to a weak formulation and then apply the machinery of geometric measure theory; for that approach see the coarea formula. In this article, we instead use a more elementary definition, based on the fact that a boundary can be discerned for full-dimensional subsets of $\mathbb{R}^2$.



Let $\gamma: [a, b] \to \mathbb{R}^2$ be a piecewise smooth Jordan plane curve. The Jordan curve theorem implies that $\gamma$ divides $\mathbb{R}^2$ into two components, a compact one and another that is non-compact. Let $D$ denote the compact part; then $D$ is bounded by $\gamma$. It now suffices to transfer this notion of boundary along a continuous map to our surface in $\mathbb{R}^3$. But we already have such a map: the parametrization of $\Sigma$.



Suppose $\psi: D \to \mathbb{R}^3$ is smooth, with $\Sigma = \psi(D)$. If $\Gamma$ is the space curve defined by $\Gamma(t) = \psi(\gamma(t))$, then we call $\Gamma$ the boundary of $\Sigma$, written $\partial \Sigma$.



With the above notation, if $\mathbf{F}$ is any smooth vector field on $\mathbb{R}^3$, then



$$\oint _{\partial\Sigma} \mathbf{F}\, \cdot\, \mathrm{d}{\mathbf{\Gamma}}  = \iint _{\Sigma} \nabla\times\mathbf{F}\, \cdot\, \mathrm{d}\mathbf{S}.$$



\paragraph{Proof}



The proof of the theorem consists of 4 steps. We assume Green's theorem, so what is of concern is how to boil down the three-dimensional complicated problem (Stokes' theorem) to a two-dimensional rudimentary problem (Green's theorem). When proving this theorem, mathematicians normally deduce it as a special case of a more general result, which is stated in terms of differential forms, and proved using more sophisticated machinery. While powerful, these techniques require substantial background, so the proof below avoids them, and does not presuppose any knowledge beyond a familiarity with basic vector calculus. At the end of this section, a short alternate proof of Stokes' theorem is given, as a corollary of the generalized Stokes' Theorem.



\paragraph{Elementary proof}



\paragraph{First step of the proof (parametrization of integral)}



As in § Theorem, we reduce the dimension by using the natural parametrization of the surface. Let $\boldsymbol{\psi}$ and $\gamma$ be as in that section, and note that by change of variables



$$\oint _{\partial\Sigma}{\mathbf{F}(\mathbf{x})\cdot\,\mathrm{d}\mathbf{l}}

= \oint _{\gamma}{\mathbf{F}(\boldsymbol{\psi}(\mathbf{y}))\cdot\,\mathrm{d}\boldsymbol{\psi}(\mathbf{y})}

= \oint _{\gamma}{\mathbf{F}(\boldsymbol{\psi}(\mathbf{y}))J _{\mathbf{y}}(\boldsymbol{\psi})\,\mathrm{d}\mathbf{y}}$$



where $J _\psi$ stands for the Jacobian matrix of $\psi$.



Now let $\{\mathbf{e} _u, \mathbf{e} _v\}$ be an orthonormal basis in the coordinate directions of $\mathbb{R} _2$. Recognizing that the columns of $J _{\mathbf{y}}\boldsymbol{\psi}$ are precisely the partial derivatives of $\boldsymbol{\psi}$ at $\mathbf{y}$, we can expand the previous equation in coordinates as



$$\begin{aligned} \oint _{\partial\Sigma}{\mathbf{F}(\mathbf{x})\cdot\,\mathrm{d}\mathbf{l}} &= \oint _{\gamma}{\mathbf{F}(\boldsymbol{\psi}(\mathbf{y}))J _{\mathbf{y}}(\boldsymbol{\psi})\mathbf{e} _u(\mathbf{e} _u\cdot\,\mathrm{d}\mathbf{y}) + \mathbf{F}(\boldsymbol{\psi}(\mathbf{y}))J _{\mathbf{y}}(\boldsymbol{\psi})\mathbf{e} _v(\mathbf{e} _v\cdot\,\mathrm{d}\mathbf{y})} \\\\ &=\oint _{\gamma}{\left(\left(\mathbf{F}(\boldsymbol{\psi}(\mathbf{y}))\cdot\frac{\partial\boldsymbol{\psi}}{\partial u}(\mathbf{y})\right)\mathbf{e} _u + \left(\mathbf{F}(\boldsymbol{\psi}(\mathbf{y}))\cdot\frac{\partial\boldsymbol{\psi}}{\partial v}(\mathbf{y})\right)\mathbf{e} _v\right)\cdot\,\mathrm{d}\mathbf{y}} \end{aligned}$$



\paragraph{Second step in the proof (defining the pullback)}



The previous step suggests we define the function



$$\mathbf{P}(u,v) = \left(\mathbf{F}(\boldsymbol{\psi}(u,v))\cdot\frac{\partial\boldsymbol{\psi}}{\partial u}(u,v)\right)\mathbf{e} _u + \left(\mathbf{F}(\boldsymbol{\psi}(u,v))\cdot\frac{\partial\boldsymbol{\psi}}{\partial v} \right)\mathbf{e} _v$$



This is the pullback of $\mathbf{F}$ along $\boldsymbol{\psi}$, and, by the above, it satisfies



$$\oint _{\partial\Sigma}{\mathbf{F}(\mathbf{x})\cdot\,\mathrm{d}\mathbf{l}}=\oint _{\gamma}{\mathbf{P}(\mathbf{y})\cdot\,\mathrm{d}\mathbf{l}}$$



We have successfully reduced one side of Stokes' theorem to a 2-dimensional formula; we now turn to the other side.



\paragraph{Third step of the proof (second equation)}



First, calculate the partial derivatives appearing in Green's theorem, via the product rule:



$$\begin{aligned} \frac{\partial P _1}{\partial v} &= \frac{\partial (\mathbf{F}\circ \boldsymbol{\psi})}{\partial v}\cdot\frac{\partial \boldsymbol\psi}{\partial u} + (\mathbf{F}\circ \boldsymbol\psi) \cdot\frac{\partial^2 \boldsymbol\psi}{\partial v \, \partial u} \\\\[5pt] \frac{\partial P _2}{\partial u} &= \frac{\partial (\mathbf{F}\circ \boldsymbol{\psi})}{\partial u}\cdot\frac{\partial \boldsymbol\psi}{\partial v} + (\mathbf{F}\circ \boldsymbol\psi) \cdot\frac{\partial^2 \boldsymbol\psi}{\partial u \, \partial v} \end{aligned}$$



Conveniently, the second term vanishes in the difference, by equality of mixed partials. So,



$$\begin{aligned} \frac{\partial P _1}{\partial v} - \frac{\partial P _2}{\partial u} &= \frac{\partial (\mathbf{F}\circ \boldsymbol\psi)}{\partial v}\cdot\frac{\partial \boldsymbol\psi}{\partial u} - \frac{\partial (\mathbf{F}\circ \boldsymbol\psi)}{\partial u}\cdot\frac{\partial \boldsymbol\psi}{\partial v} \\\\[5pt] &= \frac{\partial \boldsymbol\psi}{\partial u}(J _{\boldsymbol\psi(u,v)}\mathbf{F})\frac{\partial \boldsymbol\psi}{\partial v} - \frac{\partial \boldsymbol\psi}{\partial v}(J _{\boldsymbol\psi(u,v)}\mathbf{F})\frac{\partial \boldsymbol\psi}{\partial u} && \text{(chain rule)} \\\\[5pt] &= \frac{\partial \boldsymbol\psi}{\partial u}\left(J _{\boldsymbol\psi(u,v)}\mathbf{F}-{(J _{\boldsymbol\psi(u,v)}\mathbf{F})}^{\mathsf{T}}\right)\frac{\partial \boldsymbol\psi}{\partial v} \end{aligned}$$



But now consider the matrix in that quadratic form — that is, $J _{\boldsymbol\psi(u,v)}\mathbf{F} - (J _{\boldsymbol\psi(u,v)}\mathbf{F})^{\mathsf{T}}$. We claim this matrix in fact describes a cross product.



To be precise, let $A=(A _{ij}) _{ij}$ be an arbitrary $3 \times 3$ matrix and let



$\mathbf{a} = \begin{bmatrix}A _{32}-A _{23} \\\\ A _{13}-A _{31} \\\\ A _{21}-A _{12}\end{bmatrix}$



Note that $\mathbf{x} \mapsto \mathbf{a} \times \mathbf{x}$ is linear, so it is determined by its action on basis elements. But by direct calculation

$$\begin{aligned} \left(A-A^{\mathsf{T}}\right)\mathbf{e} _1 &= \begin{bmatrix} 0 \\\\ a _3 \\\\ -a _2 \end{bmatrix} = \mathbf{a}\times\mathbf{e} _1 \\\\ \left(A-A^{\mathsf{T}}\right)\mathbf{e} _2 &= \begin{bmatrix} -a _3 \\\\ 0 \\\\ a _1 \end{bmatrix} = \mathbf{a}\times\mathbf{e} _2 \\\\ \left(A-A^{\mathsf{T}}\right)\mathbf{e} _3 &= \begin{bmatrix} a _2 \\\\ -a _1 \\\\ 0 \end{bmatrix} = \mathbf{a}\times\mathbf{e} _3 \end{aligned}$$



Thus $(A - A^{\mathsf{T}})\mathbf{x} = \mathbf{a} \times \mathbf{x}$ for any $\mathbf{x}$. Substituting $J \mathbf{F}$ for $A$, we obtain



$$\left({(J _{\boldsymbol\psi(u,v)}\mathbf{F})} _{\psi(u,v)} - {(J _{\boldsymbol\psi(u,v)}\mathbf{F})}^{\mathsf{T}} \right) \mathbf{x} =(\nabla\times\mathbf{F})\times \mathbf{x}, \quad \text{for all}\, \mathbf{x}\in\R^{3}$$



We can now recognize the difference of partials as a (scalar) triple product:



$$\begin{aligned} \frac{\partial P _1}{\partial v} - \frac{\partial P _2}{\partial u} &= \frac{\partial \boldsymbol\psi}{\partial u}\cdot(\nabla\times\mathbf{F}) \times \frac{\partial \boldsymbol\psi}{\partial v} \\\\ &= \det \begin{bmatrix} (\nabla\times\mathbf{F})(\boldsymbol\psi(u,v)) & \frac{\partial \boldsymbol\psi}{\partial u}(u,v) & \frac{\partial \boldsymbol\psi}{\partial v}(u,v) \end{bmatrix} \end{aligned}$$



On the other hand, the definition of a surface integral also includes a triple product — the very same one!



$$\begin{aligned} \iint _S (\nabla\times\mathbf{F})\cdot \, d^2\mathbf{S} &=\iint _D {(\nabla\times\mathbf{F})(\boldsymbol\psi(u,v))\cdot\left(\frac{\partial \boldsymbol\psi}{\partial u}(u,v)\times \frac{\partial \boldsymbol\psi}{\partial v}(u,v)\,\mathrm{d}u\,\mathrm{d}v\right)} \\\\ &= \iint _D \det \begin{bmatrix} (\nabla\times\mathbf{F})(\boldsymbol\psi(u,v)) & \frac{\partial \boldsymbol\psi}{\partial u}(u,v) & \frac{\partial \boldsymbol\psi}{\partial v}(u,v) \end{bmatrix} \,\mathrm{d}u \,\mathrm{d}v \end{aligned}$$



So, we obtain



$$\iint _S (\nabla\times\mathbf{F})\cdot \,\mathrm{d}^2\mathbf{S} = \iint _D \left( \frac{\partial P _2}{\partial u} - \frac{\partial P _1}{\partial v} \right) \,\mathrm{d}u\,\mathrm{d}v$$



\paragraph{Fourth step of the proof (reduction to Green's theorem)}



Combining the second and third steps, and then applying Green's theorem completes the proof.



\paragraph{Proof via differential forms}



$\mathbb{R} \to \mathbb{R}^3$



can be identified with the differential 1-forms on $\mathbb{R}^3$ via the map



$F _1\mathbf{e} _1+F _2\mathbf{e} _2+F _3\mathbf{e} _3 \mapsto F _1\,\mathrm{d}x+F _2\,\mathrm{d}y+F _3\mathrm{d}z .$



Write the differential 1-form associated with a function $\mathbf{F}$ as $\omega _{\mathbf{F}}$. Then one can calculate that



$$\star\omega _{\nabla\times\mathbf{F}}=\mathrm{d}\omega _{\mathbf{F}}$$



where $\star$ is the Hodge star and $\mathrm{d}$ is the exterior derivative. Thus, by the generalized Stokes' theorem,



$$\oint _{\partial\Sigma}{\mathbf{F}\cdot\,\mathrm{d}\mathbf{l}}

=\oint _{\partial\Sigma}{\omega _{\mathbf{F}}}

=\int _{\Sigma}{\mathrm{d}\omega _{\mathbf{F}}}

=\int _{\Sigma}{\star\omega _{\nabla\times\mathbf{F}}}

=\iint _{\Sigma}{\nabla\times\mathbf{F}\cdot\,\mathrm{d}^2\mathbf{S}}$$



\paragraph{Licensing}



Content obtained and/or adapted from:



\begin{itemize}
    \item \href{https://en.wikipedia.org/wiki/Stokes\%27_theorem}{Stokes' Theorem} under a CC BY-SA license
\end{itemize}

\subsection{Learning Objectives}

\begin{enumerate}
    \item Here are three Student Learning Objectives (SLOs) related to Stokes' theorem:
    \item \textbf{Understand Stokes' Theorem:}
    \item \textbf{Outcome:} Students will be able to state and interpret Stokes' theorem in the context of vector calculus.
    \item \textbf{Details:} This includes recognizing the relationship between the curl of a vector field and the circulation of the field around the boundary of a surface, and understanding the geometric and physical implications of the theorem.
    \item \textbf{Apply Stokes' Theorem to Compute Integrals:}
    \item \textbf{Outcome:} Students will be able to apply Stokes' theorem to compute the surface integral of the curl of a vector field over a surface and relate it to the line integral of the vector field around the boundary of the surface.
    \item \textbf{Details:} This includes choosing appropriate surfaces and boundaries, setting up and evaluating both the surface and line integrals, and verifying that the two integrals are equal as per Stokes' theorem.
    \item \textbf{Parametrize Surfaces and Boundaries:}
    \item \textbf{Outcome:} Students will be able to parametrize surfaces and their boundaries in three-dimensional space to facilitate the application of Stokes' theorem.
    \item \textbf{Details:} This includes constructing appropriate parameterizations for given surfaces, identifying the correct orientation of the surface and its boundary, and using these parameterizations to set up and evaluate integrals according to the theorem.
\end{enumerate}

\subsection{Assessment Instrument}

\begin{enumerate}
    \item \textbf{Assessment Instrument: Stokes' Theorem Quiz} \textbf{Format:}
    \item \textbf{Type:} Mixed-format quiz (multiple-choice, short answer, and problem-solving questions)
    \item \textbf{Duration:} 45 minutes
    \item \textbf{Total Marks:} 30 points \textbf{Sections:}
    \item \textbf{Understanding Stokes' Theorem (10 points):}
    \item \textbf{Question 1:} (Multiple Choice) Which of the following best describes Stokes' theorem?
    \item a. It relates the surface integral of the divergence of a vector field over a surface to the volume integral over the region enclosed by the surface.
    \item b. It relates the line integral of a vector field around a closed loop to the flux of its curl through the surface bounded by the loop.
    \item c. It equates the line integral of a scalar function to the surface integral of its gradient.
    \item d. It states that the total flux of a vector field through a surface is equal to the circulation around the boundary of the surface.
    \item \textbf{Question 2:} (Short Answer) Briefly explain the geometric interpretation of Stokes' theorem.
    \item \textbf{Question 3:} (Problem-Solving) Given a vector field $\mathbf{A} = (x^2, y^2, z^2)$ and the surface $\Sigma$ defined by $z = 1 - x^2 - y^2$ in the first octant, write the expression for the curl of $\mathbf{A}$ and set up the surface integral according to Stokes' theorem.
    \item \textbf{Applying Stokes' Theorem (10 points):}
    \item \textbf{Question 4:} (Problem-Solving) Consider the vector field $\mathbf{F} = (yz, xz, xy)$. Use Stokes' theorem to evaluate the line integral $\oint_{\partial \Sigma} \mathbf{F} \cdot d\mathbf{l}$, where $\Sigma$ is the upper half of the unit sphere, and $\partial \Sigma$ is the boundary circle in the $xy$-plane.
    \item \textbf{Question 5:} (Multiple Choice) Which of the following is a necessary condition for applying Stokes' theorem?
    \item a. The vector field must be conservative.
    \item b. The surface must be closed.
    \item c. The vector field must have continuous first-order partial derivatives.
    \item d. The surface must be flat.
    \item \textbf{Parametrizing Surfaces and Boundaries (10 points):}
    \item \textbf{Question 6:} (Problem-Solving) Parametrize the surface $\Sigma$ of a cylinder given by $x^2 + y^2 = 1$ and $0 \leq z \leq 2$, and its boundary $\partial \Sigma$ at $z = 2$. Set up the integrals required to apply Stokes' theorem for a given vector field $\mathbf{G} = (2z, -x, y)$.
    \item \textbf{Question 7:} (Short Answer) Explain how the orientation of a surface and its boundary affects the application of Stokes' theorem. \textbf{Scoring Rubric:}
    \item \textbf{Understanding Stokes' Theorem:}
    \item Question 1: 2 points for the correct choice.
    \item Question 2: 3 points for a clear and accurate explanation.
    \item Question 3: 5 points for correctly setting up the curl and the integral.
    \item \textbf{Applying Stokes' Theorem:}
    \item Question 4: 6 points for correctly applying Stokes' theorem and evaluating the line integral.
    \item Question 5: 4 points for the correct choice.
    \item \textbf{Parametrizing Surfaces and Boundaries:}
    \item Question 6: 6 points for correctly parametrizing the surface and boundary and setting up the required integrals.
    \item Question 7: 4 points for a clear and accurate explanation. The quiz assesses students' understanding and application of Stokes' theorem, ensuring they can correctly state, interpret, and apply the theorem to compute integrals and work with parametrized surfaces and boundaries.
\end{enumerate}

%%===============================================================
\section{Taylor´s Formula for Two Variables (old)}
\label{sec:taylor-s-formula-for-two-variables-old}
%%===============================================================

Taylor´s Formula for Two Variables
